% -*- coding: utf-8 -*-
\documentclass[spanish,11pt,a4paper,openany]{book}

% -*- coding: utf-8 -*-
% Preámbulo común para la asignatura Lógica (Grado en Filosofía, Universidad de Sevilla)
% Compatible tanto para compilación independiente (article) como en el manual conjunto (book/report)

\usepackage[T1]{fontenc}
\usepackage[utf8]{inputenc}
\usepackage[spanish,es-tabla]{babel}
\usepackage{amsmath,amssymb,amstext,amsthm,array}
\usepackage[top=20mm,bottom=20mm,left=20mm,right=20mm,headsep=5mm,footskip=10mm]{geometry}
\usepackage{enumerate}
\usepackage{tikz}
\usetikzlibrary{trees}

% Paquete Fitch y reglas canónicas de la asignatura
% Altura de línea algo menor que la predeterminada (4.5ex/3.5ex/1.5ex) para que
% las derivaciones sean más compactas y den menos problemas de paginación.
\usepackage[justsep=1em,rules=myrules,justformat=myjust,height=3.8ex,topheight=3.2ex,depth=1.3ex]{fitch}
\newcommand{\myjust}[2]{#1 #2}
\newcommand{\myrules}{
 \ndrules
 \def\ci{\by{$\to$I}}
 \def\ce{\by{$\to$E}}
 \def\x{\by{X}}
 \def\ip{\by{IP}}
 \def\ds{\by{DS}}
 \def\mt{\by{MT}}
 \def\lem{\by{LEM}}
 \def\dne{\by{DNE}}
 \def\dem{\by{DeM}}
 \def\cq{\by{CQ}}
 \def\bi{\by{$\leftrightarrow$I}}
 \def\be{\by{$\leftrightarrow$E}}
 \def\idi{\by{$=$I}}
 \def\ide{\by{$=$E}}}

% Enlaces hipertexto
\usepackage[hidelinks,colorlinks=true,linkcolor=blue!60!black,urlcolor=blue!60!black]{hyperref}

% Aplicaciones web de apoyo (GitHub Pages del repositorio fersoler/AplicacionesLogica).
% Si cambia la dirección, basta cambiarla aquí.
\newcommand{\urlaplicaciones}{https://fersoler.github.io/AplicacionesLogica/}
\newcommand{\webaplicaciones}{\expandafter\url\expandafter{\urlaplicaciones}}
\newcommand{\enlaceapp}[2]{\href{\urlaplicaciones #1/}{\emph{#2}}}
\newcommand{\appsintaxis}{\enlaceapp{sintaxis}{Sintaxis de las fórmulas}}
\newcommand{\appdeduccion}{\enlaceapp{deduccion}{Deducción natural}}
\newcommand{\apptablas}{\enlaceapp{tablas}{Tablas de verdad}}
\newcommand{\appkalmar}{\enlaceapp{kalmar}{La prueba de Kalmár}}
\newcommand{\appmundo}{\enlaceapp{mundo}{El mundo de Tarski}}
\newcommand{\apptableaux}{\enlaceapp{tableaux}{Tableaux semánticos}}
% Nota breve tras una explicación: con qué aplicación se practica. Empieza con
% el icono de un portátil (fontawesome5) en el color de los enlaces.
\usepackage{fontawesome5}
\newcommand{\practica}[1]{\par\smallskip\noindent{\small{\color{blue!60!black}\faLaptop}\enspace\textit{Para practicar.} #1\par}\smallskip}

% Teoremas y entornos formales
\ifcsname chapter\endcsname
  \newtheorem{theorem}{Teorema}[chapter]
\else
  \newtheorem{theorem}{Teorema}[section]
\fi
\newtheorem{lemma}[theorem]{Lema}
\newtheorem{definition}[theorem]{Definición}
\newtheorem{definicion}[theorem]{Definición}
\newtheorem{example}[theorem]{Ejemplo}
\newtheorem{ejemplo}[theorem]{Ejemplo}
\newtheorem{corollary}[theorem]{Corolario}
\newtheorem{proposition}[theorem]{Proposición}
\newtheorem{observacion}[theorem]{Observación}
\newtheorem{nota}[theorem]{Nota}
\newtheorem{aplicacion}[theorem]{Aplicación}

\newenvironment{resumenbox}{%
  \begin{quote}\small
}{%
  \end{quote}
}

% Ejercicios: en el manual (clase book) se numeran por tema (Ejercicio 4.12).
% El título entre paréntesis va en redonda, como en los teoremas y ejemplos.
\newcounter{preg}
\ifcsname chapter\endcsname
  \makeatletter\@addtoreset{preg}{chapter}\makeatother
  \renewcommand{\thepreg}{\thechapter.\arabic{preg}}
\fi
\newenvironment{pregunta}[1][]{%
  \par\addvspace{0.4cm}\refstepcounter{preg}\noindent\textbf{Ejercicio \thepreg}\if\relax\detokenize{#1}\relax\textbf{.}\else\ (#1)\textbf{.}\fi\enspace\ignorespaces
}{%
  \par\vspace{0.4cm}%
}

% Soluciones (apéndice de soluciones)
% Dentro de una solución no se permite un salto de página entre el título y lo
% que le sigue (por ejemplo, una derivación en un entorno center), ni que un
% párrafo deje una sola línea al pie o al principio de una página.
\newenvironment{solucion}[1]{%
  \par\medskip
  \csname @beginparpenalty\endcsname=10000\relax
  \clubpenalty=10000\relax\widowpenalty=10000\relax
  \noindent\textbf{Solución del ejercicio~\ref{#1}.}\enspace\ignorespaces
}{%
  \par\medskip%
}

% Derivación dentro de un recuadro de regla (tabla con \hline): deja un poco
% de aire para que la línea vertical de la derivación no toque las líneas
% horizontales de la tabla.
\newcommand{\cajaregla}[1]{\vbox{\vskip 5pt\hbox{#1}\vskip 4pt}}

% Marca para secciones de ampliación (no necesarias para seguir el curso)
\newcommand{\ampliacion}{\texorpdfstring{\,$^{\star}$}{ (*)}}

\newcounter{pregexamen}
\newenvironment{preguntaexamen}[2]{%
  \refstepcounter{pregexamen}%
  \vspace{0.4cm}%
  \noindent\textbf{Pregunta \arabic{pregexamen}} \hfill \textbf{[#1 puntos]}\\[0.15cm]%
  \noindent\textit{#2}\\[0.2cm]%
}{%
  \vspace{0.3cm}%
}

\newtheorem{ejercicio}{Ejercicio}

\newcounter{ejerciciocnt}
\newcommand{\ejer}[1]{\vspace{1.5em}\noindent\stepcounter{ejerciciocnt}\textbf{Ejercicio \theejerciciocnt.} #1\vspace{0.5em}}

\newenvironment{resumentema}{\begin{quote}\small\textbf{Resumen.}\enspace}{\end{quote}\bigskip}

% Los capítulos se llaman «temas» en todo el texto: el encabezado dice «Tema 1»
% en lugar de «Capítulo 1» (los apéndices siguen siendo «Apéndice A»...).
\addto\captionsspanish{\renewcommand{\chaptername}{Tema}}

% Las páginas terminan donde termina su contenido, sin estirar los espacios
% verticales cuando una derivación o un árbol no cabe y pasa a la página siguiente.
\raggedbottom

\title{%
  \vspace{-1.2cm}
  {\Huge \textbf{Apuntes de Lógica}}\\[0.35cm]
  {\Large \textbf{Grado en Filosofía}}
}
\author{\textbf{Fernando Soler Toscano}}
\date{Universidad de Sevilla}

\begin{document}

\pagestyle{plain}
\frontmatter

\maketitle
\thispagestyle{empty}
\vspace*{\fill}
\begingroup
\small
\noindent \textbf{Apuntes de Lógica (Grado en Filosofía)}\\
\textbf{Autor:} Fernando Soler Toscano. Dpto. de Filosofía,
Lógica y Filosofía de la Ciencia. Universidad de Sevilla.\\
\textbf{Año:} 2026\\[0.5cm]
\noindent Esta obra está sujeta a una licencia internacional Creative Commons:\\
\textbf{\href{https://creativecommons.org/licenses/by-sa/4.0/}{Atribución-CompartirIgual 4.0 Internacional (CC BY-SA 4.0)}}.\\[0.2cm]
\noindent Enlace a la licencia: \url{https://creativecommons.org/licenses/by-sa/4.0/}\\[0.4cm]
\noindent Usted es libre de:
\begin{itemize}\setlength{\itemsep}{1pt}
    \item \textbf{Compartir:} copiar y redistribuir el material en cualquier medio o formato, para cualquier propósito, incluso comercial.
    \item \textbf{Adaptar:} remezclar, transformar y crear a partir del material, para cualquier propósito, incluso comercial.
\end{itemize}
\noindent Bajo las siguientes condiciones:
\begin{itemize}\setlength{\itemsep}{1pt}
    \item \textbf{Atribución (BY):} debe reconocer adecuadamente la autoría, proporcionar un enlace a la licencia e indicar si se han realizado cambios.
    \item \textbf{CompartirIgual (SA):} si remezcla, transforma o crea a partir del material, debe difundir sus contribuciones bajo la misma licencia que el original.
\end{itemize}
\noindent \textbf{Fuentes.} Los ficheros \LaTeX{} (\texttt{.tex}) con los que se compilan estos apuntes están en \url{https://github.com/fersoler/ApuntesLogica}, con la misma licencia.\\[0.2cm]
\noindent \textbf{Aplicaciones web.} Las que acompañan a estos apuntes (véase «Cómo usar estos apuntes») son software libre, con licencia GNU GPL versión 3. Su código está en \url{https://github.com/fersoler/AplicacionesLogica}.
\endgroup
\vspace*{1.5cm}
\clearpage

% Control fino de profundidad para que el índice sea legible y compacto
\setcounter{tocdepth}{1}
\tableofcontents

\chapter*{Cómo usar estos apuntes}
\addcontentsline{toc}{chapter}{Cómo usar estos apuntes}

Estos apuntes recorren dos veces el mismo ciclo: primero para la
lógica proposicional y después para la lógica de primer orden. En cada
ciclo se empieza por la deducción natural (cómo se \emph{demuestra}
que una conclusión se sigue de unas premisas), se continúa con la
semántica (qué significa que se \emph{siga}) y se termina con la metalógica, que relaciona ambas perspectivas. El tema~\ref{cap:tableaux} presenta un segundo método de demostración, los tableaux.

\begin{itemize}
    \item \textbf{Secciones de ampliación.} Las secciones y los ejercicios marcados con $^{\star}$ profundizan en las demostraciones metalógicas o en cuestiones complementarias. No son necesarios para seguir el resto del curso.
    \item \textbf{Ejercicios.} Al final de cada tema hay ejercicios numerados (el ejercicio 4.12 es el duodécimo del tema 4), agrupados por niveles de dificultad o por bloques temáticos. El apéndice~\ref{ap:soluciones} contiene soluciones de una selección de ellos (los del tema~\ref{cap:meta-prop} llevan su solución guiada en el propio tema), y el apéndice~\ref{ap:reglas} reúne todas las reglas en pocas páginas.
    \item \textbf{Teoría de conjuntos.} El apéndice~\ref{ap:conjuntos} repasa las nociones de conjuntos que usan la semántica y la metalógica desde el tema~\ref{cap:sem-prop}: pertenencia e inclusión, pares ordenados y relaciones, funciones y numerabilidad.
    \item \textbf{Aplicaciones web.} En \webaplicaciones{} hay seis aplicaciones para practicar en el navegador, sin instalar nada y con la notación de estos apuntes: \appsintaxis{} y \appdeduccion{} (temas~\ref{cap:dn-prop} y~\ref{cap:dn-lpo}), \apptablas{} (tema~\ref{cap:sem-prop}), \appkalmar{} (tema~\ref{cap:meta-prop}), \appmundo{} (tema~\ref{cap:sem-lpo}) y \apptableaux{} (tema~\ref{cap:tableaux}). A lo largo del texto, las notas «Para practicar» indican qué aplicación ilustra cada noción, y las instrucciones de los ejercicios de cada tema dicen con cuál se pueden comprobar.
\end{itemize}

\mainmatter

%%%%%%%%%%%%%%%%%%%%%%%%%%%%%%%%%%%%%%%%%%%%%%%%%%%%%%%%%%%%%%%%%%%%%%%%%%%%%%%
\part{Lógica proposicional}
%%%%%%%%%%%%%%%%%%%%%%%%%%%%%%%%%%%%%%%%%%%%%%%%%%%%%%%%%%%%%%%%%%%%%%%%%%%%%%%

\chapter{Deducción natural en lógica proposicional}\label{cap:dn-prop}
\begin{resumentema}
En este tema estudiamos la lógica proposicional como una teoría formal
de la inferencia válida. A través del cálculo de deducción natural,
analizamos las conectivas lógicas partiendo de sus reglas operacionales
de introducción y eliminación, y aprendemos a
construir derivaciones. 
\end{resumentema}

% -*- coding: utf-8 -*-
\section{La lógica como teoría de la inferencia formal}

El objetivo central de la lógica es el estudio de los argumentos válidos: aquellos en los que la conclusión se sigue de las premisas. En este tema estudiamos una manera de mostrarlo: obtener la conclusión a partir de las premisas mediante reglas de inferencia que solo atienden a la forma de las fórmulas. En el tema~\ref{cap:sem-prop} veremos qué significa, desde el punto de vista semántico, que una conclusión se siga de unas premisas, y en el tema~\ref{cap:meta-prop}, que las dos perspectivas coinciden. Para analizar la validez de un argumento, aislamos su estructura formal distinguiendo entre enunciados básicos o proposiciones atómicas (representadas por letras como $p, q, r, \dots$) y partículas de enlace o \textbf{conectivas lógicas} («y», «o», «no», «si\ldots{} entonces», «si y solo si»).

En 1934, el lógico Gerhard Gentzen (y, de forma independiente,
Stanisław Jaśkowski) propuso un sistema formal diseñado
para reflejar la manera en que razonamos en la
práctica: la \textbf{deducción natural}. En este marco, cada conectiva
lógica queda caracterizada por las reglas que establecen su uso en una demostración:

\begin{enumerate}
    \item \textbf{Regla de introducción (I).} Determina bajo qué condiciones o pasos previos estamos autorizados a afirmar una fórmula que contiene dicha conectiva como operador principal.
    \item \textbf{Regla de eliminación (E).} Determina qué consecuencias podemos extraer cuando disponemos de una fórmula con dicha conectiva como operador principal.
\end{enumerate}

\subsection{Sintaxis, precedencia y asociatividad de las conectivas}\label{sec:precedencia}

El lenguaje formal de la lógica proposicional se construye a partir de letras o variables proposicionales ($p, q, r, \dots$) y las conectivas lógicas: negación ($\neg$), conjunción ($\land$), disyunción ($\lor$), condicional ($\to$) y bicondicional ($\leftrightarrow$). Usaremos además el símbolo $\bot$ (lo absurdo, la contradicción), que se comporta como una fórmula siempre falsa.

Las \textbf{fórmulas} se definen inductivamente:
\begin{itemize}
    \item Las variables proposicionales y $\bot$ son fórmulas (fórmulas \textbf{atómicas}).
    \item Si $A$ y $B$ son fórmulas, también lo son $\neg A$, $(A \land B)$, $(A \lor B)$, $(A \to B)$ y $(A \leftrightarrow B)$.
    \item Solo son fórmulas las expresiones que se obtienen aplicando un número finito de veces las dos cláusulas anteriores.
\end{itemize}
La conectiva que se aplica en el último paso de la construcción es la \textbf{conectiva principal} (u operador principal) de la fórmula, y las fórmulas a las que se aplica son sus \textbf{subfórmulas inmediatas}: $A$ en $\neg A$; $A$ y $B$ en $(A \land B)$, $(A \lor B)$, $(A \to B)$ y $(A \leftrightarrow B)$. Por ejemplo, la conectiva principal de $((p \land q) \to r)$ es $\to$, y sus subfórmulas inmediatas son $(p \land q)$ y $r$. Las \textbf{subfórmulas} de una fórmula son ella misma, sus subfórmulas inmediatas, las subfórmulas inmediatas de estas, y así sucesivamente.

En la notación presentada, cada aplicación de una conectiva binaria requiere paréntesis que delimiten su alcance exacto, como en $((p \land q) \to r)$. No obstante, para simplificar la escritura y facilitar la lectura de fórmulas complejas, adoptamos las convenciones estándar de \textbf{precedencia} y \textbf{asociatividad} (análogas a las de la aritmética ordinaria, donde en $7 \times 2 + 8$ la multiplicación precede a la suma):

\begin{enumerate}
    \item \textbf{Jerarquía de precedencia (de mayor a menor fuerza de enlace):}
    \[
    \neg \quad \succ \quad \land \quad \succ \quad \lor \quad \succ \quad \to \quad \succ \quad \leftrightarrow
    \]
    \begin{itemize}
        \item El negador ($\neg$) tiene la máxima precedencia: afecta solo a la fórmula más breve que le sigue (una letra, otra negación o una expresión entre paréntesis). Así, $\neg p \land q$ representa siempre $(\neg p) \land q$, y no $\neg(p \land q)$.
        \item La conjunción ($\land$) y la disyunción ($\lor$) poseen mayor fuerza de enlace que el condicional ($\to$) y el bicondicional ($\leftrightarrow$). Por consiguiente, $p \land q \to r$ se interpreta como $(p \land q) \to r$, y $p \to q \lor r$ como $p \to (q \lor r)$. Entre ellas, la conjunción precede a la disyunción: $p \land q \lor r$ se lee $(p \land q) \lor r$.
        \item El bicondicional ($\leftrightarrow$) es el operador con menor precedencia relativa, de modo que abarca las subfórmulas completas situadas a su izquierda y derecha. Así, $p \to q \leftrightarrow \neg q \to \neg p$ se lee como $(p \to q) \leftrightarrow (\neg q \to \neg p)$.
    \end{itemize}

    \item \textbf{Asociatividad por la derecha para el condicional.}
    Cuando varios condicionales se suceden consecutivamente sin paréntesis que especifiquen el orden, se asocia por la derecha:
    \[
    p \to q \to r \quad \text{equivale a} \quad p \to (q \to r)
    \]
    Esta convención se corresponde de forma natural con la estructura deductiva: afirma que a partir del supuesto $p$ se infiere el condicional $q \to r$.
\end{enumerate}

En caso de duda, o siempre que se desee forzar un agrupamiento distinto (como en $(p \lor q) \land r$), el empleo explícito de paréntesis prevalece sobre estas convenciones abreviadas.

Cuando se encadenan varias conjunciones, o varias disyunciones, a veces se omiten los paréntesis interiores, como en $p \land q \land r$. Todas las agrupaciones tienen el mismo valor de verdad (tema~\ref{cap:sem-prop}), así que en una formalización da igual cuál se lea; pero en una derivación (y en un tableau, tema~\ref{cap:tableaux}) hay que fijar una, por ejemplo $(p \land q) \land r$, porque las reglas solo separan los dos miembros de una conjunción.

\practica{La aplicación \appsintaxis{} ilustra lo que se presenta en esta sección: da la conectiva principal, las subfórmulas y el árbol de formación de cualquier fórmula, y propone ejercicios para reconocer la conectiva principal y cómo se agrupa una fórmula según estas convenciones.}

\section{Estructura de las pruebas de Fitch: premisas, hipótesis y cajas}

En el formato de Fitch (popularizado por Frederic Fitch en 1952), una deducción consiste en una secuencia ordenada de líneas numeradas con barras verticales que delimitan el \textbf{alcance} (\emph{scope}) de cada hipótesis. A las subpruebas delimitadas por esas barras las llamaremos también \emph{cajas}:

\begin{itemize}
    \item \textbf{Premisas.} Enunciados que se asumen como punto de
      partida en una prueba.
    \item \textbf{Hipótesis provisionales (subpruebas o cajas).} Suposiciones que abrimos temporalmente («supongamos por un momento que $A$»). Se representan visualmente mediante una barra vertical subordinada.
    \item \textbf{Líneas accesibles.} En cada paso solo se pueden citar líneas anteriores que estén en la prueba principal o dentro de una caja que siga abierta en ese punto. Las líneas de una caja ya cerrada no pueden citarse sueltas; solo la caja completa (por ejemplo, «3--5») puede citarse, y únicamente por las reglas que cierran hipótesis.
    \item \textbf{Descarga o cierre de hipótesis.} Al cerrar una caja, la hipótesis interior queda sellada (no se puede volver a invocar aisladamente), pero la caja completa nos autoriza a afirmar una nueva conclusión en el nivel exterior (como un condicional o una negación).
\end{itemize}

Escribimos $\Gamma \vdash A$ para indicar que existe una derivación cuya última línea es $A$, cuyas premisas pertenecen al conjunto $\Gamma$ y en la que todas las hipótesis provisionales han sido descargadas. Cuando no hace falta ninguna premisa escribimos $\vdash A$ y decimos que $A$ es un \textbf{teorema}. Usaremos como sinónimos \emph{derivación}, \emph{deducción}, \emph{prueba} y \emph{demostración} (formal).

A la derecha de cada línea se escribe su justificación: «Prem.» en las premisas, «Hip.» en las hipótesis («Hip. IP» si se abren para aplicar la prueba indirecta, IP) y, en las demás líneas, la regla aplicada con los números de las líneas o cajas que cita. En los esquemas de las reglas, $m, n, i, j, \dots$ representan números de línea cualesquiera y los puntos verticales, pasos intermedios.

\practica{En la aplicación \appdeduccion{} puedes construir derivaciones con las reglas de este tema y del tema~\ref{cap:dn-lpo}. Comprueba cada línea mientras escribes y explica los errores. Contiene los ejercicios de deducción natural de los dos temas, sin sus soluciones.}

\section{Reglas básicas del cálculo proposicional}

\subsection{Reiteración y conjunción (\texorpdfstring{$\land$}{y})}

La conjunción une dos afirmaciones. Se introduce cuando tenemos demostradas ambas partes por separado, y se elimina extrayendo cualquiera de ellas. La regla de reiteración permite copiar una línea accesible en el contexto actual.

\vspace{0.3cm}
\begin{center}
\begin{tabular}{cc}
\textbf{Reiteración (R)} & \textbf{Introducción de la conjunción ($\land$I)} \\[0.2cm]
$\begin{nd}[justformat=myjust]
    \have [m] {1} A 
    \have [\mbox{}] {2} A \r{1}
\end{nd}$ &
$\begin{nd}
    \have [m] {1} A
    \have [n] {2} B
    \have [\mbox{}] {3} {A \land B} \ai{1,2}
\end{nd}$ \\[1.2cm]
\multicolumn{2}{c}{\textbf{Eliminación de la conjunción ($\land$E)}} \\[0.2cm]
$\begin{nd}
    \have [m] {1} {A \land B}
    \have [\mbox{}] {2} {A} \ae{1}
\end{nd}$ &
$\begin{nd}
    \have [m] {1} {A \land B}
    \have [\mbox{}] {2} {B} \ae{1}
\end{nd}$
\end{tabular}
\end{center}

\begin{ejemplo}[Demostración directa con conjunción]
Demostremos que a partir de las premisas $p$ y $q \land r$ se deduce formalmente la conclusión $r \land p$:
\[
p, \ q \land r \ \vdash \ r \land p
\]
\begin{center}
$\begin{nd}[rules=myrules]
    \hypo [1] {1} {p} \by{Prem.}{}
    \hypo [2] {2} {q \land r} \by{Prem.}{}
    \have [3] {3} {r} \ae{2}
    \have [4] {4} {r \land p} \ai{3, 1}
\end{nd}$
\end{center}
\textbf{Comentario:} 
\begin{itemize}
    \item Las líneas 1 y 2 registran las premisas iniciales del argumento.
    \item La meta es construir la conjunción $r \land p$. La regla de introducción ($\land$I) exige disponer de sus dos componentes por separado. Como la premisa 1 ya nos proporciona $p$, el único subobjetivo pendiente es aislar $r$.
    \item Aplicamos la regla de eliminación ($\land$E) sobre la premisa 2 para extraer su miembro derecho $r$ (línea 3).
    \item Con $r$ y $p$ disponibles en líneas accesibles del nivel principal, aplicamos $\land$I combinando las líneas 3 y 1 para alcanzar la conclusión en la línea 4. Nótese que la justificación en el margen derecho cita con exactitud las dos líneas utilizadas.
\end{itemize}
\end{ejemplo}

\subsection{Condicional (\texorpdfstring{$\to$}{entonces})}

El condicional representa el paso deductivo de un antecedente a un consecuente. 
\begin{itemize}
    \item \textbf{Introducción ($\to$I).} Para justificar $A \to B$, abrimos una caja asumiendo provisionalmente la hipótesis $A$. Si dentro de esa caja logramos alcanzar $B$, cerramos la caja y afirmamos fuera: $A \to B$.
    \item \textbf{Eliminación ($\to$E / \emph{modus ponens}).} Si disponemos de $A \to B$ y además tenemos $A$, obtenemos directamente $B$.
\end{itemize}

\vspace{0.3cm}
\begin{center}
\begin{tabular}{cc}
\textbf{Introducción del condicional ($\to$I)} & \textbf{Eliminación del condicional ($\to$E)} \\[0.2cm]
$\begin{nd}[rules=myrules]
    \open
    \hypo [i] {1} {A}
    \have [\mbox{}] {} {\vdots}
    \have [j] {2} {B}
    \close
    \have [\mbox{}] {3} {A \to B} \ci{1-2}
\end{nd}$ &
$\begin{nd}[rules=myrules]
    \have [m] {1} {A \to B}
    \have [n] {2} {A}
    \have [\mbox{}] {3} {B} \ce{1,2}
\end{nd}$
\end{tabular}
\end{center}

\begin{ejemplo}[Silogismo hipotético y apertura de hipótesis]
Demostremos la transitividad del condicional formalizando el silogismo hipotético:
\[
p \to q, \ q \to r \ \vdash \ p \to r
\]
\begin{center}
$\begin{nd}[rules=myrules]
    \hypo [1] {1} {p \to q} \by{Prem.}{}
    \hypo [2] {2} {q \to r} \by{Prem.}{}
    \open
        \hypo [3] {3} {p} \by{Hip.}{}
        \have [4] {4} {q} \ce{1, 3}
        \have [5] {5} {r} \ce{2, 4}
    \close
    \have [6] {6} {p \to r} \ci{3-5}
\end{nd}$
\end{center}
\textbf{Comentario:} 
\begin{itemize}
    \item La conclusión que pretendemos alcanzar es un condicional: $p \to r$. La estrategia natural para probar un condicional es aplicar su regla de introducción ($\to$I): abrimos una caja subordinada asumiendo provisionalmente su antecedente $p$ (línea 3) y nos fijamos como único objetivo dentro de la caja alcanzar el consecuente $r$.
    \item Dentro de la caja disponemos de $p$ y de la premisa 1 ($p \to q$). Aplicando la eliminación del condicional ($\to$E / \emph{modus ponens}) obtenemos $q$ (línea 4).
    \item Con $q$ recién obtenido y la premisa 2 ($q \to r$), aplicamos nuevamente $\to$E para alcanzar $r$ (línea 5).
    \item Habiendo completado el objetivo interno, cerramos la
      caja. Fuera de ella afirmamos $p \to r$,
      justificando mediante $\to$I que cita el rango continuo de la
      subprueba (3--5). A partir de la línea 6, la hipótesis $p$ queda
      descargada e inaccesible. 
\end{itemize}
\end{ejemplo}

\begin{ejemplo}[Interacción de condicional y conjunción: ley de importación]
Veamos cómo se combinan las reglas del condicional con las de la conjunción:
\[
p \to (q \to r) \ \vdash \ (p \land q) \to r
\]
\begin{center}
$\begin{nd}[rules=myrules]
    \hypo [1] {1} {p \to (q \to r)} \by{Prem.}{}
    \open
        \hypo [2] {2} {p \land q} \by{Hip.}{}
        \have [3] {3} {p} \ae{2}
        \have [4] {4} {q} \ae{2}
        \have [5] {5} {q \to r} \ce{1, 3}
        \have [6] {6} {r} \ce{5, 4}
    \close
    \have [7] {7} {(p \land q) \to r} \ci{2-6}
\end{nd}$
\end{center}
\textbf{Comentario.} Para demostrar el condicional exterior abrimos la hipótesis $p \land q$ (línea 2). Dentro de la caja, descomponemos de inmediato la conjunción mediante $\land$E para obtener tanto $p$ como $q$ (líneas 3 y 4). Esto nos permite aplicar dos veces el \emph{modus ponens}, primero con la premisa 1 y después con la línea 5, hasta alcanzar $r$, cerrando la caja con $\to$I.
\end{ejemplo}

\subsection{Negación y contradicción (\texorpdfstring{$\neg, \bot$}{no, absurdo})}

El símbolo $\bot$ representa una contradicción o absurdo lógico.
\begin{itemize}
    \item \textbf{Eliminación de la negación ($\neg$E).} De una fórmula $A$ y su negación $\neg A$ se concluye $\bot$.
    \item \textbf{Introducción de la negación ($\neg$I).} Si asumiendo $A$ llegamos a una contradicción $\bot$, cerramos la caja y concluimos $\neg A$.
    \item \textbf{Principio de explosión ($\text{X}$).} De una contradicción $\bot$ se sigue cualquier fórmula $B$.
    \item \textbf{Prueba indirecta (IP / reducción al absurdo).} Si asumiendo $\neg A$ alcanzamos $\bot$, cerramos la caja y concluimos $A$.
\end{itemize}

\vspace{0.3cm}
\begin{center}
\begin{tabular}{cc}
\textbf{Eliminación de la negación ($\neg$E)} & \textbf{Introducción de la negación ($\neg$I)} \\[0.2cm]
$\begin{nd}[rules=myrules]
    \have [m] {1} {\lnot A}
    \have [n] {2} {A}
    \have [\mbox{}] {3} {\bot} \ne{1,2}
\end{nd}$ &
$\begin{nd}[rules=myrules]
    \open
    \hypo [i] {1} {A}
    \have [\mbox{}] {} {\vdots}
    \have [j] {2} {\bot}
    \close 
    \have [\mbox{}] {3} {\lnot A} \ni{1-2}  
\end{nd}$ \\[1.2cm]
\textbf{Explosión ($\text{X}$)} & \textbf{Prueba indirecta (IP)} \\[0.2cm]
$\begin{nd}[rules=myrules]
    \have [m] {1} {\bot}
    \have [\mbox{}] {2} {B} \x{1}
\end{nd}$ &
$\begin{nd}[rules=myrules]
    \open
    \hypo [i] {1} {\lnot A}
    \have [\mbox{}] {} {\vdots}
    \have [j] {2} {\bot}
    \close 
    \have [\mbox{}] {3} {A} \ip{1-2}  
\end{nd}$
\end{tabular}
\end{center}

\begin{ejemplo}[Introducción de la negación]
Demostremos que si una fórmula implica su propia negación, podemos derivar formalmente la negación de dicha fórmula:
\[
p \to \neg p \ \vdash \ \neg p
\]
\begin{center}
$\begin{nd}[rules=myrules]
    \hypo [1] {1} {p \to \lnot p} \by{Prem.}{}
    \open
        \hypo [2] {2} {p} \by{Hip.}{}
        \have [3] {3} {\lnot p} \ce{1, 2}
        \have [4] {4} {\bot} \ne{3, 2}
    \close
    \have [5] {5} {\lnot p} \ni{2-4}
\end{nd}$
\end{center}
\textbf{Comentario:} 
\begin{itemize}
    \item La meta es demostrar la negación $\neg p$. La regla $\neg$I nos indica el camino: abrir una subprueba asumiendo temporalmente lo contrario, es decir, la afirmación positiva $p$ (línea 2).
    \item Mediante el \emph{modus ponens} con la premisa 1 obtenemos $\neg p$ (línea 3).
    \item En la línea 4 confrontamos $\neg p$ con la hipótesis $p$ mediante la eliminación de la negación ($\neg$E), produciendo el absurdo formal $\bot$.
    \item Habiendo demostrado que la hipótesis $p$ conduce a una
      contradicción, cerramos la caja y concluimos fuera $\neg p$
      mediante $\neg$I (2--4). 
\end{itemize}
\end{ejemplo}

\begin{ejemplo}[Prueba indirecta o reducción al absurdo: la regla IP]
En el cálculo deductivo clásico, si la suposición provisional de la negación de una fórmula ($\neg A$) conduce a una contradicción formal ($\bot$), la regla IP permite descargar dicha hipótesis y derivar la fórmula afirmativa ($A$). A modo de ilustración, derivamos la regla de eliminación de la doble negación ($\neg\neg p \vdash p$):
\begin{center}
$\begin{nd}[rules=myrules]
    \hypo [1] {1} {\lnot\lnot p} \by{Prem.}{}
    \open
        \hypo [2] {2} {\lnot p} \by{Hip. IP}{}
        \have [3] {3} {\bot} \ne{1, 2}
    \close
    \have [4] {4} {p} \ip{2-3}
\end{nd}$
\end{center}
\textbf{Comentario.} Nótese la diferencia entre ambas reglas: mientras que la introducción de la negación ($\neg$I) es una regla constructiva que parte de una hipótesis $A$ para establecer su negación $\neg A$, la prueba indirecta clásica (IP) tiene un carácter no constructivo, pues parte del supuesto negativo $\neg A$ para derivar la afirmación $A$. Por esta razón, reservamos el nombre tradicional de \emph{reducción al absurdo} para la regla clásica IP, distinguiéndola de la \emph{introducción de la negación} ($\neg$I).
\end{ejemplo}

\subsection{Disyunción (\texorpdfstring{$\lor$}{o})}

La disyunción expresa alternativas.
\begin{itemize}
    \item \textbf{Introducción ($\lor$I).} Si tenemos probado $A$, podemos concluir $A \lor B$ (o $B \lor A$) para cualquier fórmula $B$.
    \item \textbf{Eliminación ($\lor$E / razonamiento por casos).} Si disponemos de $A \lor B$, para alcanzar una conclusión común $C$ debemos examinar ambos supuestos por separado: en una primera caja suponemos $A$ y derivamos $C$; en una segunda caja paralela suponemos $B$ y derivamos también $C$. Al tener $C$ garantizado en ambos casos, afirmamos $C$ fuera de las cajas.
\end{itemize}

\vspace{0.3cm}
\begin{center}
\begin{tabular}{cc}
\textbf{Introducción de la disyunción ($\lor$I)} & \textbf{Eliminación de la disyunción ($\lor$E)} \\[0.2cm]
$\begin{nd}
    \have [m] {1} {A}
    \have [\mbox{}] {2} {A \lor B} \oi{1}
\end{nd}$\quad
$\begin{nd}
    \have [m] {1} {B}
    \have [\mbox{}] {2} {A \lor B} \oi{1}
\end{nd}$ &
$\begin{nd}
    \have [m] {1} {A \lor B}
    \open
    \hypo [i] {2} {A}
    \have [\mbox{}] {} {\vdots}
    \have [j] {3} {C}
    \close 
    \open
    \hypo [k] {4} {B}
    \have [\mbox{}] {} {\vdots}
    \have [l] {5} {C}
    \close 
    \have [\mbox{}] {6} {C} \oe{1,2-3, 4-5}    
\end{nd}$
\end{tabular}
\end{center}

\begin{ejemplo}[Eliminación de la disyunción en un dilema constructivo simple: $p \lor q, p \to r, q \to r \vdash r$]
Supongamos que sabemos que «llueve o nieva» ($p \lor q$), que «si llueve, el suelo se moja» ($p \to r$) y que «si nieva, el suelo se moja» ($q \to r$). Queremos concluir que «el suelo se moja» ($r$).

\begin{center}
$\begin{nd}[rules=myrules]
    \hypo [1] {1} {p \lor q} \by{Prem.}{}
    \hypo [2] {2} {p \to r} \by{Prem.}{}
    \hypo [3] {3} {q \to r} \by{Prem.}{}
    \open
        \hypo [4] {4} {p} \by{Hip.}{}
        \have [5] {5} {r} \ce{2, 4}
    \close
    \open
        \hypo [6] {6} {q} \by{Hip.}{}
        \have [7] {7} {r} \ce{3, 6}
    \close
    \have [8] {8} {r} \oe{1, 4-5, 6-7}
\end{nd}$
\end{center}
\textbf{Comentario.} La regla $\lor$E formaliza el razonamiento por casos. Nótese que:
\begin{enumerate}
    \item Para justificar la línea 8 se citan tres elementos: la disyunción de partida (línea 1) y los dos intervalos completos de subdemostración (líneas 4--5 y 6--7).
    \item Ambas cajas deben derivar \emph{exactamente la misma fórmula final} ($r$).
    \item Las hipótesis temporales $p$ y $q$ quedan descargadas al cerrar cada caja; no se pueden volver a citar fuera. Lo único que se extrae es $r$, que queda establecido sea cual sea el disyunto del que se parta.
\end{enumerate}
\end{ejemplo}

\subsection{Bicondicional (\texorpdfstring{$\leftrightarrow$}{si y solo si})}

El bicondicional $A \leftrightarrow B$ funciona como un condicional en ambos sentidos. Se introduce demostrando ambas direcciones en dos subpruebas independientes, y se elimina actuando como condicional en cualquiera de los dos sentidos.

\vspace{0.3cm}
\begin{center}
\begin{tabular}{cc}
\textbf{Introducción del bicondicional ($\leftrightarrow$I)} & \textbf{Eliminación del bicondicional ($\leftrightarrow$E)} \\[0.2cm]
$\begin{nd}[rules=myrules]
    \open
    \hypo [i] {1} {A}
    \have [\mbox{}] {} {\vdots}
    \have [j] {2} {B}
    \close 
    \open
    \hypo [k] {3} {B}
    \have [\mbox{}] {} {\vdots}
    \have [l] {4} {A}
    \close 
    \have [\mbox{}] {5} {A \leftrightarrow B} \bi{1-2,3-4}  
\end{nd}$ &
\begin{tabular}{c}
$\begin{nd}[rules=myrules]
    \have [m] {1} {A \leftrightarrow B}
    \have [n] {2} {A}
    \have [\mbox{}] {3} {B} \be{1,2}
\end{nd}$ \\[0.3cm]
$\begin{nd}[rules=myrules]
    \have [m] {1} {A \leftrightarrow B}
    \have [n] {2} {B}
    \have [\mbox{}] {3} {A} \be{1,2}
\end{nd}$
\end{tabular}
\end{tabular}
\end{center}

\begin{ejemplo}[Transitividad del bicondicional: $p \leftrightarrow q, q \leftrightarrow r \vdash p \leftrightarrow r$]
Ilustramos la regla de introducción del bicondicional, que exige construir dos cajas independientes que prueben la implicación en ambos sentidos:

\begin{center}
$\begin{nd}[rules=myrules]
    \hypo [1] {1} {p \leftrightarrow q} \by{Prem.}{}
    \hypo [2] {2} {q \leftrightarrow r} \by{Prem.}{}
    \open
        \hypo [3] {3} {p} \by{Hip.}{}
        \have [4] {4} {q} \be{1, 3}
        \have [5] {5} {r} \be{2, 4}
    \close
    \open
        \hypo [6] {6} {r} \by{Hip.}{}
        \have [7] {7} {q} \be{2, 6}
        \have [8] {8} {p} \be{1, 7}
    \close
    \have [9] {9} {p \leftrightarrow r} \bi{3-5, 6-8}
\end{nd}$
\end{center}
\textbf{Comentario.} Observamos que:
\begin{enumerate}
    \item Para alcanzar el bicondicional $p \leftrightarrow r$, la regla $\leftrightarrow$I exige probar dos derivaciones o direcciones independientes: deducir $r$ asumiendo provisionalmente $p$ (caja 3--5) y, recíprocamente, deducir $p$ asumiendo provisionalmente $r$ (caja 6--8).
    \item Dentro de cada caja, la regla de eliminación del bicondicional ($\leftrightarrow$E) se utiliza en ambos sentidos: en la línea 4 se extrae $q$ a partir de la hipótesis $p$ mediante la premisa 1 (dirección de izquierda a derecha), mientras que en la línea 8 se extrae $p$ a partir de $q$ mediante la misma premisa 1 (dirección de derecha a izquierda), pues el bicondicional autoriza el paso inferencial en cualquiera de las dos orientaciones.
\end{enumerate}
\end{ejemplo}

\section{Reglas derivadas}

Las reglas derivadas son esquemas de inferencia demostrables en el cálculo básico que sirven como herramientas abreviadas para simplificar derivaciones extensas:

\vspace{0.3cm}
\begin{center}
\begin{tabular}{cc}
\textbf{\emph{Modus tollens} (MT)} & \textbf{Silogismo disyuntivo (DS)} \\[0.2cm]
$\begin{nd}[rules=myrules]
    \have [m] {1} {A \to B}
    \have [n] {2} {\lnot B}
    \have [\mbox{}] {3} {\lnot A} \mt{1,2}
\end{nd}$ &
\begin{tabular}[t]{@{}c@{}}
$\begin{nd}[rules=myrules]
    \have [m] {1} {A \lor B}
    \have [n] {2} {\lnot A}
    \have [\mbox{}] {3} {B} \ds{1,2}
\end{nd}$ \\[0.9cm]
$\begin{nd}[rules=myrules]
    \have [m] {1} {A \lor B}
    \have [n] {2} {\lnot B}
    \have [\mbox{}] {3} {A} \ds{1,2}
\end{nd}$
\end{tabular} \\ \noalign{\vspace{0.5cm}}
\textbf{Eliminación de la doble negación (DNE)} & \textbf{Tercero excluido (LEM)} \\[0.2cm]
$\begin{nd}[rules=myrules]
    \have [m] {1} {\lnot\lnot A}
    \have [\mbox{}] {2} {A} \dne{1}
\end{nd}$ &
$\begin{nd}[justsep=.3em,rules=myrules]
    \open
    \hypo [i] {1} {A}
    \have [\mbox{}] {} {\vdots}
    \have [j] {2} {B}
    \close 
    \open
    \hypo [k] {3} {\lnot A}
    \have [\mbox{}] {} {\vdots}
    \have [l] {4} {B}
    \close 
    \have [\mbox{}] {5} {B} \lem{1-2,3-4}  
\end{nd}$ \\ \noalign{\vspace{0.7cm}}
\multicolumn{2}{c}{\textbf{Leyes de De Morgan (DeM)}} \\[0.2cm]
$\begin{nd}[rules=myrules]
    \have [m] {1} {\lnot(A \lor B)}
    \have [\mbox{}] {2} {\lnot A \land \lnot B} \dem{1}
\end{nd}$ &
$\begin{nd}[rules=myrules]
    \have [m] {1} {\lnot(A \land B)}
    \have [\mbox{}] {2} {\lnot A \lor \lnot B} \dem{1}
\end{nd}$ \\[0.6cm]
$\begin{nd}[rules=myrules]
    \have [m] {1} {\lnot A \land \lnot B}
    \have [\mbox{}] {2} {\lnot(A \lor B)} \dem{1}
\end{nd}$ &
$\begin{nd}[rules=myrules]
    \have [m] {1} {\lnot A \lor \lnot B}
    \have [\mbox{}] {2} {\lnot(A \land B)} \dem{1}
\end{nd}$
\end{tabular}
\end{center}

La regla LEM formaliza el razonamiento por casos sobre una afirmación y su negación: si asumiendo $A$ se deduce $B$, y asumiendo $\neg A$ se deduce también $B$, la regla autoriza a concluir $B$ descargando ambas hipótesis.

\begin{ejemplo}[Derivación del \emph{modus tollens} mediante reglas básicas]
Toda regla derivada puede justificarse usando solo reglas básicas:
\begin{center}
$\begin{nd}[rules=myrules]
    \hypo [1] {1} {A \to B} \by{Prem.}{}
    \hypo [2] {2} {\lnot B} \by{Prem.}{}
    \open
        \hypo [3] {3} {A} \by{Hip.}{}
        \have [4] {4} {B} \ce{1,3}
        \have [5] {5} {\bot} \ne{2,4}
    \close
    \have [6] {6} {\lnot A} \ni{3-5}
\end{nd}$
\end{center}
\end{ejemplo}

\section{Estrategias para la construcción de pruebas}

\begin{itemize}
    \item \textbf{Pensar hacia atrás desde la conclusión.} Si la meta es $A \to B$, abre una caja con $A$ y busca $B$. Si la meta es $A \land B$, busca $A$ y $B$ por separado. Si la meta es $\neg A$, abre una caja con $A$ y busca $\bot$ ($\neg$I).
    \item \textbf{Aprovechar las premisas compuestas.} Si tienes $A \lor B$, prepara dos cajas paralelas para razonar por casos. Si tienes $A \land B$, descomponlo inmediatamente con $\land$E.
    \item \textbf{El recurso del absurdo.} Si no ves un camino directo hacia tu conclusión, asume su negación y busca una contradicción para aplicar IP (si la conclusión es una negación $\neg A$, basta suponer $A$ y aplicar $\neg$I).
\end{itemize}

Una derivación tiene una propiedad notable: para saber si es correcta no hace falta entender de qué trata; basta comprobar línea a línea que cada paso se ajusta a una regla. Por eso puede comprobarla un programa, como la aplicación \appdeduccion{} que usamos en los ejercicios. En la sección~\ref{sec:asistentes} (tema~\ref{cap:dn-lpo}) veremos cómo esta idea sostiene hoy los asistentes de demostración que verifican resultados matemáticos, también cuando los propone un modelo de lenguaje.

\section*{Referencias bibliográficas}

\begin{enumerate}
    \item \textbf{Fitch, F. B.} (1952). \emph{Symbolic Logic: An Introduction}. Ronald Press, Nueva York. [Origen del formato de subdemostraciones con barras verticales que usamos en estos apuntes].
    \item \textbf{Gentzen, G.} (1934/1935). «Untersuchungen über das logische Schließen». \emph{Mathematische Zeitschrift}, vol. 39, pp. 176--210 y 405--431. [Origen de la deducción natural: cada conectiva se caracteriza por sus reglas de introducción y eliminación].
    \item \textbf{Jaśkowski, S.} (1934). «On the rules of suppositions in formal logic». \emph{Studia Logica}, n.º 1, pp. 5--32. [Propuesta independiente de un cálculo de suposiciones, con una notación gráfica de cajas que anticipa la de Fitch].
    \item \textbf{Magnus, P. D., Button, T., et al.} (2020). \emph{forall x: Calgary. An Introduction to Formal Logic}. Open Logic Project. [Parte IV: el cálculo de deducción natural proposicional que seguimos].
    \item \textbf{Pelletier, F. J.} (1999). «A brief history of natural deduction». \emph{History and Philosophy of Logic}, vol. 20, n.º 1, pp. 1--31. [Historia de los sistemas de Gentzen, Jaśkowski y Fitch y de su llegada a los manuales].
\end{enumerate}

%%% Local Variables:
%%% mode: latex
%%% TeX-master: "../ApuntesLogica"
%%% End:


\section{Ejercicios de deducción natural}
% -*- coding: utf-8 -*-
\begin{quote}
\small \textbf{Instrucciones metodológicas.} Trabaja las pruebas a lápiz o con la aplicación \appdeduccion{}, que contiene estos ejercicios y comprueba cada línea; su ayuda explica cómo escribir las fórmulas y las justificaciones. Cada nivel introduce una dificultad lógica nueva. En los ejercicios señalados como «plantilla guiada» debes completar las justificaciones que faltan (marcadas con «?») en la columna de la derecha. Salvo indicación en contra, puedes usar las reglas derivadas (MT, DS, DNE, LEM, DeM); pero en los ejercicios~\ref{e1-15} y~\ref{e1-22}, cuyo objetivo es obtener DS y el tercero excluido, usa solo reglas básicas, y en los ejercicios~\ref{e1-18} y~\ref{e1-23}, que piden derivar dos de las leyes de De Morgan, no uses DeM. Al final de los apuntes (apéndice~\ref{ap:soluciones}) hay soluciones de una selección de ejercicios.
\end{quote}

\section*{Nivel 1: reglas directas sin cajas ($\land$, $\to$E, $\lor$I)}
\textit{Objetivo: aplicar reglas básicas de descomposición y composición sin abrir supuestos provisionales.}

\begin{pregunta}[Conmutatividad de la conjunción]\label{e1-1}
$p \land q \vdash q \land p$
\end{pregunta}
\begin{pregunta}[Encadenamiento de \emph{modus ponens}]\label{e1-2}
$p, p \to q, q \to r \vdash r$
\end{pregunta}
\begin{pregunta}[Asociatividad de la conjunción]\label{e1-3}
$p \land (q \land r) \vdash (p \land q) \land r$
\end{pregunta}
\begin{pregunta}\label{e1-4}
$p \land q, p \to r, q \to s \vdash r \land s$
\end{pregunta}
\begin{pregunta}[Introducción de la disyunción]\label{e1-5}
$p \vdash (p \lor q) \land (r \lor p)$
\end{pregunta}
\begin{pregunta}\label{e1-6}
$(p \to q) \land p, (q \to r) \land s \vdash r \land s$
\end{pregunta}

\section*{Nivel 2: cajas simples e introducción del condicional ($\to$I)}
\textit{Objetivo: abrir una hipótesis provisional para alcanzar una meta, cerrar la caja y justificar con $\to$I a la derecha.}

\begin{pregunta}[Monotonía condicional con conjunción]\label{e1-7}
$p \to q \vdash (p \land r) \to (q \land r)$

\textbf{Plantilla guiada} (completa las justificaciones marcadas con «?»):
\begin{center}
$\begin{nd}[rules=myrules]
    \hypo [1] {1} {p \to q} \by{Prem.}{}
    \open
        \hypo [2] {2} {p \land r} \by{Hip.}{}
        \have [3] {3} {p} \by{$\land$E}{2}
        \have [4] {4} {r} \by{?}{}
        \have [5] {5} {q} \by{?}{}
        \have [6] {6} {q \land r} \by{$\land$I}{5, 4}
    \close
    \have [7] {7} {(p \land r) \to (q \land r)} \ci{2-6}
\end{nd}$
\end{center}

\end{pregunta}
\begin{pregunta}\label{e1-8}
$p \vdash q \to (p \land q)$
\end{pregunta}
\begin{pregunta}[Distributividad condicional]\label{e1-9}
$(p \to q) \land (p \to r) \vdash p \to (q \land r)$
\end{pregunta}
\begin{pregunta}[Exportación: cajas anidadas]\label{e1-10}
$(p \land q) \to r \vdash p \to (q \to r)$
\end{pregunta}
\begin{pregunta}[Teorema sin premisas]\label{e1-11}
$\vdash p \to p$
\end{pregunta}
\begin{pregunta}[\emph{Verum ex quolibet}]\label{e1-12}
$\vdash p \to (q \to p)$
\end{pregunta}

\section*{Nivel 3: negación, contradicción y razonamiento por casos ($\neg\text{I}, \lor\text{E}, \text{DS}$)}
\textit{Objetivo: manejar la contradicción ($\bot$), la introducción de la negación ($\neg\text{I}$) y la prueba por casos con dos cajas paralelas.}

\begin{pregunta}[Equivalencia del condicional]\label{e1-13}
$p \to q \vdash \neg(p \land \neg q)$

\textbf{Plantilla guiada} (completa la justificación marcada con «?»):
\begin{center}
$\begin{nd}[rules=myrules]
    \hypo [1] {1} {p \to q} \by{Prem.}{}
    \open
        \hypo [2] {2} {p \land \lnot q} \by{Hip.}{}
        \have [3] {3} {p} \by{$\land$E}{2}
        \have [4] {4} {\lnot q} \by{$\land$E}{2}
        \have [5] {5} {q} \ce{1, 3}
        \have [6] {6} {\bot} \by{?}{}
    \close
    \have [7] {7} {\lnot(p \land \lnot q)} \ni{2-6}
\end{nd}$
\end{center}

\end{pregunta}
\begin{pregunta}[Razonamiento disyuntivo por casos]\label{e1-14}
$p \lor q, p \to r \vdash r \lor q$
\end{pregunta}
\begin{pregunta}[Derivación del silogismo disyuntivo mediante casos y $\text{X}$]\label{e1-15}
$p \lor q, \neg p \vdash q$
\end{pregunta}
\begin{pregunta}[Esquema de la paradoja del mentiroso]\label{e1-16}
$p \leftrightarrow \neg p \vdash \bot$
\end{pregunta}
\begin{pregunta}[Conmutatividad de la disyunción]\label{e1-17}
$p \lor q \vdash q \lor p$
\end{pregunta}
\begin{pregunta}[Ley de De Morgan]\label{e1-18}
$\neg(p \lor q) \vdash \neg p \land \neg q$
\end{pregunta}
\begin{pregunta}[Principio de explosión]\label{e1-19}
$p \land \neg p \vdash q$
\end{pregunta}
\begin{pregunta}\label{e1-20}
$p \to (q \lor r), \neg q \vdash p \to r$
\end{pregunta}

\section*{Nivel 4: reducción al absurdo clásica (IP) y teoremas avanzados}
\textit{Objetivo: dominar la prueba indirecta clásica (suponer $\neg A$ para deducir $A$), la doble negación y los teoremas con supuestos múltiples.}

\begin{pregunta}[Principio de no contradicción]\label{e1-21}
$\vdash \neg(p \land \neg p)$
\end{pregunta}
\begin{pregunta}[Principio del tercero excluido]\label{e1-22}
$\vdash p \lor \neg p$

\end{pregunta}
\begin{pregunta}[Ley de De Morgan clásica]\label{e1-23}
$\neg(p \land q) \vdash \neg p \lor \neg q$
\end{pregunta}
\begin{pregunta}[Ley de Peirce]\label{e1-24}
$(p \to q) \to p \vdash p$
\end{pregunta}
\begin{pregunta}\label{e1-25}
$\vdash ((p \to q) \to p) \to p$
\end{pregunta}
\begin{pregunta}[Contraposición clásica]\label{e1-26}
$\neg p \to \neg q \vdash q \to p$
\end{pregunta}
\begin{pregunta}[Definición clásica del condicional]\label{e1-27}
$p \to q \vdash \neg p \lor q$
\end{pregunta}
\begin{pregunta}\label{e1-28}
$\neg p \lor q \vdash p \to q$
\end{pregunta}
\begin{pregunta}[Ley de linealidad o de Dummett]\label{e1-29}
$\vdash (p \to q) \lor (q \to p)$
\end{pregunta}
\begin{pregunta}\label{e1-30}
$p \leftrightarrow q \vdash (p \land q) \lor (\neg p \land \neg q)$
\end{pregunta}

%%% Local Variables:
%%% mode: latex
%%% TeX-master: "../ApuntesLogica"
%%% End:


\chapter{Semántica de la lógica proposicional}\label{cap:sem-prop}
\begin{resumentema}
En el tema~\ref{cap:dn-prop} aprendimos a construir derivaciones a partir de premisas
y reglas formales en el cálculo deductivo ($\vdash$). En este tema abordamos la perspectiva complementaria: la semántica ($\models$). Introducimos las nociones de interpretación, función de verdad y tabla de verdad, analizando cuándo una fórmula es una verdad lógica (tautología), cuándo dos fórmulas son lógicamente equivalentes, cuándo un argumento es semánticamente válido (consecuencia lógica) y cómo construir contraejemplos para refutar argumentos inválidos. También estudiamos cómo formalizar enunciados del lenguaje natural y discutimos los límites del condicional material.
\end{resumentema}

% -*- coding: utf-8 -*-
\section{De la deducción a la verdad: dos perspectivas sobre la lógica}

Hasta ahora hemos trabajado con la lógica desde una perspectiva
\textbf{sintáctica o inferencial}: afirmar 
$\Gamma \vdash A$ significa que existe una prueba formal en el cálculo
de deducción natural que, partiendo de las premisas $\Gamma$, alcanza
la conclusión $A$ aplicando correctamente reglas de introducción y
eliminación. 

Sin embargo, también nos interesa responder otra pregunta: 
\emph{¿qué relación guardan nuestras demostraciones con la verdad?}
Cuando decimos que un argumento es válido, exigimos que la verdad de
las premisas asegure necesariamente la verdad de la conclusión. Para formalizar esta intuición necesitamos una \textbf{semántica formal}: una teoría que defina con precisión qué significa que una fórmula sea verdadera o falsa en una interpretación dada.

\section{Interpretaciones y funciones de verdad}

En lógica proposicional, el valor de verdad de un enunciado complejo
depende solo del valor de verdad de los enunciados que lo componen: las conectivas son \emph{veritativo-funcionales}. Es un caso particular del principio de composicionalidad, que suele atribuirse a Frege.

\begin{definition}[Interpretación o asignación de verdad]
Sea $\text{Prop} = \{p, q, r, \dots\}$ el conjunto de variables proposicionales. Una \textbf{interpretación} (o asignación de verdad) $v$ es una función que asigna a cada variable proposicional exactamente un valor de verdad del conjunto $\{1, 0\}$, donde $1$ representa la verdad y $0$ la falsedad:
\[
v: \text{Prop} \to \{1, 0\}
\]
\end{definition}

Una vez fijado el valor de verdad de las variables proposicionales mediante $v$, el valor de verdad de cualquier fórmula compleja $A$ queda determinado unívocamente mediante las \textbf{condiciones semánticas de las conectivas}:

\begin{enumerate}
    \item \textbf{Negación ($\neg A$).} Invierte el valor de verdad.
    \[
    v(\neg A) = 1 \iff v(A) = 0
    \]
    \item \textbf{Conjunción ($A \land B$).} Verdadera si y solo si ambos miembros son verdaderos.
    \[
    v(A \land B) = 1 \iff v(A) = 1 \text{ y } v(B) = 1
    \]
    \item \textbf{Disyunción ($A \lor B$).} Verdadera si y solo si al menos uno de los miembros es verdadero.
    \[
    v(A \lor B) = 1 \iff v(A) = 1 \text{ o } v(B) = 1
    \]
    \item \textbf{Condicional ($A \to B$).} Falso únicamente cuando el antecedente es verdadero y el consecuente es falso.
    \[
    v(A \to B) = 1 \iff \text{no ocurre que } (v(A) = 1 \text{ y } v(B) = 0)
    \]
    \item \textbf{Bicondicional ($A \leftrightarrow B$).} Verdadero si y solo si ambas fórmulas tienen el mismo valor de verdad.
    \[
    v(A \leftrightarrow B) = 1 \iff v(A) = v(B)
    \]
    \item \textbf{Absurdo ($\bot$).} Es falso en toda interpretación: $v(\bot) = 0$ para toda $v$.
\end{enumerate}

\begin{ejemplo}[Evaluación composicional paso a paso]
Consideremos una asignación concreta $v$ dada por $v(p) = 1$, $v(q) = 0$ y $v(r) = 1$. Deseamos calcular el valor de verdad de la fórmula compleja:
\[
A = (p \land \neg q) \to (q \lor \neg r)
\]
Para evitar confusiones, procedemos de dentro hacia afuera, evaluando las subfórmulas según su orden jerárquico de complejidad:
\begin{enumerate}
    \item \textbf{Negaciones de letras proposicionales:}
    \begin{itemize}
        \item Como $v(q) = 0$, la inversión de la negación nos da $v(\neg q) = 1$.
        \item Como $v(r) = 1$, la inversión nos da $v(\neg r) = 0$.
    \end{itemize}
    \item \textbf{Subfórmula del antecedente ($p \land \neg q$).} Disponemos de $v(p) = 1$ y $v(\neg q) = 1$. Dado que ambos miembros son verdaderos, la conjunción es verdadera: $v(p \land \neg q) = 1$.
    \item \textbf{Subfórmula del consecuente ($q \lor \neg r$).} Tenemos $v(q) = 0$ y $v(\neg r) = 0$. Al ser ambos miembros falsos, la disyunción es falsa: $v(q \lor \neg r) = 0$.
    \item \textbf{Conectiva principal ($\to$).} El antecedente toma valor $1$ y el consecuente valor $0$. Aplicando la semántica del condicional, $v(A) = 0$.
\end{enumerate}
\textbf{Comentario.} Este procedimiento ilustra el carácter veritativo-funcional de las conectivas: el valor de verdad de la fórmula completa se calcula a partir de los valores asignados a las letras proposicionales.
\end{ejemplo}

\subsection{Tablas de verdad características de las conectivas}

Estas definiciones se resumen en tablas de verdad:

\vspace{0.3cm}
\begin{center}
\begin{tabular}{c|c}
$A$ & $\neg A$ \\
\hline
1 & 0 \\
0 & 1
\end{tabular}
\qquad
\begin{tabular}{cc|c}
$A$ & $B$ & $A \land B$ \\
\hline
1 & 1 & 1 \\
1 & 0 & 0 \\
0 & 1 & 0 \\
0 & 0 & 0
\end{tabular}
\qquad
\begin{tabular}{cc|c}
$A$ & $B$ & $A \lor B$ \\
\hline
1 & 1 & 1 \\
1 & 0 & 1 \\
0 & 1 & 1 \\
0 & 0 & 0
\end{tabular}
\qquad
\begin{tabular}{cc|c}
$A$ & $B$ & $A \to B$ \\
\hline
1 & 1 & 1 \\
1 & 0 & 0 \\
0 & 1 & 1 \\
0 & 0 & 1
\end{tabular}
\\[1cm]
\begin{tabular}{cc|c}
$A$ & $B$ & $A \leftrightarrow B$ \\
\hline
1 & 1 & 1 \\
1 & 0 & 0 \\
0 & 1 & 0 \\
0 & 0 & 1
\end{tabular}
\end{center}
\vspace{0.3cm}

\subsection{El condicional material y sus «paradojas»}\label{sec:condicional-material}

La tabla del condicional suele sorprender: $A \to B$ es verdadero siempre que $A$ sea falso o $B$ sea verdadero, haya o no alguna relación entre ambos. Este condicional veritativo-funcional se llama \textbf{condicional material} (o \emph{filónico}, por Filón de Megara, quien defendió este criterio de verdad para las proposiciones condicionales).

\paragraph{¿Por qué esta tabla?} Si queremos un condicional veritativo-funcional, la tabla está prácticamente forzada:
\begin{itemize}
    \item La fila $1, 0$ debe dar $0$: si no, el \emph{modus ponens} no sería válido.
    \item Las filas $1, 1$ y $0, 0$ deben dar $1$: si no, $p \to p$ no sería siempre verdadera.
    \item La fila $0, 1$ debe dar $1$: si diera $0$, la tabla sería la del bicondicional y «si es un perro, es un mamífero» equivaldría a «si es un mamífero, es un perro».
\end{itemize}
Además, esta tabla encaja con las reglas del tema~\ref{cap:dn-prop}: con ellas se derivan $q \vdash p \to q$ (suponemos $p$, reiteramos $q$ y aplicamos $\to$I) y $\neg p \vdash p \to q$ (suponemos $p$, obtenemos $\bot$ y, por X, $q$). Un condicional verdadero siempre que el consecuente sea verdadero o el antecedente falso es justo lo que piden esas reglas.

\paragraph{Las «paradojas».} El precio de esta elección son consecuencias que chocan con el uso ordinario de «si\ldots{} entonces». Las escribimos con la notación de la consecuencia lógica, que se define más adelante en este tema ($\Gamma \models A$ indica que toda interpretación que hace verdaderas las fórmulas de $\Gamma$ hace verdadera $A$, y $\models A$, que $A$ es verdadera en toda interpretación):
\[
q \models p \to q \qquad\qquad \neg p \models p \to q \qquad\qquad \models (p \to q) \lor (q \to p)
\]
Por ejemplo, como París está en Francia, es verdadero (materialmente) que «si la Luna es de queso, París está en Francia»; y es verdadero que «o bien si llueve hace sol, o bien si hace sol llueve». No son contradicciones, sino síntomas de que el condicional material ignora toda conexión de contenido entre antecedente y consecuente.

\paragraph{Tres respuestas filosóficas.}
\begin{enumerate} 
    \item \textbf{Pragmática (H. P. Grice).} El condicional material recoge las \emph{condiciones de verdad} de «si\ldots{} entonces», es decir, lo que se \emph{dice}. Pero al hablar comunicamos más: lo que el oyente infiere suponiendo que el hablante colabora en la conversación, las \emph{implicaturas conversacionales}. Quien afirma «si llueve, se suspenderá el partido» da a entender que tiene razones que enlazan una cosa con la otra y que no sabe todavía si lloverá ni si se suspenderá; si supiera que no va a llover, o que el partido se suspenderá, lo cooperativo sería decir eso. Esa conexión no forma parte de las condiciones de verdad, sino de lo que se da a entender al usar un condicional. Los ejemplos anteriores resultan extraños porque quien los afirmara daría a entender una conexión que no existe, aunque lo que dice sea verdadero.
    \item \textbf{Otros condicionales.} C. I. Lewis propuso la \emph{implicación estricta} (es necesario que no ocurra $A$ sin $B$), y las \emph{lógicas de la relevancia} (A. R. Anderson y N. Belnap) exigen que antecedente y consecuente compartan contenido.
    \item \textbf{Los contrafácticos no son veritativo-funcionales.} «Si hubiera soltado el vaso, se habría caído» parece verdadero y «si hubiera soltado el vaso, habría salido volando» parece falso, aunque si no he soltado el vaso ambos tienen antecedente falso. Por tanto, su valor no depende solo de los valores de sus partes. Compárense también «si Oswald no mató a Kennedy, lo hizo otra persona» (verdadero) y «si Oswald no hubiera matado a Kennedy, lo habría hecho otra persona» (probablemente falso), un ejemplo de E. Adams. Para los contrafácticos se han propuesto semánticas de mundos posibles (R. Stalnaker, D. Lewis).
\end{enumerate}
En este curso usamos el condicional material: es el que corresponde a las reglas $\to$I y $\to$E y basta para analizar la validez de una gran parte de los argumentos. Pero conviene recordar estos límites al formalizar, sobre todo los condicionales en subjuntivo (véase el ejercicio~\ref{e2-cond}).

\section{Tablas de verdad y clasificación semántica de fórmulas}

Si una fórmula contiene $n$ variables proposicionales distintas, su valor solo depende de los valores de esas $n$ variables, que pueden combinarse de $2^n$ maneras. Una \textbf{tabla de verdad} calcula el valor de la fórmula en cada uno de esos $2^n$ casos evaluando progresivamente sus subfórmulas de menor a mayor complejidad.

\begin{ejemplo}[Tabla de verdad del silogismo hipotético]
Evaluemos la fórmula $((p \to q) \land (q \to r)) \to (p \to r)$, que contiene 3 variables ($2^3 = 8$ combinaciones de valores):
\begin{center}
\small
\begin{tabular}{ccc|c|c|c|c|c}
$p$ & $q$ & $r$ & $p \to q$ & $q \to r$ & $(p \to q) \land (q \to r)$ & $p \to r$ & Fórmula completa \\
\hline
1 & 1 & 1 & 1 & 1 & 1 & 1 & \textbf{1} \\
1 & 1 & 0 & 1 & 0 & 0 & 0 & \textbf{1} \\
1 & 0 & 1 & 0 & 1 & 0 & 1 & \textbf{1} \\
1 & 0 & 0 & 0 & 1 & 0 & 0 & \textbf{1} \\
0 & 1 & 1 & 1 & 1 & 1 & 1 & \textbf{1} \\
0 & 1 & 0 & 1 & 0 & 0 & 1 & \textbf{1} \\
0 & 0 & 1 & 1 & 1 & 1 & 1 & \textbf{1} \\
0 & 0 & 0 & 1 & 1 & 1 & 1 & \textbf{1}
\end{tabular}
\end{center}
\textbf{Comentario.} Como la columna de la fórmula
completa solo contiene unos, la
fórmula es verdadera en toda interpretación, con independencia de lo
que representen sus letras: es una tautología (escribiremos $\models ((p \to q) \land (q \to
r)) \to (p \to r)$). 
\end{ejemplo}

\begin{ejemplo}[Tabla de verdad de una fórmula contingente]
Analicemos la fórmula $(p \lor q) \to (p \land q)$ ($2^2 = 4$ casos):
\begin{center}
\begin{tabular}{cc|c|c|c}
$p$ & $q$ & $p \lor q$ & $p \land q$ & $(p \lor q) \to (p \land q)$ \\
\hline
1 & 1 & 1 & 1 & \textbf{1} \\
1 & 0 & 1 & 0 & \textbf{0} \\
0 & 1 & 1 & 0 & \textbf{0} \\
0 & 0 & 0 & 0 & \textbf{1}
\end{tabular}
\end{center}
\textbf{Comentario.} La columna final combina unos y ceros. Por tanto, no es una tautología sino una fórmula contingente: es verdadera cuando $p$ y $q$ tienen el mismo valor y falsa cuando lo tienen distinto (de hecho, equivale a $p \leftrightarrow q$).
\end{ejemplo}

\subsection{Clasificación de fórmulas}

Según el resultado de su evaluación en el conjunto de todas las interpretaciones posibles, toda fórmula proposicional pertenece a exactamente una de las siguientes categorías:

\begin{itemize}
    \item \textbf{Tautología (o verdad lógica).} Una fórmula $A$ es una tautología si y solo si es verdadera en \emph{todas} las interpretaciones posibles ($v(A) = 1$ para toda $v$). Se denota:
    \[
    \models A
    \]
    En su tabla de verdad, la columna final solo contiene unos.
    \item \textbf{Contradicción.} Una fórmula $A$ es contradictoria si y solo si es falsa en \emph{todas} las interpretaciones posibles ($v(A) = 0$ para toda $v$).
    \item \textbf{Fórmula contingente.} Una fórmula que no es ni tautología ni contradicción; es verdadera en al menos una interpretación y falsa en al menos otra. Su verdad depende de los hechos que representen las variables.
\end{itemize}

\begin{definition}[Satisfacibilidad]
Una fórmula $A$ es \textbf{satisfacible} si y solo si existe al menos una interpretación $v$ tal que $v(A) = 1$, e \textbf{insatisfacible} en caso contrario. Son satisfacibles las tautologías y las fórmulas contingentes; las fórmulas insatisfacibles son, por tanto, las contradicciones.
Un conjunto de fórmulas $\Gamma$ es satisfacible si existe alguna
interpretación que haga verdaderas a la vez todas las fórmulas de $\Gamma$, e insatisfacible en caso contrario.
\end{definition}

\practica{En la aplicación \apptablas{} se rellena la tabla de verdad de una o varias fórmulas y se comprueba. Después pregunta cómo es cada fórmula, si el conjunto es satisfacible y, con dos fórmulas, si son equivalentes (sección siguiente). Su pestaña «Una interpretación» calcula paso a paso el valor de una fórmula en una interpretación.}

\section{Equivalencia lógica (\texorpdfstring{$\equiv$}{equivalencia})}

\begin{definition}[Equivalencia lógica]
Dos fórmulas $A$ y $B$ son \textbf{lógicamente equivalentes} (denotado $A \equiv B$) si y solo si toman el mismo valor de verdad en toda interpretación posible:
\[
A \equiv B \iff \text{para toda } v: (v(A) = v(B)) 
\]
Equivalentemente, $A \equiv B$ si y solo si la fórmula bicondicional $A \leftrightarrow B$ es una tautología ($\models A \leftrightarrow B$).
\end{definition}

\subsection{Leyes notables de equivalencia}

\begin{itemize}
    \item \textbf{Doble negación:} $\neg\neg A \equiv A$
    \item \textbf{Leyes de De Morgan:} 
    \[
    \neg(A \land B) \equiv \neg A \lor \neg B \qquad\qquad \neg(A \lor B) \equiv \neg A \land \neg B
    \]
    \item \textbf{Definición del condicional:} $A \to B \equiv \neg A \lor B \equiv \neg(A \land \neg B)$
    \item \textbf{Contraposición:} $A \to B \equiv \neg B \to \neg A$
    \item \textbf{Distributividad:} 
    \[
    A \land (B \lor C) \equiv (A \land B) \lor (A \land C) \qquad\qquad A \lor (B \land C) \equiv (A \lor B) \land (A \lor C)
    \]
\end{itemize}

Las leyes de De Morgan se corresponden con la regla derivada DeM del tema~\ref{cap:dn-prop}, que en una prueba permite pasar de cada lado de estas equivalencias al otro.

\section{Formalización: del lenguaje natural a las fórmulas}\label{sec:formalizacion}

Para analizar un argumento del lenguaje natural hay que \textbf{formalizarlo}: traducir sus enunciados a fórmulas que conserven lo que importa para la validez.
\begin{enumerate}
    \item Identifica los enunciados simples y asigna a cada uno una letra, formulándolos en \emph{forma afirmativa} («llueve», no «no llueve»; la negación se expresa con $\neg$). Escribe este diccionario explícitamente.
    \item Identifica las conectivas y, sobre todo, su \emph{alcance}: qué parte del enunciado afecta a cada una.
    \item Escribe la fórmula y comprueba, leyéndola de nuevo en castellano, que dice lo mismo que el original.
\end{enumerate}

Las conectivas del lenguaje natural no siempre tienen una única traducción. La tabla recoge las más frecuentes:
\begin{center}
\renewcommand{\arraystretch}{1.25}
\begin{tabular}{|p{9cm}|l|}
\hline
\textbf{Expresión} & \textbf{Formalización} \\
\hline
no $A$; no es cierto que $A$; es falso que $A$ & $\neg A$ \\
$A$ y $B$; $A$ pero $B$; $A$ aunque $B$; $A$, sin embargo, $B$ & $A \land B$ \\
$A$ o $B$ (inclusiva: quizá ambas) & $A \lor B$ \\
o bien $A$ o bien $B$ (si el contexto indica que no ambas) & $A \leftrightarrow \lnot B$ \\
ni $A$ ni $B$ & $\neg A \land \neg B$ \\
no a la vez $A$ y $B$ & $\neg(A \land B)$ \\
si $A$, $B$; $B$ si $A$; cuando $A$, $B$; siempre que $A$, $B$ & $A \to B$ \\
$A$ solo si $B$; solo si $B$, $A$ & $A \to B$ \\
$A$ es condición suficiente para $B$ & $A \to B$ \\
$A$ es condición necesaria para $B$ & $B \to A$ \\
$A$ si y solo si $B$; $A$ exactamente cuando $B$ & $A \leftrightarrow B$ \\
$A$ a menos que $B$; $A$ salvo que $B$ & $\neg B \to A$ \quad (equivale a $A \lor B$) \\
\hline
\end{tabular}
\end{center}

Algunas observaciones:
\begin{itemize}
    \item \textbf{«Pero», «aunque», «sin embargo».} Todas se formalizan con $\land$: el matiz de contraste que añaden no afecta a la verdad del enunciado ni a la validez.
    \item \textbf{«Solo si».} Es la fuente de errores más frecuente. «Voy al cine (c) solo si acabo pronto (a)» ($c \to a$) dice que acabar pronto es condición \emph{necesaria} para ir al cine, no que baste con acabar pronto para ir: no equivale a $a \to c$.
    \item \textbf{«A menos que».} «Iré a la playa a menos que llueva» se formaliza $\neg l \to p$ (o $p \lor l$): si no llueve, iré. A menudo se entiende además que si llueve no iré ($p \leftrightarrow \neg l$); como esa lectura más fuerte no siempre es la intencionada, la formalización habitual es la débil.
    \item \textbf{Alcance de la negación.} «No es cierto que Juan venga y María se quede» es $\neg(j \land m)$; «Juan no viene y María se queda» es $\neg j \land m$.
    \item \textbf{Condicionales en subjuntivo.} Su formalización con $\to$ es solo aproximada (véase la sección~\ref{sec:condicional-material}).
\end{itemize}

\begin{ejemplo}[Formalización de un enunciado filosófico]
«Si el determinismo es verdadero y no hay libre albedrío, la responsabilidad moral es una ilusión, a menos que el compatibilismo sea correcto.»

\noindent Diccionario: $d$: «el determinismo es verdadero»; $l$: «hay libre albedrío»; $r$: «la responsabilidad moral es una ilusión»; $c$: «el compatibilismo es correcto». El «a menos que» afecta a todo el condicional:
\[
\neg c \to ((d \land \neg l) \to r)
\]
\end{ejemplo}

\practica{La pestaña «Formalizar» de \apptablas{} propone frases con su diccionario y compara tu fórmula con una formalización correcta en todas las interpretaciones. Si no son equivalentes, muestra una interpretación en la que difieren.}

\section{Consecuencia lógica (\texorpdfstring{$\models$}{consecuencia})}

La consecuencia lógica es el concepto central de la
semántica. Formaliza la noción de argumento deductivamente válido:  

\begin{definition}[Consecuencia lógica]
Sean $\Gamma$ un conjunto de fórmulas (premisas) y $A$ una fórmula (conclusión). Decimos que $A$ es \textbf{consecuencia lógica} de $\Gamma$, y escribimos
\[
\Gamma \models A
\]
si y solo si no existe ninguna interpretación $v$ que haga verdaderas todas las premisas de $\Gamma$ y falsa la conclusión $A$.

Dicho en términos positivos: toda interpretación que hace verdaderas
todas las premisas de $\Gamma$ hace verdadera la
conclusión $A$. Si $\Gamma$ es finito, en la tabla de verdad conjunta de todas las fórmulas
de $\Gamma$ y $A$ no hay ninguna fila que haga verdaderas todas
las fórmulas de $\Gamma$ y falsa $A$. 
\end{definition}

\subsection{Contraejemplos semánticos y el método de asignación directa}

Para comprobar si se cumple $\Gamma \models A$, construir
la tabla de verdad completa resulta largo cuando hay muchas variables proposicionales. El \textbf{método de asignación directa} (o reducción al absurdo semántica) suele ser más rápido: \emph{intentar construir un contraejemplo} y ver si el intento fracasa. Cuando una condición admite varias posibilidades (por ejemplo, una disyunción que debe ser verdadera o una conjunción que debe ser falsa), hay que examinar cada caso por separado: el argumento es válido si y solo si \emph{todos} los casos conducen a contradicción.

\begin{definition}[Contraejemplo]
Un contraejemplo para la inferencia $\Gamma \models A$ es una interpretación concreta $v$ tal que:
\begin{enumerate}
    \item $v(B) = 1$ para toda premisa $B \in \Gamma$.
    \item $v(A) = 0$ (la conclusión es falsa).
\end{enumerate}
\end{definition}


\begin{ejemplo}[Demostración de invalidez mediante contraejemplo: falacia de afirmación del consecuente]
Consideremos el argumento con premisas $p \to q$ y $q$ y conclusión $p$. Para averiguar si se cumple $p \to q, q \models p$, intentamos construir un contraejemplo:
\begin{enumerate}
    \item Forzamos la conclusión como falsa: exigimos que $v(p) = 0$.
    \item Forzamos la segunda premisa como verdadera: exigimos que $v(q) = 1$.
    \item Comprobamos si la primera premisa puede ser simultáneamente verdadera: evaluamos $v(p \to q) = 1$.
\end{enumerate}
La asignación obtenida $v = \{p: 0, q: 1\}$ hace verdaderas todas las premisas y falsa la conclusión sin generar ningún conflicto. Por consiguiente, $v$ es un contraejemplo y el argumento es inválido: $\{p \to q, q\} \not\models p$.
\end{ejemplo}

\begin{ejemplo}[Demostración de validez por reducción al absurdo semántica: \emph{modus tollens}]
Analicemos si se cumple $p \to q, \neg q \models \neg p$.
Intentamos construir un contraejemplo:
\begin{enumerate}
    \item Forzamos la conclusión a falso: $v(\neg p) = 0$, lo que obliga a que $v(p) = 1$.
    \item Forzamos la segunda premisa a verdadero: $v(\neg q) = 1$, lo que obliga a que $v(q) = 0$.
    \item Evaluamos ahora la primera premisa con los valores
      obligados: $v(p \to q) =  0$.
    Esto genera una contradicción: la primera premisa toma necesariamente el valor $0$, por lo que resulta imposible hacerla verdadera simultáneamente con las restantes condiciones.
\end{enumerate}
\textbf{Comentario.} Como no puede construirse ningún contraejemplo, la inferencia es válida ($p \to q, \neg q \models \neg p$), y lo hemos comprobado sin construir la tabla de verdad completa.
\end{ejemplo}

\practica{En la pestaña «Consecuencia lógica» de \apptablas{} se rellena la tabla conjunta de las premisas y la conclusión, se marcan los contraejemplos y se decide si el argumento es válido. También resuelve el argumento por asignación directa, paso a paso.}

\section{Propiedades fundamentales de la consecuencia lógica}

La relación $\models$ satisface las siguientes propiedades (las tres primeras se llaman \emph{estructurales}):

\begin{enumerate}
    \item \textbf{Reflexividad.} Cualquier fórmula es consecuencia
      lógica de un conjunto que la contenga:
    \[
    \Gamma \cup \{A\} \models A
    \]
    \item \textbf{Monotonía.} Si una conclusión se sigue de ciertas premisas, sigue siendo consecuencia si añadimos premisas adicionales:
    \[
    \text{Si } \Gamma \models A, \text{ entonces } \Gamma \cup \Delta \models A
    \]
    \item \textbf{Transitividad o corte.} Si $A$ se sigue de $\Gamma \cup \{B\}$ y a su vez $B$ es consecuencia de $\Delta$, entonces $A$ es consecuencia de $\Gamma \cup \Delta$:
    \[
    \frac{\Gamma \cup \{B\} \models A \qquad \Delta \models B}{\Gamma \cup \Delta \models A}
    \]
    \item \textbf{Teorema de la deducción (versión semántica).} Una
      fórmula $B$ se sigue de un conjunto ampliado con $A$ si y solo
      si el condicional $A \to B$ se sigue del conjunto original:
    \[
    \Gamma \cup \{A\} \models B \iff \Gamma \models A \to B
    \]
    En particular, $\{B_1, \dots, B_k\} \models A$ si y solo si
    $\models (B_1 \land \dots \land B_k) \to A$. (Compruébalo: ambas
    afirmaciones dicen que no hay ninguna interpretación que haga
    verdaderas las $B_i$ y falsa $A$). 
\end{enumerate}

\subsection{El principio de refutación (conexión con contradicciones)}

La consecuencia lógica puede expresarse en términos de satisfacibilidad:
\[
\Gamma \models A \iff \Gamma \cup \{\neg A\} \text{ es insatisfacible}
\]
Un argumento es válido si y solo si suponer que las premisas son
verdaderas y la conclusión es falsa genera una contradicción. Esta
propiedad es la contrapartida semántica de la regla de reducción al
absurdo (IP) del cálculo y la base del método de tableaux (tema~\ref{cap:tableaux}) y
de muchos métodos automáticos de demostración. 

\section{Hacia la metalógica: el diálogo entre sintaxis y semántica}

Al término de estos dos temas disponemos de dos herramientas independientes para evaluar argumentos:
\begin{itemize}
    \item \textbf{Deducibilidad:} $\Gamma \vdash A$ (demostrabilidad paso a paso mediante reglas).
    \item \textbf{Consecuencia lógica:} $\Gamma \models
      A$ (preservación de la verdad en todas las interpretaciones posibles).
\end{itemize}

Surgen así dos preguntas que constituyen el núcleo de la metalógica:
\begin{enumerate}
    \item \textbf{¿Es nuestro cálculo correcto?} Si derivamos $A$ a partir de $\Gamma$ en el cálculo ($\Gamma \vdash A$), ¿tenemos la garantía de que $A$ es consecuencia lógica de $\Gamma$ ($\Gamma \models A$)?
    \item \textbf{¿Es nuestro cálculo completo?} Si $A$ es consecuencia lógica de $\Gamma$ ($\Gamma \models A$), ¿existe siempre una derivación de $A$ a partir de $\Gamma$ en el cálculo ($\Gamma \vdash A$)?
\end{enumerate}
Estas dos cuestiones son el contenido del
siguiente tema.

\section*{Referencias bibliográficas}

\begin{enumerate}
    \item \textbf{Adams, E. W.} (1970). «Subjunctive and indicative conditionals». \emph{Foundations of Language}, vol. 6, n.º 1, pp. 89--94. [El ejemplo de Oswald y Kennedy].
    \item \textbf{Anderson, A. R. y Belnap, N. D.} (1975). \emph{Entailment: The Logic of Relevance and Necessity}, vol. I. Princeton University Press, Princeton. [Obra de referencia de las lógicas de la relevancia].
    \item \textbf{Bennett, J.} (2003). \emph{A Philosophical Guide to Conditionals}. Clarendon Press, Oxford. [Introducción general a los debates sobre condicionales indicativos y contrafácticos].
    \item \textbf{Grice, H. P.} (1975). «Logic and conversation». En P. Cole y J. L. Morgan (eds.), \emph{Syntax and Semantics 3: Speech Acts}. Academic Press, Nueva York, pp. 41--58. [Las implicaturas conversacionales y la defensa pragmática del condicional material].
    \item \textbf{Lewis, C. I.} (1912). «Implication and the algebra of logic». \emph{Mind}, vol. 21, n.º 84, pp. 522--531. [Crítica de la implicación material y primera propuesta de la implicación estricta].
    \item \textbf{Lewis, C. I.} (1918). \emph{A Survey of Symbolic Logic}. University of California Press, Berkeley. [Primera presentación sistemática de la implicación estricta].
    \item \textbf{Lewis, D.} (1973). \emph{Counterfactuals}. Blackwell, Oxford. [Semántica de los contrafácticos basada en la semejanza entre mundos posibles].
    \item \textbf{Stalnaker, R. C.} (1968). «A theory of conditionals». En N. Rescher (ed.), \emph{Studies in Logical Theory}. Blackwell, Oxford, pp. 98--112. [Semántica de mundos posibles para los condicionales].
\end{enumerate}

%%% Local Variables:
%%% mode: latex
%%% TeX-master: "../ApuntesLogica"
%%% End:


\section{Ejercicios de semántica de la lógica proposicional}
% -*- coding: utf-8 -*-
\begin{quote}
\small \textbf{Instrucciones.} Puedes comprobar las tablas de verdad, las equivalencias, las consecuencias y las formalizaciones con la aplicación \apptablas{}. Rellena antes la tabla a mano: la aplicación te dirá si es correcta.
\end{quote}

\section*{Bloque 1: tablas de verdad y clasificación de fórmulas}

\begin{pregunta}\label{e2-tablas}
Construye la tabla de verdad de las siguientes fórmulas y clasifícalas en tautología ($\models A$), contradicción o fórmula contingente:
\begin{enumerate}[(a)]
    \item $(p \to q) \leftrightarrow (\neg q \to \neg p)$ \hfill (Principio de contraposición)
    \item $(p \to q) \land (p \land \neg q)$
    \item $((p \to q) \to p) \to p$ \hfill (Ley de Peirce)
    \item $(p \to q) \to (q \to p)$ \hfill (Recíproca del condicional)
    \item $(p \to q) \lor (q \to p)$
\end{enumerate}
\end{pregunta}

\begin{pregunta}Demuestra mediante tablas de verdad que las siguientes fórmulas son contingentes proponiendo, para cada una de ellas, una interpretación que la haga verdadera y otra que la haga falsa:
\begin{enumerate}[(a)]
    \item $(p \land q) \to r$
    \item $p \to (q \lor r)$
\end{enumerate}
\end{pregunta}

\section*{Bloque 2: equivalencia lógica (\texorpdfstring{$\equiv$}{equivalencia})}

\begin{pregunta}Demuestra mediante tablas de verdad las siguientes
  leyes de De Morgan: 
\begin{enumerate}[(a)]
    \item $\neg(p \lor q) \equiv \neg p \land \neg q$
    \item $\neg(p \land q) \equiv \neg p \lor \neg q$
\end{enumerate}
\end{pregunta}

\begin{pregunta}\label{e2-negcond}
Encuentra una fórmula lógicamente equivalente a $p \land q$ en la que únicamente se utilicen la negación ($\neg$), el condicional ($\to$) y paréntesis. Justifica semánticamente la respuesta demostrando que ambas fórmulas tienen la misma tabla de verdad.
\end{pregunta}

\begin{pregunta}\label{e2-negdisy}
Encuentra una fórmula equivalente a $p \leftrightarrow q$ en la que únicamente aparezcan las conectivas $\{\neg, \lor\}$.
\end{pregunta}

\section*{Bloque 3: consecuencia lógica semántica (\texorpdfstring{$\models$}{consecuencia}) y contraejemplos}

\begin{pregunta}\label{e2-consec}
Determina si se cumplen o no las siguientes relaciones de consecuencia lógica ($\models$). Si la relación es válida, justifícalo semánticamente y construye su derivación en deducción natural ($\vdash$); si es inválida, proporciona un contraejemplo explícito (una asignación $v$ que satisfaga las premisas y haga falsa la conclusión):
\begin{enumerate}[(a)]
    \item $p \to q, \neg p \models \neg q$ \hfill (Falacia de negación del antecedente)
    \item $p \to q, r \to q \models p \to r$ \hfill (Análogo proposicional del término medio no distribuido)
    \item $p \lor q, \neg p \models q$ \hfill (Silogismo disyuntivo)
    \item $p \to (q \to r), p \land \neg r \models \neg q$
    \item $p \to q, r \to q \models (p \lor r) \to q$
\end{enumerate}
\end{pregunta}

\section*{Bloque 4: formalización, tablas de verdad y puente hacia la deducción}

\begin{pregunta}[Conectivas del castellano]\label{e2-formal}
Formaliza los siguientes enunciados (consulta la sección~\ref{sec:formalizacion}), indicando el diccionario de letras que utilizas:
\begin{enumerate}[(a)]
    \item Iré a la fiesta solo si me invitan.
    \item Aprobar el examen es condición necesaria para aprobar la asignatura (pero no es condición suficiente).
    \item Saldremos de excursión a menos que llueva.
    \item Ni Kant ni Hume aceptan el argumento ontológico.
    \item El conocimiento no puede ser a la vez innato y adquirido por la experiencia.
    \item Hume es empirista si y solo si rechaza las ideas innatas.
\end{enumerate}
\end{pregunta}

\begin{pregunta}[Argumentos del lenguaje ordinario]\label{e2-arg}
Formaliza las premisas y la conclusión de los siguientes argumentos. A continuación:
\begin{enumerate}[1.]
    \item Comprueba mediante una tabla de verdad si el argumento es semánticamente válido ($\Gamma \models A$) o no. Si no lo es, indica qué fila de la tabla actúa como contraejemplo.
    \item Para aquellos argumentos que resulten ser válidos,
      construye su demostración formal en el cálculo de
        deducción natural ($\Gamma \vdash A$) utilizando las reglas
      aprendidas en el tema~\ref{cap:dn-prop}.
\end{enumerate}

\begin{enumerate}[(a)]
    \item Voy al cine solo si acabo pronto de limpiar la casa. Si me acompañas, voy al cine. Por tanto, si me acompañas, acabo pronto de limpiar la casa.
    \item Si hace buen tiempo, iremos a la playa. Iremos a la playa. Por tanto, hace buen tiempo.
    \item Si el mundo es redondo y tengo un barco, entonces navego alrededor del mundo. Tengo un barco, pero no navego alrededor del mundo. Por tanto, el mundo no es redondo.
    \item Cuando las condiciones climáticas son las adecuadas, salen setas. Si salen setas, los bosques se llenan de buscadores de setas. Los bosques están llenos de buscadores de setas. En consecuencia, las condiciones climáticas son las adecuadas.
    \item Si el presentador fuese competente, habría mantenido la audiencia, y si lo hubiese hecho, la empresa lo habría promocionado. Sin embargo, la empresa no lo ha promocionado. Así pues, el presentador no es competente.
    \item Si el espacio es euclídeo, se cumple el quinto postulado. Pero si el espacio es elíptico, no se cumple el quinto postulado. Por tanto, el espacio no es a la vez euclídeo y elíptico.
\end{enumerate}
\end{pregunta}

\begin{pregunta}[Argumentos filosóficos clásicos]\label{e2-filo}
Formaliza los siguientes argumentos filosóficos y determina su validez semántica mediante tablas de verdad o búsqueda de contraejemplos. Si son válidos, demuestra la conclusión mediante deducción natural:
\begin{enumerate}[(a)]
    \item Si el determinismo es verdadero, entonces las acciones humanas son previsibles. Si las acciones humanas son previsibles, no hay libre albedrío. Por tanto, si el determinismo es verdadero, no hay libre albedrío.
    \item Si el conocimiento requiere certeza absoluta, entonces no tenemos conocimiento del mundo empírico. Pero tenemos conocimiento del mundo empírico. En consecuencia, el conocimiento no requiere certeza absoluta.
\end{enumerate}
\end{pregunta}

\section*{Bloque 5: razonamiento metateórico sobre verdad y validez}

\begin{pregunta}\label{e2-vf1}
Indica si la siguiente afirmación es verdadera o falsa y justifica detalladamente tu respuesta:
\begin{quote}
«Sea $A$ una fórmula y $\Gamma$ un conjunto de fórmulas. Si $A$ no es consecuencia lógica de $\Gamma$ ($\Gamma \not\models A$), entonces la negación de $A$ es consecuencia lógica de $\Gamma$ ($\Gamma \models \neg A$)».
\end{quote}
\end{pregunta}

\begin{pregunta}\label{e2-vf2}
Indica si la siguiente afirmación es verdadera o falsa y justifica detalladamente tu respuesta:
\begin{quote}
«Sea $A$ una fórmula y $\Gamma$ un conjunto de premisas. Si $\Gamma \models A$, entonces necesariamente $\Gamma \not\models \neg A$».
\end{quote}
\end{pregunta}

\begin{pregunta}\label{e2-insat}
Sean $A$ y $B$ dos fórmulas tales que el conjunto $\{\neg A, B\}$ no es satisfacible. ¿Qué podemos afirmar necesariamente acerca de la fórmula $B \to A$?
\begin{enumerate}[(a)]
    \item Es una contradicción (insatisfacible).
    \item Es contingente.
    \item Es una tautología ($\models B \to A$).
    \item No se puede determinar sin conocer las variables de $A$ y $B$.
\end{enumerate}
Justifica tu elección demostrando la relación a partir de la definición de satisfacibilidad.
\end{pregunta}

\begin{pregunta}
Sea $A \to B$ una fórmula contradictoria (no satisfacible). ¿Qué podemos afirmar de la fórmula $B \to A$? Justifica la respuesta evaluando qué valores de verdad deben tener necesariamente $A$ y $B$ para que $A \to B$ sea falsa en toda interpretación.
\end{pregunta}

\begin{pregunta}[Cuestiones sobre consecuencia y satisfacibilidad]\label{e2-cuest}
Responde razonadamente a las siguientes cuestiones conceptuales:
\begin{enumerate}[(a)]
    \item Supongamos que la fórmula $A$ es una contradicción y la fórmula $B$ es una tautología. Determina justificadamente si las siguientes relaciones de consecuencia lógica son verdaderas o falsas:
    \begin{itemize}
        \item $\{A\} \models B$
        \item $\{B\} \models A$
    \end{itemize}
    \item Sea $\{A, B, \neg C\}$ un conjunto insatisfacible de fórmulas. Determina justificadamente cuáles de las siguientes relaciones de consecuencia se cumplen necesariamente:
    \begin{itemize}
        \item $\{A, B\} \models C$
        \item $\{A, \neg C\} \models \neg B$
        \item $\{A, B\} \models \neg C$
    \end{itemize}
\end{enumerate}
\end{pregunta}

\begin{pregunta}[El condicional material]\label{e2-cond}
Responde razonadamente (consulta la sección~\ref{sec:condicional-material}):
\begin{enumerate}[(a)]
    \item Comprueba mediante tablas de verdad que $q \models p \to q$ y que $\neg p \models p \to q$.
    \item Según el condicional material, ¿es verdadero el enunciado «si $2 + 2 = 5$, entonces Madrid es la capital de Francia»? ¿Y «si Madrid es la capital de España, entonces $2 + 2 = 4$»? ¿Te parece que estos resultados muestran que el condicional material es inadecuado, o que con «si\ldots{} entonces» comunicamos algo más que sus condiciones de verdad?
    \item En el argumento (e) del ejercicio~\ref{e2-arg} los condicionales están en subjuntivo. ¿Qué se pierde al formalizarlos con $\to$? ¿Crees que eso afecta a la validez de ese argumento concreto?
\end{enumerate}
\end{pregunta}

%%% Local Variables:
%%% mode: latex
%%% TeX-master: "../ApuntesLogica"
%%% End:


\chapter{Metalógica proposicional: corrección y completud}\label{cap:meta-prop}
\begin{resumentema}
En este tema cerramos el ciclo de la lógica proposicional abordando
sus dos resultados metateóricos principales: el teorema de corrección y el
teorema de completud. Demostramos que nuestro cálculo de deducción
natural ($\vdash$) y la semántica ($\models$)
coinciden: todo lo que se deduce de unas premisas es consecuencia
lógica de ellas y, a la inversa, toda consecuencia lógica puede
deducirse. Para la completud
presentamos la demostración constructiva de László Kalmár
(1935), con ejemplos que muestran paso a paso
cómo transformar una tabla de verdad en una deducción formal. 
\end{resumentema}

% -*- coding: utf-8 -*-
\section{El problema de la adecuación de un sistema formal}

A lo largo de los dos temas precedentes hemos desarrollado dos herramientas distintas para analizar los argumentos:
\begin{enumerate}
    \item \textbf{La perspectiva sintáctica:}
      $\Gamma \vdash A$  (tema~\ref{cap:dn-prop}). Existe una derivación formal en el cálculo de
      deducción natural que desde las premisas $\Gamma$ alcanza la
      conclusión $A$ mediante la aplicación de reglas sintácticas de
      introducción y eliminación de conectivas. 
    \item \textbf{La perspectiva semántica:} $\Gamma \models A$ (tema~\ref{cap:sem-prop}). En toda asignación de verdad $v$, si todas las premisas de $\Gamma$ son verdaderas, la conclusión $A$ es verdadera (no existen contraejemplos).
\end{enumerate}

La pregunta central de la metalógica es: \emph{¿qué relación
  guardan entre sí $\vdash$ y $\models$?} Para que un sistema formal (en nuestro caso, el cálculo de deducción natural, cuya relación de derivabilidad es $\vdash$)
sea adecuado, debe tener dos propiedades: la
\textbf{corrección} (no demostrar nada que no se siga lógicamente de las premisas) y la \textbf{completud} (poder demostrar todo lo que se sigue lógicamente de ellas).

En ambos casos conviene distinguir entre una versión \textbf{débil} (referida a fórmulas sin premisas, es decir, a teoremas y tautologías) y una versión \textbf{fuerte} (referida a inferencias a partir de un conjunto arbitrario de premisas $\Gamma$).

\subsection{Corrección débil frente a corrección fuerte}
\begin{itemize}
    \item \textbf{Corrección débil.} Si una fórmula se demuestra formalmente sin premisas en el cálculo, entonces es una verdad lógica (tautología):
    \[
    \vdash A \quad\implies\quad \models A
    \]
    Garantiza que nuestro cálculo no permite demostrar sin premisas ninguna fórmula que no sea una tautología.
    
    \item \textbf{Corrección fuerte.} Para cualquier conjunto de premisas $\Gamma$ (finito o infinito) y cualquier fórmula $A$:
    \[
    \Gamma \vdash A\quad \implies\quad \Gamma \models A
    \]
    Garantiza que el cálculo preserva la verdad paso a paso en presencia de cualquier conjunto de premisas o hipótesis: si partimos de supuestos verdaderos, cualquier conclusión derivada en deducción natural será verdadera.
\end{itemize}
La corrección fuerte implica trivialmente la corrección débil tomando
el conjunto vacío de premisas ($\Gamma = \varnothing$).

\subsection{Completud débil frente a completud fuerte}
\begin{itemize}
    \item \textbf{Completud débil.} Si una fórmula es una tautología
      (verdadera bajo toda asignación de verdad), entonces es
      demostrable sin premisas como un teorema del cálculo:
    \[
    \models A \quad\implies\quad \vdash A
    \]
    Garantiza que el cálculo alcanza todas las verdades lógicas: cualquier tautología puede derivarse con nuestras reglas de inferencia.
    
    \item \textbf{Completud fuerte.} Para cualquier conjunto de premisas $\Gamma$ y cualquier fórmula $A$:
    \[
    \Gamma \models A \quad\implies\quad \Gamma \vdash A
    \]
    Garantiza que siempre que un argumento sea semánticamente válido
    (la conclusión es consecuencia lógica de las premisas, es decir,
    la conclusión es verdadera en toda interpretación donde las
    premisas sean verdaderas), existe una prueba
    formal que lo justifica paso a paso en deducción natural.
\end{itemize}

Las versiones débiles, por sí solas, se quedan cortas: en la práctica
rara vez nos limitamos a
constatar tautologías (aunque veremos que, para conjuntos finitos de
premisas, la versión fuerte se obtiene fácilmente de la débil). Casi siempre razonamos a partir de
conjuntos de premisas $\Gamma$ (hipótesis empíricas, principios
éticos, axiomas). La completud fuerte asegura que la
deducción natural permite derivar cualquier fórmula que sea
consecuencia lógica de cualquier conjunto de premisas. 

Cuando un sistema formal satisface corrección y completud (en sus versiones fuertes), decimos que es \textbf{adecuado}, pues entonces coinciden exactamente la noción sintáctica de derivabilidad y la noción semántica de consecuencia lógica:
\[
\Gamma \vdash A \quad\iff\quad \Gamma \models A
\]

\subsection{La inducción como método de demostración}\label{sec:induccion}

Los teoremas de este tema hablan de \emph{todas} las derivaciones o de \emph{todas} las fórmulas, y hay infinitas: no se pueden comprobar una a una. Que una propiedad se cumpla en todos los casos que hemos probado no garantiza que se cumpla siempre. Para demostrarlo se usa la \textbf{inducción}.

El modelo es la inducción sobre los números naturales. Para probar que una propiedad $\varphi(n)$ vale para todo $n \in \mathbb{N}$ basta probar dos cosas:
\begin{itemize}
    \item \textbf{Caso base.} Se cumple $\varphi(0)$.
    \item \textbf{Paso inductivo.} Para cualquier $n$, si se cumple $\varphi(n)$, también se cumple $\varphi(n+1)$. Lo que se supone, $\varphi(n)$, es la \textbf{hipótesis de inducción} (HI).
\end{itemize}
El caso base da $\varphi(0)$, y el paso inductivo lleva de $\varphi(0)$ a $\varphi(1)$, de $\varphi(1)$ a $\varphi(2)$, y así hasta cualquier $n$. El caso base no tiene por qué ser el $0$: si la propiedad se quiere probar a partir de otro número, por ejemplo el $1$, el caso base es $\varphi(1)$ y la conclusión vale para todo $n \geq 1$. Es lo que ocurre en el ejemplo siguiente.

\begin{ejemplo}[Una fórmula que detecta la paridad]
Definimos $\Phi_1 = p_1$ y $\Phi_{n+1} = \neg\Phi_n \leftrightarrow p_{n+1}$; por ejemplo, $\Phi_2 = \neg p_1 \leftrightarrow p_2$. Queremos probar que, para todo $n$, $\Phi_n$ es verdadera exactamente cuando es verdadero un número impar de las variables $p_1, \dots, p_n$. Con tablas de verdad haría falta una tabla de $2^n$ filas para cada $n$; por inducción se prueba para todos a la vez. Sea $\varphi(n)$ la propiedad: «para toda interpretación $v$, $v(\Phi_n) = 1$ si y solo si $v$ hace verdadero un número impar de variables entre $p_1, \dots, p_n$».
\begin{itemize}
    \item \textbf{Caso base.} $v(\Phi_1) = 1$ si y solo si $v(p_1) = 1$, es decir, si es verdadera una de las variables de $\{p_1\}$, y uno es impar.
    \item \textbf{Paso inductivo.} Supongamos $\varphi(n)$ (HI) y sea $v$ una interpretación cualquiera. Hay dos casos.
    \begin{itemize}
        \item Si $v(\Phi_n) = 1$, por la HI es impar el número de variables verdaderas entre $p_1, \dots, p_n$. Como $v(\neg\Phi_n) = 0$, $v(\Phi_{n+1}) = 1$ si y solo si $v(p_{n+1}) = 0$, es decir, si al añadir $p_{n+1}$ el número sigue siendo impar.
        \item Si $v(\Phi_n) = 0$, por la HI ese número es par, y $v(\Phi_{n+1}) = 1$ si y solo si $v(p_{n+1}) = 1$, es decir, si al añadir $p_{n+1}$ pasa a ser impar.
    \end{itemize}
    En los dos casos se cumple $\varphi(n+1)$.
\end{itemize}
\end{ejemplo}

El mismo esquema sirve para objetos que se construyen paso a paso a partir de otros más simples. En este tema se usa de dos formas, y en las dos la hipótesis de inducción supone la propiedad para \emph{todos} los casos más simples, no solo para el anterior:
\begin{itemize}
    \item \textbf{Inducción sobre la longitud de las derivaciones} (teorema de corrección, sección~\ref{sec:correccion}). Se prueba que la propiedad vale para cada línea suponiendo que vale para todas las anteriores. El caso base son las premisas y las hipótesis, que no se obtienen de otras líneas; el paso inductivo se comprueba regla por regla.
    \item \textbf{Inducción sobre la complejidad de las fórmulas} (lema de Kalmár, \ref{lema_kalmar}). La complejidad de una fórmula es su número de conectivas. El caso base son las fórmulas atómicas; en el paso inductivo se supone la propiedad para las subfórmulas inmediatas (sección~\ref{sec:precedencia}) y se comprueba, conectiva por conectiva, para la fórmula compuesta.
\end{itemize}
Como cada derivación tiene un número finito de líneas y cada fórmula se construye en un número finito de pasos, así quedan cubiertas todas las derivaciones y todas las fórmulas.

\section{El teorema de corrección}\label{sec:correccion}

Demostraremos en primer lugar que el cálculo de deducción natural proposicional es fuertemente correcto:

\begin{theorem}[Corrección fuerte]
Para cualquier conjunto de premisas $\Gamma$ (finito o infinito) y cualquier fórmula proposicional $A$:
\[
\text{Si } \Gamma \vdash A, \text{ entonces } \Gamma \models A.
\]
\end{theorem}
Como consecuencia inmediata, tomando $\Gamma = \varnothing$, obtenemos la corrección débil: todo teorema del cálculo es una tautología ($\vdash A \implies \models A$). El teorema sirve además para demostrar que una fórmula \emph{no} es derivable, algo que no prueba el mero hecho de no encontrar una derivación: si hay un contraejemplo, es decir, si $\Gamma \not\models A$, entonces $\Gamma \nvdash A$ («$A$ no es derivable de $\Gamma$»). Así se razona en el ejercicio~\ref{e3-4}.

\begin{proof}
La demostración procede por inducción sobre el número de pasos (longitud) de la deducción. 

Demostraremos que en cada paso de una prueba formal se cumple la siguiente \textbf{propiedad invariante}:
\begin{quote}
\textit{Cualquier interpretación $v$ que satisfaga todas las hipótesis no descargadas activas en la línea $k$ (lo que incluye las premisas) hace también verdadera a la fórmula de la línea $k$.}
\end{quote}

\textbf{Caso base.} Si la línea $k$ es una premisa o una hipótesis provisional recién abierta, la propiedad se cumple trivialmente, pues $v$ hace verdadera la fórmula por propia definición de satisfacer dicha hipótesis.

\textbf{Paso inductivo.} Asumimos como hipótesis de inducción (HI) que
todas las líneas anteriores satisfacen la propiedad
invariante. Obsérvese que, si una línea $j$ es accesible desde la
línea $k$, las hipótesis activas en $j$ también lo están en $k$; por
eso, toda interpretación que satisface las hipótesis activas en $k$
hace verdaderas las líneas citadas. (Este es el motivo de que no se
puedan citar líneas de cajas ya cerradas). Si lo que se cita es una
caja $i$--$j$, las hipótesis activas en su última línea $j$ son, como
mucho, las activas en $k$ más la hipótesis de la propia caja (la de la
línea $i$); así se razona en los casos de $\to$I, $\neg$I e IP, y del
mismo modo en $\lor$E y $\leftrightarrow$I. Analizamos qué ocurre según
la regla que se aplica en la línea $k$:

\begin{itemize}
    \item \textbf{Regla $\land$I.} La línea $k$ es $A \land B$, justificada a partir de líneas previas $A$ y $B$. Por HI, cualquier interpretación $v$ que satisfaga las hipótesis activas hace verdaderas a $A$ y a $B$. Por la semántica de la conjunción, $v(A) = 1$ y $v(B) = 1$ implica $v(A \land B) = 1$.
    
    \item \textbf{Regla $\to$E (\emph{modus ponens}).} La línea $k$ es
      $B$, obtenida a partir de $A \to B$ y $A$. Por HI, bajo
      cualquier interpretación $v$ que satisfaga las hipótesis activas, $v(A \to B) = 1$ y $v(A) = 1$. Por la tabla del condicional, la única forma de que $A \to B$ sea verdadero siendo $A$ verdadero es que $B$ sea verdadero. Luego $v(B) = 1$.

    \item \textbf{Regla $\to$I (descarga de hipótesis).} Supongamos
      que en la línea $k$ cerramos una caja que se abrió con la
      hipótesis $A$ y alcanzó la conclusión $B$, deduciendo $A \to B$
      en el nivel exterior. Sea $\Delta$ el conjunto de hipótesis
      activas fuera de la caja. Por HI aplicada a la última línea
      dentro de la caja, toda interpretación $v$ que satisfaga $\Delta \cup \{A\}$ satisface $B$. Esto significa que
    \[
    \text{Si } v \text{ satisface } \Delta \text{ y } v(A) = 1, \text{ entonces } v(B) = 1.
    \]
    Por consiguiente, si $v$ satisface $\Delta$, no puede ocurrir simultáneamente que $v(A) = 1$ y $v(B) = 0$. Por definición semántica del condicional, esto garantiza que $v(A \to B) = 1$. La regla $\to$I preserva la corrección al descargar la hipótesis.

    \item \textbf{Regla $\neg$I (introducción de la negación).}
      Asumimos la hipótesis $A$ y alcanzamos $\bot$. Por HI, toda
      interpretación $v$ que satisfaga las hipótesis exteriores
      $\Delta$ y además satisfaga $A$ satisface $\bot$. Pero
      ninguna interpretación posible satisface $\bot$ ($v(\bot) = 0$
      siempre). Por tanto, ninguna interpretación que satisfaga
      $\Delta$ puede hacer verdadera a $A$, de donde toda
      interpretación que satisfaga $\Delta$ debe hacer falsa a $A$, es
      decir, $v(\neg A) = 1$. 

    \item \textbf{Reglas $\neg$E y X.} Si la línea $k$ es $\bot$, obtenida de $A$ y $\neg A$, ninguna interpretación puede satisfacer las hipótesis activas (haría verdaderas a la vez $A$ y $\neg A$), de modo que la propiedad invariante se cumple trivialmente. Por la misma razón, la regla X (que parte de una línea $\bot$) también la conserva.

    \item \textbf{Regla IP.} Se razona como en $\neg$I: si toda interpretación que satisface $\Delta \cup \{\neg A\}$ satisface $\bot$, ninguna interpretación que satisfaga $\Delta$ hace falsa a $A$; luego todas ellas hacen verdadera a $A$.

    \item El resto de reglas básicas (R, $\land$E, $\lor$I, $\lor$E, $\leftrightarrow$I, $\leftrightarrow$E) se justifican análogamente comprobando que respetan las condiciones semánticas de las conectivas definidas en el tema~\ref{cap:sem-prop}.
\end{itemize}

Las reglas derivadas (MT, DS, DNE, LEM, DeM) no necesitan un tratamiento aparte, porque cada aplicación de una de ellas puede sustituirse por una derivación que solo usa reglas básicas. Puesto que cada regla conserva la propiedad invariante y al finalizar la prueba todas las hipótesis provisionales han sido descargadas (quedando únicamente las premisas originales $\Gamma$), se concluye que si $\Gamma \vdash A$, entonces $\Gamma \models A$.
\end{proof}

\section{El teorema de completud: la prueba de Kalmár (1935)}

En su versión débil, el teorema de completud proposicional establece que toda tautología es demostrable en el cálculo:
\[
\text{Si } \models A, \text{ entonces } \vdash A.
\]

La primera demostración publicada de la completud de la lógica proposicional es la de Emil Post (1921), a partir de su tesis doctoral de 1920
(Paul Bernays había obtenido el resultado en 1918, en su tesis de
habilitación, que no se publicó entonces; véase Zach, 1999). En 1935, el matemático húngaro László Kalmár
publicó otra demostración, especialmente transparente, que es
constructiva y algorítmica.

Kalmár la formuló para un sistema axiomático; aquí la adaptamos a la
deducción natural. Frente a otras demostraciones más abstractas, como
la de Henkin (tema~\ref{cap:sem-lpo} y apéndice~\ref{ap:henkin}), el método de Kalmár proporciona un procedimiento explícito para traducir la tabla de verdad de una tautología en una derivación.

\subsection{Notación y conceptos previos}

Sea $A$ una fórmula cuyas variables proposicionales distintas son $p_1, p_2, \dots, p_n$.
Dada una asignación de verdad $v$, definimos para cualquier fórmula $B$ la fórmula $B^v$ (la «fórmula evaluada» según $v$) como:
\[
B^v = \begin{cases}
B & \text{si } v(B) = 1 \\
\neg B & \text{si } v(B) = 0
\end{cases}
\]
En particular, para cada letra proposicional $p_i$, el literal $p_i^v$ (la propia letra o su negación) recoge el valor que toma $p_i$ en la fila de la tabla que corresponde a $v$:
\[
p_i^v = \begin{cases}
p_i & \text{si } v(p_i) = 1 \\
\neg p_i & \text{si } v(p_i) = 0
\end{cases}
\]

\begin{example}[Fórmulas evaluadas en una fila concreta]
Si consideramos la asignación $v$ donde $v(p) = 1$ y $v(q) = 0$:
\begin{itemize}
    \item $p^v = p$
    \item $q^v = \neg q$
    \item Para la fórmula $A = p \to q$, como $v(p \to q) = 0$, tenemos $A^v = \neg(p \to q)$.
\end{itemize}
\end{example}

\subsection{El lema fundamental de Kalmár}

El núcleo de la demostración es el siguiente lema, que afirma que cada fila de una tabla de verdad puede reproducirse como una deducción formal:

\begin{lemma}[Lema de Kalmár, 1935]\label{lema_kalmar}
Sea $A$ una fórmula cuyas variables proposicionales están entre $p_1, \dots, p_n$. Para cualquier asignación de verdad $v$:
\[
p_1^v, p_2^v, \dots, p_n^v \vdash A^v
\]
\end{lemma}
\noindent{}(Enunciamos el lema para fórmulas cuyas variables están \emph{entre} $p_1, \dots, p_n$, y no exactamente para las que tienen esas variables, porque en el paso inductivo lo aplicaremos a subfórmulas de $A$, que pueden contener menos variables).


\begin{proof}
La demostración procede por inducción sobre la complejidad sintáctica de la fórmula $A$ (número de conectivas):

\textbf{Caso base ($A$ es una variable atómica $p_i$).}  
La fórmula $A^v$ es idéntica a $p_i^v$. La deducción $p_1^v, \dots, p_n^v \vdash p_i^v$ es inmediata, pues $p_i^v$ forma parte del propio conjunto de premisas. (Como nuestro lenguaje incluye $\bot$, también es un caso base: $v(\bot) = 0$, así que $\bot^v = \neg\bot$, que se deriva sin premisas suponiendo $\bot$ y cerrando la caja con $\neg$I).


\textbf{Paso inductivo.} Asumimos como hipótesis de inducción (HI) que el lema se cumple para las subfórmulas inmediatas de $A$. Examinamos los casos según la conectiva principal:

\begin{enumerate}
    \item \textbf{Caso de la negación ($A = \neg B$):}
    \begin{itemize}
        \item Si $v(B) = 1$, entonces $v(A) = 0$. Por tanto, $B^v = B$ y $A^v = \neg A = \neg\neg B$. Por HI, tenemos $p_1^v, \dots, p_n^v \vdash B$. De $B$ se deduce $\neg\neg B$: suponemos $\neg B$, obtenemos $\bot$ por $\neg$E con $B$ y cerramos la caja con $\neg$I. Luego $p_1^v, \dots, p_n^v \vdash A^v$.
        \item Si $v(B) = 0$, entonces $v(A) = 1$. Por tanto, $B^v = \neg B$ y $A^v = A = \neg B$. Por HI, $p_1^v, \dots, p_n^v \vdash \neg B$, que es exactamente $A^v$.
    \end{itemize}

    \item \textbf{Caso del condicional ($A = B \to C$):}
    \begin{itemize}
        \item \textbf{Subcaso 1 ($v(B) = 0$).} Entonces $v(B \to C) = 1$, con lo que $B^v = \neg B$ y $A^v = B \to C$. Por HI, tenemos una prueba de $\neg B$. A partir de $\neg B$, derivamos $B \to C$ mediante deducción natural: abrimos una caja con hipótesis $B$, alcanzamos $\bot$ por $\neg$E con $\neg B$, aplicamos explosión $\text{X}$ para obtener $C$, y cerramos por $\to$I.
        \item \textbf{Subcaso 2 ($v(C) = 1$).} Entonces $v(B \to C) = 1$, con lo que $C^v = C$ y $A^v = B \to C$. Por HI, tenemos una prueba de $C$. A partir de $C$, abrimos una caja con $B$, reiteramos $C$, y cerramos por $\to$I para obtener $B \to C$.
        \item \textbf{Subcaso 3 ($v(B) = 1$ y $v(C) = 0$).} Entonces $v(B \to C) = 0$, con lo que $B^v = B$, $C^v = \neg C$ y $A^v = \neg(B \to C)$. Por HI tenemos pruebas de $B$ y de $\neg C$. Derivamos $\neg(B \to C)$ abriendo una caja con $B \to C$, aplicando $\to$E con $B$ para obtener $C$, generando $\bot$ con $\neg C$, y cerrando por $\neg$I.
    \end{itemize}

    \item \textbf{Caso de la disyunción ($A = B \lor C$):}
    \begin{itemize}
        \item \textbf{Subcaso 1 ($v(B) = 1$).} Entonces $A^v = B \lor C$. Por HI tenemos una prueba de $B$, y por $\lor$I obtenemos $B \lor C$.
        \item \textbf{Subcaso 2 ($v(B) = 0$ y $v(C) = 1$).} Entonces $A^v = B \lor C$. Por HI tenemos una prueba de $C$, y por $\lor$I obtenemos $B \lor C$.
        \item \textbf{Subcaso 3 ($v(B) = 0$ y $v(C) = 0$).} Entonces $A^v = \neg(B \lor C)$. Por HI tenemos pruebas de $\neg B$ y de $\neg C$. Abrimos una caja con $B \lor C$ y aplicamos $\lor$E: suponiendo $B$ obtenemos $\bot$ con $\neg B$, y suponiendo $C$ obtenemos $\bot$ con $\neg C$. Así llegamos a $\bot$ y cerramos la caja con $\neg$I.
    \end{itemize}

    \item \textbf{Casos de la conjunción ($A = B \land C$) y del bicondicional ($A = B \leftrightarrow C$).}  
    Se demuestran de modo análogo, aplicando las reglas I/E correspondientes a cada combinación de valores de verdad de las partes (ejercicio~\ref{e3-6}).
\end{enumerate}
En todos los casos, $p_1^v, \dots, p_n^v \vdash A^v$.
\end{proof}

\subsection{El colapso de premisas: cómo eliminar las hipótesis}

El lema de Kalmár nos proporciona $2^n$ deducciones (una por cada fila de la tabla de verdad). Si la fórmula $A$ es una tautología ($\models A$), entonces en \emph{todas} las interpretaciones $v$ se cumple $v(A) = 1$, lo que significa que
\[
A^v = A \quad \text{para toda asignación } v.
\]

Por tanto, para cada combinación de valores de $p_1, \dots, p_{n-1}$, tendremos dos deducciones que solo difieren en el signo de la última variable, $p_n$:
\begin{enumerate}
    \item $p_1^v, \dots, p_{n-1}^v, p_n \vdash A$
    \item $p_1^v, \dots, p_{n-1}^v, \neg p_n \vdash A$
\end{enumerate}

El paso siguiente consiste en eliminar la variable $p_n$ aplicando el principio de prueba por casos:

\begin{lemma}[Lema de eliminación de premisas]\label{lema_eliminacion}
Si $\Gamma, p \vdash A$ y $\Gamma, \neg p \vdash A$, entonces $\Gamma \vdash A$.
\end{lemma}

\begin{proof}
Disponemos de una derivación de $A$ bajo la hipótesis $p$ (Derivación~1), y de otra derivación de $A$ bajo la hipótesis $\neg p$ (Derivación~2). La regla derivada del tercero excluido (LEM) permite descargar simultáneamente ambas hipótesis y concluir directamente $A$:
\begin{center}
$\begin{nd}[rules=myrules]
    \hypo [\mbox{}] {} {\text{Premisas } \Gamma} \by{Prem.}{}
    \open
        \hypo [i] {i} {p} \by{Hip.}{}
        \have [\mbox{}] {} {\vdots}
        \have [j] {j} {A} \by{Deriv.~1}{}
    \close
    \open
        \hypo [k] {k} {\lnot p} \by{Hip.}{}
        \have [\mbox{}] {} {\vdots}
        \have [l] {l} {A} \by{Deriv.~2}{}
    \close
    \have [\mbox{}] {} {A} \lem{i-j, k-l}
\end{nd}$
\end{center}
La variable $p$ queda eliminada en un único paso inferencial, conservando únicamente las premisas $\Gamma$.
\end{proof}

\subsection{Demostración final del teorema de completud}

\begin{theorem}[Completud proposicional]
Si $\models A$, entonces $\vdash A$.
\end{theorem}

\begin{proof}
Sea $A$ una tautología con variables $p_1, \dots, p_n$. (Si $A$ no tiene variables, como $\neg\bot$, el lema de Kalmár da directamente $\vdash A$).
\begin{enumerate}
    \item Por el lema de Kalmár (\ref{lema_kalmar}), puesto que $A$ es verdadera en las $2^n$ filas de su tabla, tenemos:
    \[
    p_1^{\pm}, p_2^{\pm}, \dots, p_{n-1}^{\pm}, p_n \vdash A \qquad\text{y}\qquad p_1^{\pm}, p_2^{\pm}, \dots, p_{n-1}^{\pm}, \neg p_n \vdash A
    \]
    donde cada $p_i^{\pm}$ representa indistintamente $p_i$ o $\neg p_i$, para cada una de las $2^{n-1}$ combinaciones posibles de signos de las primeras $n-1$ variables.
    \item Aplicando el lema de eliminación (\ref{lema_eliminacion}) a cada pareja de pruebas, eliminamos la variable $p_n$, obteniendo $2^{n-1}$ pruebas de la forma:
    \[
    p_1^{\pm}, p_2^{\pm}, \dots, p_{n-1}^{\pm} \vdash A.
    \]
    \item Repetimos el mismo procedimiento eliminando sucesivamente $p_{n-1}, p_{n-2}, \dots$ hasta eliminar $p_1$.
    \item Tras $n$ etapas de reducción se han descargado todas las premisas y se obtiene:
    \[
    \vdash A.
    \qedhere
    \]
\end{enumerate}
\end{proof}

\subsection{De la completud débil a la completud fuerte: premisas finitas y compacidad}

La demostración constructiva de Kalmár resuelve el problema de la completud débil: establece que toda tautología es deducible ($\models A \implies \vdash A$). ¿Cómo se traslada este resultado a la completud fuerte ($\Gamma \models A \implies \Gamma \vdash A$)? La respuesta requiere distinguir si el conjunto de premisas $\Gamma$ es finito o infinito.

\subsubsection{Caso 1: conjuntos finitos de premisas (reducción a la completud débil)}

Si el conjunto de premisas es finito, $\Gamma = \{B_1, B_2, \dots, B_k\}$, la completud fuerte se obtiene de modo elemental a partir de la completud débil:

\begin{corollary}[Completud para conjuntos finitos de premisas]\label{corolario_completud_fuerte_finita}
Sean $B_1, \dots, B_k$ y $A$ fórmulas proposicionales. Si $\{B_1, \dots, B_k\} \models A$, entonces $\{B_1, \dots, B_k\} \vdash A$.
\end{corollary}
\begin{proof}
La deducción se construye en tres pasos:
\begin{enumerate}
    \item Por la versión semántica del teorema de la deducción (tema~\ref{cap:sem-prop}), $\{B_1, \dots, B_k\} \models A$ equivale a afirmar que la fórmula condicional que reúne las premisas es una tautología:
    \[
    \models (B_1 \land \dots \land B_k) \to A.
    \]
    \item Por el teorema de completud de Kalmár, al ser dicha fórmula una tautología, existe una prueba formal sin premisas en deducción natural:
    \[
    \vdash (B_1 \land \dots \land B_k) \to A.
    \]
    \item Supongamos ahora disponibles las premisas originales $\{B_1, \dots, B_k\}$. Aplicando de forma sucesiva la regla de introducción de la conjunción ($\land$I), formamos la conjunción $B_1 \land \dots \land B_k$. Aplicando entonces el \emph{modus ponens} ($\to$E) entre esta conjunción y el teorema demostrado en el paso 2, obtenemos la conclusión deseada:
    \[
    \{B_1, \dots, B_k\} \vdash A.
    \qedhere
    \]
\end{enumerate}
\end{proof}

\subsubsection{Caso 2: conjuntos infinitos de premisas\ampliacion}

Si $\Gamma$ es infinito, la completud fuerte también se cumple. Basta saber que, si $\Gamma \models A$, entonces $\Gamma_0 \models A$ para algún subconjunto finito $\Gamma_0 \subseteq \Gamma$ (\textbf{teorema de compacidad}; en el apéndice~\ref{ap:henkin} se enuncia, para primer orden, en la forma equivalente «un conjunto es satisfacible si lo son todos sus subconjuntos finitos»): el corolario~\ref{corolario_completud_fuerte_finita} da entonces $\Gamma_0 \vdash A$ y, como $\Gamma_0 \subseteq \Gamma$, también $\Gamma \vdash A$.

Queda por demostrar la compacidad, que puede hacerse directamente, sin
pasar por el cálculo. Por ejemplo, si el lenguaje tiene una cantidad
numerable de variables, mediante el lema de König. Otra posibilidad es
el método de Henkin del apéndice~\ref{ap:henkin}, que funciona igual
en lógica proposicional y da directamente la completud fuerte. Con ese
método el orden se invierte: primero se obtiene la completud fuerte y
de ella se deduce la compacidad. En efecto, si $\Gamma \models A$,
entonces $\Gamma \vdash A$; la derivación solo usa un subconjunto
finito $\Gamma_0$ de premisas, luego $\Gamma_0 \vdash A$ y, por el
teorema de corrección, $\Gamma_0 \models A$. 

\section{Ejemplos completos del método de Kalmár}

\practica{La aplicación \appkalmar{} sigue estos pasos para cualquier fórmula de hasta tres variables: la tabla de verdad, la derivación de cada fila según el lema y el colapso de las premisas con LEM, paso a paso y con un color para cada fila. Si la fórmula no es tautología ni contradicción, muestra dónde se bloquea el colapso.}

\begin{example}[Una variable: el principio del tercero excluido]
Demostremos mediante el método de Kalmár que $\models p \lor \neg p$ implica $\vdash p \lor \neg p$.
La tabla de verdad tiene $2^1 = 2$ filas:
\begin{center}
\begin{tabular}{c|c}
$p$ & $p \lor \neg p$ \\
\hline
1 & 1 \\
0 & 1
\end{tabular}
\end{center}

\textbf{Paso 1: aplicación del lema de Kalmár a cada fila.}
\begin{itemize}
    \item En la fila 1 ($v(p) = 1$) demostramos $p \vdash p \lor \neg p$:
    \begin{center}
    $\begin{nd}[rules=myrules]
        \hypo [1] {1} {p} \by{Prem.}{}
        \have [2] {2} {p \lor \lnot p} \oi{1}
    \end{nd}$
    \end{center}
    \item En la fila 2 ($v(p) = 0$) demostramos $\neg p \vdash p \lor \neg p$:
    \begin{center}
    $\begin{nd}[rules=myrules]
        \hypo [1] {1} {\lnot p} \by{Prem.}{}
        \have [2] {2} {p \lor \lnot p} \oi{1}
    \end{nd}$
    \end{center}
\end{itemize}

\textbf{Paso 2: colapso de la variable $p$ mediante LEM.}  
Combinamos ambas subpruebas aplicando directamente la regla del tercero excluido (LEM):
\begin{center}
$\begin{nd}[rules=myrules]
    \open
        \hypo [1] {1} {p} \by{Hip.}{}
        \have [2] {2} {p \lor \lnot p} \oi{1}
    \close
    \open
        \hypo [3] {3} {\lnot p} \by{Hip.}{}
        \have [4] {4} {p \lor \lnot p} \oi{3}
    \close
    \have [5] {5} {p \lor \lnot p} \lem{1-2, 3-4}
\end{nd}$
\end{center}
Se obtiene directamente $\vdash p \lor \neg p$ sin ninguna premisa.
\end{example}

\begin{example}[Dos variables: la ley de linealidad]\label{ej:kalmar-linealidad}
Consideremos la fórmula $A = (p \to q) \lor (q \to p)$, que ya apareció en el ejercicio~\ref{e1-29} del tema~\ref{cap:dn-prop}. Su tabla de verdad es:
\begin{center}
\begin{tabular}{cc|c|c|c}
$p$ & $q$ & $p \to q$ & $q \to p$ & $(p \to q) \lor (q \to p)$ \\
\hline
1 & 1 & 1 & 1 & \textbf{1} \\
1 & 0 & 0 & 1 & \textbf{1} \\
0 & 1 & 1 & 0 & \textbf{1} \\
0 & 0 & 1 & 1 & \textbf{1}
\end{tabular}
\end{center}

\textbf{Paso 1: las cuatro deducciones del lema de Kalmár.} Seguimos el caso de la disyunción del lema: si $v(p \to q) = 1$, derivamos $p \to q$ y aplicamos $\lor$I; si no, derivamos $q \to p$. Cada condicional se obtiene con el subcaso correspondiente del caso del condicional:
\begin{enumerate}
    \item Fila 1 ($p, q$): $p \to q$ por el subcaso 2 (suponemos $p$ y reiteramos $q$).
    \item Fila 2 ($p, \neg q$): $p \to q$ es falso, así que derivamos $q \to p$ por el subcaso 1 (suponemos $q$, obtenemos $\bot$ con $\neg q$ y, por X, $p$).
    \item Filas 3 ($\neg p, q$) y 4 ($\neg p, \neg q$): $p \to q$ por el subcaso 1 (suponemos $p$, obtenemos $\bot$ con $\neg p$ y, por X, $q$).
\end{enumerate}

\textbf{Paso 2: colapso de $q$ y de $p$ mediante LEM.} Colocamos cada derivación de fila dentro de las cajas de hipótesis correspondientes y aplicamos LEM dos veces sobre $q$ (dentro de las cajas de $p$ y de $\neg p$) y una vez sobre $p$:
\begin{center}
$\begin{nd}[rules=myrules]
    \open
        \hypo [1] {1} {p} \by{Hip.}{}
        \open
            \hypo [2] {2} {q} \by{Hip.}{}
            \open
                \hypo [3] {3} {p} \by{Hip.}{}
                \have [4] {4} {q} \r{2}
            \close
            \have [5] {5} {p \to q} \ci{3-4}
            \have [6] {6} {(p \to q) \lor (q \to p)} \oi{5}
        \close
        \open
            \hypo [7] {7} {\lnot q} \by{Hip.}{}
            \open
                \hypo [8] {8} {q} \by{Hip.}{}
                \have [9] {9} {\bot} \ne{7, 8}
                \have [10] {10} {p} \x{9}
            \close
            \have [11] {11} {q \to p} \ci{8-10}
            \have [12] {12} {(p \to q) \lor (q \to p)} \oi{11}
        \close
        \have [13] {13} {(p \to q) \lor (q \to p)} \lem{2-6, 7-12}
    \close
    \open
        \hypo [14] {14} {\lnot p} \by{Hip.}{}
        \open
            \hypo [15] {15} {q} \by{Hip.}{}
            \open
                \hypo [16] {16} {p} \by{Hip.}{}
                \have [17] {17} {\bot} \ne{14, 16}
                \have [18] {18} {q} \x{17}
            \close
            \have [19] {19} {p \to q} \ci{16-18}
            \have [20] {20} {(p \to q) \lor (q \to p)} \oi{19}
        \close
        \open
            \hypo [21] {21} {\lnot q} \by{Hip.}{}
            \open
                \hypo [22] {22} {p} \by{Hip.}{}
                \have [23] {23} {\bot} \ne{14, 22}
                \have [24] {24} {q} \x{23}
            \close
            \have [25] {25} {p \to q} \ci{22-24}
            \have [26] {26} {(p \to q) \lor (q \to p)} \oi{25}
        \close
        \have [27] {27} {(p \to q) \lor (q \to p)} \lem{15-20, 21-26}
    \close
    \have [28] {28} {(p \to q) \lor (q \to p)} \lem{1-13, 14-27}
\end{nd}$
\end{center}
\textbf{Comentario.} La estructura de la derivación reproduce la de la tabla: el nivel exterior distingue $p$ y $\neg p$, el siguiente $q$ y $\neg q$, y cada una de las cuatro cajas de $q$ o de $\neg q$ es una fila. Las líneas 3--6, 8--12, 16--20 y 22--26 son las derivaciones del paso 1. El método no busca la prueba más corta (compárala con la que hayas encontrado en el ejercicio~\ref{e1-29}), sino una que existe siempre y se construye mecánicamente.
\end{example}

% \section{Significado conceptual del método de Kalmár}

% El procedimiento constructivo de László Kalmár constituye uno de los resultados más lúcidos de la metalógica clásica. Frente a métodos no constructivos, proporciona un algoritmo transparente que conecta de manera directa la semántica de la verdad con el cálculo de la prueba:
% \begin{enumerate}
%     \item Cada fila de la tabla de verdad codifica un lema del cálculo: una garantía de que, bajo esas premisas atómicas, la fórmula (o su negación) se deduce formalmente en Fitch (lema de Kalmár).
%     \item Si una fórmula es una tautología, el cálculo de deducción
%       natural permite fusionar sistemáticamente las filas de dos en
%       dos mediante la regla del tercero excluido (LEM), descargando
%       sucesivamente las hipótesis hasta alcanzar una prueba sin
%       premisas ($\vdash A$).
% \end{enumerate}
% La tabla de verdad de una tautología no es, por tanto, una matriz estática; contiene en su propia estructura finita toda la información necesaria para construir paso a paso su demostración formal.

\section{Ejercicios para afianzar las ideas de la prueba de completud}

Los siguientes ejercicios están diseñados para asimilar los mecanismos clave de la demostración de Kalmár y comprender la interacción entre sintaxis y semántica. Después de resolverlos, puedes compararlos con lo que construye la aplicación \appkalmar{} para las mismas fórmulas.

\begin{pregunta}[Notación de Kalmár y fórmulas evaluadas: $p_i^v$ y $A^v$]\label{e3-1}
Considera la fórmula $A = (p \lor q) \to \neg p$.
\begin{enumerate}
    \item Para la asignación $v_1$ dada por $v_1(p) = 1$ y $v_1(q) = 0$:
    \begin{itemize}
        \item Determina los literales $p^{v_1}$ y $q^{v_1}$.
        \item Evalúa el valor de verdad $v_1(A)$ según las tablas de verdad.
        \item Escribe la fórmula evaluada $A^{v_1}$.
        \item Escribe la relación de derivabilidad que afirma el lema de Kalmár para esta fila.
    \end{itemize}
    \item Repite el análisis para la asignación $v_2$ dada por $v_2(p) = 0$ y $v_2(q) = 1$.
\end{enumerate}

\paragraph{Solución guiada:}
\begin{enumerate}
    \item Para $v_1$:
    \begin{itemize}
        \item Como $v_1(p) = 1$, $p^{v_1} = p$. Como $v_1(q) = 0$, $q^{v_1} = \neg q$.
        \item Evaluamos la fórmula: $v_1(p \lor q) = 1$ y $v_1(\neg p) = 0$. Por tanto, $v_1((p \lor q) \to \neg p) = 0$.
        \item Como $v_1(A) = 0$, la fórmula evaluada es su negación: $A^{v_1} = \neg((p \lor q) \to \neg p)$.
        \item El lema de Kalmár garantiza que $p, \neg q \vdash \neg((p \lor q) \to \neg p)$.
    \end{itemize}
    \item Para $v_2$:
    \begin{itemize}
        \item Como $v_2(p) = 0$, $p^{v_2} = \neg p$. Como $v_2(q) = 1$, $q^{v_2} = q$.
        \item Evaluamos la fórmula: $v_2(p \lor q) = 1$ y $v_2(\neg p) = 1$. Por tanto, $v_2((p \lor q) \to \neg p) = 1$.
        \item Como $v_2(A) = 1$, la fórmula evaluada es ella misma: $A^{v_2} = (p \lor q) \to \neg p$.
        \item El lema de Kalmár garantiza que $\neg p, q \vdash (p \lor q) \to \neg p$.
    \end{itemize}
\end{enumerate}

\end{pregunta}

\begin{pregunta}[Reproducción en deducción natural de una fila de la tabla]\label{e3-2}
Considera la fórmula condicional $A = p \to q$ y la fila $v$ donde $v(p) = 1$ y $v(q) = 0$. 
Puesto que en esta fila el condicional es falso ($v(A) = 0$), el lema de Kalmár afirma que existe una prueba formal en deducción natural de:
\[
p, \neg q \vdash \neg(p \to q)
\]
Construye la derivación formal completa en el cálculo de Fitch utilizando únicamente reglas de introducción y eliminación.

\paragraph{Solución guiada.}
La estrategia consiste en abrir una hipótesis provisional con $p \to
q$, obtener $q$ mediante el \emph{modus ponens} con la premisa $p$,
usar $\neg q$ para alcanzar una contradicción, y cerrar mediante $\neg$I:
\begin{center}
$\begin{nd}[rules=myrules]
    \hypo [1] {1} {p} \by{Prem.}{}
    \hypo [2] {2} {\lnot q} \by{Prem.}{}
    \open
        \hypo [3] {3} {p \to q} \by{Hip.}{}
        \have [4] {4} {q} \ce{3, 1}
        \have [5] {5} {\bot} \by{$\neg$E}{2, 4}
    \close
    \have [6] {6} {\lnot(p \to q)} \by{$\neg$I}{3-5}
\end{nd}$
\end{center}
Esta derivación reproduce el motivo por el que la combinación $p = 1$ y $q = 0$ hace falso el condicional.

\end{pregunta}

\begin{pregunta}[El colapso inductivo con la regla LEM]\label{e3-3}
Supongamos que para una fórmula tautológica con dos variables $A(p, q)$ ya disponemos de las cuatro derivaciones formales proporcionadas por el lema de Kalmár (una por cada fila de la tabla):
\begin{align*}
\text{(1) } & p, q \vdash A & \text{(2) } & p, \neg q \vdash A \\
\text{(3) } & \neg p, q \vdash A & \text{(4) } & \neg p, \neg q \vdash A
\end{align*}
Describe con precisión cómo la regla derivada del tercero excluido (LEM) permite reducir estas cuatro derivaciones hasta obtener una prueba sin premisas ($\vdash A$).

\paragraph{Solución guiada.}
El procedimiento de colapso se desarrolla en dos niveles:
\begin{enumerate}
    \item \textbf{Nivel 1 (eliminación de la variable $q$):}
    \begin{itemize}
        \item Tomamos las derivaciones (1) y (2). En ambas, la premisa $p$ es común. Abrimos dos subpruebas: una que supone $q$ y concluye $A$, y otra que supone $\neg q$ y concluye $A$. Aplicando la regla LEM sobre las dos subpruebas, descargamos $q$ y obtenemos una nueva derivación: $p \vdash A$.
        \item Tomamos las derivaciones (3) y (4). La premisa $\neg p$ es común. Procediendo de modo idéntico con dos subpruebas bajo $q$ y $\neg q$, la regla LEM descarga $q$ y produce: $\neg p \vdash A$.
    \end{itemize}
    \item \textbf{Nivel 2 (eliminación de la variable $p$).}
    Disponemos ahora de dos deducciones: una que alcanza $A$ bajo la hipótesis $p$, y otra que alcanza $A$ bajo la hipótesis $\neg p$. Abriendo una subprueba bajo $p$ y otra bajo $\neg p$ y aplicando LEM sobre ambas, descargamos la última variable $p$, obteniendo la derivación final sin premisas:
    \[
    \vdash A.
    \]
\end{enumerate}

\end{pregunta}

\begin{pregunta}[¿Por qué falla el colapso si la fórmula no es tautológica?]\label{e3-4}
Considera la fórmula $B = p \lor q$, que no es una tautología (es una contingencia que resulta falsa cuando $p$ y $q$ son ambas falsas).
\begin{enumerate}
    \item Escribe las derivaciones que garantiza el lema de Kalmár para:
    \begin{itemize}
        \item La fila $v_1$: $v_1(p) = 0$ y $v_1(q) = 1$.
        \item La fila $v_2$: $v_2(p) = 0$ y $v_2(q) = 0$.
    \end{itemize}
    \item Explica por qué es imposible aplicar la regla LEM para eliminar la variable $q$ entre estas dos derivaciones y concluir $\neg p \vdash p \lor q$.
    \item ¿Qué lección metalógica nos enseña este bloqueo?
\end{enumerate}

\paragraph{Solución guiada:}
\begin{enumerate}
    \item Evaluando la fórmula en ambas filas:
    \begin{itemize}
        \item En $v_1$, $v_1(p \lor q) = 1$. Por el lema de Kalmár, tenemos: $\neg p, q \vdash p \lor q$.
        \item En $v_2$, $v_2(p \lor q) = 0$. Por el lema de Kalmár, tenemos: $\neg p, \neg q \vdash \neg(p \lor q)$.
    \end{itemize}
    \item La regla LEM exige que las dos subpruebas bajo $q$ y bajo $\neg q$ lleguen \emph{a la misma fórmula}:
    \[
    \text{Si } \Delta, q \vdash C \quad\text{y}\quad \Delta, \neg q \vdash C, \quad\text{entonces } \Delta \vdash C.
    \]
    En nuestro caso, la primera derivación concluye $p \lor q$, mientras que la segunda concluye su negación, $\neg(p \lor q)$. Al ser conclusiones distintas (de hecho, contradictorias entre sí), la regla LEM no es aplicable y el colapso queda bloqueado. Y ninguna otra derivación podría darnos $\neg p \vdash p \lor q$: la fila $v_2$ es un contraejemplo ($v_2(\neg p) = 1$ y $v_2(p \lor q) = 0$), así que $\neg p \not\models p \lor q$ y, por el teorema de corrección, $\neg p \nvdash p \lor q$.
  \item \textbf{Lección metalógica.} El método de Kalmár solo produce pruebas de tautologías. Solo cuando una fórmula $A$ es verdadera en \emph{todas} las filas de la tabla de verdad, las derivaciones que da el lema, una por fila, concluyen todas en $A$. Basta una fila en la que la fórmula sea falsa para que su derivación concluya $\neg A$, y entonces la regla LEM no puede descargar todas las premisas.
\end{enumerate}

\end{pregunta}

\begin{pregunta}[Paso constructivo de la completud débil a la fuerte]\label{e3-5}
Supongamos demostrado que toda tautología es derivable sin premisas ($\models C \implies \vdash C$).
Demuestra formalmente que si $\{B_1, B_2\} \models A$, entonces $\{B_1, B_2\} \vdash A$.
Escribe el esquema general de la demostración formal en deducción natural.

\paragraph{Solución guiada:}
\begin{enumerate}
    \item \textbf{Paso semántico.} Por definición de consecuencia lógica, en toda asignación donde $B_1$ y $B_2$ sean verdaderos, $A$ es verdadero. Por consiguiente, es imposible que $(B_1 \land B_2)$ sea verdadero y $A$ sea falso. Esto significa que la fórmula condicional es una tautología:
    \[
    \models (B_1 \land B_2) \to A.
    \]
    \item \textbf{Paso de completud débil.} Como dicha fórmula es una tautología, el teorema de Kalmár nos asegura la existencia de una demostración $\mathcal{D}$:
    \[
    \vdash (B_1 \land B_2) \to A.
    \]
    \item \textbf{Esquema formal en Fitch.} Tomando ahora $B_1$ y $B_2$ como premisas explícitas, podemos completar la derivación formal en muy pocas líneas:
    \begin{center}
    $\begin{nd}[rules=myrules]
        \hypo [1] {1} {B_1} \by{Prem.}{}
        \hypo [2] {2} {B_2} \by{Prem.}{}
        \have [3] {3} {(B_1 \land B_2) \to A} \by{Teorema $\mathcal{D}$}{}
        \have [4] {4} {B_1 \land B_2} \by{$\land$I}{1, 2}
        \have [5] {5} {A} \ce{3, 4}
    \end{nd}$
    \end{center}
    La línea 3 abrevia la demostración $\mathcal{D}$: como no usa premisas, sus líneas pueden copiarse en esta derivación. Esto demuestra de forma constructiva que la completud débil es suficiente para recuperar la completud fuerte en cualquier argumento con un número finito de premisas.
\end{enumerate}

\end{pregunta}

\begin{pregunta}[Los casos de la conjunción y el bicondicional]\label{e3-6}
Completa la demostración del lema de Kalmár para los casos que se dejaron indicados:
\begin{enumerate}
    \item $A = B \land C$. Distingue los casos según los valores $v(B)$ y $v(C)$ y describe cómo se obtiene $A^v$ a partir de las derivaciones de $B^v$ y $C^v$ que proporciona la hipótesis de inducción.
    \item $A = B \leftrightarrow C$, del mismo modo.
\end{enumerate}

\paragraph{Solución guiada:}
\begin{enumerate}
    \item \textbf{Conjunción.}
    \begin{itemize}
        \item Si $v(B) = v(C) = 1$, entonces $A^v = B \land C$: de $B$ y $C$ se obtiene $B \land C$ por $\land$I.
        \item Si $v(B) = 0$, entonces $A^v = \neg(B \land C)$: con la prueba de $\neg B$, suponemos $B \land C$, obtenemos $B$ por $\land$E, $\bot$ por $\neg$E y cerramos con $\neg$I.
        \item Si $v(B) = 1$ y $v(C) = 0$, se procede igual con $\neg C$ (extrayendo $C$ por $\land$E).
    \end{itemize}
    \item \textbf{Bicondicional.}
    \begin{itemize}
        \item Si $v(B) = v(C) = 1$, entonces $A^v = B \leftrightarrow C$: aplicamos $\leftrightarrow$I con dos cajas; en la primera suponemos $B$ y reiteramos $C$, en la segunda suponemos $C$ y reiteramos $B$.
        \item Si $v(B) = v(C) = 0$, también $A^v = B \leftrightarrow C$: en la caja de $B$ obtenemos $\bot$ con $\neg B$ y, por X, $C$; en la caja de $C$ obtenemos $\bot$ con $\neg C$ y, por X, $B$.
        \item Si $v(B) = 1$ y $v(C) = 0$, entonces $A^v = \neg(B \leftrightarrow C)$: suponemos $B \leftrightarrow C$, obtenemos $C$ por $\leftrightarrow$E con $B$, $\bot$ con $\neg C$, y cerramos con $\neg$I.
        \item Si $v(B) = 0$ y $v(C) = 1$, es simétrico: de $B \leftrightarrow C$ y $C$ obtenemos $B$, y $\bot$ con $\neg B$.
    \end{itemize}
\end{enumerate}
\end{pregunta}

\section*{Referencias bibliográficas}

\begin{enumerate}
    \item \textbf{Church, A.} (1956). \emph{Introduction to Mathematical Logic}, vol. I. Princeton University Press. [Capítulo I, §15, pp. 94--106: lema *151 (p. 98) y teorema *152, «si $B$ es una tautología, $\vdash B$» (p. 99), con el método de Kalmár adaptado al sistema P$_1$. La atribución a Kalmár está en la nota histórica del §29, pp. 163--164].
    \item \textbf{Garrido, M.} (1974). \emph{Lógica simbólica}. Tecnos, Madrid. [Manual de referencia en español para la metalógica elemental].
    \item \textbf{Gentzen, G.} (1934/1935). «Untersuchungen über das logische Schließen». \emph{Mathematische Zeitschrift}, vol. 39, pp. 176--210 y 405--431. [Nacimiento de la deducción natural].
    \item \textbf{Kalmár, L.} (1935). «Über die Axiomatisierbarkeit des Aussagenkalküls». \emph{Acta Scientiarum Mathematicarum (Szeged)}, vol. 7, n.º 4, pp. 222--243. [Artículo fundacional de la prueba constructiva de completud].
    \item \textbf{Kleene, S. C.} (1967). \emph{Mathematical Logic}. John Wiley \& Sons. [Capítulo I: análisis de la completud del cálculo proposicional].
    \item \textbf{Mendelson, E.} (1997). \emph{Introduction to Mathematical Logic} (4.ª ed.). Chapman \& Hall. [§1.4: lema 1.13 (lema de Kalmár, p. 41) y proposición 1.14, teorema de completud (p. 42), cuya demostración Mendelson atribuye a Kalmár (1935)].
    \item \textbf{Post, E. L.} (1921). «Introduction to a general theory of elementary propositions». \emph{American Journal of Mathematics}, vol. 43, n.º 3, pp. 163--185. [Primera demostración publicada de la completud proposicional].
    \item \textbf{Zach, R.} (1999). «Completeness before Post: Bernays, Hilbert, and the development of propositional logic». \emph{The Bulletin of Symbolic Logic}, vol. 5, n.º 3, pp. 331--366. [Sobre la demostración de completud de Bernays de 1918].
\end{enumerate}

%%% Local Variables:
%%% mode: latex
%%% TeX-master: "../ApuntesLogica"
%%% End:


%%%%%%%%%%%%%%%%%%%%%%%%%%%%%%%%%%%%%%%%%%%%%%%%%%%%%%%%%%%%%%%%%%%%%%%%%%%%%%%
\part{Lógica de primer orden}
%%%%%%%%%%%%%%%%%%%%%%%%%%%%%%%%%%%%%%%%%%%%%%%%%%%%%%%%%%%%%%%%%%%%%%%%%%%%%%%

\chapter{Deducción natural en lógica de primer orden}\label{cap:dn-lpo}
\begin{resumentema}
En este tema ampliamos el cálculo de deducción natural para dar cuenta
de la estructura interna de los enunciados mediante la lógica de
predicados de primer orden (LPO). Presentamos las reglas de
introducción y eliminación para los cuantificadores universal ($\forall$)
y existencial ($\exists$). Analizamos las restricciones
sobre las constantes (constantes arbitrarias y testigos), mostrando cómo previenen la deducción de
falacias formales clásicas. Después introducimos la regla derivada
de conversión de cuantificadores (CQ), las reglas de la identidad y
heurísticas para la construcción de pruebas con relaciones binarias y
múltiples cuantificadores. El tema termina con una breve presentación
de los asistentes de demostración y de su uso para verificar pruebas
propuestas por modelos de lenguaje. 
\end{resumentema}

% -*- coding: utf-8 -*-
\section{La insuficiencia de la lógica proposicional}

En los tres primeros temas hemos analizado el razonamiento desde la óptica de la lógica proposicional. Sin embargo, consideremos este conocido argumento:
\begin{quote}
\textit{Todos los seres humanos son mortales.}\\
\textit{Sócrates es un ser humano.}\\
\textit{Por consiguiente, Sócrates es mortal.}
\end{quote}
Cualquier persona reconoce intuitivamente que este argumento es válido: la conclusión se sigue necesariamente de las premisas. Ahora bien, si intentamos traducirlo al lenguaje de la lógica proposicional, cada enunciado simple debe representarse con una variable proposicional distinta: premisas $p$ y $q$, conclusión $r$. Y como
\[
p, \ q \ \not\models \ r
\]
(basta $v(p) = v(q) = 1$ y $v(r) = 0$), desde el punto de vista proposicional el argumento resulta inválido: la lógica proposicional trata los enunciados simples como unidades indivisibles y es incapaz de analizar su estructura interna. La validez del silogismo no depende de cómo se conectan enunciados completos mediante conjunciones o condicionales, sino de la articulación interna entre sujetos (individuos particulares como Sócrates), propiedades (ser humano, mortal) y cantidades (todos, alguno).

Para formalizar y derivar estos argumentos necesitamos enriquecer nuestro lenguaje formal:
\begin{enumerate}
    \item \textbf{Constantes individuales ($a, b, c, s, \dots$).} Nombres propios que designan objetos o individuos específicos de un dominio (por ejemplo, $s$: Sócrates).
    \item \textbf{Variables individuales ($x, y, z, \dots$).} Pronombres indefinidos que actúan como posiciones abiertas («cualquiera», «algo»).
    \item \textbf{Predicados ($P, Q, R, H, M, \dots$).} Expresan propiedades de individuos o relaciones entre varios individuos. Por ejemplo, $H(s)$ simboliza «Sócrates es humano», $M(s)$ simboliza «Sócrates es mortal» y $A(a,b)$ representa «Antonio ama a Belén».
    \item \textbf{Cuantificadores:}
    \begin{itemize}
        \item Cuantificador universal ($\forall x$): «para todo $x$», «cualquiera que sea $x$».
        \item Cuantificador existencial ($\exists x$): «existe al menos un $x$ tal que», «hay algún $x$».
    \end{itemize}
\end{enumerate}

Estos recursos permiten formalizar el silogismo: 
\[
\forall x (H(x) \to M(x)), \ H(s) \ \vdash \ M(s)
\]
El objetivo de este tema es extender nuestro cálculo de Fitch con reglas de inferencia que permitan derivar conclusiones como esta. Las reglas del tema~\ref{cap:dn-prop}, básicas y derivadas (MT, DS, DNE, LEM y DeM), siguen valiendo cuando las fórmulas son de primer orden; lo que añadimos son las reglas de los cuantificadores y de la identidad.

\section{Sintaxis básica, variables libres e instancias}\label{sec:sintaxis-lpo}

Antes de formular las reglas de inferencia para los cuantificadores, debemos precisar la sintaxis formal de la lógica de primer orden y la operación básica que vincula una afirmación general con individuos concretos: la sustitución.

\subsection{Términos, fórmulas y alcance}

El lenguaje de primer orden distingue dos categorías de expresiones:
\begin{enumerate}
    \item \textbf{Términos.} Expresiones que designan o refieren a objetos del universo. En nuestro cálculo básico, los términos son las constantes individuales ($a, b, c, s, \dots$), que nombran individuos fijos, y las variables individuales ($x, y, z, \dots$), que sirven como marcadores de posición o pronombres indefinidos.
    \item \textbf{Fórmulas.} Expresiones que dicen algo de los objetos y tienen un valor de verdad (o pueden adquirirlo al asignar valores a sus variables).
\end{enumerate}

\noindent Las fórmulas se definen inductivamente. Al mismo tiempo, definimos qué apariciones (\emph{ocurrencias}) de las variables en una fórmula son \textbf{libres} y cuáles están \textbf{ligadas}:
\begin{itemize}
    \item \textbf{Fórmulas atómicas.} Si $P$ es una letra de predicado $n$-aria y $t_1, \dots, t_n$ son términos, entonces $P(t_1, \dots, t_n)$ es una fórmula atómica (por ejemplo, $H(s)$, $R(x, y)$ o $P(c)$). Además, si $t_1$ y $t_2$ son términos, $t_1 = t_2$ es una fórmula atómica (identidad; véase la sección~\ref{sec:identidad}). En una fórmula atómica, todas las ocurrencias de variables son libres.
    \item \textbf{Conectivas.} Si $A$ y $B$ son fórmulas, también lo son $\neg A$, $(A \land B)$, $(A \lor B)$, $(A \to B)$ y $(A \leftrightarrow B)$. Cada ocurrencia de una variable es libre o ligada en ellas según lo sea en $A$ o en $B$.
    \item \textbf{Cuantificadores.} Si $A$ es una fórmula y $x$ es una variable individual, también lo son $\forall x A$ y $\exists x A$. La fórmula $A$ es el \textbf{alcance} (\emph{scope}) del cuantificador. Las ocurrencias de $x$ que eran libres en $A$ quedan \textbf{ligadas} por el cuantificador; las demás ocurrencias conservan su carácter. (La variable que acompaña al cuantificador, como la $x$ de $\forall x$, no se cuenta como una ocurrencia).
\end{itemize}
Por ejemplo, en $P(x) \land \forall x Q(x)$ la primera ocurrencia de $x$ es libre: está en la fórmula atómica $P(x)$ y no queda dentro del alcance de ningún cuantificador $\forall x$ o $\exists x$. La de $Q(x)$ está ligada: era libre en $Q(x)$, y $\forall x$ la liga. Una misma variable puede, por tanto, aparecer libre y ligada en distintas partes de una fórmula. En $\forall x (P(x) \to \exists x Q(x))$, el cuantificador $\forall x$ solo liga la $x$ de $P(x)$: la de $Q(x)$ ya no estaba libre en $P(x) \to \exists x Q(x)$, porque la había ligado $\exists x$.

\subsection{Fórmulas abiertas y sentencias}

La distinción entre ocurrencias libres y ligadas es necesaria para aplicar correctamente las reglas de inferencia. Decimos que una variable está \textbf{libre} en una fórmula si tiene en ella al menos una ocurrencia libre.

Una fórmula que contiene al menos una variable libre se denomina fórmula \textbf{abierta} (por ejemplo, $H(x) \to M(x)$ o $R(x, c)$); no expresa una proposición completa mientras no se especifique a qué individuo refiere $x$. Una fórmula sin ninguna variable libre se denomina fórmula \textbf{cerrada}, \textbf{enunciado} o \textbf{sentencia} (por ejemplo, $\forall x (H(x) \to M(x))$ o $H(s)$); solo las sentencias tienen por sí mismas un valor de verdad en un modelo (para evaluar una fórmula abierta hay que asignar valores a sus variables libres, como veremos en el tema~\ref{cap:sem-lpo}). En las derivaciones de este tema todas las líneas son sentencias: las reglas de los cuantificadores operan con constantes, no con variables libres.

\subsection{Sustitución e instancias}

Con estos conceptos podemos definir con precisión la sustitución:

\begin{definition}[Sustitución e instancia]
Sea $A$ una fórmula y $x$ una variable libre en $A$. Si $c$ es una constante individual, denotamos por $A(c)$ (o $A[c/x]$) la fórmula que resulta de reemplazar todas las apariciones libres de la variable $x$ en $A$ por la constante $c$. La fórmula resultante $A(c)$ se denomina una \textbf{instancia} de $A$.
\end{definition}

\begin{ejemplo}[Ejemplos de sustitución]
Consideremos los siguientes casos para ilustrar cómo opera la sustitución únicamente sobre las apariciones libres:
\begin{itemize}
    \item Si $A(x)$ es $H(x) \to M(x)$, entonces $A(s)$ es $H(s) \to M(s)$.
    \item Si $A(x)$ es $P(x) \land \forall y R(x, y)$, entonces $A(c)$ es $P(c) \land \forall y R(c, y)$ (la variable ligada $y$ no cambia).
    \item Si $A(x)$ es $R(x,c) \land \forall x P(x)$, entonces $A(b)$
      es
      $R(b,c) \land \forall x P(x)$ (la $x$ de $P(x)$ no
      se sustituye porque está ligada por el cuantificador).
\end{itemize}
\end{ejemplo}

\practica{En \appsintaxis{} puedes ver, para cualquier fórmula de primer orden, el alcance de cada cuantificador, qué apariciones de las variables son libres o ligadas y el resultado de sustituir variables libres por constantes. También propone ejercicios de variables libres y ligadas y de sustitución.}

\section{Las reglas del cuantificador universal \texorpdfstring{($\forall$)}{(para todo)}}

El cuantificador universal afirma que una propiedad se cumple para
\emph{todos y cada uno} de los individuos del universo. Sus dos reglas
de deducción natural se corresponden con la \textbf{instanciación} (ir de lo general a
lo particular) y la \textbf{generalización} (ir de un caso genérico
arbitrario a lo general). 

\subsection{Eliminación del universal \texorpdfstring{($\forall\text{E}$)}{(Eliminacion del universal)}}

La regla más intuitiva es la eliminación: si algo es verdadero para todos los individuos, entonces tiene que ser verdadero para cualquier individuo particular que elijamos.

\begin{center}
\begin{tabular}{|c|}
\hline
\textbf{Regla de eliminación del universal ($\forall\text{E}$)} \\
\hline
\cajaregla{$\begin{nd}[rules=myrules]
    \have [m] {1} {\forall x A(x)} \by{}{}
    \have [\mbox{}] {2} {A(c)} \Ae{1}
\end{nd}$} \\ 
\hline
\end{tabular}
\end{center}

\noindent \textbf{Condiciones de aplicación:}
\begin{itemize}
    \item La constante $c$ puede ser cualquier constante individual: una que ya aparecía antes en las premisas o una nueva.
    \item Se sustituyen todas las apariciones libres de $x$ en $A(x)$ por la misma constante $c$.
\end{itemize}

\begin{ejemplo}[Resolución formal del argumento socrático]
Veamos cómo la regla $\forall\text{E}$ resuelve en solo dos pasos el silogismo clásico:
\begin{center}
$\begin{nd}[rules=myrules]
    \hypo [1] {1} {\forall x (H(x) \to M(x))} \by{Prem.}{}
    \hypo [2] {2} {H(s)} \by{Prem.}{}
    \have [3] {3} {H(s) \to M(s)} \Ae{1}
    \have [4] {4} {M(s)} \ce{3, 2}
\end{nd}$
\end{center}
En la línea 3 instanciamos el principio universal al caso particular de Sócrates ($s$), y en la línea 4 aplicamos un simple \emph{modus ponens} proposicional.
\end{ejemplo}

\subsection{Introducción del universal \texorpdfstring{($\forall\text{I}$)}{(Introduccion del universal)}}\label{sec:forall-i}

¿Cómo podemos llegar a demostrar que una propiedad vale para \emph{todos} los objetos? En un universo infinito o desconocido, no podemos revisar los individuos uno por uno.

En lógica se utiliza el método del \textbf{caso arbitrario}: tomamos
un individuo genérico $c$ del cual no asumimos ninguna propiedad
particular, y demostramos $A(c)$. Puesto que no hemos
utilizado ninguna peculiaridad de $c$, lo que hemos demostrado para
$c$ vale para cualquiera. 

\begin{center}
\begin{tabular}{|c|}
\hline
\textbf{Regla de introducción del universal ($\forall\text{I}$)} \\
\hline
\cajaregla{$\begin{nd}[rules=myrules]
    \have [m] {1} {A(c)} \by{}{}
    \have [\mbox{}] {2} {\forall x A(x)} \Ai{1}
\end{nd}$} \\
\hline
\end{tabular}
\end{center}

\noindent Las dos restricciones de la constante propia o parámetro arbitrario $c$ son:
\begin{enumerate}[(1)]
    \item $c$ no debe aparecer en ninguna premisa ni en ninguna hipótesis activa (no descargada) en el momento de aplicar la regla;
    \item $c$ no debe aparecer en la conclusión $\forall x A(x)$: hay que sustituir \emph{todas} sus apariciones por $x$.
\end{enumerate}
La restricción (1) es necesaria. Si $c$ apareciera en una premisa
(por ejemplo, si una premisa dijera «Carlos es tenista», $T(c)$), entonces $c$ no sería un individuo arbitrario, sino uno con propiedades particulares asumidas. Si permitiéramos generalizar sobre $c$, deduciríamos la falsedad de que «todo individuo es tenista».

La restricción (2) también es necesaria: de $\forall x R(x, x)$ obtenemos $R(c, c)$ por $\forall$E; si pudiéramos generalizar sustituyendo solo la primera aparición de $c$, obtendríamos $\forall x R(x, c)$, y de ahí $R(d, c)$ para cualquier $d$, lo que no se sigue de la premisa. Por último (y lo mismo valdrá para $\exists$I), la variable $x$ elegida no debe aparecer ya en $A(c)$, para que no quede capturada por otro cuantificador.

\begin{ejemplo}[Demostración de $\forall x (P(x) \land Q(x)) \vdash \forall x P(x) \land \forall x Q(x)$]
Queremos probar que si todo individuo tiene las propiedades $P$ y $Q$, entonces todos tienen $P$ y todos tienen $Q$.
\begin{center}
$\begin{nd}[rules=myrules]
    \hypo [1] {1} {\forall x (P(x) \land Q(x))} \by{Prem.}{}
    \have [2] {2} {P(c) \land Q(c)} \Ae{1}
    \have [3] {3} {P(c)} \by{$\land$E}{2}
    \have [4] {4} {Q(c)} \by{$\land$E}{2}
    \have [5] {5} {\forall x P(x)} \Ai{3}
    \have [6] {6} {\forall x Q(x)} \Ai{4}
    \have [7] {7} {\forall x P(x) \land \forall x Q(x)} \by{$\land$I}{5, 6}
\end{nd}$
\end{center}
\textbf{Comprobación de las restricciones.} La constante $c$ introducida en la línea 2 no aparece en la premisa 1 (donde solo está la variable ligada $x$), ni en las conclusiones de las líneas 5 y 6. Por tanto, $c$ es arbitraria y las aplicaciones de $\forall\text{I}$ en las líneas 5 y 6 son correctas.
\end{ejemplo}

\begin{ejemplo}[Silogismo clásico Barbara]
Formalizamos la deducción aristotélica «todos los mamíferos son vertebrados ($M \to P$); todos los gatos son mamíferos ($S \to M$); por tanto, todos los gatos son vertebrados ($S \to P$)»:
\[
\forall x (M(x) \to P(x)), \ \forall x (S(x) \to M(x)) \ \vdash \ \forall x (S(x) \to P(x))
\]
\begin{center}
$\begin{nd}[rules=myrules]
    \hypo [1] {1} {\forall x (M(x) \to P(x))} \by{Prem.}{}
    \hypo [2] {2} {\forall x (S(x) \to M(x))} \by{Prem.}{}
    \open
        \hypo [3] {3} {S(c)} \by{Hip.}{}
        \have [4] {4} {S(c) \to M(c)} \Ae{2}
        \have [5] {5} {M(c)} \ce{4, 3}
        \have [6] {6} {M(c) \to P(c)} \Ae{1}
        \have [7] {7} {P(c)} \ce{6, 5}
    \close
    \have [8] {8} {S(c) \to P(c)} \ci{3-7}
    \have [9] {9} {\forall x (S(x) \to P(x))} \Ai{8}
\end{nd}$
\end{center}
\textbf{Comentario.} Nótese la articulación de la prueba:
\begin{enumerate}
    \item La meta es un enunciado universal con un condicional ($\forall x (S(x) \to P(x))$). La estrategia consiste en derivar primero el condicional $S(c) \to P(c)$ para una constante arbitraria $c$ y aplicar después $\forall$I.
    \item Para probar dicho condicional, abrimos una subprueba asumiendo el antecedente $S(c)$ (línea 3) y buscamos el consecuente $P(c)$.
    \item Las premisas universales se instancian sobre la constante $c$ que está en juego en la subprueba (líneas 4 y 6).
    \item Al cerrar la caja en la línea 8, la hipótesis $S(c)$ queda descargada. Puesto que $c$ no aparece en ninguna de las premisas de partida (líneas 1 y 2) ni en la conclusión, la regla $\forall$I en la línea 9 es legítima.
\end{enumerate}
\end{ejemplo}

\section{Las reglas del cuantificador existencial \texorpdfstring{($\exists$)}{(existe)}}

El cuantificador existencial afirma que \emph{existe al menos un individuo} que satisface la propiedad en cuestión.

\subsection{Introducción del existencial \texorpdfstring{($\exists\text{I}$)}{(Introduccion del existencial)}}

Si sabemos que un individuo concreto $c$ posee la propiedad $A$, se sigue inmediatamente que existe al menos alguien que la posee.

\begin{center}
\begin{tabular}{|c|}
\hline
\textbf{Regla de introducción del existencial ($\exists\text{I}$)} \\
\hline
\cajaregla{$\begin{nd}[rules=myrules]
    \have [m] {1} {A(c)} \by{}{}
    \have [\mbox{}] {2} {\exists x A(x)} \Ei{1}
\end{nd}$} \\
\hline
\end{tabular}
\end{center}

\noindent \textbf{Condiciones de aplicación:}
\begin{itemize}
    \item No tiene restricciones sobre la constante $c$. (Solo se exige que la variable $x$ elegida no aparezca ya en $A(c)$, para que no quede capturada por otro cuantificador: de $\forall x R(x, c)$ no podemos pasar a $\exists x \forall x R(x, x)$).
    \item No es obligatorio sustituir todas las apariciones de $c$. Por ejemplo, de $R(c, c)$ podemos derivar tanto $\exists x R(x, x)$ («alguien se relaciona consigo mismo») como $\exists x R(x, c)$ («alguien se relaciona con $c$»).
\end{itemize}

\subsection{Eliminación del existencial \texorpdfstring{($\exists\text{E}$)}{(Eliminacion del existencial)}}

La regla de eliminación del existencial es la más delicada del cálculo de primer orden. Captura el razonamiento mediante un \textbf{testigo provisional anónimo}.

Si sabemos que «existe un culpable» ($\exists x C(x)$), no sabemos quién es con exactitud. Por ello, no podemos asumir que el culpable es Juan ($j$) o Pedro ($p$). Lo que hacemos es asignarle un nombre en clave provisional: «llamémosle $c$» ($C(c)$). Si, sin suponer de $c$ nada más que cometió el crimen, logramos deducir una conclusión $B$ que no menciona a $c$, entonces la conclusión $B$ queda establecida sin depender del nombre ficticio utilizado.

\begin{center}
\begin{tabular}{|c|}
\hline
\textbf{Regla de eliminación del existencial ($\exists\text{E}$)} \\
\hline
\cajaregla{$\begin{nd}[rules=myrules]
    \have [m] {1} {\exists x A(x)} \by{}{}
    \open
        \hypo [i] {2} {A(c)} \by{Hip. testigo $c$}{}
        \have [\mbox{}] {3} {\vdots}
        \have [j] {4} {B} \by{}{}
    \close
    \have [\mbox{}] {5} {B} \Ee{1, 2-4}
\end{nd}$} \\
\hline
\end{tabular}
\end{center}

\noindent Las tres restricciones canónicas del testigo $c$ son:
\begin{enumerate}
    \item \textbf{Frescura.} La constante testigo $c$ no debe aparecer en las premisas ni en ninguna hipótesis activa (no descargada) previa al inicio de la subprueba. El propósito de esta exigencia es asegurar que sobre $c$ no recaiga ninguna asunción previa no justificada.
    
    \textit{Matiz formal sobre la frescura.} En muchos textos introductorios se dice informalmente que $c$ debe ser un nombre «completamente nuevo» en la demostración. Esa pauta práctica evita errores, pero la restricción lógica precisa solo prohíbe que $c$ figure en premisas o hipótesis activas:
    \begin{itemize}
        \item Si en líneas previas se aplicó $\forall\text{E}$ sobre $c$ a partir de premisas universales (por ejemplo, derivando $P(c) \to Q(c)$), esa línea depende solo de premisas en las que $c$ no aparece, de modo que vale para cualquier individuo; usar luego $c$ como testigo de un existencial no compromete la corrección. No obstante, para evitar confusiones es buena práctica elegir una constante no utilizada previamente.
        \item Análogamente, si una constante aparecía en una hipótesis de una subprueba que ya se ha cerrado, esa hipótesis ya no está activa, y reutilizar la constante en un bloque posterior no es incorrecto.
    \end{itemize}
    \item \textbf{Independencia del supuesto.} La constante $c$ no debe aparecer en la propia fórmula cuantificada $\exists x A(x)$.
    \item \textbf{No escape.} La constante $c$ no debe aparecer en la conclusión $B$. El nombre provisional del testigo no puede «fugarse» fuera de la caja de hipótesis.
\end{enumerate}

\begin{ejemplo}[Demostración de $\forall x (P(x) \to Q(x)), \ \exists x P(x) \ \vdash \ \exists x Q(x)$]
Veamos cómo interactúan el universal y el existencial:
\begin{center}
$\begin{nd}[rules=myrules]
    \hypo [1] {1} {\forall x (P(x) \to Q(x))} \by{Prem.}{}
    \hypo [2] {2} {\exists x P(x)} \by{Prem.}{}
    \open
        \hypo [3] {3} {P(c)} \by{Hip. testigo $c$}{}
        \have [4] {4} {P(c) \to Q(c)} \Ae{1}
        \have [5] {5} {Q(c)} \ce{4, 3}
        \have [6] {6} {\exists x Q(x)} \Ei{5}
    \close
    \have [7] {7} {\exists x Q(x)} \Ee{2, 3-6}
\end{nd}$
\end{center}
\textbf{Análisis de las restricciones:}
\begin{itemize}
    \item $c$ es una constante fresca: no aparece ni en la premisa 1 ni en la premisa 2.
    \item En la línea 4, instanciamos el universal de la premisa 1 sobre la constante testigo $c$ activa en la subprueba.
    \item En la línea 6, generalizamos existencialmente sobre $c$ obteniendo $\exists x Q(x)$.
    \item La conclusión de la subprueba (línea 6) es $\exists x Q(x)$, en la cual la constante $c$ ya no aparece. Por tanto, no hay escape de testigo y la aplicación de $\exists\text{E}$ en la línea 7 es correcta.
\end{itemize}
\end{ejemplo}

\begin{ejemplo}[Silogismo clásico Darii: $\forall x (M(x) \to P(x)), \ \exists x (S(x) \land M(x)) \ \vdash \ \exists x (S(x) \land P(x))$]
\begin{center}
$\begin{nd}[rules=myrules]
    \hypo [1] {1} {\forall x (M(x) \to P(x))} \by{Prem.}{}
    \hypo [2] {2} {\exists x (S(x) \land M(x))} \by{Prem.}{}
    \open
        \hypo [3] {3} {S(c) \land M(c)} \by{Hip. testigo $c$}{}
        \have [4] {4} {S(c)} \by{$\land$E}{3}
        \have [5] {5} {M(c)} \by{$\land$E}{3}
        \have [6] {6} {M(c) \to P(c)} \Ae{1}
        \have [7] {7} {P(c)} \ce{6, 5}
        \have [8] {8} {S(c) \land P(c)} \by{$\land$I}{4, 7}
        \have [9] {9} {\exists x (S(x) \land P(x))} \Ei{8}
    \close
    \have [10] {10} {\exists x (S(x) \land P(x))} \Ee{2, 3-9}
\end{nd}$
\end{center}
\textbf{Comentario.} Este ejemplo ilustra una secuencia habitual para combinar existenciales y universales:
\begin{enumerate}
    \item \textbf{Prioridad de apertura.} Conviene abrir en primer lugar la subprueba del existencial (línea 3) con una constante testigo fresca $c$, antes de operar con los universales.
    \item \textbf{Instanciación dirigida.} Una vez fijado el testigo provisional $c$, instanciamos la premisa universal sobre esa misma constante (línea 6).
    \item \textbf{Generalización previa al cierre.} Antes de cerrar la caja, generalizamos existencialmente (línea 9). De este modo, la constante de trabajo $c$ no aparece en la conclusión de la caja, y en la línea 10 no hay escape del testigo.
\end{enumerate}
\end{ejemplo}

\practica{En \appdeduccion{} puedes practicar las reglas de los cuantificadores con los ejercicios de este tema. Al comprobar cada línea, la aplicación vigila las restricciones de $\forall$I y $\exists$E y dice cuál se incumple.}

\section{Catálogo de falacias formales por violación de restricciones}

Comprender las restricciones de las reglas de los cuantificadores no consiste en memorizar prohibiciones arbitrarias, sino en advertir qué falacias lógicas se producirían si las ignorásemos.

\subsection{Falacia 1: generalización apresurada \texorpdfstring{(violación de $\forall\text{I}$)}{(violacion de forall I)}}
Supongamos que sabemos que Juan es filósofo ($F(j)$). Si aplicáramos $\forall\text{I}$ sobre $j$, tendríamos la siguiente pseudoderivación:
\begin{center}
$\begin{nd}[rules=myrules]
    \hypo [1] {1} {F(j)} \by{Prem.}{}
    \have [2] {2} {\forall x F(x)} \by{¡Error $\forall$I! ($j$ está en
      la premisa)}{1}
\end{nd}$
\end{center}
\textbf{Error:} de que un individuo concreto posea un atributo no se sigue que todos lo posean. La restricción de que la constante de $\forall\text{I}$ no aparezca en las premisas bloquea esta falacia.

\subsection{Falacia 2: fuga del testigo provisional a la conclusión \texorpdfstring{(violación de $\exists\text{E}$)}{(violacion de exists E)}}
Supongamos que sabemos que «alguien es un ladrón» ($\exists x L(x)$):
\begin{center}
$\begin{nd}[rules=myrules]
    \hypo [1] {1} {\exists x L(x)} \by{Prem.}{}
    \open
        \hypo [2] {2} {L(c)} \by{Hip. testigo $c$}{}
    \close
    \have [3] {3} {L(c)} \by{¡Error $\exists$E! ($c$ aparece en la conclusión)}{1, 2-2}
\end{nd}$
\end{center}
\textbf{Error:} en la línea 3 se habría concluido categóricamente que una constante específica $c$ satisface $L$. Si la constante testigo $c$ pudiera transferirse a la conclusión exterior, a partir de la mera existencia indeterminada de alguien que cumple una propiedad se atribuiría dicha propiedad a un individuo particular. La regla de que el testigo no figure en la conclusión impide esta atribución arbitraria.

\subsection{Falacia 3: fusión de testigos independientes \texorpdfstring{(violación de frescura en $\exists\text{E}$)}{(violacion de frescura en exists E)}}
Consideremos las premisas «alguien es sabio» ($\exists x S(x)$) y «alguien es ignorante» ($\exists x I(x)$):
\begin{center}
$\begin{nd}[rules=myrules]
    \hypo [1] {1} {\exists x S(x)} \by{Prem.}{}
    \hypo [2] {2} {\exists x I(x)} \by{Prem.}{}
    \open
        \hypo [3] {3} {S(c)} \by{Hip. testigo $c$}{}
        \open
            \hypo [4] {4} {I(c)} \by{Hip. testigo $c$}{}
            \have [5] {5} {S(c) \land I(c)} \by{$\land$I}{3, 4}
            \have [6] {6} {\exists x (S(x) \land I(x))} \Ei{5}
        \close
        \have [7] {7} {\exists x (S(x) \land I(x))} \by{¡Error $\exists$E! ($c$ está en la hipótesis 3)}{2, 4-6}
    \close
    \have [8] {8} {\exists x (S(x) \land I(x))} \Ee{1, 3-7}
\end{nd}$
\end{center}
\textbf{Error:} habríamos «demostrado» que existe una persona que es simultáneamente sabia e ignorante. El error está en la elección del testigo de la línea 4, que hace ilegítima la aplicación de $\exists$E de la línea 7: la constante $c$ ya aparecía en la hipótesis 3, todavía activa. Dos existenciales distintos afirman la existencia de individuos que pueden ser diferentes; por ello, cada caja de existencial debe introducir un testigo fresco (por ejemplo, $c$ para la primera y $d$ para la segunda).

\section{Relaciones binarias y la asimetría de los cuantificadores}\label{sec:asimetria}

Cuando formalizamos enunciados que involucran relaciones entre varios individuos (como el amor $A(x, y)$, el conocimiento o la causalidad), el orden de los cuantificadores importa.

\subsection{Una inferencia válida: de \texorpdfstring{$\exists\forall$}{EA} a \texorpdfstring{$\forall\exists$}{AE}}
El siguiente principio es lógicamente válido y demostrable en deducción natural:
\[
\exists x \forall y R(x, y) \ \vdash \ \forall y \exists x R(x, y)
\]
\textit{Lectura intuitiva:} «Si existe alguien que ama a todo el mundo, entonces para cualquier persona existe alguien que la ama».
\begin{center}
$\begin{nd}[rules=myrules]
    \hypo [1] {1} {\exists x \forall y R(x, y)} \by{Prem.}{}
    \open
        \hypo [2] {2} {\forall y R(c, y)} \by{Hip. testigo $c$}{}
        \have [3] {3} {R(c, d)} \Ae{2}
        \have [4] {4} {\exists x R(x, d)} \Ei{3}
        \have [5] {5} {\forall y \exists x R(x, y)} \Ai{4}
    \close
    \have [6] {6} {\forall y \exists x R(x, y)} \Ee{1, 2-5}
\end{nd}$
\end{center}
\textbf{Examen de las constantes:}
\begin{itemize}
    \item En la línea 2, $c$ es un testigo fresco para el existencial exterior de la premisa 1.
    \item En la línea 3, instanciamos el universal $\forall y$ con un parámetro arbitrario $d$.
    \item En la línea 4, generalizamos existencialmente sobre $c$, obteniendo $\exists x R(x, d)$.
    \item En la línea 5, generalizamos universalmente sobre $d$ ($\forall\text{I}$). Esto es correcto porque $d$ no aparece en la hipótesis 2 ni en las premisas.
    \item En la línea 6, cerramos la caja de $\exists\text{E}$. Como $c$ no aparece en la fórmula de la línea 5, la aplicación de $\exists$E es correcta.
\end{itemize}

\begin{ejemplo}[Conmutatividad de cuantificadores homogéneos: $\forall x \forall y R(x, y) \vdash \forall y \forall x R(x, y)$]
A diferencia del cruce entre existenciales y universales, dos cuantificadores del mismo tipo pueden permutar libremente:
\begin{center}
$\begin{nd}[rules=myrules]
    \hypo [1] {1} {\forall x \forall y R(x, y)} \by{Prem.}{}
    \have [2] {2} {\forall y R(a, y)} \Ae{1}
    \have [3] {3} {R(a, b)} \Ae{2}
    \have [4] {4} {\forall x R(x, b)} \Ai{3}
    \have [5] {5} {\forall y \forall x R(x, y)} \Ai{4}
\end{nd}$
\end{center}
\textbf{Comentario:}
\begin{enumerate}
    \item En las líneas 2 y 3 instanciamos sucesivamente con dos constantes arbitrarias distintas $a$ y $b$.
    \item En la línea 4 generalizamos sobre $a$ ($\forall$I): esto es legal porque $a$ no figura en la premisa 1.
    \item En la línea 5 generalizamos sobre $b$ ($\forall$I): igualmente legal porque $b$ tampoco aparece en las premisas.
    \item La independencia de los dos parámetros arbitrarios explica por qué el orden de dos universales consecutivos no altera el significado. El caso de dos existenciales ($\exists\exists$) es análogo (véase el ejercicio~\ref{e4-15}).
\end{enumerate}
\end{ejemplo}

\subsection{La falacia del orden inverso: ¿por qué la conversa no es derivable?}
Consideremos ahora el argumento inverso:
\[
\forall y \exists x R(x, y) \ \nvdash \ \exists x \forall y R(x, y)
\]
\textit{Lectura intuitiva:} «Todo ser humano tiene una madre ($\forall y \exists x M(x, y)$); por tanto, existe una mujer que es madre de todos los seres humanos ($\exists x \forall y M(x, y)$)».

Esta inferencia es una falacia. En nuestro cálculo de deducción natural, si intentáramos derivarla, las restricciones de las reglas bloquearían el intento:
\begin{itemize}
    \item Para eliminar el existencial interior de $\forall y \exists x R(x, y)$, primero tendríamos que instanciar $y$ con una constante $d$, obteniendo $\exists x R(x, d)$.
    \item Al abrir la caja de testigo para este existencial, suponemos $R(c, d)$: el testigo $c$ se ha elegido para \emph{ese} $d$.
    \item Si intentásemos generalizar con $\forall\text{I}$ sobre $d$ para formar $\forall y R(c, y)$, violaríamos la restricción (1) de $\forall$I, porque $d$ aparece en la hipótesis activa $R(c, d)$. Y si cerrásemos antes la caja de $\exists$E, el testigo $c$ no podría salir de ella.
\end{itemize}
Que este intento fracase no demuestra por sí solo que la inferencia sea inválida (podría haber otra derivación). La invalidez se demuestra semánticamente, con un contramodelo, como haremos en el tema~\ref{cap:sem-lpo}; y el teorema de corrección garantiza entonces que no existe ninguna derivación.

\section{Reglas derivadas: conversión de cuantificadores (CQ)}\label{sec:cq}

Al igual que en la lógica proposicional disponíamos de las leyes de De Morgan para transformar conjunciones en disyunciones mediante la negación, en la lógica de primer orden disponemos de la regla de \textbf{conversión de cuantificadores (CQ)}:

\begin{center}
\begin{tabular}{|c|c|}
\hline
\textbf{Premisa} & \textbf{Conclusión por CQ} \\
\hline
$\neg \forall x A(x)$ & $\exists x \neg A(x)$ \\
$\exists x \neg A(x)$ & $\neg \forall x A(x)$ \\
$\neg \exists x A(x)$ & $\forall x \neg A(x)$ \\
$\forall x \neg A(x)$ & $\neg \exists x A(x)$ \\
\hline
\end{tabular}
\end{center}

\begin{theorem}[Derivabilidad de CQ en el cálculo básico]
La regla CQ es derivable: no añade nada que no pueda probarse con las reglas básicas de Fitch.
\end{theorem}

\begin{proof}
A modo de ilustración, demostramos que de $\neg \forall x P(x)$ se deduce $\exists x \neg P(x)$ usando solo reglas básicas:
\begin{center}
$\begin{nd}[rules=myrules]
    \hypo [1] {1} {\lnot \forall x P(x)} \by{Prem.}{}
    \open
        \hypo [2] {2} {\lnot \exists x \lnot P(x)} \by{Hip. IP}{}
        \open
            \hypo [3] {3} {\lnot P(c)} \by{Hip. IP}{}
            \have [4] {4} {\exists x \lnot P(x)} \Ei{3}
            \have [5] {5} {\bot} \by{$\neg$E}{2, 4}
        \close
        \have [6] {6} {P(c)} \ip{3-5}
        \have [7] {7} {\forall x P(x)} \Ai{6}
        \have [8] {8} {\bot} \by{$\neg$E}{1, 7}
    \close
    \have [9] {9} {\exists x \lnot P(x)} \ip{2-8}
\end{nd}$
\end{center}
En la línea 7 la constante $c$ es elegible para $\forall\text{I}$ porque la hipótesis local 3 fue descargada en la línea 6, y ni la premisa 1 ni la hipótesis 2 mencionan a $c$. Las conversiones entre $\neg \exists x P(x)$ y $\forall x \neg P(x)$ son los ejercicios~\ref{e4-16} y~\ref{e4-17}; la que falta, de $\exists x \neg P(x)$ a $\neg \forall x P(x)$, se obtiene suponiendo $\forall x P(x)$ y abriendo una caja de $\exists$E con $\neg P(c)$. Al comprobar que la regla CQ es derivable, queda legitimado su uso como atajo en derivaciones complejas.
\end{proof}

\section{La identidad}\label{sec:identidad}

Muchos razonamientos dependen de que dos nombres designen el mismo individuo: «Cicerón es Tulio; Cicerón era orador; luego Tulio era orador». Para analizarlos añadimos al lenguaje el símbolo lógico $=$: si $t_1$ y $t_2$ son términos, $t_1 = t_2$ es una fórmula atómica. Escribimos $t_1 \neq t_2$ como abreviatura de $\neg(t_1 = t_2)$. Su significado es fijo: $t_1 = t_2$ es verdadera exactamente cuando $t_1$ y $t_2$ designan el mismo objeto (tema~\ref{cap:sem-lpo}).

\begin{center}
\begin{tabular}{cc}
\textbf{Introducción de la identidad ($=$I)} & \textbf{Eliminación de la identidad ($=$E)} \\[0.2cm]
$\begin{nd}[rules=myrules]
    \have [m] {1} {c = c} \idi{}
\end{nd}$ &
$\begin{nd}[rules=myrules]
    \have [m] {1} {a = b}
    \have [n] {2} {A(\dots a \dots a \dots)}
    \have [\mbox{}] {3} {A(\dots b \dots a \dots)} \ide{1, 2}
\end{nd}$
\end{tabular}
\end{center}

\begin{itemize}
    \item \textbf{$=$I} permite escribir $c = c$ en cualquier momento, para cualquier constante $c$, sin citar ninguna línea: todo objeto es idéntico a sí mismo.
    \item \textbf{$=$E} permite, a partir de $a = b$ y de una fórmula en la que aparece $a$, sustituir \emph{una, varias o todas} las apariciones de $a$ por $b$. También se permite la sustitución en sentido inverso: de $a = b$ y una fórmula con $b$, sustituir apariciones de $b$ por $a$.
\end{itemize}
La regla $=$E recoge el principio de \textbf{indiscernibilidad de los idénticos} (a menudo llamado \emph{ley de Leibniz}): si $a$ y $b$ son el mismo objeto, todo lo que vale de uno vale del otro. El principio recíproco, la \emph{identidad de los indiscernibles} (si $a$ y $b$ tienen exactamente las mismas propiedades, son el mismo objeto), no es una ley de la lógica de primer orden.

\begin{ejemplo}[Simetría de la identidad: $a = b \vdash b = a$]\mbox{}
\begin{center}
$\begin{nd}[rules=myrules]
    \hypo [1] {1} {a = b} \by{Prem.}{}
    \have [2] {2} {a = a} \idi{}
    \have [3] {3} {b = a} \ide{1, 2}
\end{nd}$
\end{center}
\textbf{Comentario.} En la línea 3 sustituimos la \emph{primera} aparición de $a$ en $a = a$ por $b$, como permite $=$E.
\end{ejemplo}

\begin{ejemplo}[Lo que distingue a dos individuos: $P(a), \neg P(b) \vdash a \neq b$]\mbox{}
\begin{center}
$\begin{nd}[rules=myrules]
    \hypo [1] {1} {P(a)} \by{Prem.}{}
    \hypo [2] {2} {\lnot P(b)} \by{Prem.}{}
    \open
        \hypo [3] {3} {a = b} \by{Hip.}{}
        \have [4] {4} {P(b)} \ide{3, 1}
        \have [5] {5} {\bot} \ne{2, 4}
    \close
    \have [6] {6} {\lnot(a = b)} \ni{3-5}
\end{nd}$
\end{center}
\textbf{Comentario.} Es la forma habitual de probar que dos individuos son distintos: encontrar algo que vale de uno y no del otro (indiscernibilidad de los idénticos, en forma contrapositiva).
\end{ejemplo}

\section{Guía estratégica para construir pruebas en LPO}

Para abordar con éxito las deducciones con cuantificadores, conviene seguir las siguientes heurísticas:

\begin{enumerate}
    \item \textbf{Estrategia con premisas existenciales (apertura temprana de subderivaciones).}  
    Si tienes una premisa o línea disponible de la forma $\exists x A(x)$, es recomendable abrir la subprueba de $\exists\text{E}$ con una constante testigo $c$ lo antes posible. Aunque, como hemos precisado, un paso previo de $\forall\text{E}$ sobre $c$ no vulnera la frescura mientras $c$ no aparezca en premisas ni en hipótesis activas, abrir la caja existencial de inmediato ofrece una doble ventaja heurística: asegura la elección de un testigo limpio y revela sobre qué constante exacta convendrá instanciar a continuación las premisas universales mediante $\forall\text{E}$.
    
    \item \textbf{Estrategia con premisas universales (instanciación dirigida al objetivo).}  
    Las fórmulas de la forma $\forall x A(x)$ son recursos universales que puedes instanciar tantas veces como necesites. No las instancies a ciegas: espera a ver qué constantes individuales han aparecido en las premisas, cuál has elegido como constante arbitraria para una meta universal o cuáles han introducido las cajas de los existenciales, y aplica $\forall\text{E}$ sobre esas constantes.
    
    \item \textbf{Estrategia con metas universales.}  
    Si tu objetivo final (o intermedio) es demostrar $\forall x B(x)$, planifica derivar primero $B(c)$ para una constante $c$ arbitraria, que no aparezca en las premisas ni en las hipótesis activas, para rematar en el último paso aplicando $\forall\text{I}$.
    
    \item \textbf{Estrategia con metas existenciales.}  
    Si tu objetivo es demostrar $\exists x B(x)$, busca construir $B(c)$ para alguna constante concreta disponible y aplica $\exists\text{I}$.
    
    \item \textbf{Gestión de negaciones y cuantificadores.}  
    Si una fórmula comienza con la negación de un cuantificador ($\neg \forall x$ o $\neg \exists x$), aplica la regla derivada CQ para llevar la negación al interior, o bien recurre a una reducción al absurdo clásica (IP) asumiendo la negación de la meta.
\end{enumerate}

\section{Para saber más: de la deducción natural a los asistentes de demostración}\label{sec:asistentes}

Como vimos al final del tema~\ref{cap:dn-prop}, \emph{comprobar} una derivación es una tarea mecánica: basta verificar que cada línea se ajusta a una regla y que se respetan las restricciones (accesibilidad, constantes frescas, etc.). \emph{Encontrar} una derivación, en cambio, exige ingenio, y en la lógica de primer orden no existe ningún algoritmo que decida en todos los casos si una fórmula es demostrable. En esta asimetría entre encontrar y comprobar se basan los asistentes de demostración.

\paragraph{Asistentes de demostración.} Programas como \emph{Lean}, \emph{Isabelle} o \emph{Coq} (hoy llamado \emph{Rocq}) permiten escribir demostraciones en un lenguaje formal. Un núcleo relativamente pequeño del programa acepta una demostración solo si cada paso se ajusta a las reglas de su sistema lógico, que contiene análogos de las reglas que hemos estudiado. En Lean, por ejemplo, el silogismo hipotético del tema~\ref{cap:dn-prop} ($p \to q, q \to r \vdash p \to r$) se escribe así:
\begin{verbatim}
example (p q r : Prop) (h1 : p -> q) (h2 : q -> r) : p -> r := by
  intro hp               -- suponemos p (introducción del condicional)
  have hq : q := h1 hp   -- modus ponens con la premisa h1
  exact h2 hq            -- modus ponens con la premisa h2: obtenemos r
\end{verbatim}
La instrucción \texttt{intro} abre la caja de $\to$I, y aplicar una hipótesis a otra (\texttt{h1 hp}) es $\to$E. Con herramientas de este tipo se han verificado por completo resultados matemáticos de gran complejidad, como el teorema de los cuatro colores (G. Gonthier, en Coq, 2005) o la conjetura de Kepler sobre el empaquetamiento de esferas (proyecto \emph{Flyspeck}, dirigido por T. Hales, 2014).

Puedes comprobar este ejemplo, sin instalar nada, en el editor web de Lean.\footnote{\url{https://live.lean-lang.org/}} Basta copiarlo y ver si Lean lo acepta. Hay que respetar la sangría: las líneas que siguen a \texttt{by} deben empezar todas con el mismo número de espacios. Un buen ejercicio es pedir a un modelo de lenguaje que escriba en Lean demostraciones de otros ejercicios del tema~\ref{cap:dn-prop} y comprobar en el editor si las acepta; cuando no lo hace, Lean indica la línea en la que falla.

\paragraph{Inteligencia artificial y verificación.} Los modelos de
lenguaje generan textos que \emph{parecen} razonamientos correctos,
pero pueden contener pasos falsos expresados con total
seguridad. Combinarlos con un asistente de demostración separa dos
tareas: el modelo (o la persona) \emph{propone} la demostración, y el
verificador la \emph{comprueba} paso a paso, rechazando cualquier
inferencia que no se ajuste a las reglas. Así, por ejemplo, el sistema
AlphaProof (Google DeepMind, 2024) resolvió problemas de la Olimpiada
Internacional de Matemáticas produciendo demostraciones comprobadas
por Lean. En septiembre de 2026, OpenAI anunció que un sistema de
agentes basado en un modelo de lenguaje había obtenido una
demostración, acompañada de una formalización en Lean, de que las
ecuaciones de Navier-Stokes tridimensionales para fluidos
incompresibles, con una fuerza externa, pueden desarrollar
singularidades en un tiempo finito.\footnote{Anuncio:
  \url{https://openai.com/es-ES/index/navier-stokes-solution/}.} Según
el enunciado oficial, eso bastaría para resolver uno de los siete
problemas del Milenio del Instituto Clay. Por analogía con el
\emph{vibe coding} (programar describiendo informalmente lo que se
quiere a una IA), se habla a veces de \emph{vibe proving}: esbozar la
idea de una prueba, dejar que una IA proponga los detalles y confiar
el veredicto final al verificador formal. 

\paragraph{Preguntas filosóficas.} Este esquema traslada la confianza: no es tan importante quién propone la prueba, sea un humano o una máquina, como el verificador y las reglas que aplica. Recuerda el ideal de Leibniz de resolver las disputas calculando (\emph{calculemus}), y plantea diversas preguntas filosóficas: qué cuenta como demostración, qué papel tiene la comprensión frente a la mera verificación, o qué reglas aceptamos como legítimas (¿lógica clásica o intuicionista?, por ejemplo). Casos como el de Navier-Stokes añaden otras: quién es el autor de una demostración obtenida así, si basta la comprobación de un programa para aceptar un resultado que ninguna persona ha leído, o qué garantiza que el enunciado formalizado dice lo mismo que el problema original. 

\section*{Referencias bibliográficas}

\begin{enumerate}
    \item \textbf{Garrido, M.} (1974). \emph{Lógica simbólica}. Tecnos, Madrid. [Capítulo sobre el cálculo de predicados de primer orden].
    \item \textbf{Gentzen, G.} (1934/1935). «Untersuchungen über das logische Schließen». \emph{Mathematische Zeitschrift}, vol. 39, pp. 176--210 y 405--431. [Introducción original de las reglas para cuantificadores].
    \item \textbf{Magnus, P. D., Button, T., et al.} (2020). \emph{forall x: Calgary. An Introduction to Formal Logic}. Open Logic Project. [Partes IV (deducción natural proposicional) y VII (deducción natural de primer orden)].
    \item \textbf{Prawitz, D.} (1965). \emph{Natural Deduction: A Proof-Theoretical Study}. Almqvist \& Wiksell, Estocolmo. [Normalización de las derivaciones y restricciones sobre las variables].
\end{enumerate}

%%% Local Variables:
%%% mode: latex
%%% TeX-master: "../ApuntesLogica"
%%% End:


\section{Ejercicios de deducción natural en LPO}
% -*- coding: utf-8 -*-
\begin{quote}
\small \textbf{Instrucciones.} Salvo indicación en contra, puedes usar las reglas derivadas del tema~\ref{cap:dn-prop} (MT, DS, DNE, LEM, DeM) y CQ. Puedes hacerlas y comprobarlas en la aplicación \appdeduccion{}, que contiene estos ejercicios; su ayuda indica qué letras son constantes y cuáles variables.
\end{quote}

\section*{Nivel 1: cuantificador universal y silogística clásica}
\textit{En este nivel practicamos la eliminación del universal ($\forall\text{E}$, instanciando en las constantes relevantes) y la introducción del universal ($\forall\text{I}$ sobre parámetros arbitrarios que no aparecen en premisas).}

\begin{pregunta}\label{e4-1}
Demuestra el modo silogístico Celarent (ningún $M$ es $P$, y todo $S$ es $M$; por tanto, ningún $S$ es $P$):
\[
\forall x (M(x) \to \neg P(x)), \ \forall x (S(x) \to M(x)) \ \vdash \ \forall x (S(x) \to \neg P(x))
\]
\textit{Pista.} Abre una caja con hipótesis $S(c)$ para un parámetro arbitrario $c$. Instancia las dos premisas sobre $c$ mediante $\forall\text{E}$ y encadena los condicionales con $\to$E para derivar $\neg P(c)$. Cierra la subprueba con $\to$I y concluye con $\forall\text{I}$.
\end{pregunta}

\begin{pregunta}\label{e4-2}
Demuestra el modo silogístico Cesare (ningún $P$ es $M$, y todo $S$ es $M$; por tanto, ningún $S$ es $P$):
\[
\forall x (P(x) \to \neg M(x)), \ \forall x (S(x) \to M(x)) \ \vdash \ \forall x (S(x) \to \neg P(x))
\]
\textit{Pista.} Abre una caja con $S(c)$ y una subcaja con $P(c)$ para buscar contradicción con $M(c)$ y $\neg M(c)$.
\end{pregunta}

\begin{pregunta}\label{e4-3}
Demuestra la distributividad del cuantificador universal respecto a la conjunción:
\[
\forall x P(x) \land \forall x Q(x) \ \vdash \ \forall x (P(x) \land Q(x))
\]
\textit{Pista.} Extrae $\forall x P(x)$ y $\forall x Q(x)$ por $\land$E, instáncialos a la misma constante arbitraria $c$, junta por $\land$I y aplica $\forall\text{I}$.
\end{pregunta}

\begin{pregunta}\label{e4-4}
Demuestra el siguiente argumento, en el que el consecuente del condicional no contiene la variable $x$:
\[
\forall x (P(x) \to Q(a)), \ P(b) \ \vdash \ Q(a)
\]
\textit{Pista.} Como el universal afirma la implicación para cualquier individuo, basta instanciarlo mediante $\forall\text{E}$ sobre la constante $b$ de la segunda premisa y aplicar $\to$E.
\end{pregunta}

\begin{pregunta}\label{e4-5}
Demuestra la siguiente combinación de $\forall$E y $\leftrightarrow$E:
\[
\forall x (P(x) \leftrightarrow Q(x)), \ P(c) \ \vdash \ Q(c)
\]
\textit{Pista.} Aplica $\forall\text{E}$ sobre la constante $c$ y luego la eliminación del bicondicional ($\leftrightarrow$E) junto con la premisa $P(c)$.
\end{pregunta}

\section*{Nivel 2: cuantificador existencial y gestión de testigos}
\textit{En este nivel introducimos la caja de hipótesis con testigo provisional ($\exists\text{E}$). Recuerda: conviene abrir pronto las cajas de los existenciales con una constante fresca y asegurarse de que el testigo no aparezca en la conclusión.}

\begin{pregunta}\label{e4-6}
Demuestra el modo silogístico Ferio (ningún $M$ es $P$, y algún $S$ es $M$; por tanto, algún $S$ no es $P$):
\[
\forall x (M(x) \to \neg P(x)), \ \exists x (S(x) \land M(x)) \ \vdash \ \exists x (S(x) \land \neg P(x))
\]
\textit{Pista.} Abre una caja de $\exists\text{E}$ con un testigo fresco $c$ asumiendo $S(c) \land M(c)$. Descompón la conjunción con $\land\text{E}$, instancia la premisa universal sobre $c$ para obtener $\neg P(c)$, forma $S(c) \land \neg P(c)$ por $\land\text{I}$, generaliza con $\exists\text{I}$ y cierra la caja de $\exists\text{E}$.
\end{pregunta}

\begin{pregunta}\label{e4-7}
Demuestra el modo silogístico Datisi (todo $M$ es $P$, y algún $M$ es $S$; por tanto, algún $S$ es $P$):
\[
\forall x (M(x) \to P(x)), \ \exists x (M(x) \land S(x)) \ \vdash \ \exists x (S(x) \land P(x))
\]
\textit{Pista.} Abre una caja de $\exists\text{E}$ con testigo fresco $c$ ($M(c) \land S(c)$), extrae las partes, usa $\forall\text{E}$ sobre la primera premisa y $\to$E para obtener $P(c)$ y recombina como $S(c) \land P(c)$ para aplicar $\exists\text{I}$.
\end{pregunta}

\begin{pregunta}\label{e4-8}
Demuestra la distributividad del existencial respecto a la disyunción:
\[
\exists x (P(x) \lor Q(x)) \ \vdash \ \exists x P(x) \lor \exists x Q(x)
\]
\textit{Pista.} Abre la caja de testigo para el existencial ($P(c) \lor Q(c)$). Dentro de la caja, haz una prueba por casos sobre la disyunción ($\lor$E) para deducir $\exists x P(x) \lor \exists x Q(x)$.
\end{pregunta}

\begin{pregunta}\label{e4-9}
Demuestra que un universal y un existencial se combinan existencialmente:
\[
\forall x P(x), \ \exists x Q(x) \ \vdash \ \exists x (P(x) \land Q(x))
\]
\textit{Pista.} Abre la caja de $\exists\text{E}$ para la premisa existencial con testigo $c$ asumiendo $Q(c)$. Dentro de ella, aplica $\forall\text{E}$ sobre $c$ a la primera premisa, construye la conjunción por $\land\text{I}$, aplica $\exists\text{I}$ y cierra la caja.
\end{pregunta}

\begin{pregunta}\label{e4-10}
Demuestra la independencia de dos existenciales con cajas anidadas:
\[
\exists x P(x), \ \exists y Q(y) \ \vdash \ \exists x \exists y (P(x) \land Q(y))
\]
\textit{Advertencia.} Cada caja de testigo debe usar una constante distinta (por ejemplo, $c$ para la primera y $d$ para la segunda).
\end{pregunta}

\section*{Nivel 3: relaciones binarias y cuantificadores múltiples}
\textit{En este nivel trabajamos con predicados diádicos (relaciones entre individuos $R(x, y)$). Presta especial atención al orden de eliminación e introducción de los cuantificadores.}

\begin{pregunta}\label{e4-11}
Demuestra la transmisión existencial de una relación universalmente implicada:
\[
\forall x \forall y (R(x, y) \to S(x, y)), \ \exists x \forall y R(x, y) \ \vdash \ \forall y \exists x S(x, y)
\]
\textit{Pista.} Abre una subprueba de $\exists\text{E}$ para la segunda premisa con testigo fresco $c$ asumiendo $\forall y R(c, y)$. Dentro de la subprueba, toma un parámetro arbitrario $d$ para alcanzar la meta universal; instancia las premisas sobre $c$ y $d$ para derivar $S(c, d)$ por $\to$E, aplica $\exists\text{I}$ sobre $c$ para obtener $\exists x S(x, d)$ y luego $\forall\text{I}$ sobre $d$ para obtener $\forall y \exists x S(x, y)$. Por último, cierra la caja de $\exists\text{E}$.
\end{pregunta}

\begin{pregunta}\label{e4-12}
Demuestra que si una relación $R$ es simétrica, transitiva y serial (todo individuo se relaciona con al menos uno), entonces $R$ es reflexiva:
\[
\forall x \forall y (R(x, y) \to R(y, x)), \ \forall x \forall y \forall z ((R(x, y) \land R(y, z)) \to R(x, z)), \ \forall x \exists y R(x, y) \ \vdash \ \forall x R(x, x)
\]
\textit{Pista.} Para probar la meta universal $\forall x R(x, x)$, fija un parámetro arbitrario $a$. Instancia la premisa de serialidad sobre $a$ para obtener $\exists y R(a, y)$ y abre una subprueba de $\exists\text{E}$ con un testigo fresco $b$ asumiendo $R(a, b)$. Aplica la simetría sobre $a$ y $b$ para obtener $R(b, a)$, forma la conjunción $R(a, b) \land R(b, a)$ y aplica la transitividad sobre $a, b, a$ para derivar $R(a, a)$. Cierra la subprueba de $\exists\text{E}$ (el testigo $b$ no escapa) y concluye con $\forall\text{I}$ sobre $a$.
\end{pregunta}

\begin{pregunta}\label{e4-13}
Demuestra que si una relación es asimétrica, entonces es irreflexiva (nadie se relaciona consigo mismo):
\[
\forall x \forall y (R(x, y) \to \neg R(y, x)) \ \vdash \ \forall x \neg R(x, x)
\]
\textit{Pista.} Abre una caja con la hipótesis $R(c, c)$ y busca una contradicción.
\end{pregunta}

\begin{pregunta}\label{e4-14}
Demuestra el paso de existencial condicional a antecedente universal:
\[
\exists x (P(x) \to Q(a)) \ \vdash \ \forall x P(x) \to Q(a)
\]
\textit{Pista.} Abre una caja para la introducción del condicional ($\to$I) asumiendo $\forall x P(x)$. Dentro de ella, abre la caja de $\exists\text{E}$ para la premisa con un testigo fresco $c$ asumiendo $P(c) \to Q(a)$. Instancia el supuesto $\forall x P(x)$ sobre $c$ para obtener $P(c)$, deduce $Q(a)$ por $\to$E, cierra la caja de $\exists\text{E}$ (ya que $Q(a)$ no contiene a $c$) y finalmente cierra la caja de $\to$I.
\end{pregunta}

\begin{pregunta}\label{e4-15}
Demuestra la conmutatividad del doble cuantificador existencial mediante cajas anidadas:
\[
\exists x \exists y R(x, y) \ \vdash \ \exists y \exists x R(x, y)
\]
\textit{Pista.} Abre una caja de $\exists\text{E}$ con testigo $c$ para $\exists x$; dentro de ella, abre otra caja de $\exists\text{E}$ con testigo $d$ para $\exists y$. A partir de $R(c, d)$, aplica dos veces $\exists\text{I}$ para obtener $\exists y \exists x R(x, y)$ y cierra ambas cajas en orden inverso.
\end{pregunta}

\section*{Nivel 4: cuantificadores con negación, CQ y reducción al absurdo}
\textit{En este nivel combinamos cuantificadores con negación utilizando la regla de conversión de cuantificadores (CQ) y demostraciones clásicas por reducción al absurdo (IP).}

\begin{pregunta}\label{e4-16}
Demuestra mediante reglas básicas (sin usar la regla derivada CQ) que:
\[
\neg \exists x P(x) \ \vdash \ \forall x \neg P(x)
\]
\textit{Pista.} Para una constante arbitraria $c$, abre una hipótesis $P(c)$ con el fin de derivar $\neg P(c)$ por $\neg\text{I}$. Dentro de ella, aplica $\exists\text{I}$ sobre $c$ para obtener $\exists x P(x)$, lo que genera contradicción inmediata ($\bot$) con la premisa $\neg \exists x P(x)$. Cierra la subprueba con $\neg\text{I}$ y concluye con $\forall\text{I}$ sobre $c$.
\end{pregunta}

\begin{pregunta}\label{e4-17}
Demuestra el paso recíproco, también con reglas básicas:
\[
\forall x \neg P(x) \ \vdash \ \neg \exists x P(x)
\]
\textit{Pista.} Procede por introducción de la negación ($\neg\text{I}$) abriendo la hipótesis $\exists x P(x)$. Dentro de esa caja, abre una subprueba de $\exists\text{E}$ con testigo fresco $c$ asumiendo $P(c)$. Instancia la premisa universal sobre $c$ mediante $\forall\text{E}$ para obtener $\neg P(c)$ y obtén $\bot$; cierra la caja de $\exists$E (pues $\bot$ no contiene a $c$) y después la de $\neg$I.
\end{pregunta}

\begin{pregunta}\label{e4-18}
Demuestra el siguiente argumento (basta usar \emph{modus tollens} dentro de una caja de $\exists$E; no hace falta CQ):
\[
\forall x (P(x) \to Q(x)), \ \exists x \neg Q(x) \ \vdash \ \exists x \neg P(x)
\]
\end{pregunta}

\begin{pregunta}\label{e4-19}
Demuestra, mediante introducción de la negación ($\neg$I):
\[
\neg \exists x (P(x) \land Q(x)), \ P(a) \ \vdash \ \neg Q(a)
\]
\end{pregunta}

\begin{pregunta}\label{e4-20}
Demuestra el principio clásico conocido como la paradoja del bebedor (en todo bar hay alguien tal que, si él bebe, todos beben):
\[
\vdash \ \exists x (B(x) \to \forall y B(y))
\]
\textit{Pista.} Aplica la regla LEM con las hipótesis $\forall y B(y)$ y $\neg \forall y B(y)$. En la primera caja, para cualquier constante $c$, supón $B(c)$, reitera $\forall y B(y)$, cierra con $\to$I y aplica $\exists$I. En la segunda, obtén $\exists y \neg B(y)$ por CQ y abre una caja de $\exists$E con testigo $d$ y $\neg B(d)$; supón $B(d)$, obtén $\bot$ y, por X, $\forall y B(y)$; cierra con $\to$I para obtener $B(d) \to \forall y B(y)$ y aplica $\exists$I antes de cerrar la caja de $\exists$E. (Obsérvese que la validez de este principio depende de que el dominio no sea vacío: en un bar vacío no hay nadie que cumpla la condición).
\end{pregunta}

\section*{Nivel 5: análisis de falacias formales}
\textit{A continuación se presentan cuatro demostraciones aparentemente correctas, pero que contienen un error de razonamiento por violar alguna de las restricciones de $\forall$I o de $\exists$E. Para cada una de ellas:}
\begin{enumerate}
    \item Identifica la línea en la que se aplica una regla sin respetar sus restricciones.
    \item Enuncia con precisión qué restricción ha sido violada.
    \item Proporciona un contraejemplo informal en lenguaje natural que demuestre que el argumento es inválido.
\end{enumerate}

\begin{pregunta}[Falacia de generalización indebida]\label{e4-21}

\begin{center}
$\begin{nd}[rules=myrules]
    \hypo [1] {1} {P(a)} \by{Prem.}{}
    \hypo [2] {2} {\forall x (P(x) \to Q(x))} \by{Prem.}{}
    \have [3] {3} {P(a) \to Q(a)} \Ae{2}
    \have [4] {4} {Q(a)} \ce{3, 1}
    \have [5] {5} {\forall x Q(x)} \Ai{4}
\end{nd}$
\end{center}
\end{pregunta}

\begin{pregunta}[Falacia de fuga de testigo]\label{e4-22}

\begin{center}
$\begin{nd}[rules=myrules]
    \hypo [1] {1} {\exists x P(x)} \by{Prem.}{}
    \hypo [2] {2} {\forall x (P(x) \to Q(x))} \by{Prem.}{}
    \open
        \hypo [3] {3} {P(c)} \by{Hip. testigo $c$}{}
        \have [4] {4} {P(c) \to Q(c)} \Ae{2}
        \have [5] {5} {Q(c)} \ce{4, 3}
    \close
    \have [6] {6} {Q(c)} \Ee{1, 3-5}
\end{nd}$
\end{center}
\end{pregunta}

\begin{pregunta}[Falacia de elección de testigo no fresco: confusión con un individuo conocido]\label{e4-23}

\begin{center}
$\begin{nd}[rules=myrules]
    \hypo [1] {1} {P(a)} \by{Prem.}{}
    \hypo [2] {2} {\exists x Q(x)} \by{Prem.}{}
    \open
        \hypo [3] {3} {Q(a)} \by{Hip. testigo $a$}{}
        \have [4] {4} {P(a) \land Q(a)} \by{$\land$I}{1, 3}
        \have [5] {5} {\exists x (P(x) \land Q(x))} \Ei{4}
    \close
    \have [6] {6} {\exists x (P(x) \land Q(x))} \Ee{2, 3-5}
\end{nd}$
\end{center}
\end{pregunta}

\begin{pregunta}[Falacia de la inversión ilícita de cuantificadores]\label{e4-24}

\begin{center}
$\begin{nd}[rules=myrules]
    \hypo [1] {1} {\forall y \exists x R(x, y)} \by{Prem.}{}
    \have [2] {2} {\exists x R(x, d)} \Ae{1}
    \open
        \hypo [3] {3} {R(c, d)} \by{Hip. testigo $c$}{}
        \have [4] {4} {\forall y R(c, y)} \Ai{3}
        \have [5] {5} {\exists x \forall y R(x, y)} \Ei{4}
    \close
    \have [6] {6} {\exists x \forall y R(x, y)} \Ee{2, 3-5}
\end{nd}$
\end{center}
\end{pregunta}

\section*{Nivel 6: identidad}
\textit{En este nivel usamos las reglas $=$I y $=$E (sección~\ref{sec:identidad}). Recuerda que $=$E permite sustituir una, varias o todas las apariciones de un término por otro idéntico, en cualquiera de los dos sentidos.}

\begin{pregunta}[Transitividad y simetría]\label{e4-25}
$a = b, \ b = c \ \vdash \ c = a$

\textit{Pista.} Obtén primero $a = c$ con $=$E y después aplica la simetría como en el ejemplo de la simetría de la identidad (sección~\ref{sec:identidad}).
\end{pregunta}

\begin{pregunta}\label{e4-26}
$\forall x (x = a \to P(x)) \ \vdash \ P(a)$
\end{pregunta}

\begin{pregunta}[Unicidad]\label{e4-27}
Si hay exactamente un individuo con la propiedad $P$, y tanto $a$ como $b$ la tienen, entonces $a$ y $b$ son el mismo individuo:
\[
\exists x (P(x) \land \forall y (P(y) \to y = x)), \ P(a), \ P(b) \ \vdash \ a = b
\]
\textit{Pista.} Abre una caja de $\exists$E con un testigo fresco $c$, obtén $a = c$ y $b = c$, y combínalas con $=$E.
\end{pregunta}

\begin{pregunta}[Simetría como teorema]\label{e4-28}
$\vdash \forall x \forall y (x = y \to y = x)$
\end{pregunta}

%%% Local Variables:
%%% mode: latex
%%% TeX-master: "../ApuntesLogica"
%%% End:


\chapter{Semántica de la lógica de primer orden}\label{cap:sem-lpo}
\begin{resumentema}
  En este tema completamos la lógica de primer orden abordando su
  perspectiva semántica. Definimos formalmente las nociones de
  estructura, dominio, asignación y las relaciones tarskianas de
  satisfacción y verdad en un modelo, y vemos cómo formalizar en primer orden
  enunciados del lenguaje natural. En el plano metateórico, mostramos que la
  completud equivale a que todo conjunto consistente sea satisfacible,
  motivamos la necesidad del enfoque de Leon Henkin (1949) para
  demostrarla mediante la construcción del modelo canónico y contrastamos
  su naturaleza con el método constructivo de Kalmár.
\end{resumentema}

% -*- coding: utf-8 -*-
\section{De las tablas de verdad a las estructuras semánticas}

En la lógica proposicional, la semántica (tema~\ref{cap:sem-prop}) resultaba
muy simple: saber si una fórmula es válida requería únicamente
inspeccionar las combinaciones de ceros y unos en una tabla finita de
$2^n$ filas (si bien esto supone un gran coste en términos de complejidad computacional). En la lógica de primer orden, sin embargo, nos enfrentamos a cuantificadores universales («para todo $x$») y existenciales («existe un $x$») que afirman propiedades sobre dominios que pueden ser arbitrariamente grandes e incluso infinitos (como los números naturales $\mathbb{N}$ o los puntos del continuo $\mathbb{R}$). 

Por consiguiente, ya no podemos apoyarnos en tablas de verdad. Necesitamos una teoría formal que defina qué significa que un enunciado de primer orden sea verdadero en una situación, contexto o universo determinado. Esa teoría es la \textbf{semántica de modelos}, cuyas bases estableció Alfred Tarski con su definición de verdad (1933) y que daría lugar a la teoría de modelos.

\subsection{Componentes de una estructura de primer orden}

Las nociones de conjuntos que usamos a partir de aquí (subconjunto, par ordenado, $n$-tupla, producto cartesiano $D^n$, función) se repasan en el apéndice~\ref{ap:conjuntos}. Para interpretar un lenguaje de primer orden necesitamos una \textbf{estructura} (también llamada \textbf{modelo}):

\begin{definition}[Estructura de primer orden]
Una \textbf{estructura} $\mathcal{M}$ para un lenguaje de primer orden $\mathcal{L}$ es un par ordenado $\mathcal{M} = \langle D, I \rangle$ compuesto por:
\begin{enumerate}
    \item \textbf{El dominio o universo de discurso ($D$).} Un conjunto \emph{no vacío} de objetos o individuos sobre los cuales versan las variables y cuantificadores.
    \item \textbf{La función de interpretación ($I$).} Una función que asigna significado a los símbolos no lógicos del lenguaje:
    \begin{itemize}
        \item A cada constante individual $c$, $I$ le asigna un objeto particular del dominio: $I(c) \in D$.
        \item A cada símbolo de predicado $n$-ario $R$, $I$
          le asigna un conjunto de $n$-tuplas de elementos del
          dominio: $I(R) \subseteq D^n$.
       \begin{itemize}
        \item En el caso particular de un predicado
            unario $P$, $I$ le asigna el conjunto de individuos
          que poseen dicha propiedad: $I(P) \subseteq D$.
        \item Si se incluye el símbolo de identidad ($=$), su
          interpretación en toda estructura es la relación de identidad: $I(=) = \{(d, d) \mid d \in D\}$. 
       \end{itemize}
    \end{itemize}
\end{enumerate}
\end{definition}

\subsection{Asignaciones de variables y denotación de términos}

Mientras que las constantes individuales poseen una referencia fija
determinada por la función de interpretación $I$, las variables
individuales ($x, y, z, \dots$) no tienen un significado fijo: actúan
como pronombres o huecos abiertos. Para dotar de valor a las variables
acudimos al concepto de \textbf{asignación}:

\begin{definition}[Asignación de variables]
Dada una estructura $\mathcal{M} = \langle D, I \rangle$, una \textbf{asignación de variables} en $\mathcal{M}$ es una función $g: \text{Var} \to D$, donde $\text{Var}$ es el conjunto de las variables individuales, que asocia a cada variable un elemento del dominio $D$.
\end{definition}

\noindent A partir de la interpretación $I$ y de la asignación $g$, podemos definir la \textbf{denotación} de cualquier término individual $t$ en $\mathcal{M}$ bajo $g$, que escribimos $[\![t]\!]^\mathcal{M}_g$:
\[
[\![t]\!]^\mathcal{M}_g = \begin{cases}
I(c) & \text{si } t \text{ es una constante individual } c, \\
g(x) & \text{si } t \text{ es una variable individual } x.
\end{cases}
\]

\begin{definition}[Variante de una asignación]
Si $g$ es una asignación, $x$ una variable y $d \in D$ un elemento del dominio, denotamos por $g[x \mapsto d]$ la asignación que coincide exactamente con $g$ en todas las variables excepto, a lo sumo, en $x$, a la cual asigna el objeto $d$:
\[
(g[x \mapsto d])(y) = \begin{cases}
d & \text{si } y = x, \\
g(y) & \text{si } y \neq x.
\end{cases}
\]
\end{definition}

\subsection{Definición inductiva de satisfacción (Tarski)}

Con estas nociones, la relación básica de la semántica de primer orden no es directamente la verdad, sino la relación de \textbf{satisfacción}: una estructura $\mathcal{M}$ satisface una fórmula $A$ bajo una asignación $g$, lo cual se escribe $\mathcal{M}, g \models A$.

\begin{definition}[Satisfacción de fórmulas al estilo Tarski]
Dada una estructura $\mathcal{M} = \langle D, I \rangle$ y una asignación $g$, la relación $\mathcal{M}, g \models A$ se define por inducción sobre la estructura sintáctica de $A$:
\begin{enumerate}
    \item \textbf{Fórmulas atómicas:}
    \begin{itemize}
        \item $\mathcal{M}, g \models P(t_1, \dots, t_n) \iff ([\![t_1]\!]^\mathcal{M}_g, \dots, [\![t_n]\!]^\mathcal{M}_g) \in I(P)$.
        \item $\mathcal{M}, g \models t_1 = t_2 \iff [\![t_1]\!]^\mathcal{M}_g = [\![t_2]\!]^\mathcal{M}_g$.
    \end{itemize}
    \item \textbf{Conectivas proposicionales:}
    \begin{itemize}
        \item $\mathcal{M}, g \not\models \bot$ ($\bot$ es falsa en toda estructura, como en el tema~\ref{cap:sem-prop}).
        \item $\mathcal{M}, g \models \neg A \iff \mathcal{M}, g \not\models A$.
        \item $\mathcal{M}, g \models A \land B \iff \mathcal{M}, g \models A \text{ y } \mathcal{M}, g \models B$.
        \item $\mathcal{M}, g \models A \lor B \iff \mathcal{M}, g \models A \text{ o } \mathcal{M}, g \models B$.
        \item $\mathcal{M}, g \models A \to B \iff (\text{si } \mathcal{M}, g \models A \text{ entonces } \mathcal{M}, g \models B)$.
        \item $\mathcal{M}, g \models A \leftrightarrow B \iff (\mathcal{M}, g \models A \iff \mathcal{M}, g \models B)$.
    \end{itemize}
    \item \textbf{Cuantificador universal:}
    \[
    \mathcal{M}, g \models \forall x A \iff \text{para todo elemento } d \in D, \ \mathcal{M}, g[x \mapsto d] \models A.
    \]
    \item \textbf{Cuantificador existencial:}
    \[
    \mathcal{M}, g \models \exists x A \iff \text{existe al menos un elemento } d \in D \text{ tal que } \mathcal{M}, g[x \mapsto d] \models A.
    \]
\end{enumerate}
\end{definition}

\subsection{De la satisfacción a la verdad en una estructura, consecuencia lógica y validez}

Fijadas la estructura $\mathcal{M}$ y la fórmula $A$, que se cumpla o no $\mathcal{M}, g \models A$ solo depende de los valores que la asignación $g$ da a las variables libres de $A$: si dos asignaciones $g$ y $g'$ coinciden en las variables libres de $A$, entonces $\mathcal{M}, g \models A$ si y solo si $\mathcal{M}, g' \models A$. Este resultado se conoce como \emph{lema de coincidencia}. En particular:
\begin{quote}
Si $A$ es una \textbf{sentencia} (fórmula cerrada, sin variables libres), entonces para cualesquiera dos asignaciones $g$ y $g'$ se cumple que $\mathcal{M}, g \models A \iff \mathcal{M}, g' \models A$.
\end{quote}
Esto nos permite definir la noción de
\textbf{verdad en una estructura o modelo}:

\begin{definition}[Verdad en un modelo]
Sea $\mathcal{M} = \langle D, I \rangle$ una estructura y $A$ una fórmula de primer orden:
\begin{enumerate}
    \item Si $A$ es una \textbf{sentencia}, decimos que $A$ es \textbf{verdadera en $\mathcal{M}$} (y escribimos simplemente $\mathcal{M} \models A$) si $\mathcal{M}, g \models A$ para alguna (o equivalentemente, para toda) asignación $g$. En caso contrario, decimos que $A$ es \textbf{falsa en $\mathcal{M}$} ($\mathcal{M} \not\models A$).
    \item Si $A$ es una \textbf{fórmula abierta} (con variables libres), decimos que $A$ es verdadera en $\mathcal{M}$ ($\mathcal{M} \models A$) si $\mathcal{M}, g \models A$ para \emph{toda} asignación $g$ (lo cual equivale a la verdad de su clausura universal $\forall x_1 \dots \forall x_n A$).
\end{enumerate}
\end{definition}

\noindent A partir de aquí se definen de manera unificada las nociones semánticas principales. Salvo indicación en contra, en lo que sigue $\Gamma$ es un conjunto de sentencias y $A$ una sentencia (como en las derivaciones del tema~\ref{cap:dn-lpo}):
\begin{itemize}
    \item \textbf{Modelo de un conjunto de fórmulas.} $\mathcal{M} \models \Gamma$ significa que $\mathcal{M} \models B$ para toda fórmula $B \in \Gamma$.
    \item \textbf{Consecuencia lógica semántica ($\Gamma \models A$).} En \emph{toda} estructura $\mathcal{M}$, si $\mathcal{M} \models \Gamma$, entonces $\mathcal{M} \models A$. Decimos que el argumento es semánticamente válido.
    \item \textbf{Fórmula universalmente válida o verdad lógica ($\models A$).} La fórmula $A$ es verdadera en \emph{toda} estructura $\mathcal{M}$.
    \item \textbf{Satisfacibilidad.} Un conjunto $\Gamma$ es \textbf{satisfacible} si existe al menos una estructura $\mathcal{M}$ tal que $\mathcal{M} \models \Gamma$.
    \item \textbf{Equivalencia lógica ($A \equiv B$).} $A$ y $B$ son verdaderas en las mismas estructuras. Como en el tema~\ref{cap:sem-prop}, esto equivale a $\models A \leftrightarrow B$.
\end{itemize}


\subsection{La identidad}

El símbolo $=$ es un símbolo \emph{lógico}: en toda estructura se interpreta como la relación de identidad (véanse la definición de estructura y la cláusula de satisfacción de $t_1 = t_2$). Escribimos $t_1 \neq t_2$ como abreviatura de $\neg(t_1 = t_2)$. La identidad permite expresar cantidades y unicidad:
\begin{itemize}
    \item «Hay al menos dos $P$»: $\exists x \exists y (P(x) \land P(y) \land x \neq y)$.
    \item «Hay como mucho un $P$»: $\forall x \forall y ((P(x) \land P(y)) \to x = y)$.
    \item «Hay exactamente un $P$»: $\exists x (P(x) \land \forall y (P(y) \to y = x))$.
\end{itemize}
La sección~\ref{sec:cuantificadores-numericos} generaliza estas construcciones a cualquier número. Las reglas de deducción natural para la identidad ($=$I y $=$E) se presentaron en el tema~\ref{cap:dn-lpo} (sección~\ref{sec:identidad}). Aquí nos interesa sobre todo su uso en la formalización y en la semántica (por ejemplo, en los ejercicios sobre el tamaño del dominio).

\subsection{Ejemplos de evaluación semántica: flexibilidad notacional}

Veamos cómo se aplican estas definiciones en la práctica a través de ejemplos:

\begin{ejemplo}[Una estructura aritmética simple]
Consideremos un lenguaje con constantes individuales $a, b$ y predicados $P$ (unario) y $R$ (binario). Definimos la estructura aritmética $\mathcal{M}_1 = \langle D, I \rangle$:
\begin{itemize}
    \item $D = \mathbb{N} = \{0, 1, 2, 3, \dots\}$ (el conjunto de los números naturales).
    \item $I(a) = 0$, \quad $I(b) = 2$.
    \item $I(P) = \{n \in \mathbb{N} \mid n \text{ es par}\} = \{n  \mid \exists m\in \mathbb{N} (n = 2\times m)\} = \{0, 2, 4, 6, \dots\}$.
    \item $I(R) = \{(m, n) \in \mathbb{N} \times \mathbb{N} \mid m < n\}$ (orden menor estricto).
\end{itemize}
Evaluemos el valor de verdad de varias sentencias en $\mathcal{M}_1$:
\begin{itemize}
    \item La sentencia atómica $P(a)$ es verdadera ($\mathcal{M}_1 \models P(a)$) porque $I(a) = 0 \in I(P)$ (el cero es par).
    \item La sentencia atómica $R(b, a)$ es falsa ($\mathcal{M}_1 \not\models R(b, a)$) porque $(I(b), I(a)) = (2, 0) \notin I(R)$ (es falso que $2 < 0$).
    \item La sentencia $\exists x R(a, x)$ es verdadera porque existe al menos un natural $d \in \mathbb{N}$ (por ejemplo, $d = 1$) tal que $\mathcal{M}_1, g[x \mapsto 1] \models R(a, x)$, ya que $(0, 1) \in I(R)$.
\end{itemize}
\end{ejemplo}

\begin{ejemplo}[Evaluación en una estructura finita y convenciones de notación]
Consideremos una estructura $\mathcal{M}_2$ con un dominio finito $D = \{1, 2, 3\}$, propiedades $P, Q$ y relación binaria $R$:
\begin{itemize}
    \item $I(P) = \{2, 3\}$.
    \item $I(Q) = \{1, 3\}$.
    \item $I(R) = \{(1, 1), (1, 2), (1, 3), (2, 2), (2, 3), (3, 3)\}$.
\end{itemize}
Evaluemos el valor de verdad de tres enunciados cuantificados:
\begin{enumerate}
    \item \textbf{Fórmula existencial: $\exists x (P(x) \land Q(x))$ («Existe algún elemento que es a la vez $P$ y $Q$»).}  
    Formalmente, debemos comprobar si existe algún $d \in D$ tal que $\mathcal{M}_2, g[x \mapsto d] \models P(x) \land Q(x)$:
    \begin{itemize}
        \item Para $d = 1$: $1 \notin I(P)$ (falso), luego la conjunción no se satisface.
        \item Para $d = 2$: $2 \in I(P)$ (verdadero), pero $2 \notin I(Q)$ (falso).
        \item Para $d = 3$: tanto $3 \in I(P)$ como $3 \in I(Q)$ son verdaderos.
    \end{itemize}
    Al existir un elemento en el dominio ($d = 3$) que satisface la conjunción, concluimos formalmente que $\mathcal{M}_2 \models \exists x (P(x) \land Q(x))$.
    
    \item \textbf{Fórmula universal: $\forall x (P(x) \to Q(x))$ («Todo elemento que es $P$ es también $Q$»).}  
    Para que sea verdadera, la condición debe cumplirse para todos los elementos de $D$:
    \begin{itemize}
        \item Para $d = 2$: el antecedente $2 \in I(P)$ es verdadero, pero el consecuente $2 \in I(Q)$ es falso. Por la semántica del condicional, la implicación no se satisface para $d = 2$.
    \end{itemize}
    Al fallar para al menos un individuo ($d=2$), la propiedad no es universal. Por tanto, $\mathcal{M}_2 \not\models \forall x (P(x) \to Q(x))$.
    
    \item \textbf{Fórmula relacional: $\forall x \exists y R(x, y)$ («Todo elemento se relaciona con alguno, quizá consigo mismo»).}  
    Para cada elemento $d_1 \in \{1, 2, 3\}$, comprobamos si existe algún $d_2 \in D$ tal que $(d_1, d_2) \in I(R)$. Vemos que para $d_1 = 1$ sirve $d_2 = 1$; para $d_1 = 2$ sirve $d_2 = 2$; y para $d_1 = 3$ sirve $d_2 = 3$ (por la reflexividad de $R$). Consecuentemente, $\mathcal{M}_2 \models \forall x \exists y R(x, y)$.
\end{enumerate}

\paragraph{Nota metodológica: rigor formal frente a flexibilidad notacional en la práctica.}
Conviene detenerse en una distinción. Desde un punto de vista formal, los números $1, 2, 3$ son objetos semánticos del dominio $D$, no símbolos del lenguaje formal $\mathcal{L}$. Por ello:
\begin{itemize}
    \item Expresiones como «evaluar para $x = 1$» o «la fórmula $A(1)$» constituyen, en rigor, un \textbf{abuso de notación}: la forma exacta es decir que $\mathcal{M}, g[x \mapsto 1] \models A(x)$ (o, para un predicado atómico $P$, que $1 \in I(P)$).
    \item Del mismo modo, escribiremos cosas como «$\text{V} \to \text{F} = \text{F}$» para indicar que un condicional con antecedente verdadero y consecuente falso es falso. En los temas de primer orden usaremos \textbf{V} y \textbf{F} para los valores de verdad (en lugar del $1$ y el $0$ del tema~\ref{cap:sem-prop}), para no confundirlos con los elementos de dominios como $\{1, 2, 3\}$.
\end{itemize}
Estos abusos de notación no introducen ningún error y simplifican mucho la resolución de ejercicios con dominios finitos. Por eso, una vez comprendida su base formal, los usaremos en los ejemplos y ejercicios.
\end{ejemplo}

\begin{ejemplo}[Cómo se demuestra semánticamente una consecuencia]
Para mostrar que un argumento no es válido basta una estructura que haga verdaderas las premisas y falsa la conclusión. Para mostrar que es válido hay que razonar sobre una estructura \emph{cualquiera}. Veamos que $\forall x (P(x) \to Q(x)), \ \exists x P(x) \models \exists x Q(x)$. Sea $\mathcal{M} = \langle D, I \rangle$ una estructura cualquiera en la que las dos premisas son verdaderas. Por la segunda, hay algún $d \in D$ con $d \in I(P)$. Por la primera, el condicional se cumple para todo elemento del dominio, y en particular para $d$: como $d \in I(P)$, también $d \in I(Q)$. Luego $\mathcal{M} \models \exists x Q(x)$. Como de $\mathcal{M}$ solo hemos supuesto que hace verdaderas las premisas, la conclusión es verdadera en toda estructura que las hace verdaderas.
\end{ejemplo}

\practica{\appmundo{} es un tablero con piezas de distintas formas, tamaños y colores que funciona como una estructura: muestra su dominio y su interpretación, y evalúa cualquier fórmula con un árbol que sigue la definición de satisfacción, con las asignaciones de cada paso.}

\section{Formalización en lógica de primer orden}\label{sec:formalizacion-lpo}

En la sección~\ref{sec:formalizacion} formalizábamos enunciados
completos con letras proposicionales. En primer orden hay que analizar
también su estructura interna: de qué individuos se habla
(constantes), qué se dice de ellos (predicados) y cuántos cumplen lo
que se dice (cuantificadores). El diccionario es ahora más rico:
indica el dominio, el significado de cada constante y de cada
predicado y, en los predicados de varios argumentos, el orden de
estos (no es lo mismo «$T(x, y)$: $x$ entrena a $y$» que «$T(x, y)$: $y$ entrena a $x$»).

\subsection{Dos esquemas básicos}

Casi todo lo que formalizaremos se construye a partir de dos esquemas, que corresponden a las formas A e I de la silogística (sección~\ref{sec:importacion}):
\begin{center}
\renewcommand{\arraystretch}{1.3}
\begin{tabular}{|l|l|}
\hline
«Todo $S$ es $P$» & $\forall x (S(x) \to P(x))$ \\
«Algún $S$ es $P$» & $\exists x (S(x) \land P(x))$ \\
\hline
\end{tabular}
\end{center}
El universal va con el condicional y el existencial con la conjunción. «Todo $S$ es $P$» no dice que todo sea $S$: dice de cada individuo que, \emph{si} es $S$, entonces es $P$. «Algún $S$ es $P$» dice que hay al menos un individuo que es $S$ \emph{y} es $P$. Intercambiar las conectivas lleva a dos errores frecuentes:
\begin{itemize}
    \item $\forall x (S(x) \land P(x))$ dice que todo individuo es $S$
      y es $P$. Con $F(x)$: «$x$ es filósofo» y $B(x)$: «$x$ es
      sabio», $\forall x (F(x) \land B(x))$ no dice que todos los
      filósofos sean sabios, sino que todo el mundo es filósofo y
      sabio. Afirma mucho más de lo que habitualmente queremos
      expresar. 
    \item $\exists x (S(x) \to P(x))$ dice demasiado poco: es verdadera en cuanto exista algo que no sea $S$. Con $U(x)$: «$x$ es un unicornio» y $W(x)$: «$x$ es blanco», $\exists x (U(x) \to W(x))$ es verdadera en el mundo real, simplemente porque hay cosas que no son unicornios. «Algún unicornio es blanco», en cambio, es falsa.
\end{itemize}

\subsection{Negaciones y otras variantes}

Con estos dos esquemas y la negación se formalizan otras expresiones de cantidad:
\begin{center}
\renewcommand{\arraystretch}{1.3}
\begin{tabular}{|p{7.2cm}|l|l|}
\hline
\textbf{Expresión} & \textbf{Formalización} & \textbf{Equivalente} \\
\hline
Todo $S$ es $P$; todos los $S$ son $P$; cualquier $S$ es $P$ & $\forall x (S(x) \to P(x))$ & \\
Algún $S$ es $P$; hay $S$ que son $P$ & $\exists x (S(x) \land P(x))$ & \\
Ningún $S$ es $P$; no hay $S$ que sean $P$ & $\forall x (S(x) \to \neg P(x))$ & $\neg \exists x (S(x) \land P(x))$ \\
Algún $S$ no es $P$; no todo $S$ es $P$; no todos los $S$ son $P$ & $\exists x (S(x) \land \neg P(x))$ & $\neg \forall x (S(x) \to P(x))$ \\
Solo los $S$ son $P$ & $\forall x (P(x) \to S(x))$ & \\
Todos los $S$, y solo ellos, son $P$ & $\forall x (S(x) \leftrightarrow P(x))$ & \\
Todo es $P$; todo el mundo es $P$ & $\forall x P(x)$ & \\
Algo es $P$; alguien es $P$ & $\exists x P(x)$ & \\
Nada es $P$; nadie es $P$ & $\neg \exists x P(x)$ & $\forall x \neg P(x)$ \\
\hline
\end{tabular}
\end{center}

Algunas observaciones:
\begin{itemize}
    \item \textbf{Equivalencias.} Las dos columnas de la derecha son equivalentes, por la conversión de cuantificadores (CQ, sección~\ref{sec:cq}) y por equivalencias proposicionales como $\neg(A \to B) \equiv A \land \neg B$ o $\neg(A \land B) \equiv A \to \neg B$. Por ejemplo:
    \[
    \neg \forall x (S(x) \to P(x)) \ \equiv\ \exists x \neg (S(x) \to P(x)) \ \equiv\ \exists x (S(x) \land \neg P(x)).
    \]
    Cualquiera de las dos formas es una formalización
    correcta. 
    \item \textbf{«Solo».} «Solo los socios pueden votar» es $\forall x (V(x) \to S(x))$: ser socio es condición \emph{necesaria} para votar, como en el «solo si» de la sección~\ref{sec:formalizacion}. No dice que todos los socios puedan votar.
    \item \textbf{Adjetivos y oraciones de relativo.} $S$ y $P$ pueden ser condiciones complejas. «Todo filósofo griego es sabio» es $\forall x ((F(x) \land G(x)) \to B(x))$, y «algún filósofo griego es sabio», $\exists x (F(x) \land G(x) \land B(x))$. En «quien estudia aprueba», la oración de relativo hace de $S$: $\forall x (E(x) \to A(x))$.
    \item \textbf{«Todos\ldots{} no».} «Todos los alumnos no aprobaron» es ambiguo: puede significar que ninguno aprobó, $\forall x (L(x) \to \neg A(x))$, o que no todos aprobaron, $\neg \forall x (L(x) \to A(x))$. El refrán «no es oro todo lo que reluce» tiene la segunda forma: $\neg \forall x (R(x) \to O(x))$, es decir, $\exists x (R(x) \land \neg O(x))$. Dice que algunas cosas que relucen no son oro, no que nada de lo que reluce lo sea.
\end{itemize}

\subsection{Anidamiento de esquemas}

En los dos esquemas, el lugar de $P(x)$ puede ocuparlo una condición
compleja que contenga, a su vez, otro cuantificador. Así se forman los
enunciados con varios cuantificadores. Conviene formalizarlos por
capas, de fuera adentro: primero se identifica el esquema principal,
dejando en lenguaje natural la condición interior, y después se formaliza esa condición, en la que la variable del cuantificador exterior aparece libre.

\begin{ejemplo}[Un existencial dentro de un universal]
«Todos los deportistas tienen un entrenador.» Diccionario: $J(x)$: «$x$ es deportista»; $N(x)$: «$x$ es entrenador»; $T(x, y)$: «$x$ entrena a $y$».
\begin{enumerate}
    \item El esquema principal es universal: $\forall x (J(x) \to [x \text{ tiene un entrenador}])$.
    \item «$x$ tiene un entrenador» significa «algún entrenador entrena a $x$»: $\exists y (N(y) \land T(y, x))$.
\end{enumerate}
Resultado: $\forall x (J(x) \to \exists y (N(y) \land T(y, x)))$.
\end{ejemplo}

\begin{ejemplo}[Un universal dentro de un existencial]
«Hay un estudiante que aprueba todas las asignaturas.» Diccionario: $E(x)$: «$x$ es estudiante»; $M(x)$: «$x$ es una asignatura»; $A(x, y)$: «$x$ aprueba $y$».
\begin{enumerate}
    \item El esquema principal es existencial: $\exists x (E(x) \land [x \text{ aprueba todas las asignaturas}])$.
    \item «$x$ aprueba todas las asignaturas» significa «toda asignatura es aprobada por $x$»: $\forall y (M(y) \to A(x, y))$.
\end{enumerate}
Resultado: $\exists x (E(x) \land \forall y (M(y) \to A(x, y)))$.
\end{ejemplo}

Los dos ejemplos anidan los esquemas en sentido inverso, y el orden importa. En el primero cada deportista puede tener un entrenador distinto. Si se invierte el orden, se obtiene un enunciado más fuerte:
\begin{center}
\begin{tabular}{ll}
«Todos los deportistas tienen un entrenador» & $\forall x (J(x) \to \exists y (N(y) \land T(y, x)))$ \\
«Hay un entrenador que entrena a todos los deportistas» & $\exists y (N(y) \land \forall x (J(x) \to T(y, x)))$
\end{tabular}
\end{center}
El segundo implica el primero, pero no al revés. En una estructura con dominio $\{1, 2, 3, 4\}$, $I(J) = \{1, 2\}$, $I(N) = \{3, 4\}$ e $I(T) = \{(3, 1), (4, 2)\}$, cada deportista tiene su propio entrenador: el primer enunciado es verdadero y el segundo, falso. Es la misma asimetría entre $\exists\forall$ y $\forall\exists$ que vimos en la sección~\ref{sec:asimetria}. Del mismo modo, «toda asignatura la aprueba algún estudiante», $\forall y (M(y) \to \exists x (E(x) \land A(x, y)))$, es más débil que el enunciado del segundo ejemplo.

Al anidar conviene:
\begin{itemize}
    \item usar una variable distinta para cada cuantificador;
    \item colocar el cuantificador interior donde está la condición que traduce: en el consecuente del universal o dentro de la conjunción del existencial;
    \item leer de nuevo la fórmula en castellano, capa a capa, para comprobar que dice lo mismo que el original.
\end{itemize}

\practica{La pestaña «Formalizar» de \appmundo{} propone frases sobre el tablero y compara tu fórmula con una formalización correcta en muchos mundos. Si en alguno toman valores distintos, permite verlo en el tablero.}

\subsection{Identidad y cuantificadores numéricos\ampliacion}\label{sec:cuantificadores-numericos}

Con la identidad pueden expresarse cantidades: «hay al menos dos», «hay como mucho uno», «hay exactamente uno». Estas construcciones forman una familia que se puede describir de forma general. Escribiremos $\exists_{\geq n} x\, P(x)$ para «hay al menos $n$ individuos que son $P$», $\exists_{\leq n} x\, P(x)$ para «hay como mucho $n$» y $\exists_{= n} x\, P(x)$ para «hay exactamente $n$». No son símbolos nuevos, sino abreviaturas de fórmulas que solo usan cuantificadores, conectivas y $=$.

\paragraph{Al menos $n$.} Hay $n$ individuos que son $P$ y son distintos entre sí:
\begin{align*}
\exists_{\geq 1} x\, P(x) &:= \exists x P(x)\\
\exists_{\geq 2} x\, P(x) &:= \exists x \exists y (P(x) \land P(y) \land x \neq y)\\
\exists_{\geq 3} x\, P(x) &:= \exists x \exists y \exists z (P(x) \land P(y) \land P(z) \land x \neq y \land x \neq z \land y \neq z)
\end{align*}

\paragraph{Como mucho $n$.} Entre $n + 1$ individuos que sean $P$, al menos dos coinciden:
\begin{align*}
\exists_{\leq 1} x\, P(x) &:= \forall x \forall y ((P(x) \land P(y)) \to x = y)\\
\exists_{\leq 2} x\, P(x) &:= \forall x \forall y \forall z ((P(x) \land P(y) \land P(z)) \to (x = y \lor x = z \lor y = z))
\end{align*}
Equivalentemente, $\exists_{\leq n} x\, P(x)$ es $\neg \exists_{\geq n+1} x\, P(x)$.

\paragraph{Exactamente $n$.} Puede escribirse como $\exists_{\geq n} x\, P(x) \land \exists_{\leq n} x\, P(x)$, o de forma más compacta: hay $n$ individuos distintos que son $P$, y todo lo que es $P$ es uno de ellos.
\begin{align*}
\exists_{= 0} x\, P(x) &:= \neg \exists x P(x)\\
\exists_{= 1} x\, P(x) &:= \exists x (P(x) \land \forall y (P(y) \to y = x))\\
\exists_{= 2} x\, P(x) &:= \exists x \exists y (P(x) \land P(y) \land x \neq y \land \forall z (P(z) \to (z = x \lor z = y)))
\end{align*}
Con la identidad también se dice que un individuo es el único que cumple una condición: «solo $a$ es $P$» es $P(a) \land \forall x (P(x) \to x = a)$.

En lugar de $P$ puede ir cualquier condición. Por ejemplo, «Antonio ama exactamente a dos personas» (ejercicio~\ref{e5-2}) es $\exists_{= 2} x\, A(a, x)$, que se desarrolla como en la última línea.

\begin{observacion}[«Tres» en el lenguaje natural]
Los numerales no siempre significan «exactamente». «Hay tres
estudiantes en el aula» se entiende como «exactamente tres». En
cambio, «para aprobar hay que resolver bien tres ejercicios» se
entiende como «al menos tres»: quien resuelve bien cuatro también
aprueba. Al formalizar hay que decidir, según el contexto, qué lectura
se quiere.
\end{observacion}

\begin{observacion}[Límites]
Para cada número $n$ fijo, «al menos $n$», «como mucho $n$» y «exactamente $n$» se expresan en primer orden, aunque las fórmulas crecen con $n$. En cambio, «hay infinitos $P$» o «la mayoría de los $S$ son $P$» no pueden expresarse con ninguna fórmula de primer orden. «Hay infinitos $P$» sí se expresa con un conjunto infinito de fórmulas, $\{\exists_{\geq n} x\, P(x) : n \geq 1\}$, y que no baste una sola se sigue del teorema de compacidad (apéndice~\ref{ap:henkin}). Para «la mayoría», véase Barwise y Cooper (1981).
\end{observacion}

\subsection{Términos complejos: functores y descripciones\ampliacion}

Hasta ahora, los únicos términos son las constantes y las variables (sección~\ref{sec:sintaxis-lpo}). Pero el lenguaje natural tiene términos singulares complejos, como «el padre de Sócrates», «el autor de la \emph{República}» o «el actual rey de Francia», y también descripciones indefinidas, como «un ministro francés». Veamos tres maneras de ampliar el lenguaje para recogerlos. No las usaremos en el resto del curso: nos interesan por lo que muestran sobre la relación entre el lenguaje y la lógica.

\paragraph{Functores.} Además de predicados, el lenguaje puede tener \textbf{símbolos de función} o \textbf{functores} ($f, h, k, \dots$), cada uno con un número fijo de argumentos.
\begin{itemize}
    \item \textbf{Sintaxis.} Si $f$ es un functor de $n$ argumentos y $t_1, \dots, t_n$ son términos, $f(t_1, \dots, t_n)$ es un término. Los términos pueden anidarse; por ejemplo, con $p(x)$: «el padre de $x$», $p(p(s))$ es «el abuelo paterno de Sócrates».
    \item \textbf{Semántica.} La interpretación de un functor de $n$ argumentos es una función $I(f): D^n \to D$, que asigna a cada $n$-tupla de elementos del dominio un único elemento del dominio. La denotación de $f(t_1, \dots, t_n)$ es el valor de $I(f)$ para las denotaciones de $t_1, \dots, t_n$. Por ejemplo, en el dominio de los números naturales, con la constante $0$ interpretada como el cero y $f(x)$ como el sucesor de $x$, el término $f(f(0))$ denota el número $2$.
\end{itemize}
La función tiene que ser \emph{total}: toda $n$-tupla de elementos del dominio debe
tener un valor, y ese valor debe estar en el dominio. Sin embargo,
muchas expresiones funcionales del lenguaje natural son parciales: «la
capital de $x$» solo tiene sentido si $x$ es un país. En esos casos es
mejor usar una relación, como hace el ejercicio~\ref{e5-3} con «la
asignatura favorita de $x$ es $y$». 

Los functores no permiten decir nada que no pueda decirse sin
ellos. Un functor $f$ de un argumento puede sustituirse por un predicado binario $F(x,
y)$, «$y$ es $f(x)$», junto con la condición de que sea funcional,
$\forall x\, \exists_{= 1} y\, F(x, y)$. Entonces $P(f(a))$ se
reescribe como $\exists y (F(a, y) \land P(y))$.

\paragraph{Descripciones definidas: el operador $\iota$.} Una descripción definida, como «el autor de la \emph{República}», designa un objeto mediante una condición que solo él cumple.
\begin{itemize}
    \item \textbf{Sintaxis.} Si $A(x)$ es una fórmula y $x$ una variable, $\iota x A(x)$ («el $x$ tal que $A(x)$») es un término. Como los cuantificadores, $\iota$ liga la variable $x$. Con $E(x, y)$: «$x$ escribió $y$», $r$: la \emph{República}, $U(x, y)$: «$x$ fue discípulo de $y$» y $s$: Sócrates, «el autor de la \emph{República} fue discípulo de Sócrates» es $U(\iota x E(x, r), s)$.
    \item \textbf{Semántica.} Si exactamente un elemento $d$ del dominio satisface $A(x)$, el término $\iota x A(x)$ denota $d$.
\end{itemize}
La dificultad está en las descripciones \emph{impropias}: «el actual rey de Francia» no describe a nadie, y «el hijo de Sócrates» no describe a uno solo, porque Sócrates tuvo tres. Hay tres respuestas clásicas:
\begin{itemize}
    \item \textbf{Frege} (1892) sostuvo que un enunciado con un término que no se refiere a nada carece de valor de verdad, y propuso que, en un lenguaje lógico, se estipulara una referencia para esos términos.
    \item \textbf{Russell} (1905) negó que las descripciones sean auténticos términos: desaparecen al analizar el enunciado. Con $R(x)$: «$x$ es actualmente rey de Francia» y $C(x)$: «$x$ es calvo», «el actual rey de Francia es calvo» dice que hay un rey de Francia, que solo hay uno y que es calvo:
    \[
    \exists x (R(x) \land \forall y (R(y) \to y = x) \land C(x)).
    \]
    Como no hay rey de Francia, el enunciado es falso. Respecto a la
    negación, caben dos opciones: 
    \[
    \exists x (R(x) \land \forall y (R(y) \to y = x) \land \neg C(x)) \qquad \neg \exists x (R(x) \land \forall y (R(y) \to y = x) \land C(x))
    \]
    La primera lectura («hay un único rey de Francia y no es calvo»)
    es falsa; la segunda («no es cierto que haya un único rey de
    Francia y sea calvo»), verdadera. Así se cumple el tercero
    excluido: de «el actual rey de Francia es calvo» y su negación en
    la segunda lectura, una es verdadera. 
    \item \textbf{Strawson} (1950) objetó que Russell confunde lo que dice un enunciado con lo que presupone quien lo usa. Quien dice «el actual rey de Francia es calvo» da por supuesto que hay un rey de Francia; si no lo hay, no dice nada verdadero ni falso. Es la misma idea que veremos en la sección~\ref{sec:importacion} a propósito de la importación existencial.
\end{itemize}
La lógica de primer orden clásica adopta el análisis de Russell: cualquier enunciado con una descripción definida puede formalizarse sin $\iota$, con $\exists$, $\forall$ y $=$. Por ejemplo, «Belén no aprueba la asignatura favorita de Antonio» (ejercicio~\ref{e5-3}, apartado 6), en su lectura más natural, es
\[
\exists x (F(a, x) \land \forall y (F(a, y) \to y = x) \land \neg A(b, x)).
\]

\paragraph{Descripciones indefinidas: el operador $\varepsilon$ de
  Hilbert.} Russell analizó también las descripciones indefinidas («un
tal y cual») con el cuantificador existencial (Russell, 1919, cap.~16). En los años veinte, Hilbert introdujo un término $\varepsilon x A(x)$ que puede leerse como una descripción indefinida, «un $x$ tal que $A(x)$». Lo hizo con otro fin: eliminar los cuantificadores en sus pruebas de consistencia, dentro de su programa de fundamentación de las matemáticas (véase Hilbert y Bernays, 1939).
\begin{itemize}
    \item \textbf{Sintaxis.} Como en $\iota$: si $A(x)$ es una fórmula, $\varepsilon x A(x)$ es un término que liga la variable $x$.
    \item \textbf{Semántica.} Cada estructura incluye una \emph{función de elección}, que escoge un elemento de cada conjunto no vacío de objetos del dominio. El término $\varepsilon x A(x)$ denota el elemento escogido entre los que satisfacen $A(x)$. Si ninguno lo satisface, denota un elemento cualquiera del dominio, también fijado por la estructura.
\end{itemize}
La diferencia entre $\iota$ y $\varepsilon$ está en lo que exigen. «El rey de Francia», $\iota x R(x)$, solo designa algo si hay exactamente un rey. «Un ministro francés», $\varepsilon x M(x)$, solo necesita que haya alguno, y entonces designa a uno de ellos, sin decir cuál.

El principio básico de $\varepsilon$ es que, si algo cumple $A(x)$, el objeto escogido lo cumple: $A(t) \to A(\varepsilon x A(x))$ para cualquier término $t$. De ahí resulta que los cuantificadores pueden definirse con $\varepsilon$:
\[
\exists x A(x) \leftrightarrow A(\varepsilon x A(x)) \qquad\qquad \forall x A(x) \leftrightarrow A(\varepsilon x \neg A(x))
\]
La segunda dice que $A(x)$ vale para todo si vale para el objeto escogido entre los que no la cumplen: si hubiera alguno que no la cumpliera, el escogido sería uno de ellos.

Hay que tener cuidado al formalizar con $\varepsilon$. «Un ministro francés visitó Sevilla» (con $e$: Sevilla) se formaliza con el término $t$ formado con toda la condición, $\varepsilon x (M(x) \land V(x, e))$, como $M(t) \land V(t, e)$, que equivale a $\exists x (M(x) \land V(x, e))$. En cambio, $V(\varepsilon x M(x), e)$ dice que el ministro escogido visitó Sevilla. Si hay ministros franceses, esto implica lo anterior, pero no al revés; si no los hay, $\varepsilon x M(x)$ denota un objeto cualquiera y la fórmula ya no habla de un ministro. Por último, el testigo de la regla $\exists$E del tema~\ref{cap:dn-lpo} se parece a un término $\varepsilon x A(x)$: es una constante de la que solo sabemos que cumple $A(x)$, y por eso no podemos suponer de ella nada más. A diferencia de $\varepsilon x A(x)$, el testigo no puede salir de su caja, y cada uso de $\exists$E exige una constante fresca.

\section{Contraejemplos semánticos: cómo refutar un argumento inválido}

En el tema~\ref{cap:dn-lpo} aprendimos que para demostrar que un
argumento es válido construimos una deducción formal en Fitch ($\Gamma
\vdash A$). Ahora bien, para demostrar la invalidez
($\Gamma\not\models A$) debemos exhibir un \textbf{contramodelo} (contraejemplo semántico): una estructura concreta $\mathcal{M} = \langle D, I \rangle$ que haga verdaderas todas las premisas de $\Gamma$ y falsa la conclusión $A$.

\begin{ejemplo}[Refutación de la falacia de inversión de cuantificadores]
En el tema~\ref{cap:dn-lpo} señalamos que pasar de $\forall y \exists x R(x, y)$ a $\exists x \forall y R(x, y)$ es una falacia formal. Demostremos ahora semánticamente su invalidez mediante un contramodelo minimal con solo dos elementos:
\[
\forall y \exists x R(x, y) \ \not\models \ \exists x \forall y R(x, y)
\]
\textbf{Construcción del contramodelo minimal $\mathcal{M}$:}
\begin{itemize}
    \item Elegimos un dominio finito de dos individuos: $D = \{1, 2\}$.
    \item Interpretamos la relación binaria $R$ como la identidad en $D$:
    \[
    I(R) = \{(1, 1), \ (2, 2)\}.
    \]
\end{itemize}
\textbf{Evaluación en $\mathcal{M}$:}
\begin{enumerate}
    \item La premisa $\forall y \exists x R(x, y)$, que afirma que «para todo elemento $y$ existe al menos un elemento $x$ que se relaciona con él», es verdadera:
    \begin{itemize}
        \item Para $y = 1$: eligiendo $x = 1$, tenemos $(1, 1) \in I(R)$.
        \item Para $y = 2$: eligiendo $x = 2$, tenemos $(2, 2) \in I(R)$.
    \end{itemize}
    Como se satisface para todos los individuos de $D$, la premisa es verdadera en $\mathcal{M}$.
    \item La conclusión $\exists x \forall y R(x, y)$, que afirma que «existe algún elemento $x$ que se relaciona con \emph{todos} los elementos $y$ del dominio», es falsa:
    \begin{itemize}
        \item Para $x = 1$: $(1, 1) \in I(R)$, pero $(1, 2) \notin I(R)$. Luego el $1$ no se relaciona con todos.
        \item Para $x = 2$: $(2, 2) \in I(R)$, pero $(2, 1) \notin I(R)$. Luego el $2$ no se relaciona con todos.
    \end{itemize}
    Dado que ningún individuo de $D$ se relaciona simultáneamente con $1$ y con $2$, la conclusión es falsa en $\mathcal{M}$.
\end{enumerate}
Puesto que en la estructura $\mathcal{M}$ la premisa es verdadera y la conclusión es falsa, queda demostrado formalmente que el argumento es inválido.
\end{ejemplo}

\begin{ejemplo}[Construcción de un contramodelo finito: falacia del término medio no distribuido]
Deseamos demostrar semánticamente que la siguiente inferencia no es válida:
\[
\forall x (P(x) \to M(x)), \ \forall x (S(x) \to M(x)) \ \not\models \ \forall x (S(x) \to P(x))
\]
\textbf{Estrategia de diseño del contramodelo.}
Buscamos una estructura con el menor número posible de elementos donde las dos premisas sean verdaderas y la conclusión falle. Para que la conclusión $\forall x (S(x) \to P(x))$ sea falsa, necesitamos que exista al menos un individuo que sea $S$ pero no sea $P$.
Definimos la estructura $\mathcal{M}$ con un dominio minimal de un único elemento $D = \{1\}$:
\begin{itemize}
    \item $I(S) = \{1\}$ (el individuo $1$ posee la propiedad $S$).
    \item $I(P) = \varnothing$ (el individuo $1$ no posee la propiedad $P$).
    \item $I(M) = \{1\}$ (el individuo $1$ posee la propiedad $M$).
\end{itemize}
\textbf{Comprobación paso a paso:}
\begin{enumerate}
    \item \textbf{Premisa 1 ($\forall x (P(x) \to M(x))$):} para $x=1$, el antecedente $1 \in I(P)$ es falso. Como un condicional con antecedente falso es verdadero ($\text{F} \to \text{V} = \text{V}$), la premisa es verdadera.
    \item \textbf{Premisa 2 ($\forall x (S(x) \to M(x))$):} para $x=1$, tanto $1 \in I(S)$ como $1 \in I(M)$ son verdaderos ($\text{V} \to \text{V} = \text{V}$). La premisa es verdadera.
    \item \textbf{Conclusión ($\forall x (S(x) \to P(x))$):} para $x=1$, el antecedente $1 \in I(S)$ es verdadero, pero el consecuente $1 \in I(P)$ es falso. El condicional toma valor $\text{V} \to \text{F} = \text{F}$. Por consiguiente, la conclusión es falsa.
\end{enumerate}
\textbf{Comentario.} Hemos demostrado la invalidez del argumento. Obsérvese que un contramodelo no necesita ser complicado: una estructura con un solo individuo basta para refutar la inferencia, porque su validez exigiría que la conclusión fuera verdadera en todas las estructuras que hacen verdaderas las premisas.
\end{ejemplo}

\practica{En la pestaña «Contramodelos» de \appmundo{} puedes construir en el tablero un contramodelo para un argumento y ver si lo es. Tiene también retos: argumentos para los que hay que buscar un contramodelo, si lo tienen.}

\section{Silogística e importación existencial}\label{sec:importacion}

En el tema~\ref{cap:dn-lpo} formalizamos varios modos silogísticos (Barbara, Celarent, Darii, Ferio\ldots). Las cuatro formas de enunciado de la silogística tradicional se formalizan así:
\begin{center}
\renewcommand{\arraystretch}{1.25}
\begin{tabular}{|c|l|l|}
\hline
\textbf{Forma} & \textbf{Enunciado} & \textbf{Formalización} \\
\hline
A & Todo $S$ es $P$ & $\forall x (S(x) \to P(x))$ \\
E & Ningún $S$ es $P$ & $\forall x (S(x) \to \neg P(x))$ \\
I & Algún $S$ es $P$ & $\exists x (S(x) \land P(x))$ \\
O & Algún $S$ no es $P$ & $\exists x (S(x) \land \neg P(x))$ \\
\hline
\end{tabular}
\end{center}

En la lógica de primer orden, $\forall x (S(x) \to P(x))$ es verdadera cuando no hay ningún $S$: el condicional se cumple trivialmente para cada individuo. En cambio, la silogística tradicional suponía que los términos no son vacíos: daba por hecho que «todo $S$ es $P$» implica «algún $S$ es $P$» (\emph{subalternación}). Esta suposición se llama \textbf{importación existencial}, y la lógica de primer orden no la hace:
\[
\forall x (S(x) \to P(x)) \not\models \exists x (S(x) \land P(x))
\]
Basta un contramodelo con $D = \{1\}$ e $I(S) = I(P) = \varnothing$: la premisa es verdadera (no hay ningún $S$) y la conclusión falsa.

\paragraph{Consecuencias.} Algunos modos considerados válidos por la tradición dejan de serlo. Por ejemplo, \emph{Darapti}:
\[
\forall x (M(x) \to P(x)), \ \forall x (M(x) \to S(x)) \ \not\models \ \exists x (S(x) \land P(x))
\]
En una estructura en la que no hay ningún $M$ las dos premisas son verdaderas, y la conclusión es falsa si no hay nada que sea a la vez $S$ y $P$. Lo mismo ocurre con otros modos como \emph{Felapton} o \emph{Bramantip}. Todos ellos vuelven a ser válidos si se añade como premisa que el término en cuestión no es vacío (en Darapti, $\exists x M(x)$). En el cuadro de oposición tradicional, las relaciones entre contradictorias (A y O, E e I) se mantienen, pero A y E pueden ser ambas verdaderas, I y O pueden ser ambas falsas, y la subalternación falla.

\paragraph{¿Quién tiene razón?} No es una cuestión puramente técnica. La lógica moderna separa la afirmación general («todo $S$ es $P$») de la afirmación de existencia («hay $S$»), lo que permite enunciar leyes sobre casos que quizá no existan («todo cuerpo libre de fuerzas se mueve con velocidad constante»). Otros autores, como P. F. Strawson, han defendido que «todos los $S$ son $P$» \emph{presupone} que hay $S$: si no los hay, el enunciado no sería ni verdadero ni falso. Al formalizar conviene preguntarse si el argumento depende de esa suposición y, si es así, añadirla como premisa explícita.

\section{Los teoremas metalógicos de primer orden}

Al igual que en lógica proposicional, la metalógica de primer orden investiga la correspondencia exacta entre la derivabilidad sintáctica ($\vdash$) y la consecuencia semántica ($\models$):

\begin{theorem}[Corrección de primer orden]
Para cualquier conjunto de sentencias $\Gamma$ y cualquier sentencia $A$ de LPO:
\[
\text{Si } \Gamma \vdash A, \text{ entonces } \Gamma \models A.
\]
\end{theorem}
El cálculo de deducción natural nunca lleva de premisas verdaderas a una conclusión falsa. La demostración sigue la del tema~\ref{cap:meta-prop}: cada línea es consecuencia lógica de las premisas e hipótesis activas en ella. Para las reglas $\forall$I y $\exists$E esto solo es cierto gracias a las restricciones sobre constantes propias y testigos; los ejemplos de falacias del tema~\ref{cap:dn-lpo} muestran lo que ocurriría sin ellas.

\begin{proof}[Idea de la demostración para los cuantificadores\ampliacion]
Como en el tema~\ref{cap:meta-prop}, se prueba por inducción que cada línea es consecuencia lógica del conjunto $\Delta$ de premisas e hipótesis activas en ella: toda estructura que hace verdaderas las fórmulas de $\Delta$ hace verdadera la línea. Las reglas proposicionales se tratan igual que allí; $\forall$E y $\exists$I son inmediatas, y también las de la identidad ($=$E es correcta porque si $\mathcal{M} \models a = b$, entonces $I(a) = I(b)$). Veamos los dos casos interesantes. En ambos usamos, dada una estructura $\mathcal{M}$ y un elemento $d$ de su dominio, la estructura $\mathcal{M}_d$, que coincide con $\mathcal{M}$ salvo en que interpreta la constante $c$ como $d$.
\begin{itemize}
    \item \textbf{$\forall$I.} La línea es $\forall x A(x)$, obtenida de $A(c)$, donde $c$ no aparece en $\Delta$ ni en $\forall x A(x)$. Sea $\mathcal{M} \models \Delta$ y sea $d$ cualquier elemento del dominio. Como $c$ no aparece en $\Delta$, también $\mathcal{M}_d \models \Delta$; por la hipótesis de inducción, $\mathcal{M}_d \models A(c)$, es decir, $d$ satisface $A(x)$ en $\mathcal{M}_d$. Como $c$ no aparece en $A(x)$, $d$ satisface $A(x)$ también en $\mathcal{M}$. Al ser $d$ arbitrario, $\mathcal{M} \models \forall x A(x)$. (Sin la restricción (1) de $\forall$I del tema~\ref{cap:dn-lpo}, $\mathcal{M}_d$ podría no satisfacer $\Delta$; sin la (2), cambiar la interpretación de $c$ cambiaría lo que dice $A(x)$).
    \item \textbf{$\exists$E.} La línea es $B$, obtenida de $\exists x A(x)$ y de una caja que va de $A(c)$ a $B$, donde $c$ no aparece en $\Delta$, ni en $\exists x A(x)$, ni en $B$. Sea $\mathcal{M} \models \Delta$. Por la hipótesis de inducción, $\mathcal{M} \models \exists x A(x)$, luego algún $d$ satisface $A(x)$. Entonces $\mathcal{M}_d \models \Delta$ (porque $c$ no aparece en $\Delta$) y $\mathcal{M}_d \models A(c)$ (porque $c$ no aparece en $A(x)$). Por la hipótesis de inducción aplicada a la última línea de la caja, $\mathcal{M}_d \models B$, y como $c$ no aparece en $B$, también $\mathcal{M} \models B$. \qedhere
\end{itemize}
\end{proof}

\begin{theorem}[Completud de primer orden: Gödel, 1930; Henkin, 1949]
Para cualquier conjunto de sentencias $\Gamma$ y cualquier sentencia $A$ de LPO:
\[
\text{Si } \Gamma \models A, \text{ entonces } \Gamma \vdash A.
\]
\end{theorem}
Toda inferencia semánticamente válida puede ser demostrada formalmente mediante las reglas de deducción natural.

\subsection{Consistencia, satisfacibilidad y completud}\label{sec:consistencia-satisfacibilidad}

La completud de primer orden no puede demostrarse siguiendo el camino de Kalmár (sección~\ref{sec:no-kalmar}). Las demostraciones habituales, como la de Henkin, parten de una reformulación del teorema que relaciona una noción sintáctica con una semántica. La semántica es la satisfacibilidad, que ya hemos definido: $\Gamma$ es satisfacible si existe una estructura $\mathcal{M}$ tal que $\mathcal{M} \models \Gamma$. La sintáctica es la consistencia:

\begin{definition}[Consistencia]
Un conjunto de sentencias $\Gamma$ es \textbf{consistente} si $\Gamma \nvdash \bot$, es decir, si no hay ninguna derivación de $\bot$ cuyas premisas pertenezcan a $\Gamma$. En caso contrario, $\Gamma$ es \textbf{inconsistente}.
\end{definition}

El teorema de corrección garantiza que todo conjunto satisfacible es consistente: si fuera $\Gamma \vdash \bot$, tendríamos $\Gamma \models \bot$, y ninguna estructura podría hacer verdaderas todas las fórmulas de $\Gamma$, porque ninguna hace verdadera $\bot$. La completud equivale a la afirmación recíproca:

\begin{lemma}[Lema de reducción]\label{lema_reduccion}
Las dos afirmaciones siguientes son equivalentes:
\begin{enumerate}
    \item \textbf{Completud fuerte:} para todo conjunto de sentencias $\Gamma$ y toda sentencia $A$, si $\Gamma \models A$, entonces $\Gamma \vdash A$.
    \item \textbf{Teorema del modelo:} todo conjunto consistente de sentencias es satisfacible.
\end{enumerate}
\end{lemma}

\begin{proof}
$(2) \Rightarrow (1)$. Supongamos (2) y que $\Gamma \models A$. Si
fuera $\Gamma \nvdash A$, el conjunto $\Gamma \cup \{\neg A\}$ sería
consistente: si se tuviera $\Gamma \cup \{\neg A\} \vdash \bot$,
bastaría tomar $\neg A$ como hipótesis de una caja que llegue a $\bot$
y aplicar IP para obtener $\Gamma \vdash A$. Por (2), existiría una
estructura $\mathcal{M}$ con $\mathcal{M} \models \Gamma$ y
$\mathcal{M} \models \neg A$, es decir, un contramodelo de $\Gamma
\models A$. Como esto no puede ocurrir (por $\Gamma\models A$),
$\Gamma \vdash A$. 

$(1) \Rightarrow (2)$. Supongamos (1) y que $\Gamma$ no es satisfacible. Entonces toda estructura que hace verdaderas todas las fórmulas de $\Gamma$ (no hay ninguna) hace verdadera $\bot$, es decir, $\Gamma \models \bot$. Por (1), $\Gamma \vdash \bot$: $\Gamma$ es inconsistente. Por contraposición, todo conjunto consistente es satisfacible.
\end{proof}

Junto con los teoremas de corrección y completud, el lema muestra que la consistencia y la satisfacibilidad, aunque se definen de modo distinto (una en el cálculo, otra en la semántica), coinciden: un conjunto de sentencias es consistente si y solo si es satisfacible. Lo mismo vale en lógica proposicional, con interpretaciones (asignaciones de verdad) en lugar de estructuras. Demostrar la completud equivale, por tanto, a demostrar que todo conjunto consistente tiene un modelo. Es lo que hizo Henkin: dado un conjunto consistente, construyó una estructura que lo satisface (sección~\ref{sec:estrategia-henkin} y apéndice~\ref{ap:henkin}).

\section{¿Por qué no se puede usar la prueba de Kalmár en primer orden?\ampliacion}\label{sec:no-kalmar}

Tras estudiar la completud proposicional (tema~\ref{cap:meta-prop}) podemos preguntarnos
por qué no aplicar la demostración constructiva de Kalmár a la lógica
de primer orden. Veamos: 

\begin{enumerate}
    \item \textbf{Imposibilidad de tablas de verdad finitas.}  
    El método de Kalmár se basa en que una fórmula con $n$ variables proposicionales tiene exactamente $2^n$ combinaciones posibles de verdad. Cada fila de la tabla genera una prueba finita, y luego el principio del tercero excluido (LEM) colapsa sucesivamente las $2^n$ pruebas de dos en dos hasta llegar a cero premisas.  
    En primer orden, una fórmula cuantificada como $\forall x P(x)$ habla sobre un dominio $D$ que puede ser infinito. No existe ninguna tabla de verdad de tamaño finito que contemple todas las interpretaciones posibles de una relación, ni un número finito de filas que podamos colapsar sucesivamente mediante LEM.
    
    \item \textbf{La barrera de la indecidibilidad (Church y Turing, 1936).}  
    El método de Kalmár parte de la tabla de verdad, que es a la vez un \emph{procedimiento de decisión}: revisando las $2^n$ filas sabemos, en un número finito de pasos, si la fórmula es válida o no (y, si lo es, construimos su deducción). Un análogo en primer orden nos daría un algoritmo para decidir la validez de cualquier fórmula.  
    Sin embargo, Alonzo Church y Alan Turing demostraron en 1936 que la lógica de primer orden es \textbf{indecidible}: no existe ningún algoritmo que determine, para cualquier fórmula de primer orden, si es válida o no. Sí existen procedimientos que, dada una fórmula \emph{válida}, acaban encontrando una deducción de ella; pero ante una fórmula no válida pueden no terminar nunca. Volveremos sobre ello en el tema~\ref{cap:tableaux}.
    
    \item \textbf{La necesidad de una prueba conjuntista no constructiva.}  
    Puesto que ningún algoritmo finitista puede inspeccionar todos los modelos posibles de primer orden, la demostración de completud debe cambiar de estrategia: en lugar de ir de las tablas de verdad a la prueba (Kalmár), debemos probar que si un conjunto de sentencias es consistente, existe un modelo que lo satisface (Henkin).
\end{enumerate}

\section{La estrategia de Leon Henkin (1949)\ampliacion}\label{sec:estrategia-henkin}

En 1949, el lógico estadounidense Leon Henkin publicó una
nueva demostración de completud. En lugar de intentar demostrar $\Gamma \models A \implies \Gamma \vdash A$ de forma directa, demostró el teorema del modelo, que equivale a la completud fuerte (lema~\ref{lema_reduccion}):
\[
\Gamma \nvdash \bot \implies \exists \mathcal{M} \ (\mathcal{M} \models \Gamma)
\]

Para fabricar ese modelo, Henkin introduce dos conceptos:

\begin{definition}[Conjunto de Henkin]
Un conjunto de sentencias $\Delta$ en un lenguaje de primer orden es un \textbf{conjunto de Henkin} si satisface tres propiedades:
\begin{enumerate}
    \item \textbf{Consistencia sintáctica.} $\Delta \nvdash \bot$ (a partir de $\Delta$ no se deriva $\bot$).
    \item \textbf{Maximalidad (Lindenbaum).} Para cualquier sentencia $A$ del lenguaje, o bien $A \in \Delta$ o bien $\neg A \in \Delta$. 
    \item \textbf{Propiedad de los testigos (Henkin).} Para cada fórmula existencial $\exists x B(x) \in \Delta$, existe una constante individual concreta $c$ en el lenguaje tal que $B(c) \in \Delta$.
\end{enumerate}
\end{definition}

El lema de Lindenbaum-Henkin (apéndice~\ref{ap:henkin}) garantiza que todo conjunto consistente $\Gamma$ puede ampliarse a un conjunto de Henkin $\Gamma^* \supseteq \Gamma$ en un lenguaje enriquecido con nuevas constantes (los testigos).

\subsection{La estructura del modelo canónico de términos}
¿Dónde encuentra Henkin un modelo para satisfacer ese conjunto maximal de fórmulas? La clave del método radica en construir el dominio a partir del propio material sintáctico: las constantes del lenguaje formal.
\begin{itemize}
    \item Toma como dominio $D$ el conjunto de las constantes del lenguaje (ampliado con nuevas constantes que sirven de testigos): $D = \{c_1, c_2, c_3, \dots\}$, y cada constante se interpreta a sí misma. Si el lenguaje contiene la identidad, el dominio no puede estar formado por las constantes mismas, sino por clases de constantes que $\Gamma^*$ declara idénticas. 
    \item Interpreta cada predicado según las fórmulas atómicas que pertenecen a $\Gamma^*$:
    \[
    (c_1, \dots, c_n) \in I(P) \iff P(c_1, \dots, c_n) \in \Gamma^*.
    \]
\end{itemize}
En el modelo canónico de Henkin para $\Gamma^*$, denotado por $\mathcal{M}_{\Gamma^*}$, la verdad de cada sentencia coincide exactamente con su pertenencia a $\Gamma^*$:
\[
\mathcal{M}_{\Gamma^*} \models A \iff A \in \Gamma^* \qquad (\text{para toda sentencia } A).
\]
En particular, como $\Gamma \subseteq \Gamma^*$, el modelo $\mathcal{M}_{\Gamma^*}$ satisface $\Gamma$. La demostración completa se desarrolla en el apéndice~\ref{ap:henkin}.

\section{Cuadro comparativo: el método de Kalmár frente al método de Henkin\ampliacion}

La siguiente tabla resume las diferencias entre los dos métodos de demostración de la completud:

\begin{center}
\renewcommand{\arraystretch}{1.4}
\begin{tabular}{|>{\raggedright\arraybackslash}p{3.2cm}|>{\raggedright\arraybackslash}p{6cm}|>{\raggedright\arraybackslash}p{6.4cm}|}
\hline
\textbf{Dimensión} & \textbf{Método de László Kalmár (1935)} & \textbf{Método de Leon Henkin (1949)} \\
\hline
\textbf{Ámbito de aplicación} & Lógica proposicional clásica. & Lógica proposicional y de primer orden (y extensiones de orden superior o modal). \\
\hline
\textbf{Naturaleza metodológica} & \textbf{Constructiva y algorítmica:} proporciona un procedimiento explícito para traducir tablas de verdad en deducciones. & \textbf{No constructiva y conjuntista:} prueba la existencia de un modelo mediante el lema de Lindenbaum. \\
\hline
\textbf{Estructura de la prueba} & Inducción finita sobre la complejidad sintáctica y colapso de premisas con la regla LEM. & Extensión a un conjunto maximalmente consistente con testigos y construcción del modelo canónico. \\
\hline
\textbf{Tratamiento del infinito} & Limitado a
                                    \textbf{conjuntos finitos de
                                    premisas}. & Trabaja con
                                                 \textbf{conjuntos
                                                 infinitos maximalmente
                                                 consistentes}. \\
\hline
\textbf{Punto de partida semántico} & Tablas de verdad finitas ($2^n$ filas). & Estructuras y modelos con dominios arbitrarios (finitos o infinitos). \\
\hline
\textbf{Relación con la computación} & Estrechamente ligada a la decidibilidad proposicional. & Compatible con la indecidibilidad de primer orden (Church-Turing). \\
\hline
\end{tabular}
\end{center}


\section*{Referencias bibliográficas}

\begin{enumerate}
    \item \textbf{Barwise, J., \& Cooper, R.} (1981). «Generalized quantifiers and natural language». \emph{Linguistics and Philosophy}, vol. 4, n.º 2, pp. 159--219.
    \item \textbf{Boolos, G., Burgess, J. P., \& Jeffrey, R.} (2007). \emph{Computability and Logic} (5.ª ed.). Cambridge University Press. [Capítulos dedicados a modelos, satisfacción y completud].
    \item \textbf{Frege, G.} (1892). «Über Sinn und Bedeutung». \emph{Zeitschrift für Philosophie und philosophische Kritik}, vol. 100, pp. 25--50.
    \item \textbf{Garrido, M.} (1974). \emph{Lógica simbólica}. Tecnos, Madrid. [Capítulos sobre la semántica del cálculo de predicados].
    \item \textbf{Gödel, K.} (1930). «Die Vollständigkeit der Axiome des logischen Funktionenkalküls». \emph{Monatshefte für Mathematik und Physik}, vol. 37, pp. 349--360. [Primera prueba de la completud de primer orden].
    \item \textbf{Henkin, L.} (1949). «The completeness of the first-order functional calculus». \emph{Journal of Symbolic Logic}, vol. 14, n.º 3, pp. 159--166. [El artículo fundacional del método de modelos canónicos].
    \item \textbf{Hilbert, D., \& Bernays, P.} (1939). \emph{Grundlagen der Mathematik}, vol. II. Springer, Berlín.
    \item \textbf{Russell, B.} (1905). «On denoting». \emph{Mind}, vol. 14, n.º 56, pp. 479--493.
    \item \textbf{Russell, B.} (1919). \emph{Introduction to Mathematical Philosophy}. George Allen \& Unwin, Londres. [Cap. 16: «Descriptions»].
    \item \textbf{Strawson, P. F.} (1950). «On referring». \emph{Mind}, vol. 59, n.º 235, pp. 320--344.
    \item \textbf{Tarski, A.} (1936). «Der Wahrheitsbegriff in den formalisierten Sprachen». \emph{Studia Philosophica}, vol. 1, pp. 261--405. [Versión alemana del original polaco de 1933, \emph{Pojęcie prawdy w językach nauk dedukcyjnych}; definición fundacional de verdad semántica y satisfacción].
\end{enumerate}

%%% Local Variables:
%%% mode: latex
%%% TeX-master: "../ApuntesLogica"
%%% End:


\section{Ejercicios de formalización y semántica en LPO}
% -*- coding: utf-8 -*-
\begin{quote}
\small \textbf{Instrucciones.} Estos ejercicios usan sus propios predicados, así que se hacen a lápiz. Para practicar más la formalización, la evaluación de fórmulas y los contramodelos, usa la aplicación \appmundo{}, con su diccionario de formas, tamaños, colores y posiciones.
\end{quote}

\section*{Bloque 1: formalización de enunciados del lenguaje natural en LPO}
\textit{En este bloque practicamos la formalización de enunciados al
  lenguaje de la lógica de predicados.}

\begin{pregunta}[Enunciados categóricos y propiedades individuales]\label{e5-1}
Formaliza los siguientes enunciados utilizando los símbolos de predicado indicados entre paréntesis:
\begin{enumerate}
    \item Las manzanas y las naranjas son gustosas y nutritivas.  
    ($M(x)$: «$x$ es una manzana»; $N(x)$: «$x$ es una naranja»; $G(x)$: «$x$ es gustosa»; $U(x)$: «$x$ es nutritiva»).
    \item Hay alimentos que solo se pueden comer si han sido cocinados.  
    ($A(x)$: «$x$ es un alimento»; $C(x)$: «$x$ se puede comer»; $K(x)$: «$x$ ha sido cocinado»).
    \item Sin frenos, no hay ningún coche seguro.  
    ($F(x)$: «$x$ tiene frenos»; $C(x)$: «$x$ es un coche»; $S(x)$: «$x$ es seguro»).
    \item No todo el mundo es rico, culto y educado, ni todos los ricos son educados y cultos.  
    ($R(x)$: «$x$ es rico»; $C(x)$: «$x$ es culto»; $E(x)$: «$x$ es educado». Dominio: todas las personas).
    \item No todas las cosas compradas a bajo precio son delicadas y quebradizas.  
    ($C(x)$: «$x$ es una cosa»; $B(x)$: «$x$ ha sido comprada a bajo precio»; $L(x)$: «$x$ es delicada»; $Q(x)$: «$x$ es quebradiza»).
\end{enumerate}
\end{pregunta}

\begin{pregunta}[Relaciones, cuantificación e identidad]\label{e5-2}
Dado el predicado binario $A(x, y)$ («$x$ ama a $y$») y las constantes individuales $a$ (Antonio) y $b$ (Belén), formaliza los siguientes enunciados:
\begin{enumerate}
    \item Antonio ama a Belén pero Belén no ama a Antonio.
    \item Todo el mundo ama a Belén.
    \item Belén ama a alguien que también la ama a ella.
    \item Belén ama a alguien que no es Antonio (lo cual no implica que no ame a Antonio).
    \item Antonio siempre es correspondido en su amor.
    \item Belén ama a una sola persona que no es ella misma.
    \item Antonio ama exactamente a dos personas (véase la sección~\ref{sec:cuantificadores-numericos}).
    \item Todo el mundo se ama a sí mismo.
    \item Todo el que ama es correspondido.
    \item Todo el que ama es amado.
    \item Hay quien no es correspondido en su amor.
    \item Todo el mundo ama a alguien distinto de sí mismo.
\end{enumerate}
\end{pregunta}

\begin{pregunta}\label{e5-3}
Utilizando los predicados:
\begin{itemize}
    \item $A(x, y)$: «$x$ aprueba la asignatura $y$»
    \item $E(x)$: «$x$ estudia filosofía»
    \item $F(x, y)$: «la asignatura favorita de $x$ es $y$»
    \item $S(x)$: «$x$ es una asignatura»
\end{itemize}
y las constantes individuales $a$ (Antonio), $b$ (Belén) y $l$ (Lógica), formaliza los siguientes enunciados:
\begin{enumerate}
    \item La lógica no es la asignatura favorita de ningún estudiante de filosofía.
    \item Hay al menos dos estudiantes de filosofía distintos que tienen la misma asignatura favorita.
    \item La asignatura favorita de Antonio no es la misma que la de Belén.
    \item Antonio y Belén aprueban exactamente las mismas asignaturas.
    \item Todos los estudiantes de filosofía cuya asignatura favorita es la lógica, la aprueban.
    \item Belén no aprueba la asignatura favorita de Antonio.
    \item Belén aprueba todas las asignaturas que Antonio suspende (entiende que suspender es no aprobar).
    \item La lógica no es la asignatura favorita de ningún estudiante de filosofía que la suspenda.
\end{enumerate}
\end{pregunta}

\begin{pregunta}[Anidamiento de cuantificadores]\label{e5-anid}
Formaliza los siguientes enunciados, indicando el diccionario:
\begin{enumerate}
    \item Todo filósofo ha leído a algún clásico.
    \item Hay un clásico que todos los filósofos han leído.
    \item Ningún estudiante aprueba todas las asignaturas.
    \item Toda asignatura tiene algún estudiante que no la aprueba.
    \item Hay un profesor al que no entiende ninguno de sus alumnos.
\end{enumerate}
¿Cuál de los enunciados 1 y 2 implica al otro? Da una estructura que muestre que el otro no lo implica.
\end{pregunta}

\begin{pregunta}[Cantidad y descripciones\ampliacion]\label{e5-cant}
Formaliza sin usar functores ni los operadores $\iota$ y $\varepsilon$:
\begin{enumerate}
    \item Hay al menos tres filósofos presocráticos.
    \item Belén suspende como mucho una asignatura.
    \item Hay exactamente dos asignaturas que Belén no aprueba.
    \item El autor de la \emph{Ética a Nicómaco} fue discípulo de Platón.
    \item La montaña de oro no existe (según el análisis de Russell).
\end{enumerate}
\end{pregunta}

\section*{Bloque 2: evaluación de fórmulas y verdad en estructuras}

\begin{pregunta}[Diseño de modelos]\label{e5-4}
Considera las oraciones formalizadas en el ejercicio~\ref{e5-2} sobre el predicado $A(x, y)$ («amar»). Diseña una estructura concreta $\mathcal{M} = \langle D, I \rangle$ con un dominio finito de cuatro personas $D = \{1, 2, 3, 4\}$ con $I(a) = 1$ e $I(b) = 2$ donde se verifiquen simultáneamente:
\begin{enumerate}
    \item Belén ama a alguien que también la ama a ella ($\exists x (A(b, x) \land A(x, b))$).
    \item Antonio siempre es correspondido en su amor ($\forall x (A(a, x) \to A(x, a))$).
    \item Hay quien no es correspondido en su amor ($\exists x \exists y (A(x, y) \land \neg A(y, x))$).
\end{enumerate}
Define explícitamente el conjunto de pares $I(A) \subseteq D \times D$ y justifica por qué satisface las tres fórmulas.
\end{pregunta}

\begin{pregunta}[Justificación semántica del silogismo clásico]\label{e5-5}
Demuestra formalmente, utilizando únicamente las definiciones semánticas de verdad y modelo, que el siguiente argumento es válido:
\begin{quote}
\textit{Los caracoles son gasterópodos.}\\
\textit{Los gasterópodos son moluscos.}\\
\textit{Por tanto, los caracoles son moluscos.}
\end{quote}
\textit{Pista.} Formaliza las premisas como $\forall x (C(x) \to
G(x))$ y $\forall x (G(x) \to M(x))$, y la conclusión como $\forall x
(C(x) \to M(x))$. Supón una estructura arbitraria $\mathcal{M} =
\langle D, I \rangle$ que hace verdaderas las premisas y razona para
cualquier elemento $d \in D$ observando que $I(C) \subseteq I(G)
\subseteq I(M)$. 
\end{pregunta}

\begin{pregunta}[Evaluación en una estructura aritmética infinita]\label{e5-6}
Sea la estructura $\mathcal{M} = \langle D, I \rangle$, donde $D = \mathbb{N} = \{0, 1, 2, 3, \dots\}$, $I(P)$ son los números pares, $I(Q)$ los números primos, e $I(R) = \le$. Determina si son verdaderas o falsas en $\mathcal{M}$:
\begin{enumerate}
    \item $\forall x \exists y R(x, y)$
    \item $\exists y \forall x R(y, x)$
    \item $\forall x \forall y (R(x, y) \lor R(y, x))$
    \item $\forall x (P(x) \to \exists y (Q(y) \land R(x, y)))$
    \item $\forall x \forall y ((R(x, y) \land x \neq y) \to \exists z (R(x, z) \land R(z, y) \land z \neq x \land z \neq y))$
\end{enumerate}
\end{pregunta}

\section*{Bloque 3: construcción de contramodelos y comparación semántica}

\begin{pregunta}[Diferenciación semántica de prefijos de cuantificadores]\label{e5-7}
Considera las siguientes cinco sentencias sobre una relación binaria $R$:
\begin{enumerate}
    \item $\exists x \exists y R(x, y)$
    \item $\forall x \forall y R(x, y)$
    \item $\exists x \forall y R(x, y)$
    \item $\forall x \exists y R(x, y)$
    \item $\exists y \forall x R(x, y)$
\end{enumerate}
Proporciona estructuras con dominio $D =
\{1, 2\}$ que muestren que ninguna de ellas es equivalente a otra (por ejemplo,
una estructura donde la fórmula 4 sea verdadera pero la 3 sea
falsa). ¿Hay alguna(s) que sea consecuencia lógica de otra(s)?\\
\textit{Pista.} Presta atención a la diferencia entre $\forall x \exists y R(x, y)$ (donde para cada $x$ puede elegirse un $y$ potencialmente distinto) y $\exists y \forall x R(x, y)$ (donde un mismo $y$ debe servir para todos los $x$).
\end{pregunta}

\begin{pregunta}[Refutación de falacias formales mediante contramodelos]\label{e5-8}
Construye explícitamente estructuras $\mathcal{M} = \langle D, I \rangle$ que hagan verdaderas todas las premisas y falsa la conclusión para refutar:
\begin{enumerate}
    \item $\forall x (P(x) \lor Q(x)) \ \not\models \ \forall x P(x) \lor \forall x Q(x)$
    \item $\forall x (P(x) \to Q(x)), \ \exists x \neg P(x) \ \not\models \ \exists x \neg Q(x)$
    \item $\forall x (P(x) \to Q(x)), \ \exists x Q(x) \ \not\models \ \exists x P(x)$
    \item $\forall x (P(x) \to Q(x)), \ \neg P(a) \ \not\models \ \neg Q(a)$
\end{enumerate}
\end{pregunta}

\section*{Bloque 4: requisitos de cardinalidad del dominio}

\begin{pregunta}[Condiciones de tamaño sobre el universo de discurso]\label{e5-9}
Justifica qué requisitos impone cada una de las siguientes fórmulas a la cardinalidad del dominio $D$ para que una estructura $\mathcal{M} = \langle D, I \rangle$ pueda satisfacerlas:
\begin{enumerate}
    \item $\exists x \exists y (x \neq y)$
    \item $\exists x \exists y (x \neq y \land \forall z (z = x \lor z = y))$
    \item $\forall x \exists y R(x, y) \ \land \ \forall x \forall y (R(x, y) \to \neg R(y, x)) \ \land \ \forall x \forall y \forall z ((R(x, y) \land R(y, z)) \to R(x, z))$
\end{enumerate}
\end{pregunta}

\begin{pregunta}[Estructuras de cardinalidad mínima para conjuntos de fórmulas]\label{e5-10}
Para cada uno de los siguientes tres conjuntos de sentencias de primer orden, propón una estructura $\mathcal{M} = \langle D, I \rangle$ que lo satisfaga procurando que el dominio tenga la menor cardinalidad posible. Justifica en cada caso por qué no es posible hacerlo con un número menor de elementos:
\begin{enumerate}
    \item $\Gamma_1 = \{P(a), \ P(b), \ P(c), \ P(d)\}$.  
    \textit{Pista.} Observa si el lenguaje exige explícitamente que las constantes se refieran a individuos distintos o si pueden coincidir.
    \item $\Gamma_2 = \{\forall x (\neg P(x) \to R(x, x)), \ \forall x (P(x) \leftrightarrow \neg Q(x)), \ P(a), \ \exists x Q(x)\}$.
    \item $\Gamma_3 = \{\exists x \exists y \exists z (x \neq y \land y \neq z \land x \neq z), \ \forall x R(x, x), \ \forall x \forall y (R(x, y) \to R(y, x)), \ \forall x \forall y \forall z ((R(x, y) \land R(y, z)) \to R(x, z))\}$.
\end{enumerate}
\end{pregunta}

\section*{Bloque 5: importación existencial}

\begin{pregunta}[Darapti]\label{e5-11}
Considera el modo silogístico Darapti: «todo $M$ es $P$; todo $M$ es $S$; por tanto, algún $S$ es $P$».
\begin{enumerate}
    \item Construye una estructura que muestre que $\forall x (M(x) \to P(x)), \ \forall x (M(x) \to S(x)) \not\models \exists x (S(x) \land P(x))$.
    \item Demuestra en deducción natural que el argumento es válido si se añade la premisa $\exists x M(x)$.
\end{enumerate}
\end{pregunta}

\begin{pregunta}[El cuadro de oposición]\label{e5-12}
Sean A: $\forall x (S(x) \to P(x))$, E: $\forall x (S(x) \to \neg P(x))$, I: $\exists x (S(x) \land P(x))$ y O: $\exists x (S(x) \land \neg P(x))$.
\begin{enumerate}
    \item Muestra que A y O no pueden ser ambas verdaderas ni ambas falsas en ninguna estructura (son contradictorias). ¿Ocurre lo mismo con E e I?
    \item Construye una estructura en la que A y E sean ambas verdaderas, y otra en la que I y O sean ambas falsas.
    \item ¿Qué premisa habría que añadir para recuperar la subalternación (de A a I)?
\end{enumerate}
\end{pregunta}

%%% Local Variables:
%%% mode: latex
%%% TeX-master: "../ApuntesLogica"
%%% End:


\chapter{El cálculo de tableaux}\label{cap:tableaux}
\begin{resumentema}
  En este tema introducimos el cálculo de árboles semánticos o
  tableaux. El carácter analítico y refutativo de los
  tableaux contrasta con el carácter sintético y generativo de la deducción
  natural. Presentamos la clasificación uniforme de fórmulas de
  Smullyan ($\alpha, \beta, \gamma, \delta$), así como el algoritmo
  para construir árboles, cerrar ramas contradictorias y
  extraer contramodelos explícitos de las ramas abiertas. Finalmente,
  exponemos la idea de las pruebas de corrección y
  completud (conjuntos de Hintikka) y explicamos la semidecidibilidad de la
  lógica de primer orden y su indecidibilidad (Church y Turing).
\end{resumentema}

% -*- coding: utf-8 -*-
\section{Dos estilos de cálculo: sintético frente a analítico}

A lo largo del curso hemos utilizado la deducción natural (temas~\ref{cap:dn-prop} y~\ref{cap:dn-lpo}) como cálculo de demostración. Sin embargo, en la historia de la lógica existen dos familias de cálculos:

\begin{enumerate}
    \item \textbf{Cálculos sintéticos (o generativos).}  
    Tanto los sistemas axiomáticos al estilo de Frege y Hilbert como
    la deducción natural de Gentzen son cálculos sintéticos. Parten de
    axiomas o de premisas iniciales y proceden hacia
      adelante, generando fórmulas paso a paso 
    mediante reglas de inferencia. La deducción natural modela
    el proceso de demostración directa: «tengo estas
    premisas, aplico una secuencia de reglas y llego a la conclusión». 
    
    \item \textbf{Cálculos analíticos (o de descomposición).}  
    Desarrollados por Evert Beth (1955), Jaakko Hintikka (1955) y
    Raymond Smullyan (1968), los \textbf{tableaux semánticos} proceden
    en sentido inverso: parten de las fórmulas complejas que queremos
    examinar y las descomponen en sus partes
      constitutivas. Su estrategia no es «construir una prueba» paso
    a paso, sino buscar sistemáticamente una estructura
      semántica o modelo que haga verdadero un conjunto de fórmulas;
    si la búsqueda fracasa en todas las ramas, ese fracaso es la
    prueba. También suele contarse entre los cálculos analíticos el de resolución, propuesto
    por John Alan Robinson (1965): como los tableaux, trabaja por
    refutación y es la base de muchos demostradores automáticos y del
    lenguaje de programación lógica Prolog.
\end{enumerate}

Mientras que la deducción natural responde a la pregunta sintética \emph{«¿Cómo justifico este enunciado paso a paso?»}, el cálculo de tableaux responde a la pregunta analítica \emph{«¿Es posible concebir una situación o mundo donde estas premisas sean verdaderas y la conclusión falsa?»}.

Esta diferencia tiene consecuencias prácticas. En deducción natural hay que decidir qué hipótesis abrir, qué fórmulas intermedias buscar o cuándo recurrir a la reducción al absurdo, y nada garantiza que una elección sea la buena. En un tableau no hay que inventar nada: las reglas solo descomponen las fórmulas que ya están en la rama, de modo que su aplicación puede mecanizarse por completo. En lógica proposicional el proceso siempre termina y proporciona un procedimiento de decisión; en primer orden, una aplicación sistemática de las reglas encuentra un tableau cerrado siempre que lo haya, aunque puede no terminar cuando no lo hay (sección~\ref{sec:ramas-infinitas}). Por eso los tableaux están también en la base de muchos demostradores automáticos, que usan variables libres y unificación para no tener que probar a ciegas una constante tras otra en la regla $\gamma$ (Fitting, 1996), y de los razonadores para lógicas modales y para las lógicas de descripción con las que se construyen ontologías. No son, sin embargo, el método más eficiente en todos los casos: hay fórmulas proposicionales cuyos tableaux cerrados son mucho más largos que su tabla de verdad (D'Agostino, 1992).

En este tema trabajamos con lenguajes sin el símbolo de identidad ni símbolos de función. Si el tableau encuentra un contramodelo (una rama abierta y completa), el argumento queda refutado; si el tableau descarta todas las posibilidades por contradicción (todas las ramas se cierran), queda demostrado que la conclusión se sigue de las premisas.

\section{Clasificación uniforme de fórmulas (Raymond Smullyan)}

Un tableau es un árbol. En su raíz se escriben las fórmulas de
partida (para comprobar si $\Gamma \models A$, las de $\Gamma$ y
$\neg A$), que luego se descomponen en otras más simples. Si una
fórmula exige que otras dos sean verdaderas, ambas se añaden a la
misma rama; si basta con que lo sea una de ellas, la rama se bifurca en
dos, una para cada posibilidad. Cada rama describe una situación
posible; si contiene una fórmula y su negación, la situación es
imposible y la rama se cierra (sección~\ref{sec:reglas-tableau}).

Para evitar tener que formular decenas de reglas \emph{ad hoc} para cada combinación de conectivas y negaciones, el lógico estadounidense Raymond Smullyan introdujo una clasificación uniforme en cuatro categorías:

\subsection{Fórmulas proposicionales: tipos \texorpdfstring{$\alpha$}{alpha} y \texorpdfstring{$\beta$}{beta}}

\begin{itemize}
    \item \textbf{Fórmulas de tipo $\alpha$ (conjuntivas, no ramifican).}  
    Son fórmulas que son verdaderas si y solo si lo son sus dos componentes. Se descomponen linealmente en la misma rama:
    \begin{center}
    \begin{tabular}{|c|c|c|}
    \hline
    $\alpha$ & $\alpha_1$ & $\alpha_2$ \\
    \hline
    $A \land B$ & $A$ & $B$ \\
    $\neg (A \lor B)$ & $\neg A$ & $\neg B$ \\
    $\neg (A \to B)$ & $A$ & $\neg B$ \\
    \hline
    \end{tabular}
    \end{center}
    
    \item \textbf{Fórmulas de tipo $\beta$ (disyuntivas, ramifican).}  
    Son fórmulas que son verdaderas si y solo si lo es al menos uno de sus componentes. Exigen bifurcar el árbol en dos ramas alternativas:
    \begin{center}
    \begin{tabular}{|c|c|c|}
    \hline
    $\beta$ & $\beta_1$ & $\beta_2$ \\
    \hline
    $A \lor B$ & $A$ & $B$ \\
    $\neg (A \land B)$ & $\neg A$ & $\neg B$ \\
    $A \to B$ & $\neg A$ & $B$ \\
    $A \leftrightarrow B$ & $A \land B$ & $\neg A \land \neg B$ \\
    $\neg (A \leftrightarrow B)$ & $A \land \neg B$ & $\neg A \land B$ \\
    \hline
    \end{tabular}
    \end{center}
\end{itemize}

\subsection{Fórmulas cuantificadas: tipos \texorpdfstring{$\gamma$}{gamma} y \texorpdfstring{$\delta$}{delta}}

\begin{itemize}
    \item \textbf{Fórmulas de tipo $\gamma$ (universales).}  
    Equivalen a una afirmación universal ($\forall x A(x)$ o $\neg \exists x A(x) \equiv \forall x \neg A(x)$). Se pueden instanciar con cualquier constante $c$ que aparezca en la rama:
    \begin{center}
    \begin{tabular}{|c|c|}
    \hline
    $\gamma$ & $\gamma(c)$ \\
    \hline
    $\forall x A(x)$ & $A(c)$ \\
    $\neg \exists x A(x)$ & $\neg A(c)$ \\
    \hline
    \end{tabular}
    \end{center}
    
    \item \textbf{Fórmulas de tipo $\delta$ (existenciales).} Equivalen a una afirmación existencial, como $\exists x A(x)$ o bien $\neg \forall x A(x)$. Exigen introducir una \textbf{constante testigo fresca} $a$ que no haya aparecido previamente en la rama:
    \begin{center}
    \begin{tabular}{|c|c|}
    \hline
    $\delta$ & $\delta(a)$ (con $a$ nueva) \\
    \hline
    $\exists x A(x)$ & $A(a)$ \\
    $\neg \forall x A(x)$ & $\neg A(a)$ \\
    \hline
    \end{tabular}
    \end{center}
\end{itemize}

\subsection{Fórmulas literales y doble negación (\texorpdfstring{$\sigma$}{sigma})}
\begin{itemize}
    \item \textbf{Literales.} Fórmulas atómicas ($P(c_1, \dots, c_n)$)
      o sus negaciones ($\neg P(c_1, \dots, c_n)$). No se
      descomponen más; constituyen los nodos básicos de evaluación.
    \item \textbf{Regla $\sigma$ (doble negación).} De $\neg\neg A$ se añade $A$ a la rama. (Smullyan trata $\neg\neg A$ como una fórmula $\alpha$ con $\alpha_1 = \alpha_2 = A$; aquí la presentamos aparte por claridad).
\end{itemize}

\section{Reglas de construcción y cierre de un tableau}\label{sec:reglas-tableau}

El árbol, que se dibuja con la raíz arriba, se construye de acuerdo con las siguientes pautas:

\begin{enumerate}
    \item \textbf{Inicio.} Se colocan las fórmulas del conjunto de partida como nodos iniciales en una columna vertical (el tronco del árbol).
    \item \textbf{Expansión no ramificada ($\alpha$, $\sigma$ y $\delta$).} Se añade el resultado al final de todas las ramas abiertas que pasen por la fórmula.
    \item \textbf{Expansión ramificada ($\beta$).} Cada rama abierta que pasa por la fórmula se bifurca en dos, con $\beta_1$ en una y $\beta_2$ en la otra.
    \item \textbf{Instanciación universal ($\gamma$).} Se añade $\gamma(c)$ para cada constante $c$ ya existente en la rama; si en la rama no hay todavía ninguna constante, se introduce una nueva. Si posteriormente aparecen nuevas constantes, la fórmula $\gamma$ debe volver a utilizarse con ellas (las fórmulas $\gamma$ nunca se «agotan»).
    \item \textbf{Regla de cierre.} Una rama queda \textbf{cerrada} en cuanto contiene una contradicción explícita: una fórmula $A$ y su negación $\neg A$. Lo más frecuente es que se cierre con un literal y su negación ($P(a)$ y $\neg P(a)$), pero no hace falta descomponer más $A$ si ya aparecen $A$ y $\neg A$ en la rama. Se marca el final de la rama con el símbolo $\bigotimes$. Una rama cerrada deja de expandirse.
\end{enumerate}

\subsection{Estado final de un tableau}
\begin{itemize}
    \item \textbf{Tableau cerrado.} Todas las ramas del árbol terminan en contradicción ($\bigotimes$). No existe ningún camino coherente para hacer verdaderas las fórmulas. Demuestra que el conjunto de partida es insatisfacible.
    \item \textbf{Tableau abierto y completo.} Al menos una rama permanece sin contradicción y se han aplicado todas las reglas posibles sobre sus fórmulas. Demuestra que el conjunto de partida es satisfacible.
\end{itemize}

\subsection{Estrategia heurística: orden de expansión}

Para evitar un crecimiento innecesario del árbol, conviene seguir una heurística habitual de prioridad (no afecta a la corrección del método, solo a su eficiencia, siempre que ninguna fórmula quede sin expandir indefinidamente; Smullyan ya recomendaba aplicar las reglas $\alpha$ antes que las $\beta$):
\begin{center}
\textbf{Prioridad de expansión: } $\alpha$ (y $\sigma$) $\ \longrightarrow \ \delta \ \longrightarrow \ \gamma \ \longrightarrow \ \beta$
\end{center}
\begin{enumerate}
    \item \textbf{Prioridad 1 ($\alpha$ y $\sigma$).} Descomponer primero todas las fórmulas que no ramifican. 
    \item \textbf{Prioridad 2 ($\delta$).} Descomponer los existenciales introduciendo una constante fresca $a$.
    \item \textbf{Prioridad 3 ($\gamma$).} Instanciar los universales con las constantes ya generadas por las reglas $\delta$ o presentes previamente en la rama (y solo con una nueva si no hay ninguna).
    \item \textbf{Prioridad 4 ($\beta$).} Posponer las bifurcaciones disyuntivas hasta que no quede otra opción. Si ramificamos prematuramente, cualquier fórmula pendiente tendrá que descomponerse por duplicado en cada una de las nuevas ramas.
\end{enumerate}

\begin{ejemplo}[Tableau proposicional: silogismo disyuntivo]
Comprobamos la inferencia $(p \lor q) \land \neg p \models q$.
Colocamos la premisa y la negación de la conclusión en la raíz, dado
que estamos buscando un modo de hacer verdadera la premisa y falsa la
conclusión:
\begin{center}
\begin{tikzpicture}[scale=0.9, -, thick, level/.style={sibling distance=4cm,level distance=.9cm}]
  \node[label=left:{1.},label=right:{Premisa}] (r1) {$(p \lor q) \land \neg p$}
  child {
    node[label=left:{2.},label=right:{Negación de la conclusión}] {$\neg q$}
    child {
      node[label=left:{3.},label=right:{$\alpha$, 1}] {$p \lor q$}
      child {
        node[label=left:{4.},label=right:{$\alpha$, 1}] {$\neg p$}
        child {
          node[label=left:{5.},label=right:{$\beta$, 3}] {$p$}
          child { node[label=below:{4, 5}] {$\bigotimes$} }
        }
        child {
          node[label=left:{6.},label=right:{$\beta$, 3}] {$q$}
          child { node[label=below:{2, 6}] {$\bigotimes$} }
        }
      }
    }
  };
\end{tikzpicture}
\end{center}
\textbf{Comentario.} Siguiendo la prioridad heurística, descomponemos
primero el nodo 1 mediante la regla lineal $\alpha$, obteniendo los
nodos 3 y 4. Al disponer ya de $\neg p$ (nodo 4) y de $\neg q$ (nodo
2), en cuanto se aplica la bifurcación $\beta$ sobre el nodo 3, ambas
ramas se cierran de inmediato. Por tanto, no podemos hacer simultáneamente verdadera
la premisa y falsa la conclusión. 
\end{ejemplo}

\practica{La aplicación \apptableaux{} construye el tableau paso a paso con estas reglas y esta prioridad, y explica cada paso. También permite construirlo uno mismo: se elige la fórmula que se descompone y se cierran las ramas señalando dos nodos contradictorios. De cada rama abierta y completa da el modelo (sección siguiente).}

\section{Extracción de contramodelos a partir de ramas abiertas}

Una ventaja importante de los tableaux frente a la deducción natural es que, cuando un argumento es inválido, el cálculo no solo diagnostica la falta de consecuencia, sino que suministra el contramodelo a través de sus ramas abiertas.

\begin{theorem}[Extracción del modelo]
Sea $\mathcal{B}$ una rama completa y abierta de un tableau, y sea $\Omega$ el conjunto de literales que aparecen en dicha rama. La estructura $\mathcal{M} = \langle D, I \rangle$ definida por:
\begin{enumerate}
    \item \textbf{Dominio ($D$).} Conjunto de todas las constantes que aparecen en la rama $\mathcal{B}$:
    \[
    D = \{c \mid c \text{ ocurre en } \mathcal{B}\}
    \]
    (Si no hay ninguna constante en la rama, se toma un dominio de un solo elemento, $D = \{1\}$).
    \item \textbf{Interpretación de constantes.} Para cada constante
      $c$, $I(c) = c$.
    \item \textbf{Interpretación de predicados.} Para cada símbolo de predicado $P$:
    \[
    (c_1, \dots, c_n) \in I(P) \iff P(c_1, \dots, c_n) \in \Omega.
    \]
\end{enumerate}
es un modelo que satisface simultáneamente todas las fórmulas de la rama $\mathcal{B}$ y, por tanto, satisface las fórmulas de la raíz del tableau.
\end{theorem}

\noindent En lógica proposicional, la rama da la interpretación $v$ con $v(p) = 1$ si $p$ está en la rama y $v(p) = 0$ en otro caso.

\section{Aplicaciones del método: validez, consecuencia y satisfacibilidad}

El cálculo de tableaux es un procedimiento analítico de búsqueda de modelos
cuya utilidad depende de qué fórmulas situemos en la raíz del árbol. A
continuación se detallan sus cuatro aplicaciones
principales, indicando qué fórmulas constituyen el punto de partida
y cómo debe interpretarse tanto el cierre del árbol como la presencia
de ramas abiertas. En todas ellas, «tableau abierto» quiere decir que
al menos una rama está abierta y completa.

\subsection{Comprobación de consecuencia lógica (\texorpdfstring{$\Gamma \models A$}{Gamma |= A})}
\begin{itemize}
    \item \textbf{Fórmulas iniciales en la raíz:} las premisas de partida y la negación de la conclusión: $\Gamma \cup \{\neg A\}$.
    \item \textbf{Tableau cerrado (todas las ramas terminan en $\bigotimes$):} es imposible satisfacer a la vez las premisas y la negación de la conclusión. Por tanto, el argumento es válido: $\Gamma \models A$.
    \item \textbf{Tableau abierto (al menos una rama completa no se cierra):} la rama abierta proporciona, de acuerdo con el teorema de extracción, un contramodelo $\mathcal{M}$ en el que todas las premisas son verdaderas ($\mathcal{M} \models \Gamma$) y la conclusión es falsa ($\mathcal{M} \not\models A$). El argumento queda refutado: $\Gamma \not\models A$.
\end{itemize}

\subsection{Comprobación de validez universal (\texorpdfstring{$\models A$}{|= A})}
\begin{itemize}
    \item \textbf{Fórmulas iniciales en la raíz:} la negación de la fórmula: $\{\neg A\}$.
    \item \textbf{Tableau cerrado:} no hay ninguna estructura en la que la fórmula sea falsa. Por consiguiente, la fórmula es universalmente válida: $\models A$.
    \item \textbf{Tableau abierto:} la rama abierta describe un contramodelo $\mathcal{M}$ donde la fórmula resulta falsa ($\mathcal{M} \not\models A$). La fórmula no es válida: $\not\models A$.
\end{itemize}

\subsection{Comprobación de satisfacibilidad (\texorpdfstring{$\Gamma$}{Gamma} satisfacible)}
\begin{itemize}
    \item \textbf{Fórmulas iniciales en la raíz:} el conjunto de fórmulas que se desea evaluar, sin negar nada: $\Gamma = \{A_1, \dots, A_n\}$.
    \item \textbf{Tableau abierto:} el conjunto es satisfacible. En este caso, la estructura $\mathcal{M}$ extraída de la rama abierta no actúa como un contramodelo, sino como un modelo confirmatorio: una situación concreta que hace simultáneamente verdaderas a todas las fórmulas examinadas ($\mathcal{M} \models \Gamma$).
    \item \textbf{Tableau cerrado:} el conjunto es insatisfacible (contradictorio): no existe ninguna estructura donde todas las fórmulas sean simultáneamente verdaderas ($\Gamma \models \bot$).
\end{itemize}

\subsection{Comprobación de equivalencia lógica (\texorpdfstring{$A \equiv B$}{A equiv B})}
\begin{itemize}
    \item \textbf{Fórmulas iniciales en la raíz:} la negación del bicondicional: $\{\neg (A \leftrightarrow B)\}$.
    \item \textbf{Tableau cerrado:} las fórmulas coinciden semánticamente en toda estructura: son lógicamente equivalentes ($A \equiv B$).
    \item \textbf{Tableau abierto:} la rama abierta genera un contramodelo de equivalencia donde una de las fórmulas es verdadera y la otra es falsa.
\end{itemize}

\begin{center}
\renewcommand{\arraystretch}{1.4}
\begin{tabular}{|>{\raggedright\arraybackslash}p{3.6cm}|>{\raggedright\arraybackslash}p{3.2cm}|>{\raggedright\arraybackslash}p{4.2cm}|>{\raggedright\arraybackslash}p{4.2cm}|}
\hline
\textbf{Objetivo lógico} & \textbf{Raíz del árbol} & \textbf{Tableau cerrado ($\bigotimes$)} & \textbf{Tableau abierto} \\
\hline
\textbf{Consecuencia} ($\Gamma \models A$) & $\Gamma \cup \{\neg A\}$ & Argumento válido. & Inválido: se extrae un contramodelo. \\
\hline
\textbf{Validez} ($\models A$) & $\{\neg A\}$ & Fórmula válida (verdad lógica). & No válida: se extrae un contramodelo. \\
\hline
\textbf{Satisfacibilidad} ($\Gamma\not\models\bot$) & $\Gamma$ (sin negar) & Conjunto insatisfacible / contradictorio. & Satisfacible: se extrae un modelo confirmatorio. \\
\hline
\textbf{Equivalencia} ($A \equiv B$) & $\{\neg(A \leftrightarrow B)\}$ & Fórmulas equivalentes. & No equivalentes: contramodelo de equivalencia. \\
\hline
\end{tabular}
\end{center}

\section{Ejemplos completos de tableaux}

\begin{ejemplo}[Demostración de consecuencia lógica mediante tableau cerrado]
Queremos comprobar la validez de la inferencia:
\[
\forall x (P(x) \to Q(x)), \ \exists x P(x) \ \models \ \exists x Q(x)
\]
Para verificar la consecuencia semántica, colocamos las premisas y la negación de la conclusión en la raíz del árbol:
\[
\{\forall x (P(x) \to Q(x)), \ \exists x P(x), \ \neg \exists x Q(x)\}
\]

\begin{center}
\begin{tikzpicture}[scale=0.9, -, thick, level/.style={sibling distance=5cm,level distance=.9cm}]
  \node[label=left:{1.},label=right:{Premisa}] (r1) {$\forall x (P(x) \to Q(x))$}
  child {
    node[label=left:{2.},label=right:{Premisa}] {$\exists x P(x)$}
    child {
      node[label=left:{3.},label=right:{Negación de la conclusión}] {$\neg \exists x Q(x)$}
      child {
        node[label=left:{4.},label=right:{$\delta$, 2 ($a$ constante fresca)}] {$P(a)$}
        child {
          node[label=left:{5.},label=right:{$\gamma$, 1}] {$P(a) \to Q(a)$}
          child {
            node[label=left:{6.},label=right:{$\gamma$, 3}] {$\neg Q(a)$}
            child {
              node[label=left:{7.},label=right:{$\beta$, 5}] {$\neg P(a)$}
              child { node[label=below:{4, 7}] {$\bigotimes$} }
            }
            child {
              node[label=left:{8.},label=right:{$\beta$, 5}] {$Q(a)$}
              child { node[label=below:{6, 8}] {$\bigotimes$} }
            }
          }
        }
      }
    }
  };
\end{tikzpicture}
\end{center}

\textbf{Comentario.} Seguimos el orden de prioridad de expansión:
\begin{enumerate}
    \item En el nodo 4 aplicamos primero la regla existencial $\delta$ sobre el nodo 2, generando la constante fresca $a$.
    \item Solo entonces instanciamos con esa constante $a$ las fórmulas $\gamma$ de los nodos 1 y 3 (nodos 5 y 6).
    \item Finalmente aplicamos la regla disyuntiva $\beta$ sobre el condicional (nodo 5). Ambas ramas chocan inmediatamente con literales previos: la rama izquierda ($\neg P(a)$) choca con el nodo 4 ($P(a)$), y la rama derecha ($Q(a)$) con el nodo 6 ($\neg Q(a)$).
    \item Todas las ramas están cerradas por contradicción ($\bigotimes$). Por consiguiente, el conjunto inicial es insatisfacible y la consecuencia semántica $\forall x (P(x) \to Q(x)), \exists x P(x) \models \exists x Q(x)$ queda demostrada.
\end{enumerate}
\end{ejemplo}

\begin{ejemplo}[Detección de invalidez y extracción sistemática de contramodelos]
Analicemos la inferencia (inválida) que «fusiona» dos existenciales en uno:
\[
\exists x P(x), \ \exists x Q(x) \ \stackrel{?}{\models} \ \exists x (P(x) \land Q(x))
\]
Colocamos las premisas y la negación de la conclusión en la raíz del árbol:

\begin{center}
\begin{tikzpicture}[scale=0.9, -, thick, level/.style={sibling distance=5cm,level distance=.9cm}]
  \node[label=left:{1.},label=right:{Premisa}] (r1) {$\exists x P(x)$}
  child {
    node[label=left:{2.},label=right:{Premisa}] {$\exists x Q(x)$}
    child {
      node[label=left:{3.},label=right:{Negación de la conclusión}] {$\neg \exists x (P(x) \land Q(x))$}
      child {
        node[label=left:{4.},label=right:{$\delta$, 1 ($a$ fresca)}] {$P(a)$}
        child {
          node[label=left:{5.},label=right:{$\delta$, 2 ($b$ fresca, distinta de $a$)}] {$Q(b)$}
          child {
            node[label=left:{6.},label=right:{$\gamma$, 3}] {$\neg(P(a) \land Q(a))$}
            child {
              node[label=left:{7.},label=right:{$\gamma$, 3}] {$\neg(P(b) \land Q(b))$}
              child {
                node[label=left:{8.},label=right:{$\beta$, 6}] {$\neg P(a)$}
                child { node[label=below:{4, 8}] {$\bigotimes$} }
              }
              child {
                node[label=left:{9.},label=right:{$\beta$, 6}] {$\neg Q(a)$}
                child {
                  node[label=left:{10.},label=right:{$\beta$, 7}] {$\neg P(b)$}
                  child { node {\textbf{ABIERTA}} }
                }
                child {
                  node[label=left:{11.},label=right:{$\beta$, 7}] {$\neg Q(b)$}
                  child { node[label=below:{5, 11}] {$\bigotimes$} }
                }
              }
            }
          }
        }
      }
    }
  };
\end{tikzpicture}
\end{center}

\textbf{Comentario y extracción del contramodelo.} 
La rama abierta que desciende por los nodos 1--7, 9 y 10 no es un fallo del método, sino la construcción algorítmica de un contramodelo:
\begin{enumerate}
    \item \textbf{Dominio ($D$):} recogemos las constantes que ocurren en la rama abierta, $D = \{a, b\}$.
    \item \textbf{Literales de la rama:} $\Omega = \{P(a), Q(b), \neg Q(a), \neg P(b)\}$.
    \item \textbf{Extensiones de los predicados:} asignamos a cada predicado los elementos que aparecen afirmados:
    \[
    I(P) = \{a\}, \qquad I(Q) = \{b\}
    \]
    \item \textbf{Evaluación:} en esta estructura, el individuo $a$ cumple $P$ pero no $Q$, y el individuo $b$ cumple $Q$ pero no $P$. Las premisas son verdaderas ($\exists x P(x)$ por el individuo $a$, y $\exists x Q(x)$ por el individuo $b$), pero la conclusión es falsa (ningún individuo cumple simultáneamente $P$ y $Q$). El tableau ha construido un contramodelo formal que refuta la inferencia.
\end{enumerate}
\end{ejemplo}

\section{Los tableaux como cálculo\ampliacion}\label{sec:tableaux-calculo}

El nombre de «tableaux semánticos» tiene una razón histórica. Gentzen
había introducido en 1934 el cálculo de secuentes con una motivación
sintáctica: analizar la estructura de las pruebas. Beth y Hintikka,
que propusieron el método de los tableaux de forma independiente en
1955, partían en cambio de una preocupación semántica. Beth lo
concebía como un intento de construir un contraejemplo,
que se encontrará si y solo si el argumento no es válido. Esa búsqueda
debe llevarse a cabo de forma sistemática, y eso es lo que hace un tableau semántico (Beth, 1955). Su tabla tenía dos columnas, una para las fórmulas que deben ser verdaderas y otra para las que deben ser falsas. Hintikka entendía del mismo modo las pruebas de verdad lógica, como pruebas de que es imposible un contraejemplo, y caracterizó los conjuntos que hoy llevan su nombre (él los llamó \emph{conjuntos modelo}), que describen parcialmente un modelo (Hintikka, 1955). Más tarde, Smullyan (1968) habló de \emph{tableaux analíticos}, para destacar que en ellos solo intervienen subfórmulas de las fórmulas de partida. Fitting (1999) cuenta esta historia con detalle.

Esa motivación explica que a menudo se vean como un procedimiento semántico. Pero son, como la deducción natural, un cálculo de demostración en sentido estricto. Sus reglas son sintácticas: indican qué fórmulas se añaden a una rama según la forma de las que ya contiene ($\alpha$, $\beta$, $\gamma$, $\delta$), y una rama se cierra cuando contiene una fórmula y su negación. Para aplicarlas, y para comprobar que un tableau está bien construido, no hace falta interpretar ningún símbolo, igual que ocurre con una derivación de Fitch. Por eso el cálculo tiene su propia relación de derivabilidad: escribimos $\Gamma \vdash_T A$ cuando hay un tableau cerrado para $\Gamma \cup \{\neg A\}$.

La semántica interviene en dos lugares, pero fuera del cálculo. Primero, en la motivación: cada regla está pensada para reflejar las condiciones de verdad de una conectiva o de un cuantificador, como ocurre con las reglas de introducción y eliminación de la deducción natural. Segundo, en la lectura de los resultados: entendemos que un tableau cerrado muestra que no hay modelo y que una rama abierta y completa describe uno. Que esas lecturas sean adecuadas no es evidente; son los teoremas de corrección y completud del cálculo de tableaux (sección siguiente), que relacionan $\vdash_T$ con $\models$ como los temas~\ref{cap:meta-prop} y~\ref{cap:sem-lpo} relacionaban la derivabilidad en deducción natural con la consecuencia lógica. 

Lo que distingue a los tableaux de la deducción natural no es, por
tanto, que uno sea semántico y el otro sintáctico, sino el modo de
construir las pruebas. La deducción natural es sintética: construye la
conclusión a partir de las premisas. Los tableaux son analíticos:
descomponen las fórmulas de partida y solo añaden fórmulas obtenidas
de ellas (sus componentes, a veces negados, o sus instancias). Un
tableau no construye la conclusión; lo que puede construir, cuando una
rama queda abierta, es un modelo. En esto sí se distinguen los
tableaux de la mayoría de los demás cálculos: cuando en un sistema
axiomático o en deducción natural no se encuentra una prueba, el
intento fallido no proporciona un contramodelo, mientras que una rama
abierta y completa describe uno. Como ha señalado P. de Rouilhan
(1998-1999), los tableaux son un método sintáctico, y su aspecto
semántico es un subproducto cuando fracasa en un intento de prueba. 

\section{Metalógica de tableaux: corrección y completud\ampliacion}

Al igual que la deducción natural, el cálculo de tableaux es un sistema deductivo adecuado:

\begin{theorem}[Corrección del cálculo de tableaux]
Si existe un tableau cerrado para $\Gamma \cup \{\neg A\}$ (es decir, si $\Gamma \vdash_T A$), entonces $\Gamma \models A$.  
(En particular, si hay un tableau cerrado para $\Gamma$, $\Gamma$ es insatisfacible).
\end{theorem}

\begin{proof}[Idea de la demostración]
La prueba descansa en la noción de \textbf{satisfacibilidad de una rama}: decimos que una rama $\mathcal{B}$ es satisfacible si existe al menos una estructura $\mathcal{M}$ en la que todas las fórmulas que pertenecen a $\mathcal{B}$ son verdaderas.
\begin{enumerate}
    \item Se demuestra inductivamente que todas las reglas de expansión preservan la satisfacibilidad:
    \begin{itemize}
        \item Si una rama satisfacible se expande mediante una regla $\alpha$, la única rama resultante sigue siendo satisfacible (pues $\alpha \equiv \alpha_1 \land \alpha_2$).
        \item Si se expande mediante una regla $\beta$, al menos una de las dos ramas hijas resultantes es satisfacible (pues $\beta \equiv \beta_1 \lor \beta_2$).
        \item Si se expande mediante una regla $\delta$, al introducir una constante fresca $a$, siempre podemos extender la interpretación asignando a $a$ el testigo que garantiza la verdad del existencial, manteniendo la rama satisfacible.
        \item La regla $\gamma$ preserva la satisfacibilidad trivialmente, pues si $\forall x A(x)$ es verdadera, cualquier instancia $A(c)$ lo es. (Si $\gamma$ introduce una constante nueva porque la rama no tenía ninguna, se le asigna cualquier elemento del dominio; al ser nueva, no altera el valor de las demás fórmulas).
        \item La regla $\sigma$ también la preserva trivialmente, pues $\neg\neg A \equiv A$.
    \end{itemize}
    \item Ahora bien, si un tableau es cerrado, todas sus ramas contienen una fórmula y su negación, y ninguna rama cerrada es satisfacible.
    \item Por tanto, si el conjunto inicial $\Gamma \cup \{\neg A\}$ fuera satisfacible, el árbol debería contener al menos una rama abierta satisfacible. Al no existir ninguna, se concluye que $\Gamma \cup \{\neg A\}$ es insatisfacible, de donde $\Gamma \models A$. \qedhere
\end{enumerate}
\end{proof}

\begin{theorem}[Completud del cálculo de tableaux]
Sea $\Gamma$ finito. Si $\Gamma \models A$, entonces cualquier tableau
para $\Gamma \cup \{\neg A\}$ construido de forma sistemática (sin que
ninguna fórmula quede sin expandir indefinidamente) se termina
cerrando. En particular, $\Gamma \vdash_T A$.  
\end{theorem}

\begin{proof}[Idea de la demostración (vía conjuntos de Hintikka)]
Razonamos por el contrarrecíproco. Si el árbol sistemático completo para $\Gamma \cup \{\neg A\}$ no se cierra, contiene al menos una rama abierta $\mathcal{B}$, quizá infinita. 
\begin{enumerate}
    \item Por la sistematicidad de la aplicación de reglas, el conjunto de fórmulas de dicha rama abierta $\mathcal{B}$ forma un \textbf{conjunto de Hintikka}:
    \begin{itemize}
        \item No contiene ninguna contradicción atómica ($P \in \mathcal{B} \implies \neg P \notin \mathcal{B}$).
        \item Si $\neg\neg A \in \mathcal{B}$, entonces $A \in \mathcal{B}$.
        \item Si $\alpha \in \mathcal{B}$, entonces $\alpha_1 \in \mathcal{B}$ y $\alpha_2 \in \mathcal{B}$.
        \item Si $\beta \in \mathcal{B}$, entonces $\beta_1 \in \mathcal{B}$ o $\beta_2 \in \mathcal{B}$.
        \item Si $\delta \in \mathcal{B}$, existe una constante $a$ tal que $\delta(a) \in \mathcal{B}$.
        \item Si $\gamma \in \mathcal{B}$, para toda constante $c$ que aparezca en la rama, $\gamma(c) \in \mathcal{B}$ (y en la rama aparece al menos una constante).
    \end{itemize}
    \item Por el \textbf{lema de Hintikka}, todo conjunto de Hintikka es satisfacible: basta construir la estructura canónica sobre las constantes de la rama (exactamente como en el teorema de extracción del modelo).
    \item Dicha estructura satisface todas las fórmulas de la rama $\mathcal{B}$, incluyendo las premisas $\Gamma$ y la negación $\neg A$.
    \item La existencia de esta estructura demuestra que $\Gamma \not\models A$. Por tanto, si $\Gamma \models A$, el tableau sistemático se cierra. \qedhere
\end{enumerate}
\end{proof}

\section{Ramas infinitas y la semidecidibilidad de primer orden}\label{sec:ramas-infinitas}

En lógica proposicional, todo tableau es finito: cada paso descompone
fórmulas en subfórmulas más pequeñas o en sus negaciones (el
bicondicional, a través de una conjunción de ellas) y nunca se repite la aplicación
de una regla sobre una misma fórmula. Por ello, el proceso concluye
siempre en un número finito de pasos (decidibilidad).

En lógica de primer orden la finitud del tableau no está
garantizada. Consideremos como ejemplo el conjunto de premisas:  
\[
\Gamma = \{\forall x \exists y R(x, y), \ \forall x \neg R(x, x)\}
\]
\begin{enumerate}
    \item Empezamos con una constante $c_1$. Por $\forall x \exists y R(x, y)$, obtenemos $\exists y R(c_1, y)$.
    \item Por la regla $\delta$, el existencial exige crear una constante fresca $c_2$, obteniendo $R(c_1, c_2)$.
    \item Pero ahora $c_2$ debe ser procesada por el universal $\forall x \exists y R(x, y)$, produciendo $\exists y R(c_2, y)$.
    \item Al aplicar $\delta$, debemos crear otra constante fresca, $c_3$, obteniendo $R(c_2, c_3)$, y así indefinidamente.
\end{enumerate}

El tableau genera una rama abierta infinita:
\[
c_1 \xrightarrow{R} c_2 \xrightarrow{R} c_3 \xrightarrow{R} c_4 \xrightarrow{R} \dots
\]

\paragraph{Ramas infinitas y modelos finitos: una limitación algorítmica.}
Conviene advertir que la aparición de esta rama infinita no significa
que $\Gamma$ solo pueda satisfacerse en dominios infinitos. De hecho,
el conjunto anterior posee un modelo finito con solo dos elementos:
\[
D = \{1, 2\}, \quad I(R) = \{(1, 2), \ (2, 1)\}
\]
(donde el 1 se relaciona con el 2, el 2 con el 1, y ningún elemento se relaciona consigo mismo). Si no tuviéramos la condición de irreflexividad $\forall x \neg R(x, x)$, la fórmula $\forall x \exists y R(x, y)$ se satisfaría trivialmente en un dominio de un solo elemento con $I(R) = \{(1, 1)\}$.

¿Por qué entonces el tableau no se detiene? La causa está en la regla $\delta$ estándar, que obliga a introducir una constante \emph{fresca} para cada fórmula existencial, sin comprobar si alguna de las constantes ya disponibles en la rama (por ejemplo, $c_1$ como testigo de $\exists y R(c_2, y)$) permitiría satisfacer la fórmula sin agrandar el dominio.

Existen modificaciones del método de árboles diseñadas para paliar este problema (véase Boolos, 1984). En dichas variantes, ante una fórmula existencial en una rama donde ya existen constantes $\{c_1, \dots, c_k\}$, el árbol se bifurca explorando ramas para cada una de las constantes previas y una rama adicional con una constante nueva. Esa ramificación permite al cálculo encontrar modelos finitos (siempre que existan) y detenerse con éxito en ejemplos como el anterior.

\practica{En \apptableaux{} se ven las dos situaciones con este ejemplo: con la regla $\delta$ estándar la rama queda incompleta al llegar al límite de constantes, y con la opción «Regla $\delta$ de Boolos» aparece el modelo finito de dos elementos.}

\paragraph{La barrera de la indecidibilidad (Church-Turing).}
No obstante, ninguna variante algorítmica puede evitar las ramas infinitas en todos los casos: si pudiera, tendríamos un procedimiento de decisión para la lógica de primer orden, lo que contradice el teorema de Church y Turing. Un ejemplo concreto lo dan los conjuntos de fórmulas satisfacibles que solo tienen modelos infinitos: por ejemplo, si al conjunto $\Gamma$ le añadimos la propiedad transitiva $\forall x \forall y \forall z (R(x, y) \land R(y, z) \to R(x, z))$, se define un orden estricto serial, que carece de modelos finitos. En tal caso, ningún tableau puede terminar con una rama abierta y finita, porque de ella se extraería un modelo finito: la rama abierta ha de ser infinita.

Este comportamiento refleja las propiedades metalógicas de la lógica de primer orden:
\begin{itemize}
    \item \textbf{Semidecidibilidad de LPO (consecuencia de la completud).} Si un argumento con un número finito de premisas es válido ($\Gamma \models A$), el tableau construido de modo sistemático se cerrará siempre en un número finito de pasos.
    \item \textbf{Indecidibilidad de LPO (Church-Turing).} Si el argumento no es válido, el algoritmo de tableaux puede permanecer ejecutándose indefinidamente en una rama infinita sin que exista ningún procedimiento mecánico general que decida si el árbol terminará cerrándose o no.
\end{itemize}
Más que un defecto del método, las ramas infinitas hacen visible en el propio árbol un límite de lo computable.

\section*{Referencias bibliográficas}

\begin{enumerate}
    \item \textbf{Beth, E. W.} (1955). «Semantic entailment and formal derivability». \emph{Mededelingen der Koninklijke Nederlandse Akademie van Wetenschappen, Afd. Letterkunde}, nueva serie, vol. 18, n.º 13, pp. 309--342. [Introducción original de los tableaux semánticos].
    \item \textbf{Boolos, G.} (1984). «Trees and finite satisfiability: proof of a conjecture of Burgess». \emph{Notre Dame Journal of Formal Logic}, vol. 25, n.º 3, pp. 193--197. [Variante de la regla existencial que encuentra modelos finitos cuando existen].
    \item \textbf{D'Agostino, M.} (1992). «Are tableaux an improvement on truth-tables? Cut-free proofs and bivalence». \emph{Journal of Logic, Language and Information}, vol. 1, n.º 3, pp. 235--252.
    \item \textbf{Fitting, M.} (1996). \emph{First-Order Logic and Automated Theorem Proving} (2.ª ed.). Springer. [Exposición moderna orientada a demostración automática y conjuntos de Hintikka].
    \item \textbf{Fitting, M.} (1999). «Introduction». En D'Agostino, M., Gabbay, D. M., Hähnle, R., \& Posegga, J. (eds.), \emph{Handbook of Tableau Methods}. Kluwer Academic Publishers. [Historia de los métodos de tableaux, de Gentzen y Beth a la demostración automática].
    \item \textbf{Hintikka, J.} (1955). «Form and content in quantification theory». \emph{Acta Philosophica Fennica}, vol. 8, pp. 7--55. [Conjuntos modelo, hoy llamados conjuntos de Hintikka].
    \item \textbf{Robinson, J. A.} (1965). «A machine-oriented logic based on the resolution principle». \emph{Journal of the ACM}, vol. 12, n.º 1, pp. 23--41. [El método de resolución, base de muchos demostradores automáticos y de la programación lógica].
    \item \textbf{Rouilhan, P. de} (1998-1999). «Les tableaux de Beth: syntaxe ou sémantique?». \emph{Philosophia Scientiae}, vol. 3, n.º 4, pp. 303--322. [Sobre el carácter sintáctico de los tableaux].
    \item \textbf{Smullyan, R. M.} (1968). \emph{First-Order Logic}. Springer-Verlag. [La formulación uniforme de los tableaux].
\end{enumerate}

%%% Local Variables:
%%% mode: latex
%%% TeX-master: "../ApuntesLogica"
%%% End:


\section{Ejercicios de tableaux semánticos}
% -*- coding: utf-8 -*-
\begin{quote}
\small \textbf{Instrucciones.} Para comprobar tus árboles puedes usar la aplicación \apptableaux{}, que sigue las reglas de estos apuntes y da el modelo de cada rama abierta. Construye antes el árbol a mano, o en la aplicación con la opción «Construirlo yo».
\end{quote}

\section*{Bloque 1: tableaux en lógica proposicional}

\begin{pregunta}[Satisfacibilidad de conjuntos de fórmulas]\label{e6-1}
Construye el tableau para cada uno de los siguientes conjuntos de fórmulas. Determina si son satisfacibles o insatisfacibles. En caso de ser satisfacibles, indica una interpretación $v$ que los haga verdaderos a partir de una rama abierta:
\begin{enumerate}
    \item $\Gamma_1 = \{p \to \neg p, \ q \to p, \ p \lor q\}$
    \item $\Gamma_2 = \{p \land \neg q, \ \neg (p \to r), \ q \lor \neg r\}$
    \item $\Gamma_3 = \{(p \lor q) \to r, \ \neg r, \ p\}$
    \item $\Gamma_4 = \{p \leftrightarrow q, \ q \leftrightarrow r, \ p \land \neg r\}$
\end{enumerate}
\end{pregunta}

\begin{pregunta}[Comprobación de validez]\label{e6-2}
Utiliza el método de tableaux para determinar si las siguientes fórmulas proposicionales son lógicamente válidas (tautologías). Recuerda que para comprobar si una fórmula $A$ es válida, se construye el tableau para su negación $\neg A$:
\begin{enumerate}
    \item $((p \lor q) \land (p \lor \neg q)) \to p$
    \item $(p \to (q \to r)) \to ((p \to q) \to (p \to r))$ \quad (Autodistribución del condicional; axioma S, el segundo axioma de Frege)
    \item $(p \to q) \lor (q \to p)$ \quad (Principio de linealidad / Dummett)
    \item $((p \to q) \to p) \to p$ \quad (Ley de Peirce)
    \item $((p \to q) \land \neg p) \to \neg q$ \quad (Falacia de negación del antecedente)
\end{enumerate}
\end{pregunta}

\begin{pregunta}[Consecuencia lógica y detección de falacias proposicionales]\label{e6-3}
Determina mediante tableaux si se verifican las siguientes relaciones de consecuencia lógica. Si la relación no se cumple, extrae de una rama abierta una interpretación que sirva de contraejemplo:
\begin{enumerate}
    \item $\{p \lor q, \ p \to r, \ q \to s\} \models \neg r \to s$ \quad (Dilema constructivo)
    \item $\{p \to (q \lor r), \ \neg r\} \models \neg p$
    \item $\{p \to q, \ \neg p \to r, \ \neg q\} \models r$
    \item $\{p \to q, \ q \to r\} \models p \land r$ \quad (Los condicionales no aseguran sus antecedentes)
\end{enumerate}
\end{pregunta}

\section*{Bloque 2: tableaux en lógica de primer orden (propiedades monádicas)}

\begin{pregunta}[Silogística categórica y cuantificación monádica]\label{e6-4}
Comprueba la validez de los siguientes argumentos clásicos en lógica de predicados monádica mediante el cierre sistemático de su tableau:
\begin{enumerate}
    \item \textbf{Celarent:} $\{\forall x (M(x) \to \neg P(x)), \ \forall x (S(x) \to M(x))\} \models \forall x (S(x) \to \neg P(x))$
    \item \textbf{Ferio:} $\{\forall x (M(x) \to \neg P(x)), \ \exists x (S(x) \land M(x))\} \models \exists x (S(x) \land \neg P(x))$
    \item \textbf{Camestres:} $\{\forall x (P(x) \to M(x)), \ \forall x (S(x) \to \neg M(x))\} \models \forall x (S(x) \to \neg P(x))$
    \item $\{\forall x ((P(x) \land Q(x)) \to \neg R(x)), \ \exists x (P(x) \land Q(x) \land S(x))\} \models \exists x (P(x) \land S(x) \land \neg R(x))$
\end{enumerate}
\end{pregunta}

\begin{pregunta}[Equivalencias entre cuantificadores]\label{e6-5}
Demuestra las siguientes equivalencias mediante tableaux:
\begin{enumerate}
    \item $\forall x (P(x) \land Q(x)) \equiv \forall x P(x) \land \forall x Q(x)$
    \item $\exists x (P(x) \lor Q(x)) \equiv \exists x P(x) \lor \exists x Q(x)$
    \item $\neg \forall x P(x) \equiv \exists x \neg P(x)$
    \item $\forall x (P(x) \to Q(a)) \equiv \exists x P(x) \to Q(a)$ \quad (el consecuente no contiene la variable $x$)
    \item $\forall x (\exists y Q(y) \land P(x)) \equiv \exists y Q(y) \land \forall x P(x)$ 
\end{enumerate}
\end{pregunta}

\section*{Bloque 3: detección de falacias y extracción formal de contramodelos}

\begin{pregunta}[Inversión de cuantificadores y falacias distributivas]\label{e6-6}
Los siguientes argumentos parecen válidos a primera vista en el lenguaje natural, pero son falacias formales. Para cada uno de ellos:
\begin{itemize}
    \item Construye el tableau a partir de las premisas y la negación de la conclusión.
    \item Verifica que al menos una rama permanece abierta y sin contradicciones.
    \item Especifica el contramodelo $\mathcal{M} = \langle D, I \rangle$ extraído de dicha rama abierta (indicando el dominio $D$ y las extensiones $I(P)$ e $I(Q)$).
    \item Explica intuitivamente por qué dicho contramodelo hace verdaderas las premisas y falsa la conclusión.
\end{itemize}

\begin{enumerate}
    \item $\forall x (P(x) \lor Q(x)) \stackrel{?}{\models} \forall x P(x) \lor \forall x Q(x)$  
    \emph{(«Todos los números son pares o impares; por tanto, o todos los números son pares o todos los números son impares»).}
    
    \item $\forall x (P(x) \to Q(x)), \ \exists x \neg P(x) \stackrel{?}{\models} \exists x \neg Q(x)$  
    \emph{(«Todos los filósofos son mortales; hay personas que no son filósofos; por tanto, hay personas que no son mortales»; falacia de negación del antecedente en versión cuantificada).}
    
    \item $\forall x (P(x) \to Q(x)), \ \exists x Q(x) \stackrel{?}{\models} \exists x P(x)$  
    \emph{(Falacia de afirmación del consecuente en versión cuantificada).}
    
    \item $\exists x (P(x) \lor Q(x)) \stackrel{?}{\models} \exists x P(x) \land \exists x Q(x)$
\end{enumerate}
\end{pregunta}

\section*{Bloque 4: relaciones binarias, dependencia y ramas infinitas}

\begin{pregunta}[Permutación de cuantificadores en relaciones binarias]\label{e6-7}
Examina mediante tableaux las dos relaciones de consecuencia siguientes:
\begin{enumerate}
    \item $\exists x \forall y R(x, y) \stackrel{?}{\models} \forall y
      \exists x R(x, y)$ \quad (¿Se cierra el tableau? Si existe,
      extrae el contramodelo).
    \item $\forall x \exists y R(x, y) \stackrel{?}{\models} \exists y
      \forall x R(x, y)$ \quad (¿Se cierra el tableau? \emph{Atención:} con las reglas estándar este tableau no termina, como ocurre en el ejemplo de la sección~\ref{sec:ramas-infinitas}. Describe el patrón que se repite y propón directamente un contramodelo finito, comprobándolo con la semántica del tema~\ref{cap:sem-lpo}).
\end{enumerate}
\end{pregunta}

\begin{pregunta}[Argumento relacional complejo]\label{e6-8}
Construye el tableau para determinar si el siguiente conjunto de premisas entraña la conclusión indicada:
\[
\text{Premisas: } \begin{cases}
1. \ \forall x (\exists y R(x, y) \to \forall z R(z, x)) \\
2. \ \exists x \exists y R(x, y)
\end{cases}
\qquad \text{Conclusión: } \forall x R(x, x)
\]
\emph{Indicación.} Aplica las reglas $\delta$ sobre las premisas iniciales antes de descomponer los condicionales universales.
\end{pregunta}

\section*{Bloque 5: preguntas conceptuales y metodológicas}

\begin{pregunta}[Tableaux frente a deducción natural]\label{e6-9}
Responde razonadamente:
\begin{enumerate}
    \item \textbf{Carácter analítico frente a sintético.} En deducción natural, para demostrar $A \lor B \vdash B \lor A$ aplicamos la regla de eliminación de la disyunción ($\lor\text{E}$), suponiendo $A$ en una subderivación y $B$ en otra, para concluir en ambas $B \lor A$. ¿Cómo demuestra esa misma inferencia un tableau semántico? Explica la diferencia en la dirección del flujo deductivo.
    \item \textbf{Cálculo generativo frente a cálculo refutativo.} Cuando se intenta demostrar una inferencia inválida en deducción natural (por ejemplo, $p \lor q \vdash p \land q$), ¿qué información proporciona el cálculo al no alcanzarse la derivación? ¿Qué información proporciona en cambio el tableau semántico?
    \item \textbf{Conservación de la satisfacibilidad en la corrección\ampliacion.} En el teorema de corrección de los tableaux, ¿por qué es necesario que la regla existencial ($\delta$) utilice una constante fresca, que no haya aparecido previamente en la rama? Muestra un ejemplo de falso razonamiento (una regla $\delta$ defectuosa) que permitiría probar una fórmula que no es válida si reutilizáramos una constante vieja.
    \item \textbf{Conjuntos de Hintikka y completud\ampliacion.} Explica con tus
      propias palabras por qué una rama abierta y completa de un
      tableau no es simplemente una lista de fórmulas, sino la
      descripción de un modelo (un conjunto de Hintikka). 
\end{enumerate}
\end{pregunta}

%%% Local Variables:
%%% mode: latex
%%% TeX-master: "../ApuntesLogica"
%%% End:


%%%%%%%%%%%%%%%%%%%%%%%%%%%%%%%%%%%%%%%%%%%%%%%%%%%%%%%%%%%%%%%%%%%%%%%%%%%%%%%
\appendix
\part{Apéndices}
%%%%%%%%%%%%%%%%%%%%%%%%%%%%%%%%%%%%%%%%%%%%%%%%%%%%%%%%%%%%%%%%%%%%%%%%%%%%%%%

\chapter{Nociones elementales de teoría intuitiva de conjuntos}\label{ap:conjuntos}
\begin{resumentema}
Este apéndice proporciona las herramientas conceptuales de teoría
intuitiva de conjuntos necesarias para fundamentar la semántica y la
metalógica del curso: noción de conjunto y extensionalidad,
pertenencia e inclusión, conjunto potencia, operaciones booleanas
(unión, intersección, diferencia y complemento), pares ordenados y producto cartesiano, relaciones y clases 
de equivalencia, funciones y nociones básicas
de finitud y numerabilidad, ilustrando cada noción con su aplicación
en los temas del programa. 
\end{resumentema}

% -*- coding: utf-8 -*-
\section{La noción intuitiva de conjunto}

La lógica contemporánea, especialmente a partir de los trabajos de
Gottlob Frege, Bertrand Russell y Alfred Tarski, utiliza el lenguaje
de la teoría de conjuntos como su principal herramienta para construir
modelos semánticos y precisar las propiedades de los sistemas
formales. Para comprender la
semántica lógica no se requiere
una fundamentación axiomática 
(como la teoría formal de Zermelo-Fraenkel): basta adoptar el punto de vista de la \textbf{teoría
  intuitiva de conjuntos} (\emph{naive set theory}). 

\begin{definicion}[Conjunto y elemento]
Un \textbf{conjunto} es una colección bien delimitada y definida de
objetos. Los objetos que integran dicha colección se denominan
\textbf{elementos} o \textbf{miembros} del conjunto. Si $a$ es un
elemento del conjunto $B$, lo anotamos como $a\in B$. En otro caso,
$a\notin B$. 
\end{definicion}

La identidad de un conjunto queda determinada por el \textbf{principio
  de extensionalidad}: un conjunto se caracteriza única y
exclusivamente por cuáles son sus elementos, sin que importe el orden
en que se presenten ni el número de veces que se repitan en una
descripción. Así:
\[
\{a, b, c\} = \{c, b, a\} = \{a, a, b, c, c\}
\]
En la práctica, existen dos formas complementarias de especificar un conjunto:
\begin{enumerate}
    \item \textbf{Definición por extensión (o enumeración).} Consiste
      en enumerar explícitamente entre llaves todos y cada uno de los
      elementos. Ejemplos: 
    \begin{itemize}
        \item $\{1, 0\}$ (el conjunto de los
          valores de verdad de la lógica clásica).
        \item $D = \{\text{Sócrates},
          \text{Platón}, \text{Aristóteles}\}$ (un dominio de
          individuos que podemos usar en un modelo). 
    \end{itemize}

    \item \textbf{Definición por comprensión (o especificación).}
      Consiste en caracterizar el conjunto mediante una propiedad o
      condición que deben satisfacer sus elementos. Se utiliza la
      notación con barra vertical $\{x \in D \mid P(x)\}$, que se lee:
      «el conjunto de los elementos $x$ pertenecientes a $D$ tales que
      $x$ tiene la propiedad $P$». Ejemplos: 
    \begin{itemize}
        \item $\mathcal{F} = \{A \mid A \text{ es una fórmula bien formada del lenguaje}\}$.
        \item La extensión de un predicado en un modelo tarskiano es el conjunto de individuos del dominio de los que el predicado es verdadero:
        \[
        I(\text{Filósofo}) = \{x \in D \mid x \text{ es filósofo}\}
        \]
    \end{itemize}
\end{enumerate}

\section{Pertenencia (\texorpdfstring{$\in$}{pertenece}) frente a inclusión (\texorpdfstring{$\subseteq$}{incluido})}

Una de las confusiones conceptuales más recurrentes consiste en asimilar la relación de pertenencia a la relación de inclusión. Se trata, sin embargo, de dos relaciones diferentes:

\begin{center}
\begin{tabular}{|l|c|l|}
\hline
\textbf{Relación} & \textbf{Símbolo} & \textbf{Significado} \\
\hline
\textbf{Pertenencia} & $\in$ & Conecta un elemento con un conjunto: $x \in A$ \\
\textbf{Inclusión} & $\subseteq$ & Conecta dos conjuntos: $A \subseteq B$ \\
\hline
\end{tabular}
\end{center}

\begin{definicion}[Subconjunto]
Decimos que un conjunto $A$ es \textbf{subconjunto} de $B$ (o que $A$ está \textbf{incluido} en $B$), y escribimos $A \subseteq B$, si todo elemento que pertenece a $A$ pertenece también a $B$:
\[
A \subseteq B \iff \text{para todo } x, \text{ si } x \in A \text{ entonces } x \in B
\]
Si además existe al menos un elemento en $B$ que no está en $A$, decimos que $A$ es un subconjunto propio de $B$, escrito $A \subsetneq B$.
\end{definicion}
Por el principio de extensionalidad, si $A \subseteq B$ y $B \subseteq A$,
entonces $A = B$ (ambos conjuntos tienen los mismos elementos). Esta es la
forma habitual de demostrar que dos conjuntos son iguales. 

\subsection{La distinción entre elemento y conjunto unitario}
Un objeto o individuo $x$ no debe identificarse con el conjunto unitario (\emph{singleton}) cuyo único miembro es $x$:
\[
x \neq \{x\}
\]
Cuando un elemento $x$ pertenece a un conjunto $A$ ($x \in A$), esto implica que el conjunto unitario formado por dicho elemento es un subconjunto de $A$:
\[
x \in A \implies \{x\} \subseteq A
\]
Sin embargo, el hecho de que $x \in A$ en general no implica
que $\{x\} \in A$. Podemos apreciar esta distinción con dos ejemplos:
\begin{enumerate}
    \item \textbf{Ejemplo negativo.} En el conjunto $A = \{a, b\}$ se cumple que $a \in \{a, b\}$ y, en consecuencia, $\{a\} \subseteq \{a, b\}$. Por el contrario, la afirmación $\{a\} \in \{a, b\}$ es falsa: los elementos de $A$ son los individuos $a$ y $b$, y ninguno de ellos es el conjunto unitario $\{a\}$. 
    
    \item \textbf{Ejemplo positivo.} Consideremos el conjunto $B = \{a, \{a\}\}$, donde figuran tanto el individuo $a$ como el conjunto unitario $\{a\}$. En este conjunto se verifican simultáneamente tres relaciones:
    \begin{itemize}
        \item $a \in \{a, \{a\}\}$ \quad (el individuo $a$ es elemento de $B$).
        \item $\{a\} \subseteq \{a, \{a\}\}$ \quad (todo elemento perteneciente a $\{a\}$, a saber $a$, pertenece a $B$).
        \item $\{a\} \in \{a, \{a\}\}$ \quad (el conjunto unitario $\{a\}$ figura explícitamente como elemento de $B$).
    \end{itemize}
\end{enumerate}

\subsection{El conjunto vacío (\texorpdfstring{$\varnothing$}{vacio})}
El conjunto que carece por completo de elementos se denomina
\textbf{conjunto vacío} y se denota mediante el símbolo $\varnothing$: 
\[
\text{Para todo } x, \quad x \notin \varnothing
\]
El conjunto vacío está incluido en cualquier conjunto. Es decir, para cualquier conjunto $A$:
\[
\varnothing \subseteq A
\]
\begin{nota}[Justificación lógica por vacuidad]
  La definición de $A \subseteq B$ exige: «para todo $x$, si $x \in A$, entonces
  $x \in B$». Cuando $A = \varnothing$, el antecedente
  «$x \in \varnothing$» es falso para cualquier $x$. En
  lógica proposicional aprendimos que un condicional con
  antecedente falso es verdadero. En consecuencia,
  $\varnothing \subseteq B$ se cumple trivialmente (o por vacuidad) para todo conjunto
  $B$.
\end{nota}

\begin{ejemplo}[Aplicación en metalógica]
En el estudio de la deducción natural y la completud:
\begin{itemize}
    \item Denotamos por $\Gamma$ un conjunto de premisas o hipótesis.
    \item Decir que $A \in \Gamma$ significa que la fórmula específica $A$ es una de las premisas asumidas.
    \item Decir que $\Delta \subseteq \Gamma$ significa que $\Delta$ es un subconjunto de premisas seleccionadas (por ejemplo, el subconjunto finito de premisas efectivamente utilizadas en una deducción formal).
    \item Escribir $\varnothing \vdash A$ significa que la fórmula $A$
      es un \textbf{teorema lógico}: se demuestra formalmente sin
      necesidad de premisa alguna, a partir de un conjunto vacío de
      supuestos iniciales.
    \item Escribir $\Gamma \models \bot$ significa que $\Gamma$ no es
      satisfacible. Como ninguna interpretación hace verdadera $\bot$,
      la condición «toda interpretación que satisface $\Gamma$
      satisface $\bot$» solo puede cumplirse por vacuidad, es decir, porque
      ninguna interpretación satisface $\Gamma$.
\end{itemize}
\end{ejemplo}

\section{Conjunto potencia o de las partes (\texorpdfstring{$\mathcal{P}(A)$}{P(A)})}

\begin{definicion}[Conjunto potencia]
Dado un conjunto $A$, su \textbf{conjunto potencia} (denotado $\mathcal{P}(A)$ o $2^A$) es el conjunto cuyos elementos son todos los subconjuntos de $A$:
\[
\mathcal{P}(A) = \{X \mid X \subseteq A\}
\]
\end{definicion}

Si un conjunto finito $A$ tiene $n$ elementos, su conjunto potencia posee exactamente $2^n$ elementos, ya que para formar un subconjunto arbitrario disponemos de dos opciones independientes para cada elemento: incluirlo o excluirlo.

\begin{ejemplo}
Consideremos el conjunto de individuos $D = \{a, b\}$:
\[
\mathcal{P}(D) = \{\varnothing, \{a\}, \{b\}, \{a, b\}\}
\]
Obsérvese que:
\begin{itemize}
    \item El conjunto vacío $\varnothing$ y el propio $D$ son elementos de $\mathcal{P}(D)$.
    \item $\varnothing \subseteq D$ y a la vez $\varnothing \in \mathcal{P}(D)$.
\end{itemize}
\end{ejemplo}

\begin{aplicacion}[Conexión con la semántica de primer orden (tema~\ref{cap:sem-lpo})]
En una estructura de primer orden $\mathcal{M} = \langle D, I \rangle$, la interpretación de un predicado monádico (propiedad) $P$ consiste en asignar a $P$ el conjunto de individuos que poseen dicha propiedad. Por consiguiente, la extensión $I(P)$ es un subconjunto de $D$, lo que equivale a afirmar que pertenece al conjunto de las partes de $D$:
\[
I(P) \subseteq D \iff I(P) \in \mathcal{P}(D)
\]
Si el universo de discurso $D$ consta de 3 individuos, hay exactamente $2^3 = 8$ extensiones posibles para un predicado monádico (desde la que contiene a los tres individuos hasta la vacía).
\end{aplicacion}

\section{Operaciones booleanas sobre conjuntos y su paralelismo lógico}

A mediados del siglo XIX, George Boole descubrió que ciertas reglas de razonamiento podían formularse algebraicamente estudiando las operaciones entre clases o conjuntos. Las operaciones básicas con conjuntos $A$ y $B$ de un universo de referencia $U$ son:

\begin{itemize}
    \item \textbf{Unión ($A \cup B$).} Colección de elementos que pertenecen a $A$, a $B$ o a ambos:
    \[
    A \cup B = \{x \mid x \in A \lor x \in B\}
    \]
    Corresponde a la disyunción lógica ($\lor$).

    \item \textbf{Intersección ($A \cap B$).} Colección de elementos que pertenecen simultáneamente a $A$ y a $B$:
    \[
    A \cap B = \{x \mid x \in A \land x \in B\}
    \]
    Corresponde a la conjunción lógica ($\land$). Si $A \cap B = \varnothing$, decimos que $A$ y $B$ son \textbf{disjuntos}.

    \item \textbf{Diferencia ($A \setminus B$).} Elementos que están en $A$ pero no están en $B$:
    \[
    A \setminus B = \{x \mid x \in A \land x \notin B\}
    \]

    \item \textbf{Complemento ($A^c$ o $\overline{A}$).} Elementos del universo $U$ que no pertenecen a $A$:
    \[
    A^c = U \setminus A = \{x \in U \mid x \notin A\}
    \]
    Corresponde a la negación lógica ($\neg$).
\end{itemize}

La correspondencia es exacta si se toma como universo $U$ el conjunto
de todas las interpretaciones (asignaciones de verdad), como $A$ el
conjunto de las interpretaciones que hacen verdadera $p$ y como $B$ el
de las que hacen verdadera $q$. Entonces $A \cup B$ es el conjunto de
las interpretaciones que hacen verdadera $p \lor q$; $A \cap B$, el de
las que hacen verdadera $p \land q$; $A^c$, el de las que hacen
verdadera $\neg p$, y así en cada fila de la tabla:

\begin{center}
\begin{tabular}{|l|c|l|c|}
\hline
\textbf{Conjuntos} & \textbf{Símbolo} & \textbf{Lógica proposicional} & \textbf{Fórmula} \\
\hline
Unión & $A \cup B$ & Disyunción & $p \lor q$ \\
Intersección & $A \cap B$ & Conjunción & $p \land q$ \\
Complemento & $A^c$ & Negación & $\neg p$ \\
Unión de $A^c$ y $B$ & $A^c \cup B$ & Condicional & $p \to q$ \\
Unión de $A \cap B$ y $A^c \cap B^c$ & $(A \cap B) \cup (A^c \cap B^c)$ & Bicondicional & $p \leftrightarrow q$ \\
Conjunto vacío & $\varnothing$ & Contradicción / Absurdo & $\bot$ \\
Universo total & $U$ & Tautología & $p \lor \neg p$ \\
\hline
\end{tabular}
\end{center}

Lo mismo vale si en lugar de $p$ y $q$ se toman fórmulas cualesquiera
$F_A$ y $F_B$.
Con esta lectura, las \emph{relaciones} entre conjuntos corresponden a
relaciones entre fórmulas: la inclusión $A \subseteq B$ significa que
toda interpretación que hace verdadera $F_A$ hace verdadera $F_B$, es
decir, $F_A \models F_B$; y la igualdad $A = B$ significa que $F_A \equiv F_B$.

Gracias a este paralelismo, las leyes de De Morgan rigen de manera idéntica en ambos campos:
\[
(A \cup B)^c = A^c \cap B^c \qquad \longleftrightarrow \qquad \neg (p \lor q) \equiv \neg p \land \neg q
\]
\[
(A \cap B)^c = A^c \cup B^c \qquad \longleftrightarrow \qquad \neg (p \land q) \equiv \neg p \lor \neg q
\]

\section{Pares ordenados y producto cartesiano (\texorpdfstring{$A \times B$}{A x B})}

En un conjunto ordinario, los elementos no guardan ninguna relación de
prioridad o sucesión: $\{a, b\} = \{b, a\}$. Sin embargo, razonamos
habitualmente con relaciones en las que importa el orden: «$a$ es más antiguo que $b$», «$a$ ama a $b$» o «$a$ divide a $b$». Para representar estas relaciones necesitamos fijar el orden.

\begin{definicion}[Par ordenado]
Un \textbf{par ordenado} $(a, b)$ es un objeto que consta de un primer componente $a$ y un segundo componente $b$. Dos pares ordenados son idénticos si y solo si coinciden componente a componente:
\[
(a, b) = (c, d) \iff a = c \text{ y } b = d
\]
En particular, si $a \neq b$, entonces $(a, b) \neq (b, a)$.
\end{definicion}

\begin{definicion}[Producto cartesiano]
Dados dos conjuntos $A$ y $B$, su \textbf{producto cartesiano} $A \times B$ es el conjunto de todos los pares ordenados cuyo primer elemento proviene de $A$ y cuyo segundo elemento proviene de $B$:
\[
A \times B = \{(a, b) \mid a \in A \text{ y } b \in B\}
\]
De forma general, para $n \ge 2$, el producto cartesiano de $n$ conjuntos $A_1, \dots, A_n$ (denotado $A_1 \times \dots \times A_n$) y, en el caso particular en que todos los conjuntos coinciden, la potencia cartesiana $A^n$, son los conjuntos de las \textbf{tuplas ordenadas} (o $n$-tuplas) de longitud $n$:
\begin{align*}
  A_1 \times \dots \times A_n &= \{(a_1, \dots, a_n) \mid a_i \in A_i \text{ para cada } i \in \{1, \dots, n\}\}, \\
  A^n &= \underbrace{A \times \dots \times A}_{n \text{ veces}} = \{(a_1, \dots, a_n) \mid a_1, \dots, a_n \in A\}.
\end{align*}
\end{definicion}

\begin{ejemplo}
Sean $A = \{1, 2\}$ y $B = \{x, y\}$. El producto cartesiano consta de $2 \times 2 = 4$ pares:
\[
A \times B = \{(1, x), (1, y), (2, x), (2, y)\}
\]
\[
B \times A = \{(x, 1), (x, 2), (y, 1), (y, 2)\}
\]
Obsérvese que $A \times B \neq B \times A$.
\end{ejemplo}

\section{Relaciones binarias y relaciones de equivalencia}

\begin{definicion}[Relación binaria]
Una \textbf{relación binaria} $R$ entre los elementos de un conjunto $D$ es un subconjunto del producto cartesiano $D \times D$:
\[
R \subseteq D \times D
\]
Escribimos $a R b$ o bien $(a, b) \in R$ para expresar que el individuo $a$ está relacionado con $b$. Si una relación involucra $n$ argumentos, es una relación $n$-aria y se define formalmente como $R \subseteq D^n$.
\end{definicion}

\begin{aplicacion}[Relaciones en lógica de primer orden]
En el tema~\ref{cap:dn-lpo} y el tema~\ref{cap:sem-lpo} formalizamos enunciados relacionales como $R(x, y)$ («$x$ admira a $y$»). En la semántica de Tarski, la interpretación $I(R)$ es un conjunto de pares ordenados del dominio:
\[
\mathcal{M} \models R(a, b) \iff (I(a), I(b)) \in I(R)
\]
\end{aplicacion}

\subsection{Propiedades estructurales de las relaciones binarias}
Una relación $R \subseteq D \times D$ puede satisfacer diversas
propiedades. Entre ellas:

\begin{itemize}
    \item \textbf{Reflexividad.} Todo elemento está relacionado consigo mismo: para todo $x \in D$, $(x, x) \in R$.
    \item \textbf{Irreflexividad.} Ningún elemento está relacionado consigo mismo: para todo $x \in D$, $(x, x) \notin R$.
    \item \textbf{Simetría.} Si un elemento se relaciona con otro, la relación es mutua: si $(x, y) \in R$, entonces $(y, x) \in R$.
    \item \textbf{Asimetría.} Si $(x, y) \in R$, entonces $(y, x) \notin R$.
    \item \textbf{Antisimetría.} Si dos elementos se relacionan mutuamente, deben ser idénticos: si $(x, y) \in R$ y $(y, x) \in R$, entonces $x = y$.
    \item \textbf{Transitividad.} La relación se transfiere a través de pasos intermedios: si $(x, y) \in R$ y $(y, z) \in R$, entonces $(x, z) \in R$.
    \item \textbf{Serialidad.} Todo elemento está relacionado con alguno: para todo $x \in D$ existe $y \in D$ tal que $(x, y) \in R$.
\end{itemize}
Estas propiedades aparecen en varios ejercicios de los temas~\ref{cap:dn-lpo} y~\ref{cap:sem-lpo} (por ejemplo, toda relación asimétrica es irreflexiva). Nótese que los nombres «irreflexiva» y «asimétrica» no significan «no reflexiva» y «no simétrica»: una relación puede no ser reflexiva sin ser irreflexiva.

\subsection{Relaciones de equivalencia y clases de equivalencia}
Cuando una relación reúne simultáneamente las propiedades de ser reflexiva, simétrica y transitiva, se denomina \textbf{relación de equivalencia} (y suele denotarse con el símbolo $\sim$ o $\equiv$).

Las relaciones de equivalencia capturan formalmente la idea de «ser indistinguible respecto a cierto criterio». Su propiedad más importante es que dividen el dominio en bloques disjuntos llamados \textbf{clases de equivalencia}:
\begin{definicion}[Clase de equivalencia]
Sea $\sim$ una relación de equivalencia en $D$. Para cada $a \in D$, la \textbf{clase de equivalencia} de $a$ (escrita $[a]$) es el conjunto de todos los elementos equivalentes a $a$:
\[
[a] = \{x \in D \mid a \sim x\}
\]
El conjunto de todas las clases de equivalencia se denomina \textbf{conjunto cociente} y se denota $D / \sim$.
\end{definicion}
Las clases de equivalencia forman una \textbf{partición} de $D$: cada clase es no vacía (pues $a \in [a]$, por la reflexividad), la unión de todas las clases es $D$, y dos clases diferentes son disjuntas. En concreto, $[a] = [b]$ si $a \sim b$, y $[a] \cap [b] = \varnothing$ si no (la simetría y la transitividad garantizan este hecho).

\begin{aplicacion}[El modelo canónico con igualdad de Henkin]
En la demostración del teorema de completud de primer orden (apéndice~\ref{ap:henkin}), si el lenguaje contiene el predicado de identidad ($=$), los
individuos del modelo canónico no pueden ser simplemente las constantes
del lenguaje, porque dos constantes distintas $c_1$ y $c_2$
(como dos nombres de Venus: «el lucero del alba» y «el lucero vespertino»)
pueden designar el mismo objeto según la teoría ($(c_1 = c_2) \in \Gamma^*$). 

La solución consiste en definir una relación de equivalencia entre constantes ($c_1 \sim c_2 \iff (c_1 = c_2) \in \Gamma^*$), comprobar con las reglas de la identidad que es compatible con los predicados (si $c_1 \sim c_2$, entonces $P(c_1) \in \Gamma^*$ si y solo si $P(c_2) \in \Gamma^*$), y construir el dominio canónico mediante las clases de equivalencia:
\[
D^* = \{[c] \mid c \text{ es una constante del lenguaje ampliado}\}
\]
Cada individuo de $D^*$ es una clase de nombres que la teoría declara idénticos. El apéndice~\ref{ap:henkin} desarrolla la demostración para lenguajes sin identidad y solo esboza este caso.
\end{aplicacion}

\section{Funciones o aplicaciones (\texorpdfstring{$f: X \to Y$}{f: X -> Y})}

Una función es un caso particular de relación: podemos verla como un
dispositivo que asocia a cada entrada exactamente una salida. 

\begin{definicion}[Función]
Una \textbf{función} o \textbf{aplicación} $f$ de un conjunto $X$ (llamado \textbf{dominio}) en un conjunto $Y$ (llamado \textbf{codominio}), denotada $f: X \to Y$, es una relación $f \subseteq X \times Y$ que cumple la condición de \textbf{univalencia y totalidad}:
\[
\text{Para todo } x \in X, \text{ existe un \emph{único} } y \in Y \text{ tal que } (x, y) \in f
\]
A dicho elemento único lo llamamos la \textbf{imagen} de $x$ bajo $f$, y escribimos $f(x) = y$. El conjunto de todos los valores efectivamente alcanzados se denomina \textbf{imagen o recorrido} de la función: $\operatorname{Im}(f) = \{f(x) \mid x \in X\} \subseteq Y$.
\end{definicion}

\subsection{Clasificación de funciones}
\begin{itemize}
    \item \textbf{Inyectiva (uno a uno).} Elementos distintos del dominio producen siempre imágenes distintas en el codominio:
    \[
    x_1 \neq x_2 \implies f(x_1) \neq f(x_2) \quad (\text{o equivalentemente: } f(x_1) = f(x_2) \implies x_1 = x_2)
    \]
    Una función inyectiva no confunde individuos diferentes.

    \item \textbf{Sobreyectiva (exhaustiva).} Todo elemento del codominio es imagen de al menos un elemento del dominio:
    \[
    \text{Para todo } y \in Y, \text{ existe al menos un } x \in X \text{ tal que } f(x) = y \quad (\operatorname{Im}(f) = Y)
    \]
    Una función sobreyectiva no deja ningún elemento del codominio sin cubrir.

    \item \textbf{Biyectiva (correspondencia biunívoca).} La función es simultáneamente inyectiva y sobreyectiva. Permite emparejar de manera reversible cada elemento de $X$ con exactamente un elemento de $Y$.
\end{itemize}

\subsection{El papel de las funciones en la semántica}
Hemos encontrado funciones en tres lugares de la semántica:
\begin{enumerate}
    \item \textbf{Interpretaciones o asignaciones de verdad (tema~\ref{cap:sem-prop}).} Una interpretación $v$ es una función que asigna a cada variable proposicional $p$ un valor booleano:
    \[
    v: \text{Prop} \to \{1, 0\}
    \]
    Esta función se extiende de manera composicional a todas las fórmulas complejas del lenguaje formal.

    \item \textbf{Asignaciones de variables en LPO (tema~\ref{cap:sem-lpo}).} Para definir cuándo una fórmula abierta $P(x)$ es satisfecha por una estructura, Tarski introduce una función de asignación $g$ que asocia a cada variable individual $x$ un objeto del dominio:
    \[
    g: \text{Var} \to D
    \]

    \item \textbf{Interpretación de constantes y símbolos funcionales (tema~\ref{cap:sem-lpo}).} La interpretación semántica asigna a cada símbolo de constante $c$ un individuo concreto: $I(c) \in D$. Y si el lenguaje incorpora símbolos de función $f$ de $n$ argumentos (como la suma, de dos, o el sucesor, de uno), la semántica les asigna una función:
    \[
    I(f): D^n \to D
    \]
\end{enumerate}

\section{Finitud, infinitud y numerabilidad}

En la metalógica (como el teorema de Kalmár en el tema~\ref{cap:meta-prop} o la demostración
de Henkin del tema~\ref{cap:sem-lpo} y el apéndice~\ref{ap:henkin}), la distinción entre lo finito y los distintos
órdenes de infinito desempeña un importante papel teórico. 

\begin{definicion}[Conjuntos finitos e infinitos]
Un conjunto $A$ es \textbf{finito} si consta de un número natural $n \in \mathbb{N}$ de elementos (incluyendo el conjunto vacío, donde $n = 0$). Un conjunto es \textbf{infinito} si no es finito; es decir, contiene una cantidad ilimitada de elementos.
\end{definicion}

\subsection{El infinito numerable (\texorpdfstring{$\aleph_0$}{alef-0})}
Georg Cantor descubrió que no todos los conjuntos infinitos tienen el
mismo «tamaño» o cardinalidad. El nivel más elemental de infinitud
(y el de menor cardinalidad posible) es la numerabilidad:

\begin{definicion}[Conjunto numerable]
Un conjunto infinito $A$ es \textbf{numerable} (o contable) si puede ponerse en biyección con el conjunto de los números naturales $\mathbb{N} = \{0, 1, 2, 3, \dots\}$.

De manera intuitiva: un conjunto infinito es numerable si y solo si es posible disponer todos sus elementos en una lista sin repeticiones ni omisiones:
\[
A = \{a_0, a_1, a_2, a_3, \dots\}
\]
\end{definicion}

Son conjuntos infinitos numerables:
\begin{itemize}
    \item El conjunto de los números naturales $\mathbb{N}$.
    \item El conjunto de los números enteros $\mathbb{Z} = \{0, 1, -1, 2, -2, \dots\}$.
    \item El conjunto de los números racionales (fracciones) $\mathbb{Q}$.
    \item \textbf{El conjunto de las fórmulas de un lenguaje formal.} Toda fórmula es una cadena finita de símbolos. Si el alfabeto es finito, basta ordenar las fórmulas por longitud y, a igual longitud, alfabéticamente. Si el alfabeto es infinito numerable ($s_1, s_2, s_3, \dots$), como ocurre cuando hay infinitas constantes, ese orden no sirve (habría infinitas cadenas de longitud 1), pero podemos proceder por etapas: en la etapa $n$ listamos las fórmulas, en número finito, de longitud como mucho $n$ que solo usan los símbolos $s_1, \dots, s_n$ y que no hayan aparecido antes. Toda fórmula aparece en alguna etapa, de modo que obtenemos una lista:
    \[
    \mathcal{F} = \{A_0, A_1, A_2, A_3, \dots\}
    \]
\end{itemize}

Por el contrario, el conjunto de los números reales $\mathbb{R}$ es \textbf{no numerable} (su cardinalidad es estrictamente mayor que la de los números naturales, como muestra el argumento diagonal de Cantor).

\subsection{La numerabilidad en la metalógica del curso}
La numerabilidad interviene en varios de los resultados teóricos que
hemos estudiado:

\begin{enumerate}
    \item \textbf{Lema de Lindenbaum-Henkin.}  
    Para ampliar un conjunto consistente $\Gamma$ hasta convertirlo en un conjunto maximalmente consistente con testigos, Henkin aprovecha que las sentencias del lenguaje ampliado forman una lista infinita numerable:
    \[
    A_0, A_1, A_2, \dots
    \]
    El lema procede paso a paso por inducción: en el paso 0 toma
    $\Gamma_0 = \Gamma$; en el paso $n+1$ examina la fórmula $A_n$ y
    decide si añadirla o no. Si el lenguaje no fuera numerable, este
    procedimiento paso a paso no podría definirse del mismo modo. 

    \item \textbf{El teorema de Löwenheim-Skolem.}  
    En su versión más sencilla, el teorema afirma que todo conjunto de sentencias de un lenguaje numerable sin identidad que tenga algún modelo tiene también un modelo numerable. ¿Por qué ocurre esto? Porque el modelo canónico de Henkin utiliza como universo de discurso las propias constantes del lenguaje ampliado, y estas, al ser cadenas finitas sobre un alfabeto numerable, forman un conjunto numerable.

    \item \textbf{Ramas infinitas y semidecidibilidad.}  
    En el cálculo de tableaux de primer orden, un árbol con fórmulas existenciales y universales puede no cerrarse nunca y desarrollar una rama infinita numerable. Si el árbol se construye de modo sistemático, cada rama infinita describe un conjunto de Hintikka, que proporciona un contramodelo con dominio numerable. Por eso la validez en primer orden solo es semidecidible: si un argumento con un número finito de premisas es válido, el tableau sistemático se cierra en un número finito de pasos; si no lo es, puede no terminar (sección~\ref{sec:ramas-infinitas}).
\end{enumerate}

\section*{Referencias bibliográficas}

\begin{enumerate}
    \item \textbf{Enderton, H. B.} (1977). \emph{Elements of Set Theory}. Academic Press, Nueva York. [Introducción a la teoría axiomática de Zermelo-Fraenkel].
    \item \textbf{Halmos, P. R.} (1960). \emph{Naive Set Theory}. Van Nostrand, Princeton. Traducción española: \emph{Teoría intuitiva de los conjuntos}, CECSA, México, 1965. [Introducción breve y clásica a la teoría intuitiva de conjuntos].
    \item \textbf{Hrbacek, K., \& Jech, T.} (1999). \emph{Introduction to Set Theory} (3.ª ed.). Marcel Dekker, Nueva York. [Teoría axiomática, cardinales y ordinales].
    \item \textbf{Ivorra, C.} \emph{Lógica y teoría de conjuntos}. Universitat de València. \url{https://www.uv.es/ivorra/Libros/Logica.pdf}. [Apuntes gratuitos en castellano; versión elemental, que el autor ya no actualiza].
    \item \textbf{Mosterín, J.} (1980). \emph{Teoría axiomática de conjuntos} (2.ª ed.; 1.ª ed., 1971). Ariel, Barcelona. [Exposición en castellano de la axiomática de Zermelo-Fraenkel].
\end{enumerate}

%%% Local Variables:
%%% mode: latex
%%% TeX-master: "../ApuntesLogica"
%%% End:


\chapter{La demostración de Henkin (1949) del teorema de completud}\label{ap:henkin}
\begin{resumentema}
  Este apéndice presenta la demostración del teorema de
  completud de primer orden publicada por Leon Henkin en 1949: testigos de Henkin, lema de Lindenbaum-Henkin, modelo canónico
  de términos, lema de verdad y los corolarios de compacidad y
  Löwenheim-Skolem.
\end{resumentema}

% -*- coding: utf-8 -*-
\section{Introducción: el teorema de completud y el teorema del modelo}

El principal objetivo metateórico de la lógica clásica de primer orden es demostrar la coincidencia entre la consecuencia semántica ($\models$) y la derivabilidad sintáctica en deducción natural ($\vdash$):
\[
\Gamma \vdash A \iff \Gamma \models A
\]
El sentido directo ($\Gamma \vdash A \implies \Gamma \models A$) es el
\textbf{teorema de corrección}. Al igual que en la lógica proposicional
(donde se demostró por inducción sobre la longitud de las deducciones en el
tema~\ref{cap:meta-prop}), la corrección en primer orden garantiza que el cálculo deductivo
no lleva de premisas verdaderas a conclusiones falsas. Se demuestra
comprobando que cada línea de una derivación es consecuencia lógica de las
premisas e hipótesis activas en ella. Para las reglas $\forall\text{I}$ y
$\exists\text{E}$ esto no significa que cada paso preserve la verdad en
una estructura fija; lo que se preserva es la validez, gracias a las restricciones sobre parámetros arbitrarios y
constantes testigo analizadas en el tema~\ref{cap:dn-lpo}.

El sentido inverso ($\Gamma \models A \implies \Gamma \vdash A$) es el
\textbf{teorema de completud}, demostrado por primera vez por Kurt
Gödel en 1930 y demostrado de nuevo por Leon Henkin en 1949 mediante un método
que destaca por su generalidad, dado que es posible adaptarlo a
diversos sistemas formales. 

\paragraph{Supuestos.} Para simplificar, en este apéndice suponemos que el lenguaje $\mathcal{L}$ es numerable, que dispone de infinitas constantes individuales (como $a, b, \dots$ con subíndices) y que no contiene ni el símbolo de identidad ni símbolos de función (el tema~\ref{cap:dn-lpo} sí incluye la identidad). Además, $\Gamma$ y $A$ denotan siempre sentencias (fórmulas sin variables libres), igual que en las derivaciones del tema~\ref{cap:dn-lpo}. En la sección~\ref{sec_henkin_identidad} indicamos qué cambia si se añade la identidad.

El punto de partida de Henkin es que la completud fuerte equivale a un enunciado sobre la existencia de modelos para conjuntos consistentes (sección~\ref{sec:consistencia-satisfacibilidad}, donde se definen la consistencia y la satisfacibilidad, y lema~\ref{lema_reduccion}):
\begin{quote}
\textbf{Teorema del modelo:} si $\Gamma$ es consistente ($\Gamma \nvdash \bot$), entonces existe una estructura $\mathcal{M}$ tal que $\mathcal{M} \models \Gamma$.
\end{quote}

La estrategia de Henkin se reduce, por tanto, a un objetivo: dada una teoría consistente arbitraria $\Gamma$, construir una estructura semántica $\mathcal{M}$ que la haga verdadera.

\section{Paso 1: el enriquecimiento del lenguaje con testigos frescos}

Sea $\mathcal{L}$ un lenguaje de primer orden numerable. Para poder dotar de «nombres» a todos los individuos cuya existencia afirmen las fórmulas existenciales de la teoría, Henkin amplía el lenguaje añadiendo un conjunto infinito numerable de nuevas constantes individuales, llamadas \textbf{testigos de Henkin}:
\[
C = \{c_1, c_2, c_3, \dots\}
\]
Definimos el lenguaje enriquecido como $\mathcal{L}^* = \mathcal{L} \cup C$. Puesto que tanto los símbolos originales como las nuevas constantes son numerables, el conjunto de todas las sentencias de $\mathcal{L}^*$ es numerable (véase el apéndice~\ref{ap:conjuntos}) y puede ser ordenado en una secuencia infinita:
\[
A_0, A_1, A_2, A_3, \dots
\]

\section{Paso 2: el lema de Lindenbaum-Henkin}

\begin{definition}[Conjunto de Henkin]
Un conjunto de sentencias $\Gamma^*$ de $\mathcal{L}^*$ es un \textbf{conjunto de Henkin} si satisface:
\begin{enumerate}
    \item \textbf{Consistencia.} $\Gamma^* \nvdash \bot$.
    \item \textbf{Maximalidad.} Para toda sentencia $A$ de $\mathcal{L}^*$, $A \in \Gamma^*$ o bien $\neg A \in \Gamma^*$.
    \item \textbf{Propiedad de los testigos.} Para toda fórmula existencial $\exists x B(x)$, si $\exists x B(x) \in \Gamma^*$, entonces existe una constante $c \in C$ tal que $B(c) \in \Gamma^*$.
\end{enumerate}
\end{definition}

\begin{lemma}[Lema de Lindenbaum-Henkin]\label{lema_lindenbaum_henkin}
Todo conjunto consistente de sentencias $\Gamma$ de $\mathcal{L}$ puede extenderse a un conjunto de Henkin $\Gamma^*$ en el lenguaje enriquecido $\mathcal{L}^*$ tal que $\Gamma \subseteq \Gamma^*$.
\end{lemma}

\begin{proof}
Construimos inductivamente una cadena creciente de conjuntos consistentes $\Gamma_0 \subseteq \Gamma_1 \subseteq \Gamma_2 \subseteq \dots$ como sigue:
\begin{enumerate}
    \item \textbf{Base.} Definimos $\Gamma_0 = \Gamma$. Puesto que $\Gamma$ es consistente en $\mathcal{L}$, sigue siendo consistente en $\mathcal{L}^*$. En efecto, si hubiera una derivación de $\bot$ a partir de $\Gamma$ en $\mathcal{L}^*$, como las constantes nuevas no aparecen en $\Gamma$, podríamos sustituir cada una de las (finitas) constantes nuevas que aparecen en ella por una constante de $\mathcal{L}$ que no aparezca en la derivación (una distinta para cada constante nueva); el resultado sigue siendo una derivación correcta (las restricciones de $\forall$I y $\exists$E se conservan) y sería una derivación de $\bot$ en $\mathcal{L}$.
    \item \textbf{Paso inductivo.} Dado el conjunto $\Gamma_n$ consistente, consideramos la fórmula $A_n$ de la enumeración:
    \begin{itemize}
        \item \textbf{Caso 1.} Si $\Gamma_n \cup \{A_n\}$ es \textbf{inconsistente}, definimos:
        \[
        \Gamma_{n+1} = \Gamma_n.
        \]
        \item \textbf{Caso 2.} Si $\Gamma_n \cup \{A_n\}$ es \textbf{consistente}:
        \begin{itemize}
            \item Si $A_n$ no es una fórmula existencial, definimos:
            \[
            \Gamma_{n+1} = \Gamma_n \cup \{A_n\}.
            \]
            \item Si $A_n$ es una fórmula existencial de la forma $\exists x B(x)$, elegimos la primera constante $c \in C$ que no ocurra en $\Gamma_n$ ni en $A_n$ (lo cual siempre es posible porque $C$ es infinito, $\Gamma$ no contiene constantes de $C$ y hasta la etapa $n$ solo se han añadido a $\Gamma$ un número finito de sentencias, que contienen un número finito de constantes de $C$), y definimos:
            \[
            \Gamma_{n+1} = \Gamma_n \cup \{\exists x B(x), B(c)\}.
            \]
        \end{itemize}
    \end{itemize}
\end{enumerate}

\paragraph{Verificación de la consistencia en el caso existencial.}  
Debemos asegurarnos de que al añadir $B(c)$ para una constante fresca $c$, el conjunto $\Gamma_{n+1}$ sigue siendo consistente. Si no lo fuera, tendríamos $\Gamma_n \cup \{\exists x B(x), B(c)\} \vdash \bot$. Por deducción natural proposicional, esto equivale a:
\[
\Gamma_n \cup \{\exists x B(x)\} \vdash \neg B(c).
\]
Pero la constante $c$ no aparece en $\Gamma_n$ ni en $\exists x B(x)$ (ni, por tanto, en $\forall x \neg B(x)$). Aplicando la regla de introducción del universal ($\forall\text{I}$), obtenemos:
\[
\Gamma_n \cup \{\exists x B(x)\} \vdash \forall x \neg B(x).
\]
Por la regla de conversión de cuantificadores (CQ), $\forall x \neg B(x)$ equivale a $\neg \exists x B(x)$. Por tanto:
\[
\Gamma_n \cup \{\exists x B(x)\} \vdash \neg \exists x B(x),
\]
lo que implicaría que $\Gamma_n \cup \{\exists x B(x)\} \vdash \bot$, contradiciendo que estábamos en el Caso 2, donde $\Gamma_n \cup \{A_n\}$ era consistente. Por consiguiente, $\Gamma_{n+1}$ es consistente.

\paragraph{Definición de $\Gamma^*$.} Definimos el conjunto límite como la unión infinita:
\[
\Gamma^* = \bigcup_{n=0}^{\infty} \Gamma_n.
\]
Comprobamos que $\Gamma^*$ cumple las propiedades requeridas:
\begin{enumerate}
    \item \textbf{$\Gamma^*$ es consistente.} Toda derivación en deducción natural es un objeto finito que solo utiliza un número finito de premisas. Si $\Gamma^* \vdash \bot$, esas finitas premisas pertenecerían todas a un mismo $\Gamma_k$ (la cadena es creciente), lo que contradice que cada $\Gamma_k$ es consistente.
    \item \textbf{$\Gamma^*$ es maximal.} Consideremos cualquier sentencia $A$ de $\mathcal{L}^*$. Dado que la lista $A_0, A_1, \dots$ contiene todas las sentencias del lenguaje, tanto $A$ como su negación $\neg A$ aparecen en dicha enumeración (digamos $A = A_n$ y $\neg A = A_k$).
    
    Si $A_n \notin \Gamma^*$, es porque en el paso $n$ el conjunto $\Gamma_n \cup \{A_n\}$ era inconsistente ($\Gamma_n \cup \{A_n\} \vdash \bot$), lo que implica que $\Gamma_n \vdash \neg A_n$ y, por ende, $\Gamma^* \vdash \neg A_n$. Cuando el proceso inductivo alcanza el paso $k$ correspondiente a la fórmula $\neg A_n$, el conjunto $\Gamma_k \cup \{\neg A_n\}$ tiene que ser consistente; en efecto, si $\Gamma_k \cup \{\neg A_n\} \vdash \bot$, se tendría $\Gamma_k \vdash A_n$, y en consecuencia $\Gamma^*$ deduciría tanto $A_n$ como $\neg A_n$ ($\Gamma^* \vdash \bot$), lo cual contradiría que $\Gamma^*$ es consistente (punto 1). Como además $\neg A_n$ no es existencial, en el paso $k$ se aplicó el Caso 2 y se incorporó $\neg A_n \in \Gamma_{k+1} \subseteq \Gamma^*$. Por tanto, para toda sentencia $A$, o bien $A \in \Gamma^*$ o bien $\neg A \in \Gamma^*$.
    \item \textbf{$\Gamma^*$ tiene la propiedad de los testigos.} Supongamos que $\exists x B(x) \in \Gamma^*$ y sea $A_n = \exists x B(x)$ su lugar en la enumeración. Como $\Gamma_n \cup \{A_n\} \subseteq \Gamma^*$ y $\Gamma^*$ es consistente, $\Gamma_n \cup \{A_n\}$ es consistente; luego en el paso $n$ se aplicó el Caso 2 y se añadió $B(c) \in \Gamma_{n+1} \subseteq \Gamma^*$.
\end{enumerate}
El lema queda demostrado.
\end{proof}

\begin{observacion}[Clausura deductiva y propiedades booleanas de $\Gamma^*$]\label{obs_clausura_deductiva}
Como consecuencia directa de su consistencia y maximalidad, el conjunto de Henkin $\Gamma^*$ satisface las siguientes propiedades:
\begin{enumerate}
    \item \textbf{Clausura deductiva.} Para cualquier sentencia $A$,
    \[
    \Gamma^* \vdash A \iff A \in \Gamma^*.
    \]
    \emph{Demostración.} La implicación $(\Leftarrow)$ es inmediata (si $A \in \Gamma^*$, se deduce tomando $A$ como premisa). Para $(\Rightarrow)$, si $\Gamma^* \vdash A$ y fuera $A \notin \Gamma^*$, por la maximalidad de $\Gamma^*$ tendríamos $\neg A \in \Gamma^*$; esto implicaría que $\Gamma^* \vdash A$ y $\Gamma^* \vdash \neg A$, de modo que $\Gamma^* \vdash \bot$, contradiciendo que $\Gamma^*$ es consistente. Por consiguiente, $A \in \Gamma^*$.
    \item \textbf{Comportamiento homomórfico respecto a conectivas.}
    \begin{itemize}
        \item $\neg A \in \Gamma^* \iff A \notin \Gamma^*$.
        \item $B \land C \in \Gamma^* \iff B \in \Gamma^* \text{ y } C \in \Gamma^*$.
        \item $B \lor C \in \Gamma^* \iff B \in \Gamma^* \text{ o } C \in \Gamma^*$.
        \item $B \to C \in \Gamma^* \iff (B \in \Gamma^* \implies C \in \Gamma^*)$.
        \item $B \leftrightarrow C \in \Gamma^* \iff (B \in \Gamma^* \iff C \in \Gamma^*)$.
        \item $\bot \notin \Gamma^*$ (por la consistencia).
    \end{itemize}
    Cada equivalencia se comprueba con las reglas de introducción y eliminación de la conectiva correspondiente y de la negación, la clausura deductiva, la consistencia y la maximalidad. Por ejemplo, si $B \lor C \in \Gamma^*$ pero $B \notin \Gamma^*$ y $C \notin \Gamma^*$, entonces $\neg B, \neg C \in \Gamma^*$, y con $\lor$E se deriva $\bot$ de $\Gamma^*$, lo que es imposible.
    \item \textbf{Universales.} Si $\forall x B(x) \in \Gamma^*$, entonces $B(c) \in \Gamma^*$ para toda constante $c$ de $\mathcal{L}^*$ (por $\forall$E y la clausura deductiva).
\end{enumerate}
\end{observacion}

\section{Paso 3: construcción del modelo canónico de términos}

Disponemos ahora de un conjunto de Henkin $\Gamma^*$. Henkin construye
a partir de él una estructura $\mathcal{M}_{\Gamma^*} = \langle D, I
\rangle$ del siguiente modo:

\begin{definition}[Modelo canónico]
Sea $\Gamma^*$ un conjunto de Henkin. Definimos la estructura $\mathcal{M}_{\Gamma^*} = \langle D, I \rangle$ como:
\begin{enumerate}
    \item \textbf{Dominio de discurso ($D$).} El conjunto de todas las constantes individuales del lenguaje $\mathcal{L}^*$, es decir, las de $\mathcal{L}$ junto con los testigos de $C$. Es un conjunto infinito numerable. Como el lenguaje no tiene símbolos de función, las constantes son precisamente los términos cerrados; por eso se habla de \emph{modelo de términos}.
    \item \textbf{Interpretación ($I$).} Cada constante se interpreta a sí misma: $I(c) = c$. Para cualquier símbolo de predicado $n$-ario $P$ y constantes $c_1, \dots, c_n \in D$, la interpretación de $P$ es:
    \[
    (c_1, \dots, c_n) \in I(P) \iff P(c_1, \dots, c_n) \in \Gamma^*.
    \]
\end{enumerate}
\end{definition}

\section{Paso 4: el lema de verdad (\emph{Truth Lemma})}

El corazón de la demostración es probar que, en el modelo canónico, la verdad semántica coincide con la pertenencia sintáctica a la teoría:

\begin{lemma}[Lema de verdad]\label{lema_verdad}
Para cualquier fórmula cerrada $A$ del lenguaje $\mathcal{L}^*$:
\[
\mathcal{M}_{\Gamma^*} \models A \iff A \in \Gamma^*.
\]
\end{lemma}

\begin{proof}
Procedemos por inducción sobre la complejidad sintáctica de la fórmula $A$ (número de conectivas y cuantificadores):

\paragraph{Caso base ($A$ es una fórmula atómica $P(c_1, \dots, c_n)$).}  
Por definición tarskiana de satisfacción en $\mathcal{M}_{\Gamma^*}$:
\[
\mathcal{M}_{\Gamma^*} \models P(c_1, \dots, c_n) \iff (I(c_1), \dots, I(c_n)) \in I(P) \iff (c_1, \dots, c_n) \in I(P).
\]
Por la definición de $I(P)$ en el modelo canónico, esto equivale a que $P(c_1, \dots, c_n) \in \Gamma^*$. Si $A$ es $\bot$, el lema también se cumple: $\mathcal{M}_{\Gamma^*} \not\models \bot$ y, por la consistencia, $\bot \notin \Gamma^*$.

\paragraph{Paso inductivo.}  
Asumimos como hipótesis de inducción (HI) que el lema se cumple para todas las sentencias de menor complejidad que $A$. No basta con las subfórmulas de $A$: en los casos de los cuantificadores aplicaremos la HI a instancias $B(c)$, que no son subfórmulas de $\exists x B(x)$ pero tienen menor complejidad.

Usaremos también un hecho sencillo: como $I(c) = c$, satisfacer $B(x)$ con una asignación que da a $x$ el valor $c$ equivale a que sea verdadera la sentencia $B(c)$; es decir, $\mathcal{M}_{\Gamma^*}, g[x \mapsto c] \models B(x) \iff \mathcal{M}_{\Gamma^*} \models B(c)$.

\begin{itemize}
    \item \textbf{Caso de la negación ($A = \neg B$):}
    \begin{align*}
    \mathcal{M}_{\Gamma^*} \models \neg B &\iff \mathcal{M}_{\Gamma^*}
                                            \not\models B &
                                                            \text{(Semántica
                                                            de $\neg$)} \\
    &\iff B \notin \Gamma^* & \text{(Por hipótesis de inducción aplicada a $B$)} \\
    &\iff \neg B \in \Gamma^* & \text{(Por consistencia y maximalidad de $\Gamma^*$)}
    \end{align*}

    \item \textbf{Caso de la conjunción ($A = B \land C$):}
    \begin{align*}
    \mathcal{M}_{\Gamma^*} \models B \land C &\iff \mathcal{M}_{\Gamma^*} \models B \text{ y } \mathcal{M}_{\Gamma^*} \models C & \text{(Semántica de $\land$)} \\
    &\iff B \in \Gamma^* \text{ y } C \in \Gamma^* & \text{(Por HI)} \\
    &\iff B \land C \in \Gamma^* & \text{(Por reglas $\land$I y $\land$E en $\Gamma^*$)}
    \end{align*}

    \item \textbf{Caso del condicional ($A = B \to C$):}
    \begin{align*}
    \mathcal{M}_{\Gamma^*} \models B \to C &\iff (\mathcal{M}_{\Gamma^*} \models B \implies \mathcal{M}_{\Gamma^*} \models C) & \text{(Semántica de $\to$)} \\
    &\iff (B \in \Gamma^* \implies C \in \Gamma^*) & \text{(Por HI)} \\
    &\iff B \to C \in \Gamma^* & \text{(Observación~\ref{obs_clausura_deductiva})}
    \end{align*}
    Los casos de la disyunción y el bicondicional se tratan igual, usando las equivalencias correspondientes de la observación~\ref{obs_clausura_deductiva}.

    \item \textbf{Caso del existencial ($A = \exists x B(x)$):}
    \begin{itemize}
        \item [$(\Rightarrow)$] Supongamos que $\mathcal{M}_{\Gamma^*} \models \exists x B(x)$. Por la semántica de Tarski, existe algún elemento del dominio (una constante $c$) tal que $\mathcal{M}_{\Gamma^*}, g[x \mapsto c] \models B(x)$, es decir, $\mathcal{M}_{\Gamma^*} \models B(c)$. Por hipótesis de inducción aplicada a la fórmula más simple $B(c)$, tenemos que $B(c) \in \Gamma^*$. Aplicando la regla de introducción del existencial ($\exists\text{I}$) en deducción natural, de $B(c)$ se deduce formalmente $\exists x B(x)$ ($\Gamma^* \vdash \exists x B(x)$). Dado que $\Gamma^*$ es cerrado bajo deducción (observación~\ref{obs_clausura_deductiva}), concluimos que $\exists x B(x) \in \Gamma^*$.
        \item [$(\Leftarrow)$] Supongamos que $\exists x B(x) \in \Gamma^*$. Por la propiedad de los testigos que caracteriza al conjunto de Henkin $\Gamma^*$, existe una constante $c \in C$ tal que $B(c) \in \Gamma^*$. Por hipótesis de inducción aplicada a $B(c)$, tenemos que $\mathcal{M}_{\Gamma^*} \models B(c)$. Dado que $c$ es un elemento del dominio $D$, por la semántica de Tarski para el existencial se concluye inmediatamente que $\mathcal{M}_{\Gamma^*} \models \exists x B(x)$.
    \end{itemize}

    \item \textbf{Caso del universal ($A = \forall x B(x)$):} 
    \begin{itemize}
        \item [$(\Leftarrow)$] Si $\forall x B(x) \in \Gamma^*$, entonces $B(c) \in \Gamma^*$ para toda constante $c$ (observación~\ref{obs_clausura_deductiva}). Por HI, $\mathcal{M}_{\Gamma^*} \models B(c)$ para toda constante $c$. Como todo elemento del dominio es una constante, $\mathcal{M}_{\Gamma^*} \models \forall x B(x)$.
        \item [$(\Rightarrow)$] Si $\forall x B(x) \notin \Gamma^*$, por maximalidad $\neg \forall x B(x) \in \Gamma^*$, y por CQ y la clausura deductiva $\exists x \neg B(x) \in \Gamma^*$. Por la propiedad de los testigos, $\neg B(c) \in \Gamma^*$ para alguna constante $c$, luego $B(c) \notin \Gamma^*$ (por consistencia). Por HI, $\mathcal{M}_{\Gamma^*} \not\models B(c)$, y por tanto $\mathcal{M}_{\Gamma^*} \not\models \forall x B(x)$.
    \end{itemize}
    Obsérvese que el caso universal, igual que el existencial, necesita la propiedad de los testigos.
\end{itemize}
El lema de verdad queda así demostrado para cualquier fórmula cerrada de $\mathcal{L}^*$.
\end{proof}

\begin{observacion}[Fórmulas cerradas y el método canónico de Henkin]\label{obs_cerradas_henkin}
Conviene resaltar por qué el lema de verdad se formula específicamente
para fórmulas cerradas (enunciados o sentencias sin variables
libres). La construcción canónica de Henkin introduce constantes
testigo ($c \in C$) para las sentencias existenciales ($\exists
x B(x)$), y el conjunto maximalmente consistente $\Gamma^*$
solo contiene sentencias del lenguaje ampliado
$\mathcal{L}^*$. Ninguna fórmula abierta pertenece a $\Gamma^*$, aunque algunas, como $P(x) \lor \neg P(x)$, son verdaderas en $\mathcal{M}_{\Gamma^*}$; por eso el lema solo puede valer para sentencias. En el modelo canónico, el universo está formado por
las propias constantes y la noción primaria es la verdad de enunciados
cerrados. Como veremos a continuación en el teorema de completud, este resultado se formula para sentencias, que constituyen el objeto propio de las teorías lógicas y matemáticas.
\end{observacion}

\section[Corolarios de la completud: compacidad y
Löwenheim-Skolem]{Conclusión de la completud y corolarios: compacidad
  y\\ Löwenheim-Skolem}

A partir del lema de verdad, la demostración del teorema de completud
fuerte para sentencias es inmediata:

\begin{theorem}[Teorema de completud fuerte para sentencias]
El cálculo de deducción natural de primer orden es fuertemente completo para enunciados cerrados: para cualquier conjunto de sentencias (fórmulas cerradas) $\Gamma$ y cualquier sentencia $A$,
\[
\Gamma \models A \implies \Gamma \vdash A.
\]
\end{theorem}

\begin{proof}
Procedemos por reducción al absurdo. Supongamos que $\Gamma \models A$ y $\Gamma \nvdash A$:
\begin{enumerate}
    \item El conjunto $\Gamma \cup \{\neg A\}$ es sintácticamente consistente. En efecto, si tuviéramos $\Gamma \cup \{\neg A\} \vdash \bot$, bastaría tomar $\neg A$ como hipótesis de una caja que llegue a $\bot$ y aplicar IP para deducir $\Gamma \vdash A$, en contra de lo supuesto.
    \item Por el lema de Lindenbaum-Henkin (\ref{lema_lindenbaum_henkin}), $\Gamma \cup \{\neg A\}$ puede extenderse a un conjunto de Henkin $\Gamma^*$ en el lenguaje ampliado $\mathcal{L}^*$, con $\Gamma \cup \{\neg A\} \subseteq \Gamma^*$.
    \item Por el lema de verdad (\ref{lema_verdad}), el modelo
      canónico de términos $\mathcal{M}_{\Gamma^*}$ satisface todas
      las sentencias de $\Gamma^*$ (y lo mismo hace la
      $\mathcal{L}$-estructura que no interpreta las constantes nuevas, pues $\Gamma$ y $A$ no las contienen):
    \[
    \mathcal{M}_{\Gamma^*} \models \Gamma \quad \text{y} \quad \mathcal{M}_{\Gamma^*} \models \neg A.
    \]
    \item La semántica de la negación establece que $\mathcal{M}_{\Gamma^*} \models \neg A \iff \mathcal{M}_{\Gamma^*} \not\models A$.
    \item Por consiguiente, $\mathcal{M}_{\Gamma^*}$ es un modelo que
      hace verdaderas todas las sentencias de $\Gamma$ y falsa la
      sentencia $A$, de modo que $\Gamma
      \not\models A$. 
\end{enumerate}
Esto contradice la suposición $\Gamma \models A$; por tanto, $\Gamma \models A \implies \Gamma \vdash A$.
\end{proof}

\begin{observacion}[Fórmulas abiertas: anomalías semánticas y deductivas]\label{obs_formulas_abiertas}
Veamos por qué el teorema de completud se enuncia para
sentencias y qué problemas surgen al emplear
fórmulas abiertas (con variables libres) como premisas o conclusión. 

En la práctica, las teorías axiomáticas (como los axiomas de Peano
para la aritmética o ZFC para la teoría de conjuntos) y los teoremas
que se pretenden demostrar suelen ser sentencias
(enunciados sin variables libres) o se entienden como sus clausuras universales. Las fórmulas abiertas son
predicados incompletos que solo adquieren un valor de verdad definido
cuando se fija una asignación para sus variables libres, o cuando estas quedan ligadas por cuantificadores o sustituidas
por constantes. 

Si consideramos fórmulas abiertas como premisas, aparecen 
anomalías tanto en la semántica como en el cálculo: 
\begin{itemize}
    \item \textbf{Anomalía semántica (fallo del teorema de la
        deducción).} Tal como se definió en el tema~\ref{cap:sem-lpo}, una fórmula
      abierta es verdadera en una estructura ($\mathcal{M} \models B(x)$)
      si resulta satisfecha bajo \emph{toda} asignación de variables
      (se evalúa igual que $\mathcal{M} \models \forall x B(x)$). Bajo
      esta convención, la inferencia $P(x) \models P(y)$ es
      formalmente válida: si todo individuo del dominio cumple $P$,
      cualquier otro individuo también lo cumple. Sin embargo,
      falla el teorema de la deducción semántico: $P(x) \to P(y)$
      no es universalmente válida. En efecto,
      basta una estructura con $D = \{1, 2\}$, $I(P) = \{1\}$ y una
      asignación con $g(x) = 1$ y $g(y) = 2$ para que
      $\mathcal{M}, g \not\models P(x) \to P(y)$ (pues $\text{V} \to \text{F} = \text{F}$). 
    \item \textbf{Anomalía sintáctica.}
      En el cálculo del tema~\ref{cap:dn-lpo} las líneas de una derivación son siempre
      sentencias (las reglas de los cuantificadores operan con
      constantes), de modo que «$P(x) \vdash P(y)$» ni siquiera está
      definido. En los cálculos de deducción natural que sí admiten fórmulas abiertas, la
      regla $\forall\text{I}$ prohíbe generalizar sobre una variable
      libre en las premisas, y por eso $P(x) \nvdash P(y)$: la
      correspondencia con $\models$ se rompería. 
\end{itemize}
Si una premisa ha de tener fuerza universal,
debe explicitarse en el lenguaje formal mediante su cuantificación:
$\forall x P(x)$. Y a partir de ahí, se verifican tanto la
derivabilidad ($\forall x P(x) \vdash \forall y P(y)$) como el teorema
de la deducción semántico ($\models \forall x P(x) \to \forall y P(y)$). 

La formulación del teorema de completud para sentencias asegura así
que la correspondencia entre $\models$ y $\vdash$ se establece sobre
sentencias, en las que se conservan las propiedades metateóricas
habituales. 
\end{observacion}

\subsection{El teorema de compacidad como corolario inmediato}

Una de las consecuencias del teorema de completud es que el
\textbf{teorema de compacidad} no requiere una demostración
independiente; surge como corolario de la finitud sintáctica: 

\begin{corollary}[Teorema de compacidad de primer orden]
Un conjunto de sentencias $\Gamma$ (de un lenguaje numerable) es satisfacible si y solo si todo subconjunto finito $\Gamma_0 \subseteq \Gamma$ es satisfacible.
\end{corollary}

\begin{proof}\mbox{}
\begin{itemize}
    \item $(\Rightarrow)$ Inmediato: si existe una estructura $\mathcal{M}$ tal que $\mathcal{M} \models \Gamma$, entonces directamente $\mathcal{M} \models \Gamma_0$ para cualquier subconjunto $\Gamma_0 \subseteq \Gamma$.
    \item $(\Leftarrow)$ Supongamos que todo subconjunto finito
      $\Gamma_0 \subseteq \Gamma$ es satisfacible. Veamos, en primer
      lugar, que $\Gamma$ es sintácticamente consistente ($\Gamma
      \nvdash \bot$). En efecto, si $\Gamma$ fuera inconsistente,
      existiría una deducción formal $\Gamma \vdash \bot$. Pero toda
      derivación en deducción natural consta de un número finito de
      pasos y, por tanto, solo puede utilizar un conjunto finito de
      premisas $\Gamma_0 \subseteq \Gamma$. Tendríamos entonces
      $\Gamma_0 \vdash \bot$. Por el teorema de corrección, $\Gamma_0
      \models \bot$, lo que significaría que $\Gamma_0$ no tiene
      ningún modelo, contradiciendo nuestra hipótesis.   
    Por consiguiente, $\Gamma$ debe ser consistente ($\Gamma \nvdash \bot$). Por el teorema del modelo, que la construcción de Henkin demuestra (lema de Lindenbaum-Henkin y lema de verdad), $\Gamma$ es satisfacible. \qedhere
\end{itemize}
\end{proof}

\subsection{El teorema de Löwenheim-Skolem}

\begin{corollary}[Teorema de Löwenheim-Skolem]
Si un conjunto de sentencias $\Gamma$ de un lenguaje numerable (sin identidad) tiene un modelo, entonces tiene un modelo numerable.
\end{corollary}

\begin{proof}
Si $\Gamma$ tiene un modelo, por el teorema de corrección $\Gamma$ es sintácticamente consistente. Aplicando la construcción de Henkin, el modelo canónico $\mathcal{M}_{\Gamma^*}$ tiene como dominio el conjunto de las constantes de $\mathcal{L}^*$, que es infinito numerable. Por tanto, la estructura construida por Henkin es un modelo numerable de $\Gamma$.
\end{proof}

\subsection{Lenguajes con identidad}\label{sec_henkin_identidad}
Si el lenguaje contiene el símbolo $=$, el cálculo incluye las reglas
$=$I y $=$E del tema~\ref{cap:dn-lpo}, y el modelo canónico no puede
tomar como dominio las propias constantes: si $(c_1 = c_2) \in
\Gamma^*$, ambas constantes deben denotar el mismo objeto. Se define
entonces $c_1 \sim c_2 \iff (c_1 = c_2) \in \Gamma^*$, se comprueba
con las reglas de la identidad que $\sim$ es una relación de
equivalencia compatible con los predicados, y se toma como dominio el
conjunto cociente de las constantes por $\sim$ (véase el
apéndice~\ref{ap:conjuntos}). El resto de la demostración es igual, con $I(c) = [c]$ y un caso más en el lema de verdad: el de las sentencias atómicas $c_1 = c_2$. En
ese caso el dominio puede ser finito, y para obtener la versión de
Löwenheim-Skolem con un modelo \emph{infinito} numerable (a partir de
un modelo infinito) se necesita un paso adicional: añadir infinitas constantes nuevas $d_1, d_2, \dots$ con las sentencias $d_i \neq d_j$ ($i \neq j$), que el modelo infinito satisface interpretándolas como elementos distintos.

\section*{Referencias bibliográficas}

\begin{enumerate}
    \item \textbf{Henkin, L.} (1949). «The completeness of the first-order functional calculus». \emph{Journal of Symbolic Logic}, vol. 14, n.º 3, pp. 159--166.
    \item \textbf{Mendelson, E.} (1997). \emph{Introduction to Mathematical Logic} (4.ª ed.). Chapman \& Hall. [§2.7, «Completeness Theorems», pp. 84--94: lema de Lindenbaum (prop. 2.14), modelo numerable de toda teoría consistente por el método de Henkin (prop. 2.17) y teorema de completud de Gödel (cor. 2.19)].
    \item \textbf{Shoenfield, J. R.} (1967). \emph{Mathematical Logic}. Addison-Wesley. [Capítulo 4, «The Characterization Problem»; §4.2, «The completeness theorem», pp. 43--48: el teorema en sus dos formas, la demostración de Henkin y el teorema de Lindenbaum (p. 47)].
    \item \textbf{Tarski, A.} (1930). «Fundamentale Begriffe der Methodologie der deduktiven Wissenschaften. I». \emph{Monatshefte für Mathematik und Physik}, vol. 37, pp. 361--404. [§7, teorema I.56 (p. 394), que Tarski atribuye a Lindenbaum: todo conjunto consistente de enunciados puede extenderse a uno consistente y completo].
\end{enumerate}

%%% Local Variables:
%%% mode: latex
%%% TeX-master: "../ApuntesLogica"
%%% End:


\chapter{Resumen de reglas}\label{ap:reglas}
% -*- coding: utf-8 -*-
Este apéndice reúne en pocas páginas las reglas del curso, para consultarlas al resolver ejercicios. Las explicaciones y los ejemplos están en los temas~\ref{cap:dn-prop},~\ref{cap:dn-lpo} y~\ref{cap:tableaux}. En los esquemas, $m, n, i, j, \dots$ son números de línea; los puntos verticales indican pasos intermedios.

% Una regla: título y esquema, alineados por arriba con la regla vecina.
\newenvironment{regla}[1]{%
  \begin{minipage}[t]{0.48\textwidth}\centering\textbf{#1}\par\medskip
}{%
  \end{minipage}%
}
\newcommand{\finfila}{\par\bigskip\bigskip}

\section{Reglas básicas de la deducción natural proposicional}

\noindent\hfill
\begin{regla}{Reiteración (R)}
$\begin{nd}[rules=myrules]
    \have [m] {1} {A}
    \have [\mbox{}] {2} {A} \r{1}
\end{nd}$
\end{regla}\hfill\mbox{}
\finfila

\noindent
\begin{regla}{Introducción de la conjunción ($\land$I)}
$\begin{nd}[rules=myrules]
    \have [m] {1} {A}
    \have [n] {2} {B}
    \have [\mbox{}] {3} {A \land B} \ai{1, 2}
\end{nd}$
\end{regla}\hfill
\begin{regla}{Eliminación de la conjunción ($\land$E)}
$\begin{nd}[rules=myrules]
    \have [m] {1} {A \land B}
    \have [\mbox{}] {2} {A} \ae{1}
\end{nd}$
\quad
$\begin{nd}[rules=myrules]
    \have [m] {1} {A \land B}
    \have [\mbox{}] {2} {B} \ae{1}
\end{nd}$
\end{regla}
\finfila

\noindent
\begin{regla}{Introducción del condicional ($\to$I)}
$\begin{nd}[rules=myrules]
    \open
    \hypo [i] {1} {A}
    \have [\mbox{}] {} {\vdots}
    \have [j] {2} {B}
    \close
    \have [\mbox{}] {3} {A \to B} \ci{1-2}
\end{nd}$
\end{regla}\hfill
\begin{regla}{Eliminación del condicional ($\to$E)}
$\begin{nd}[rules=myrules]
    \have [m] {1} {A \to B}
    \have [n] {2} {A}
    \have [\mbox{}] {3} {B} \ce{1, 2}
\end{nd}$
\end{regla}
\finfila

\noindent
\begin{regla}{Introducción de la negación ($\neg$I)}
$\begin{nd}[rules=myrules]
    \open
    \hypo [i] {1} {A}
    \have [\mbox{}] {} {\vdots}
    \have [j] {2} {\bot}
    \close
    \have [\mbox{}] {3} {\lnot A} \ni{1-2}
\end{nd}$
\end{regla}\hfill
\begin{regla}{Eliminación de la negación ($\neg$E)}
$\begin{nd}[rules=myrules]
    \have [m] {1} {\lnot A}
    \have [n] {2} {A}
    \have [\mbox{}] {3} {\bot} \ne{1, 2}
\end{nd}$
\end{regla}
\finfila

\noindent
\begin{regla}{Explosión (X)}
$\begin{nd}[rules=myrules]
    \have [m] {1} {\bot}
    \have [\mbox{}] {2} {B} \x{1}
\end{nd}$
\end{regla}\hfill
\begin{regla}{Prueba indirecta (IP)}
$\begin{nd}[rules=myrules]
    \open
    \hypo [i] {1} {\lnot A}
    \have [\mbox{}] {} {\vdots}
    \have [j] {2} {\bot}
    \close
    \have [\mbox{}] {3} {A} \ip{1-2}
\end{nd}$
\end{regla}
\finfila

\noindent
\begin{regla}{Introducción de la disyunción ($\lor$I)}
$\begin{nd}[rules=myrules]
    \have [m] {1} {A}
    \have [\mbox{}] {2} {A \lor B} \oi{1}
\end{nd}$
\quad
$\begin{nd}[rules=myrules]
    \have [m] {1} {B}
    \have [\mbox{}] {2} {A \lor B} \oi{1}
\end{nd}$
\end{regla}\hfill
\begin{regla}{Eliminación de la disyunción ($\lor$E)}
$\begin{nd}[rules=myrules]
    \have [m] {1} {A \lor B}
    \open
    \hypo [i] {2} {A}
    \have [\mbox{}] {} {\vdots}
    \have [j] {3} {C}
    \close
    \open
    \hypo [k] {4} {B}
    \have [\mbox{}] {} {\vdots}
    \have [l] {5} {C}
    \close
    \have [\mbox{}] {6} {C} \oe{1, 2-3, 4-5}
\end{nd}$
\end{regla}
\finfila

\noindent
\begin{regla}{Introducción del bicondicional ($\leftrightarrow$I)}
$\begin{nd}[rules=myrules]
    \open
    \hypo [i] {1} {A}
    \have [\mbox{}] {} {\vdots}
    \have [j] {2} {B}
    \close
    \open
    \hypo [k] {3} {B}
    \have [\mbox{}] {} {\vdots}
    \have [l] {4} {A}
    \close
    \have [\mbox{}] {5} {A \leftrightarrow B} \bi{1-2, 3-4}
\end{nd}$
\end{regla}\hfill
\begin{regla}{Eliminación del bicondicional ($\leftrightarrow$E)}
$\begin{nd}[rules=myrules]
    \have [m] {1} {A \leftrightarrow B}
    \have [n] {2} {A}
    \have [\mbox{}] {3} {B} \be{1, 2}
\end{nd}$
\quad
$\begin{nd}[rules=myrules]
    \have [m] {1} {A \leftrightarrow B}
    \have [n] {2} {B}
    \have [\mbox{}] {3} {A} \be{1, 2}
\end{nd}$
\end{regla}
\finfila

\section{Reglas básicas de los cuantificadores y de la identidad}

\noindent
\begin{regla}{Eliminación del universal ($\forall$E)}
$\begin{nd}[rules=myrules]
    \have [m] {1} {\forall x A(x)}
    \have [\mbox{}] {2} {A(c)} \Ae{1}
\end{nd}$

\medskip\small\raggedright
$c$ puede ser cualquier constante; se sustituyen todas las apariciones libres de $x$.
\end{regla}\hfill
\begin{regla}{Introducción del universal ($\forall$I)}
$\begin{nd}[rules=myrules]
    \have [m] {1} {A(c)}
    \have [\mbox{}] {2} {\forall x A(x)} \Ai{1}
\end{nd}$

\medskip\small\raggedright
$c$ no aparece en ninguna premisa ni hipótesis activa, ni en $\forall x A(x)$ (se sustituyen \emph{todas} las apariciones de $c$). La variable $x$ no aparece en $A(c)$.
\end{regla}
\finfila

\noindent
\begin{regla}{Introducción del existencial ($\exists$I)}
$\begin{nd}[rules=myrules]
    \have [m] {1} {A(c)}
    \have [\mbox{}] {2} {\exists x A(x)} \Ei{1}
\end{nd}$

\medskip\small\raggedright
Sin restricciones sobre $c$; pueden sustituirse solo algunas de sus apariciones. La variable $x$ no aparece en $A(c)$.
\end{regla}\hfill
\begin{regla}{Eliminación del existencial ($\exists$E)}
$\begin{nd}[rules=myrules]
    \have [m] {1} {\exists x A(x)}
    \open
    \hypo [i] {2} {A(c)} \by{Hip. testigo $c$}{}
    \have [\mbox{}] {} {\vdots}
    \have [j] {3} {B}
    \close
    \have [\mbox{}] {4} {B} \Ee{1, 2-3}
\end{nd}$

\medskip\small\raggedright
$c$ no aparece en ninguna premisa ni hipótesis activa anterior a la caja, ni en $\exists x A(x)$, ni en $B$.
\end{regla}
\finfila

\noindent
\begin{regla}{Introducción de la identidad ($=$I)}
$\begin{nd}[rules=myrules]
    \have [m] {1} {c = c} \idi{}
\end{nd}$

\medskip\small\raggedright
Para cualquier constante $c$; no se cita ninguna línea.
\end{regla}\hfill
\begin{regla}{Eliminación de la identidad ($=$E)}
$\begin{nd}[rules=myrules]
    \have [m] {1} {a = b}
    \have [n] {2} {A(\dots a \dots a \dots)}
    \have [\mbox{}] {3} {A(\dots b \dots a \dots)} \ide{1, 2}
\end{nd}$

\medskip\small\raggedright
Se sustituyen una, varias o todas las apariciones de $a$ por $b$. También en sentido inverso: de $a = b$ y una fórmula con $b$, sustituir $b$ por $a$.
\end{regla}
\finfila

\section{Reglas derivadas}

\noindent
\begin{regla}{\emph{Modus tollens} (MT)}
$\begin{nd}[rules=myrules]
    \have [m] {1} {A \to B}
    \have [n] {2} {\lnot B}
    \have [\mbox{}] {3} {\lnot A} \mt{1, 2}
\end{nd}$
\end{regla}\hfill
\begin{regla}{Silogismo disyuntivo (DS)}
$\begin{nd}[rules=myrules]
    \have [m] {1} {A \lor B}
    \have [n] {2} {\lnot A}
    \have [\mbox{}] {3} {B} \ds{1, 2}
\end{nd}$ \qquad
$\begin{nd}[rules=myrules]
    \have [m] {1} {A \lor B}
    \have [n] {2} {\lnot B}
    \have [\mbox{}] {3} {A} \ds{1, 2}
\end{nd}$
\end{regla}
\finfila

\noindent
\begin{regla}{Eliminación de la doble negación (DNE)}
$\begin{nd}[rules=myrules]
    \have [m] {1} {\lnot\lnot A}
    \have [\mbox{}] {2} {A} \dne{1}
\end{nd}$
\end{regla}\hfill
\begin{regla}{Tercero excluido (LEM)}
$\begin{nd}[rules=myrules]
    \open
    \hypo [i] {1} {A}
    \have [\mbox{}] {} {\vdots}
    \have [j] {2} {B}
    \close
    \open
    \hypo [k] {3} {\lnot A}
    \have [\mbox{}] {} {\vdots}
    \have [l] {4} {B}
    \close
    \have [\mbox{}] {5} {B} \lem{1-2, 3-4}
\end{nd}$
\end{regla}
\finfila

\noindent\hfill
\begin{minipage}[t]{0.8\textwidth}\centering
\textbf{Leyes de De Morgan (DeM)}\par\medskip
\begin{tabular}{cc}
$\begin{nd}[rules=myrules]
    \have [m] {1} {\lnot(A \lor B)}
    \have [\mbox{}] {2} {\lnot A \land \lnot B} \dem{1}
\end{nd}$ &
$\begin{nd}[rules=myrules]
    \have [m] {1} {\lnot(A \land B)}
    \have [\mbox{}] {2} {\lnot A \lor \lnot B} \dem{1}
\end{nd}$ \\[0.6cm]
$\begin{nd}[rules=myrules]
    \have [m] {1} {\lnot A \land \lnot B}
    \have [\mbox{}] {2} {\lnot(A \lor B)} \dem{1}
\end{nd}$ &
$\begin{nd}[rules=myrules]
    \have [m] {1} {\lnot A \lor \lnot B}
    \have [\mbox{}] {2} {\lnot(A \land B)} \dem{1}
\end{nd}$
\end{tabular}
\end{minipage}\hfill\mbox{}
\finfila

\noindent\hfill
\begin{minipage}[t]{0.8\textwidth}\centering
\textbf{Conversión de cuantificadores (CQ)}\par\medskip
De una línea con la fórmula de la izquierda se obtiene la de la derecha, citando esa línea:\par\medskip
\renewcommand{\arraystretch}{1.3}
\begin{tabular}{ccc|ccc}
$\neg \forall x A(x)$ & $\Rightarrow$ & $\exists x \neg A(x)$ & $\exists x \neg A(x)$ & $\Rightarrow$ & $\neg \forall x A(x)$ \\
$\neg \exists x A(x)$ & $\Rightarrow$ & $\forall x \neg A(x)$ & $\forall x \neg A(x)$ & $\Rightarrow$ & $\neg \exists x A(x)$
\end{tabular}
\end{minipage}\hfill\mbox{}
\finfila

\clearpage
\section{Tableaux}

\begin{center}
\renewcommand{\arraystretch}{1.2}
\begin{tabular}[t]{|c|c|c|}
\hline
$\alpha$ & $\alpha_1$ & $\alpha_2$ \\
\hline
$A \land B$ & $A$ & $B$ \\
$\neg (A \lor B)$ & $\neg A$ & $\neg B$ \\
$\neg (A \to B)$ & $A$ & $\neg B$ \\
\hline
\end{tabular}
\qquad
\begin{tabular}[t]{|c|c|c|}
\hline
$\beta$ & $\beta_1$ & $\beta_2$ \\
\hline
$A \lor B$ & $A$ & $B$ \\
$\neg (A \land B)$ & $\neg A$ & $\neg B$ \\
$A \to B$ & $\neg A$ & $B$ \\
$A \leftrightarrow B$ & $A \land B$ & $\neg A \land \neg B$ \\
$\neg (A \leftrightarrow B)$ & $A \land \neg B$ & $\neg A \land B$ \\
\hline
\end{tabular}
\\[0.4cm]
\begin{tabular}{|c|c|}
\hline
$\gamma$ & $\gamma(c)$ \\
\hline
$\forall x A(x)$ & $A(c)$ \\
$\neg \exists x A(x)$ & $\neg A(c)$ \\
\hline
\end{tabular}
\qquad
\begin{tabular}{|c|c|}
\hline
$\delta$ & $\delta(a)$, $a$ nueva \\
\hline
$\exists x A(x)$ & $A(a)$ \\
$\neg \forall x A(x)$ & $\neg A(a)$ \\
\hline
\end{tabular}
\qquad
\begin{tabular}{|c|c|}
\hline
$\sigma$ & \\
\hline
$\neg\neg A$ & $A$ \\
\hline
\end{tabular}
\end{center}

\begin{itemize}
    \item \textbf{$\alpha$, $\delta$, $\sigma$:} se añaden a todas las ramas abiertas que pasan por la fórmula. \textbf{$\beta$:} la rama se bifurca.
    \item \textbf{$\gamma$:} se instancia con todas las constantes de la rama (con una nueva si no hay ninguna), y de nuevo cada vez que aparezca una constante nueva.
    \item \textbf{Cierre:} una rama se cierra si contiene $A$ y $\neg A$.
    \item \textbf{Prioridad orientativa:} $\alpha$ (y $\sigma$), después $\delta$, después $\gamma$ y, al final, $\beta$.
    \item \textbf{Usos:} para $\Gamma \models A$, raíz $\Gamma \cup \{\neg A\}$; para $\models A$, raíz $\{\neg A\}$; para la satisfacibilidad de $\Gamma$, raíz $\Gamma$; para $A \equiv B$, raíz $\{\neg(A \leftrightarrow B)\}$. Tableau cerrado: consecuencia válida, fórmula válida, conjunto insatisfacible o fórmulas equivalentes. Rama abierta y completa: contramodelo o modelo.
\end{itemize}

%%% Local Variables:
%%% mode: latex
%%% TeX-master: "../ApuntesLogica"
%%% End:


\chapter{Soluciones de ejercicios seleccionados}\label{ap:soluciones}
% -*- coding: utf-8 -*-
Este apéndice recoge soluciones de una selección de ejercicios de los temas~\ref{cap:dn-prop}, \ref{cap:sem-prop}, \ref{cap:dn-lpo}, \ref{cap:sem-lpo} y~\ref{cap:tableaux}. Las derivaciones no son únicas: si tu solución es distinta, comprueba con el asistente web que cada línea está bien justificada. Los ejercicios del tema~\ref{cap:meta-prop} ya incluyen soluciones guiadas en el propio tema.

\section{Tema 1: deducción natural en lógica proposicional}

\begin{solucion}{e1-10}
\begin{center}
$\begin{nd}[rules=myrules]
    \hypo [1] {1} {(p \land q) \to r} \by{Prem.}{}
    \open
        \hypo [2] {2} {p} \by{Hip.}{}
        \open
            \hypo [3] {3} {q} \by{Hip.}{}
            \have [4] {4} {p \land q} \ai{2, 3}
            \have [5] {5} {r} \ce{1, 4}
        \close
        \have [6] {6} {q \to r} \ci{3-5}
    \close
    \have [7] {7} {p \to (q \to r)} \ci{2-6}
\end{nd}$
\end{center}
\end{solucion}

\begin{solucion}{e1-15}
\begin{center}
$\begin{nd}[rules=myrules]
    \hypo [1] {1} {p \lor q} \by{Prem.}{}
    \hypo [2] {2} {\lnot p} \by{Prem.}{}
    \open
        \hypo [3] {3} {p} \by{Hip.}{}
        \have [4] {4} {\bot} \ne{2, 3}
        \have [5] {5} {q} \x{4}
    \close
    \open
        \hypo [6] {6} {q} \by{Hip.}{}
        \have [7] {7} {q} \r{6}
    \close
    \have [8] {8} {q} \oe{1, 3-5, 6-7}
\end{nd}$
\end{center}
\end{solucion}

\begin{solucion}{e1-16}
Primero derivamos $\neg p$ (si $p$, la premisa da $\neg p$) y después, con la premisa, $p$:
\begin{center}
$\begin{nd}[rules=myrules]
    \hypo [1] {1} {p \leftrightarrow \lnot p} \by{Prem.}{}
    \open
        \hypo [2] {2} {p} \by{Hip.}{}
        \have [3] {3} {\lnot p} \be{1, 2}
        \have [4] {4} {\bot} \ne{3, 2}
    \close
    \have [5] {5} {\lnot p} \ni{2-4}
    \have [6] {6} {p} \be{1, 5}
    \have [7] {7} {\bot} \ne{5, 6}
\end{nd}$
\end{center}
\textbf{Comentario.} Si $p$ es «esta oración es falsa», la premisa dice que la oración es verdadera si y solo si es falsa. La derivación muestra que esa premisa es contradictoria: es la estructura lógica de la paradoja del mentiroso.
\end{solucion}

\begin{solucion}{e1-18}
\begin{center}
$\begin{nd}[rules=myrules]
    \hypo [1] {1} {\lnot(p \lor q)} \by{Prem.}{}
    \open
        \hypo [2] {2} {p} \by{Hip.}{}
        \have [3] {3} {p \lor q} \oi{2}
        \have [4] {4} {\bot} \ne{1, 3}
    \close
    \have [5] {5} {\lnot p} \ni{2-4}
    \open
        \hypo [6] {6} {q} \by{Hip.}{}
        \have [7] {7} {p \lor q} \oi{6}
        \have [8] {8} {\bot} \ne{1, 7}
    \close
    \have [9] {9} {\lnot q} \ni{6-8}
    \have [10] {10} {\lnot p \land \lnot q} \ai{5, 9}
\end{nd}$
\end{center}
\end{solucion}

\begin{solucion}{e1-22}
Sin usar LEM, suponemos la negación de la conclusión y aplicamos IP:
\begin{center}
$\begin{nd}[rules=myrules]
    \open
        \hypo [1] {1} {\lnot(p \lor \lnot p)} \by{Hip. IP}{}
        \open
            \hypo [2] {2} {p} \by{Hip.}{}
            \have [3] {3} {p \lor \lnot p} \oi{2}
            \have [4] {4} {\bot} \ne{1, 3}
        \close
        \have [5] {5} {\lnot p} \ni{2-4}
        \have [6] {6} {p \lor \lnot p} \oi{5}
        \have [7] {7} {\bot} \ne{1, 6}
    \close
    \have [8] {8} {p \lor \lnot p} \ip{1-7}
\end{nd}$
\end{center}
\end{solucion}

\begin{solucion}{e1-24}
\begin{center}
$\begin{nd}[rules=myrules]
    \hypo [1] {1} {(p \to q) \to p} \by{Prem.}{}
    \open
        \hypo [2] {2} {\lnot p} \by{Hip. IP}{}
        \open
            \hypo [3] {3} {p} \by{Hip.}{}
            \have [4] {4} {\bot} \ne{2, 3}
            \have [5] {5} {q} \x{4}
        \close
        \have [6] {6} {p \to q} \ci{3-5}
        \have [7] {7} {p} \ce{1, 6}
        \have [8] {8} {\bot} \ne{2, 7}
    \close
    \have [9] {9} {p} \ip{2-8}
\end{nd}$
\end{center}
\textbf{Comentario.} La ley de Peirce no puede demostrarse sin una regla clásica como IP (no es válida en lógica intuicionista).
\end{solucion}

\begin{solucion}{e1-29}
Una derivación sistemática se obtiene con el método de Kalmár: véase el ejemplo~\ref{ej:kalmar-linealidad} del tema~\ref{cap:meta-prop}. Una más breve usa LEM sobre $q$: si $q$, se obtiene $p \to q$ (suponiendo $p$ y reiterando $q$); si $\neg q$, se obtiene $q \to p$ (suponiendo $q$, $\bot$ con $\neg q$ y X); en ambos casos se obtiene la disyunción por $\lor$I, y LEM permite concluirla fuera de las cajas.
\end{solucion}

\section{Tema 2: semántica de la lógica proposicional}

\begin{solucion}{e2-negcond}
$p \land q \equiv \neg(p \to \neg q)$:
\begin{center}
\begin{tabular}{cc|c|c|c|c}
$p$ & $q$ & $\neg q$ & $p \to \neg q$ & $\neg(p \to \neg q)$ & $p \land q$ \\
\hline
1 & 1 & 0 & 0 & 1 & 1 \\
1 & 0 & 1 & 1 & 0 & 0 \\
0 & 1 & 0 & 1 & 0 & 0 \\
0 & 0 & 1 & 1 & 0 & 0
\end{tabular}
\end{center}
\end{solucion}

\begin{solucion}{e2-negdisy}
$p \leftrightarrow q \equiv \neg(\neg p \lor \neg q) \lor \neg(p \lor q)$. En efecto, por De Morgan $\neg(\neg p \lor \neg q) \equiv p \land q$ y $\neg(p \lor q) \equiv \neg p \land \neg q$, y $p \leftrightarrow q$ es verdadera exactamente cuando $p$ y $q$ son ambas verdaderas o ambas falsas.
\end{solucion}

\begin{solucion}{e2-consec}
\begin{enumerate}[(a)]
    \item No es válida. Contraejemplo: $v(p) = 0$, $v(q) = 1$ (las premisas son verdaderas y $\neg q$ es falsa).
    \item No es válida. Contraejemplo: $v(p) = 1$, $v(q) = 1$, $v(r) = 0$.
    \item Es válida (si $v(\neg p) = 1$, la única forma de que $p \lor q$ sea verdadera es $v(q) = 1$). La derivación es la del ejercicio~\ref{e1-15}.
    \item Es válida. Para un contraejemplo necesitaríamos $v(q) = 1$, $v(p) = 1$ y $v(r) = 0$, pero entonces $v(q \to r) = 0$ y $v(p \to (q \to r)) = 0$. Derivación:
    \begin{center}
    $\begin{nd}[rules=myrules]
        \hypo [1] {1} {p \to (q \to r)} \by{Prem.}{}
        \hypo [2] {2} {p \land \lnot r} \by{Prem.}{}
        \have [3] {3} {p} \ae{2}
        \have [4] {4} {\lnot r} \ae{2}
        \have [5] {5} {q \to r} \ce{1, 3}
        \open
            \hypo [6] {6} {q} \by{Hip.}{}
            \have [7] {7} {r} \ce{5, 6}
            \have [8] {8} {\bot} \ne{4, 7}
        \close
        \have [9] {9} {\lnot q} \ni{6-8}
    \end{nd}$
    \end{center}
    \item Es válida. Para un contraejemplo necesitaríamos $v(p \lor r) = 1$ y $v(q) = 0$. Pero entonces $v(p) = 1$ o $v(r) = 1$, y en ambos casos una de las premisas, $p \to q$ o $r \to q$, sería falsa. Derivación:
    \begin{center}
    $\begin{nd}[rules=myrules]
        \hypo [1] {1} {p \to q} \by{Prem.}{}
        \hypo [2] {2} {r \to q} \by{Prem.}{}
        \open
            \hypo [3] {3} {p \lor r} \by{Hip.}{}
            \open
                \hypo [4] {4} {p} \by{Hip.}{}
                \have [5] {5} {q} \ce{1, 4}
            \close
            \open
                \hypo [6] {6} {r} \by{Hip.}{}
                \have [7] {7} {q} \ce{2, 6}
            \close
            \have [8] {8} {q} \oe{3, 4-5, 6-7}
        \close
        \have [9] {9} {(p \lor r) \to q} \ci{3-8}
    \end{nd}$
    \end{center}
\end{enumerate}
\end{solucion}

\begin{solucion}{e2-formal}
\begin{enumerate}[(a)]
    \item $f$: «iré a la fiesta»; $i$: «me invitan». $f \to i$.
    \item $e$: «apruebo el examen»; $a$: «apruebo la asignatura». $a \to e$. Esta fórmula dice que aprobar el examen es condición necesaria para aprobar la asignatura, y no dice que sea suficiente. La segunda parte del enunciado no puede añadirse como $\neg(e \to a)$, porque esa fórmula equivale a $e \land \neg a$, que afirma algo que el enunciado no dice: que apruebo el examen y suspendo la asignatura. Que una condición $C$ no sea suficiente para $P$ es una afirmación general (puede cumplirse $C$ sin que se cumpla $P$) que la lógica proposicional no recoge; en lógica de primer orden se expresaría con cuantificadores.
    \item $s$: «saldremos de excursión»; $l$: «llueve». $\neg l \to s$ (o $s \lor l$).
    \item $k$: «Kant acepta el argumento ontológico»; $h$: «Hume acepta el argumento ontológico». $\neg k \land \neg h$.
    \item $n$: «el conocimiento es innato»; $a$: «el conocimiento es adquirido por la experiencia». $\neg(n \land a)$. (El «no puede» modal se pierde en la formalización).
    \item $e$: «Hume es empirista»; $r$: «Hume rechaza las ideas innatas». $e \leftrightarrow r$.
\end{enumerate}
\end{solucion}

\begin{solucion}{e2-arg}
Apartados (a)--(c):
\begin{enumerate}[(a)]
    \item $c$: «voy al cine»; $l$: «acabo pronto de limpiar la casa»; $a$: «me acompañas». Premisas $c \to l$ y $a \to c$; conclusión $a \to l$. Es válido (silogismo hipotético): para un contraejemplo necesitaríamos $v(a) = 1$ y $v(l) = 0$; entonces la segunda premisa exige $v(c) = 1$, y la primera sería falsa. Derivación:
    \begin{center}
    $\begin{nd}[rules=myrules]
        \hypo [1] {1} {c \to l} \by{Prem.}{}
        \hypo [2] {2} {a \to c} \by{Prem.}{}
        \open
            \hypo [3] {3} {a} \by{Hip.}{}
            \have [4] {4} {c} \ce{2, 3}
            \have [5] {5} {l} \ce{1, 4}
        \close
        \have [6] {6} {a \to l} \ci{3-5}
    \end{nd}$
    \end{center}
    \item $b$: «hace buen tiempo»; $p$: «iremos a la playa». Premisas $b \to p$ y $p$; conclusión $b$. No es válido (afirmación del consecuente): la fila $v(b) = 0$, $v(p) = 1$ es un contraejemplo.
    \item $m$: «el mundo es redondo»; $b$: «tengo un barco»; $n$: «navego alrededor del mundo». Premisas $(m \land b) \to n$ y $b \land \neg n$; conclusión $\neg m$. Es válido: para un contraejemplo necesitaríamos $v(m) = 1$; la segunda premisa exige $v(b) = 1$ y $v(n) = 0$, y entonces la primera sería falsa. Derivación:
    \begin{center}
    $\begin{nd}[rules=myrules]
        \hypo [1] {1} {(m \land b) \to n} \by{Prem.}{}
        \hypo [2] {2} {b \land \lnot n} \by{Prem.}{}
        \have [3] {3} {b} \ae{2}
        \have [4] {4} {\lnot n} \ae{2}
        \open
            \hypo [5] {5} {m} \by{Hip.}{}
            \have [6] {6} {m \land b} \ai{5, 3}
            \have [7] {7} {n} \ce{1, 6}
            \have [8] {8} {\bot} \ne{4, 7}
        \close
        \have [9] {9} {\lnot m} \ni{5-8}
    \end{nd}$
    \end{center}
\end{enumerate}
\end{solucion}

\begin{solucion}{e2-vf2}
Es falsa. Si $\Gamma$ es insatisfacible, no hay ninguna interpretación que haga verdaderas a la vez todas sus fórmulas, de modo que $\Gamma \models A$ y $\Gamma \models \neg A$ se cumplen ambas (por vacuidad). Por ejemplo, con $\Gamma = \{p, \neg p\}$ tenemos $\Gamma \models q$ y $\Gamma \models \neg q$. La afirmación sí es verdadera si $\Gamma$ es satisfacible.
\end{solucion}

\begin{solucion}{e2-cuest}
\begin{enumerate}[(a)]
    \item $\{A\} \models B$ es verdadera: $B$ es verdadera en toda interpretación (y, además, ninguna interpretación hace verdadera a $A$). $\{B\} \models A$ es falsa: cualquier interpretación hace verdadera a $B$ y falsa a $A$.
    \item $\{A, B\} \models C$ se cumple: si una interpretación hiciera verdaderas $A$ y $B$ y falsa $C$, satisfaría $\{A, B, \neg C\}$. Por la misma razón se cumple $\{A, \neg C\} \models \neg B$. En cambio, $\{A, B\} \models \neg C$ no se cumple necesariamente: con $A = p$, $B = q$ y $C = p$, el conjunto $\{p, q, \neg p\}$ es insatisfacible, pero la interpretación $v(p) = v(q) = 1$ hace verdaderas $p$ y $q$ y falsa $\neg p$.
\end{enumerate}
\end{solucion}

\section{Tema 4: deducción natural en lógica de primer orden}

\begin{solucion}{e4-1}
\begin{center}
$\begin{nd}[rules=myrules]
    \hypo [1] {1} {\forall x (M(x) \to \lnot P(x))} \by{Prem.}{}
    \hypo [2] {2} {\forall x (S(x) \to M(x))} \by{Prem.}{}
    \open
        \hypo [3] {3} {S(c)} \by{Hip.}{}
        \have [4] {4} {S(c) \to M(c)} \Ae{2}
        \have [5] {5} {M(c)} \ce{4, 3}
        \have [6] {6} {M(c) \to \lnot P(c)} \Ae{1}
        \have [7] {7} {\lnot P(c)} \ce{6, 5}
    \close
    \have [8] {8} {S(c) \to \lnot P(c)} \ci{3-7}
    \have [9] {9} {\forall x (S(x) \to \lnot P(x))} \Ai{8}
\end{nd}$
\end{center}
\end{solucion}

\begin{solucion}{e4-8}
\begin{center}
$\begin{nd}[rules=myrules]
    \hypo [1] {1} {\exists x (P(x) \lor Q(x))} \by{Prem.}{}
    \open
        \hypo [2] {2} {P(c) \lor Q(c)} \by{Hip. testigo $c$}{}
        \open
            \hypo [3] {3} {P(c)} \by{Hip.}{}
            \have [4] {4} {\exists x P(x)} \Ei{3}
            \have [5] {5} {\exists x P(x) \lor \exists x Q(x)} \oi{4}
        \close
        \open
            \hypo [6] {6} {Q(c)} \by{Hip.}{}
            \have [7] {7} {\exists x Q(x)} \Ei{6}
            \have [8] {8} {\exists x P(x) \lor \exists x Q(x)} \oi{7}
        \close
        \have [9] {9} {\exists x P(x) \lor \exists x Q(x)} \oe{2, 3-5, 6-8}
    \close
    \have [10] {10} {\exists x P(x) \lor \exists x Q(x)} \Ee{1, 2-9}
\end{nd}$
\end{center}
\end{solucion}

\begin{solucion}{e4-12}
\begin{center}
$\begin{nd}[rules=myrules]
    \hypo [1] {1} {\forall x \forall y (R(x, y) \to R(y, x))} \by{Prem.}{}
    \hypo [2] {2} {\forall x \forall y \forall z ((R(x, y) \land R(y, z)) \to R(x, z))} \by{Prem.}{}
    \hypo [3] {3} {\forall x \exists y R(x, y)} \by{Prem.}{}
    \have [4] {4} {\exists y R(a, y)} \Ae{3}
    \open
        \hypo [5] {5} {R(a, b)} \by{Hip. testigo $b$}{}
        \have [6] {6} {\forall y (R(a, y) \to R(y, a))} \Ae{1}
        \have [7] {7} {R(a, b) \to R(b, a)} \Ae{6}
        \have [8] {8} {R(b, a)} \ce{7, 5}
        \have [9] {9} {\forall y \forall z ((R(a, y) \land R(y, z)) \to R(a, z))} \Ae{2}
        \have [10] {10} {\forall z ((R(a, b) \land R(b, z)) \to R(a, z))} \Ae{9}
        \have [11] {11} {(R(a, b) \land R(b, a)) \to R(a, a)} \Ae{10}
        \have [12] {12} {R(a, b) \land R(b, a)} \ai{5, 8}
        \have [13] {13} {R(a, a)} \ce{11, 12}
    \close
    \have [14] {14} {R(a, a)} \Ee{4, 5-13}
    \have [15] {15} {\forall x R(x, x)} \Ai{14}
\end{nd}$
\end{center}
\textbf{Comprobación de las restricciones:} el testigo $b$ no aparece en las premisas, ni en $\exists y R(a, y)$, ni en la conclusión $R(a, a)$ de la caja; y $a$ no aparece en ninguna premisa ni hipótesis activa en la línea 15.
\end{solucion}

\begin{solucion}{e4-17}
\begin{center}
$\begin{nd}[rules=myrules]
    \hypo [1] {1} {\forall x \lnot P(x)} \by{Prem.}{}
    \open
        \hypo [2] {2} {\exists x P(x)} \by{Hip.}{}
        \open
            \hypo [3] {3} {P(c)} \by{Hip. testigo $c$}{}
            \have [4] {4} {\lnot P(c)} \Ae{1}
            \have [5] {5} {\bot} \ne{4, 3}
        \close
        \have [6] {6} {\bot} \Ee{2, 3-5}
    \close
    \have [7] {7} {\lnot \exists x P(x)} \ni{2-6}
\end{nd}$
\end{center}
\end{solucion}

\begin{solucion}{e4-20}
\begin{center}
$\begin{nd}[rules=myrules]
    \open
        \hypo [1] {1} {\forall y B(y)} \by{Hip.}{}
        \open
            \hypo [2] {2} {B(a)} \by{Hip.}{}
            \have [3] {3} {\forall y B(y)} \r{1}
        \close
        \have [4] {4} {B(a) \to \forall y B(y)} \ci{2-3}
        \have [5] {5} {\exists x (B(x) \to \forall y B(y))} \Ei{4}
    \close
    \open
        \hypo [6] {6} {\lnot \forall y B(y)} \by{Hip.}{}
        \have [7] {7} {\exists y \lnot B(y)} \cq{6}
        \open
            \hypo [8] {8} {\lnot B(d)} \by{Hip. testigo $d$}{}
            \open
                \hypo [9] {9} {B(d)} \by{Hip.}{}
                \have [10] {10} {\bot} \ne{8, 9}
                \have [11] {11} {\forall y B(y)} \x{10}
            \close
            \have [12] {12} {B(d) \to \forall y B(y)} \ci{9-11}
            \have [13] {13} {\exists x (B(x) \to \forall y B(y))} \Ei{12}
        \close
        \have [14] {14} {\exists x (B(x) \to \forall y B(y))} \Ee{7, 8-13}
    \close
    \have [15] {15} {\exists x (B(x) \to \forall y B(y))} \lem{1-5, 6-14}
\end{nd}$
\end{center}
\end{solucion}

\begin{solucion}{e4-24}
La infracción está en la línea 4: se aplica $\forall$I sobre $d$, pero $d$ aparece en la hipótesis activa $R(c, d)$ (línea 3), lo que viola la restricción (1) de $\forall$I. Contraejemplo informal: «todo el mundo tiene una madre» ($\forall y \exists x R(x, y)$, con $R(x, y)$: «$x$ es madre de $y$») es verdadero, pero «hay alguien que es madre de todo el mundo» ($\exists x \forall y R(x, y)$) es falso.
\end{solucion}

\begin{solucion}{e4-27}
\begin{center}
$\begin{nd}[rules=myrules]
    \hypo [1] {1} {\exists x (P(x) \land \forall y (P(y) \to y = x))} \by{Prem.}{}
    \hypo [2] {2} {P(a)} \by{Prem.}{}
    \hypo [3] {3} {P(b)} \by{Prem.}{}
    \open
        \hypo [4] {4} {P(c) \land \forall y (P(y) \to y = c)} \by{Hip. testigo $c$}{}
        \have [5] {5} {\forall y (P(y) \to y = c)} \ae{4}
        \have [6] {6} {P(a) \to a = c} \Ae{5}
        \have [7] {7} {a = c} \ce{6, 2}
        \have [8] {8} {P(b) \to b = c} \Ae{5}
        \have [9] {9} {b = c} \ce{8, 3}
        \have [10] {10} {a = b} \ide{9, 7}
    \close
    \have [11] {11} {a = b} \Ee{1, 4-10}
\end{nd}$
\end{center}
\textbf{Comentario.} En la línea 10 usamos $b = c$ para sustituir $c$ por $b$ en $a = c$ (sustitución en sentido inverso). La conclusión no contiene el testigo $c$, de modo que la aplicación de $\exists$E es correcta.
\end{solucion}

\begin{solucion}{e4-28}
\begin{center}
$\begin{nd}[rules=myrules]
    \open
        \hypo [1] {1} {a = b} \by{Hip.}{}
        \have [2] {2} {a = a} \idi{}
        \have [3] {3} {b = a} \ide{1, 2}
    \close
    \have [4] {4} {a = b \to b = a} \ci{1-3}
    \have [5] {5} {\forall y (a = y \to y = a)} \Ai{4}
    \have [6] {6} {\forall x \forall y (x = y \to y = x)} \Ai{5}
\end{nd}$
\end{center}
\end{solucion}

\section{Tema 5: semántica de la lógica de primer orden}

\begin{solucion}{e5-2}
Algunos apartados (hay otras formalizaciones equivalentes):
\begin{itemize}
    \item[4.] $\exists x (A(b, x) \land x \neq a)$.
    \item[5.] $\forall x (A(a, x) \to A(x, a))$.
    \item[6.] $\exists x (x \neq b \land A(b, x) \land \forall y ((A(b, y) \land y \neq b) \to y = x))$.
    \item[7.] $\exists x \exists y (x \neq y \land A(a, x) \land A(a, y) \land \forall z (A(a, z) \to (z = x \lor z = y)))$.
    \item[9.] $\forall x \forall y (A(x, y) \to A(y, x))$.
    \item[10.] $\forall x (\exists y A(x, y) \to \exists z A(z, x))$.
    \item[12.] $\forall x \exists y (y \neq x \land A(x, y))$.
\end{itemize}
\end{solucion}

\begin{solucion}{e5-3}
Algunos apartados:
\begin{itemize}
    \item[1.] $\neg \exists x (E(x) \land F(x, l))$.
    \item[2.] $\exists x \exists y \exists z (E(x) \land E(y) \land x \neq y \land F(x, z) \land F(y, z))$.
    \item[4.] $\forall y (S(y) \to (A(a, y) \leftrightarrow A(b, y)))$.
    \item[7.] $\forall y ((S(y) \land \neg A(a, y)) \to A(b, y))$ (entendiendo que suspender es no aprobar).
\end{itemize}
\end{solucion}

\begin{solucion}{e5-anid}
Diccionario: $F(x)$: «$x$ es filósofo»; $C(x)$: «$x$ es un clásico»; $L(x, y)$: «$x$ ha leído a $y$»; $E(x)$: «$x$ es estudiante»; $S(x)$: «$x$ es una asignatura»; $A(x, y)$: «$x$ aprueba $y$»; $P(x)$: «$x$ es profesor»; $U(x, y)$: «$x$ es alumno de $y$»; $N(x, y)$: «$x$ entiende a $y$».
\begin{enumerate}
    \item $\forall x (F(x) \to \exists y (C(y) \land L(x, y)))$.
    \item $\exists y (C(y) \land \forall x (F(x) \to L(x, y)))$.
    \item $\neg \exists x (E(x) \land \forall y (S(y) \to A(x, y)))$, o bien $\forall x (E(x) \to \exists y (S(y) \land \neg A(x, y)))$.
    \item $\forall y (S(y) \to \exists x (E(x) \land \neg A(x, y)))$.
    \item $\exists x (P(x) \land \neg \exists y (U(y, x) \land N(y, x)))$, o bien $\exists x (P(x) \land \forall y (U(y, x) \to \neg N(y, x)))$.
\end{enumerate}
El enunciado 2 implica el 1: si hay un clásico que han leído todos los filósofos, cada filósofo ha leído al menos ese clásico. El 1 no implica el 2: en la estructura con dominio $\{1, 2, 3, 4\}$, $I(F) = \{1, 2\}$, $I(C) = \{3, 4\}$ e $I(L) = \{(1, 3), (2, 4)\}$, cada filósofo ha leído un clásico, pero no hay ningún clásico que hayan leído los dos.
\end{solucion}

\begin{solucion}{e5-cant}
Diccionario: $F(x)$: «$x$ es un filósofo presocrático»; $b$: Belén; $S(x)$: «$x$ es una asignatura»; $A(x, y)$: «$x$ aprueba $y$» (suspender es no aprobar); $E(x, y)$: «$x$ escribió $y$»; $n$: la \emph{Ética a Nicómaco}; $U(x, y)$: «$x$ fue discípulo de $y$»; $p$: Platón; $O(x)$: «$x$ es una montaña de oro».
\begin{enumerate}
    \item $\exists x \exists y \exists z (F(x) \land F(y) \land F(z) \land x \neq y \land x \neq z \land y \neq z)$.
    \item $\forall x \forall y ((S(x) \land \neg A(b, x) \land S(y) \land \neg A(b, y)) \to x = y)$.
    \item $\exists x \exists y (S(x) \land \neg A(b, x) \land S(y) \land \neg A(b, y) \land x \neq y \land \forall z ((S(z) \land \neg A(b, z)) \to (z = x \lor z = y)))$.
    \item $\exists x (E(x, n) \land \forall y (E(y, n) \to y = x) \land U(x, p))$.
    \item $\neg \exists x (O(x) \land \forall y (O(y) \to y = x))$: no hay exactamente una montaña de oro. Para Russell, decir que «la montaña de oro» no existe no es atribuir a un objeto la propiedad de no existir, sino negar que la descripción se aplique a uno y solo un objeto.
\end{enumerate}
\end{solucion}

\begin{solucion}{e5-6}
\begin{enumerate}
    \item Verdadera: para cada $n$ basta tomar $y = n$ (o $n + 1$).
    \item Verdadera: $y = 0$ es menor o igual que todo natural.
    \item Verdadera: $\le$ es un orden total.
    \item Verdadera: para todo $x$ (par o no) hay un primo mayor o igual que $x$, porque hay infinitos primos.
    \item Falsa: para $x = 0$ e $y = 1$ se cumple $0 \le 1$ y $0 \neq 1$, pero no hay ningún natural $z$ distinto de $0$ y de $1$ con $0 \le z \le 1$.
\end{enumerate}
\end{solucion}

\begin{solucion}{e5-7}
Con $D = \{1, 2\}$, estas estructuras bastan (V: verdadera; F: falsa):
\begin{center}
\begin{tabular}{l|ccccc}
$I(R)$ & 1 & 2 & 3 & 4 & 5 \\
\hline
$\{(1, 1)\}$ & V & F & F & F & F \\
$\{(1, 1), (1, 2)\}$ & V & F & V & F & F \\
$\{(1, 1), (2, 2)\}$ & V & F & F & V & F \\
$\{(1, 1), (2, 1)\}$ & V & F & F & V & V \\
$D \times D$ & V & V & V & V & V
\end{tabular}
\end{center}
No hay dos columnas iguales, de modo que ninguna fórmula es equivalente a otra. Relaciones de consecuencia: la 2 implica todas las demás; la 3, la 4 y la 5 implican la 1 (el dominio no es vacío); la 5 implica la 4 (si un mismo $y$ sirve para todos los $x$, cada $x$ tiene alguno). Entre fórmulas tomadas de una en una, la tabla muestra que no hay más: por ejemplo, la segunda fila muestra que la 3 no implica la 4 ni la 5, y la cuarta que la 5 no implica la 3. Si se toman varias premisas a la vez hay que comprobarlo aparte, y a veces dos elementos no bastan. Así, en un dominio de dos elementos la 3 y la 4 juntas hacen verdadera la 5, pero no la implican: con $D = \{1, 2, 3\}$ e $I(R) = \{(1, 1), (1, 2), (1, 3), (2, 3), (3, 2)\}$, la 3 y la 4 son verdaderas y la 5 es falsa.
\end{solucion}

\begin{solucion}{e5-8}
\begin{enumerate}
    \item $D = \{1, 2\}$, $I(P) = \{1\}$, $I(Q) = \{2\}$: todo elemento es $P$ o $Q$, pero no todos son $P$ ni todos son $Q$.
    \item $D = \{1\}$, $I(P) = \varnothing$, $I(Q) = \{1\}$: la premisa universal se cumple, $1$ no es $P$, pero todo es $Q$.
    \item La misma estructura: hay un $Q$, pero ningún $P$.
    \item La misma estructura con $I(a) = 1$: $\neg P(a)$ es verdadera y $\neg Q(a)$ es falsa.
\end{enumerate}
\end{solucion}

\begin{solucion}{e5-9}
\begin{enumerate}
    \item Al menos dos elementos.
    \item Exactamente dos elementos.
    \item Un dominio infinito. El segundo miembro de la conjunción (asimetría) implica que nadie se relaciona consigo mismo. Por el primero, partiendo de un elemento $d_0$ hay $d_1$ con $R(d_0, d_1)$, $d_2$ con $R(d_1, d_2)$, etc. Por transitividad, $R(d_i, d_j)$ siempre que $i < j$; si fuera $d_i = d_j$ con $i < j$, tendríamos $R(d_i, d_i)$, lo que es imposible. Luego todos los $d_i$ son distintos. Un modelo es $\mathbb{N}$ interpretando $R$ como la relación «menor que» ($m < n$).
\end{enumerate}
\end{solucion}

\begin{solucion}{e5-10}
\begin{enumerate}
    \item Basta un elemento: $D = \{1\}$, todas las constantes denotan $1$ e $I(P) = \{1\}$. Nada obliga a que constantes distintas denoten objetos distintos.
    \item Hacen falta dos elementos: si $d$ es $Q$, por la segunda fórmula no es $P$, mientras que $a$ sí es $P$; luego $d \neq I(a)$. Un modelo: $D = \{1, 2\}$, $I(a) = 1$, $I(P) = \{1\}$, $I(Q) = \{2\}$, $I(R) = \{(2, 2)\}$ (la primera fórmula exige $R(2, 2)$ porque $2$ no es $P$).
    \item Hacen falta tres elementos (lo exige la primera fórmula); las demás dicen que $R$ es una relación de equivalencia. Un modelo: $D = \{1, 2, 3\}$ e $I(R) = D \times D$.
\end{enumerate}
\end{solucion}

\begin{solucion}{e5-11}
\begin{enumerate}
    \item $D = \{1\}$ e $I(M) = I(S) = I(P) = \varnothing$: las premisas son verdaderas porque no hay ningún $M$, y la conclusión es falsa.
    \item Derivación:
    \begin{center}
    $\begin{nd}[rules=myrules]
        \hypo [1] {1} {\forall x (M(x) \to P(x))} \by{Prem.}{}
        \hypo [2] {2} {\forall x (M(x) \to S(x))} \by{Prem.}{}
        \hypo [3] {3} {\exists x M(x)} \by{Prem.}{}
        \open
            \hypo [4] {4} {M(c)} \by{Hip. testigo $c$}{}
            \have [5] {5} {M(c) \to P(c)} \Ae{1}
            \have [6] {6} {P(c)} \ce{5, 4}
            \have [7] {7} {M(c) \to S(c)} \Ae{2}
            \have [8] {8} {S(c)} \ce{7, 4}
            \have [9] {9} {S(c) \land P(c)} \ai{8, 6}
            \have [10] {10} {\exists x (S(x) \land P(x))} \Ei{9}
        \close
        \have [11] {11} {\exists x (S(x) \land P(x))} \Ee{3, 4-10}
    \end{nd}$
    \end{center}
\end{enumerate}
\end{solucion}

\begin{solucion}{e5-12}
\begin{enumerate}
    \item O es verdadera exactamente cuando hay algún $S$ que no es $P$, es decir, exactamente cuando A es falsa; luego A y O son contradictorias. Del mismo modo, I es verdadera exactamente cuando E es falsa.
    \item En $D = \{1\}$ con $I(S) = \varnothing$ (y $I(P)$ cualquiera), A y E son verdaderas, e I y O son falsas.
    \item La premisa $\exists x S(x)$: con ella, de A se deduce I (como en el ejercicio~\ref{e5-11}).
\end{enumerate}
\end{solucion}

\section{Tema 6: el cálculo de tableaux}

En las soluciones de este tema indicamos cada nodo con su número y, entre paréntesis, la regla aplicada y el nodo al que se aplica; una barra vertical ($\mid$) separa los dos nodos que abre una regla $\beta$ (cada uno con su número), y «la rama de $m$ se cierra con $n$» indica que las fórmulas de los nodos $m$ y $n$ son una la negación de la otra, de modo que la rama se cierra ($\bigotimes$).

\begin{solucion}{e6-1}
\textbf{Apartado 1.} Raíz: 1. $p \to \neg p$; 2. $q \to p$; 3. $p \lor q$.
\begin{itemize}
    \item 4. $\neg q \mid$ 5. $p$ ($\beta$, 2).
    \item Rama de 5: 6. $\neg p \mid$ 7. $\neg p$ ($\beta$, 1); ambas ramas se cierran con 5.
    \item Rama de 4: 8. $p \mid$ 9. $q$ ($\beta$, 3). La rama de 9 se cierra con 4. En la de 8, 10. $\neg p \mid$ 11. $\neg p$ ($\beta$, 1), y ambas se cierran con 8.
\end{itemize}
Todas las ramas se cierran: $\Gamma_1$ es insatisfacible.

\noindent\textbf{Apartado 2.} Raíz: 1. $p \land \neg q$; 2. $\neg(p \to r)$; 3. $q \lor \neg r$. 4. $p$ y 5. $\neg q$ ($\alpha$, 1); 6. $p$ y 7. $\neg r$ ($\alpha$, 2); 8. $q \mid$ 9. $\neg r$ ($\beta$, 3). La rama de 8 se cierra con 5; la de 9 queda abierta y completa. $\Gamma_2$ es satisfacible: $v(p) = 1$, $v(q) = 0$, $v(r) = 0$.
\end{solucion}

\begin{solucion}{e6-4}
\textbf{Celarent.} Raíz: 1. $\forall x (M(x) \to \neg P(x))$; 2. $\forall x (S(x) \to M(x))$; 3. $\neg \forall x (S(x) \to \neg P(x))$.
\begin{itemize}
    \item 4. $\neg(S(a) \to \neg P(a))$ ($\delta$, 3, con $a$ nueva); 5. $S(a)$ y 6. $\neg\neg P(a)$ ($\alpha$, 4); 7. $P(a)$ ($\sigma$, 6).
    \item 8. $M(a) \to \neg P(a)$ ($\gamma$, 1); 9. $S(a) \to M(a)$ ($\gamma$, 2).
    \item 10. $\neg S(a) \mid$ 11. $M(a)$ ($\beta$, 9). La rama de 10 se cierra con 5.
    \item En la rama de 11: 12. $\neg M(a) \mid$ 13. $\neg P(a)$ ($\beta$, 8). La rama de 12 se cierra con 11 y la de 13 con 7.
\end{itemize}
El tableau se cierra: el argumento es válido.
\end{solucion}

\begin{solucion}{e6-6}
\textbf{Apartado 1.} Raíz: 1. $\forall x (P(x) \lor Q(x))$; 2. $\neg(\forall x P(x) \lor \forall x Q(x))$.
\begin{itemize}
    \item 3. $\neg \forall x P(x)$ y 4. $\neg \forall x Q(x)$ ($\alpha$, 2).
    \item 5. $\neg P(a)$ ($\delta$, 3); 6. $\neg Q(b)$ ($\delta$, 4).
    \item 7. $P(a) \lor Q(a)$ y 8. $P(b) \lor Q(b)$ ($\gamma$, 1).
    \item 9. $P(a) \mid$ 10. $Q(a)$ ($\beta$, 7). La rama de 9 se cierra con 5.
    \item En la rama de 10: 11. $P(b) \mid$ 12. $Q(b)$ ($\beta$, 8). La rama de 12 se cierra con 6; la de 11 queda abierta y completa.
\end{itemize}
Contramodelo extraído de la rama abierta (literales $\neg P(a)$, $\neg Q(b)$, $Q(a)$, $P(b)$): $D = \{a, b\}$, $I(P) = \{b\}$, $I(Q) = \{a\}$. Todo elemento es $P$ o $Q$, pero no todos son $P$ ($a$ no lo es) ni todos son $Q$ ($b$ no lo es). Es la situación de los pares y los impares.
\end{solucion}

\begin{solucion}{e6-8}
Raíz: 1. $\forall x (\exists y R(x, y) \to \forall z R(z, x))$; 2. $\exists x \exists y R(x, y)$; 3. $\neg \forall x R(x, x)$.
\begin{center}
\begin{tikzpicture}[scale=0.9, -, thick, level/.style={sibling distance=6cm,level distance=1.1cm}]
  \node[label=left:{4.},label=right:{$\delta$, 2}] {$\exists y R(a, y)$}
  child {
    node[label=left:{5.},label=right:{$\delta$, 4}] {$R(a, b)$}
    child {
      node[label=left:{6.},label=right:{$\delta$, 3}] {$\neg R(c, c)$}
      child {
        node[label=left:{7.},label=right:{$\gamma$, 1}] {$\exists y R(a, y) \to \forall z R(z, a)$}
        child {
          node[label=left:{8.},label=right:{$\beta$, 7}] {$\neg \exists y R(a, y)$}
          child {
            node[label=left:{10.},label=right:{$\gamma$, 8}] {$\neg R(a, b)$}
            child { node[label=below:{5, 10}] {$\bigotimes$} }
          }
        }
        child {
          node[label=left:{9.},label=right:{$\beta$, 7}] {$\forall z R(z, a)$}
          child {
            node[label=left:{11.},label=right:{$\gamma$, 9}] {$R(c, a)$}
            child {
              node[label=left:{12.},label=right:{$\gamma$, 1}] {$\exists y R(c, y) \to \forall z R(z, c)$}
              child {
                node[label=left:{13.},label=right:{$\beta$, 12}] {$\neg \exists y R(c, y)$}
                child {
                  node[label=left:{15.},label=right:{$\gamma$, 13}] {$\neg R(c, a)$}
                  child { node[label=below:{11, 15}] {$\bigotimes$} }
                }
              }
              child {
                node[label=left:{14.},label=right:{$\beta$, 12}] {$\forall z R(z, c)$}
                child {
                  node[label=left:{16.},label=right:{$\gamma$, 14}] {$R(c, c)$}
                  child { node[label=below:{6, 16}] {$\bigotimes$} }
                }
              }
            }
          }
        }
      }
    }
  };
\end{tikzpicture}
\end{center}
Todas las ramas se cierran: las premisas implican la conclusión. (Como la regla de cierre no exige literales, la rama de 8 podría cerrarse ya con el nodo 4, pues $\neg \exists y R(a, y)$ es la negación de $\exists y R(a, y)$; el paso 10 no es necesario, aunque tampoco es incorrecto).
\end{solucion}

%%% Local Variables:
%%% mode: latex
%%% TeX-master: "../ApuntesLogica"
%%% End:


\end{document}

%%% Local Variables:
%%% mode: latex
%%% TeX-master: t
%%% End:
